% =====================================================================
% THE WORKHOUSE PROGRAM --- master edition
%
% SELF-CONTAINED EDITION.  Generated by paper/master/assemble_selfcontained.py
% from the modular sources in paper/master/.  Do not edit this file: edit
% the section files and re-assemble, or the two will disagree.
%
% Needs no other file.  Build with:
%     tectonic workhouse_master_selfcontained.tex
% or  pdflatex workhouse_master_selfcontained.tex   (twice, for cross-references)
%
% Standard packages only, no shell escape, no images, no .bib.
% =====================================================================
% ===================================================================
% THE WORKHOUSE PROGRAM --- master edition
%
% Driver for the full monograph.  Build:
%     tectonic paper/master/workhouse_master.tex
% or
%     pdflatex ... ; bibtex ... ; pdflatex ... ; pdflatex ...
%
% Part I (sections \ref{part:exact}) is written to be extractable on its
% own: paper/master/workhouse_band_paper.tex compiles exactly those files
% as a standalone paper with no dangling references.
%
% Regenerate the apparatus before building:
%     python paper/master/generate_apparatus.py
%     python paper/master/generate_sources.py
% ===================================================================
\documentclass[11pt,letterpaper]{article}
% ---- begin preamble.tex --------------------------------------------
% Shared preamble for both drivers: the monograph and the extractable Part I.
% Conventions follow paper/workhouse_publication_edition_rev5_2026-08-30.tex
% so that material can move between editions without renaming symbols.

\usepackage[T1]{fontenc}
\usepackage[utf8]{inputenc}
\usepackage{lmodern}
\usepackage[margin=1in,headheight=14pt]{geometry}
\usepackage{amsmath,amssymb,amsthm,mathtools}
\usepackage{booktabs,array,longtable,enumitem}
\usepackage{xcolor,microtype,xurl}
\usepackage{fancyhdr}
\usepackage[hidelinks]{hyperref}

\hypersetup{colorlinks=true,
  linkcolor=blue!45!black, citecolor=blue!45!black, urlcolor=blue!45!black}

\setlength{\emergencystretch}{3em}
\setlength{\parskip}{2pt}
\setlist{itemsep=2pt,topsep=4pt}
\numberwithin{equation}{section}

\theoremstyle{plain}
\newtheorem{theorem}{Theorem}[section]
\newtheorem{proposition}[theorem]{Proposition}
\newtheorem{lemma}[theorem]{Lemma}
\newtheorem{corollary}[theorem]{Corollary}
\newtheorem{conjecture}[theorem]{Conjecture}
\theoremstyle{definition}
\newtheorem{definition}[theorem]{Definition}
\newtheorem{problem}{Problem}
\theoremstyle{remark}
\newtheorem{remark}[theorem]{Remark}
\newtheorem{scopenote}[theorem]{Scope}

% ---------------------------------------------------------------- notation
\newcommand{\R}{\mathbb{R}}
\newcommand{\C}{\mathbb{C}}
\newcommand{\Q}{\mathbb{Q}}
\newcommand{\Z}{\mathbb{Z}}
\newcommand{\N}{\mathbb{N}}
\newcommand{\SU}{\mathrm{SU}}
\newcommand{\Tr}{\operatorname{Tr}}
\newcommand{\spec}{\operatorname{spec}}
\newcommand{\Ric}{\operatorname{Ric}}
\newcommand{\im}{\operatorname{im}}
\newcommand{\rank}{\operatorname{rank}}
\newcommand{\norm}[1]{\left\lVert #1 \right\rVert}
\newcommand{\abs}[1]{\left\lvert #1 \right\rvert}
\newcommand{\ip}[2]{\left\langle #1, #2 \right\rangle}
\newcommand{\Hil}{\mathcal{H}}
\newcommand{\Alg}{\mathcal{A}}
\newcommand{\gp}{\mathfrak{g}}
\newcommand{\Hph}{\Hil_{\mathrm{phys}}}

% ------------------------------------------------- verification-tier badges
% The badge records how a statement is machine-certified in the repository.
% It is INDEPENDENT of mathematical status: a T3 badge on a proven theorem
% means the proof is analytic and no machine check covers the whole
% statement.  Conflating the two is the error this notation exists to stop.
\definecolor{tierzero}{RGB}{20,90,40}
\definecolor{tierone}{RGB}{25,70,140}
\definecolor{tiertwo}{RGB}{130,90,10}
\definecolor{tierthree}{RGB}{100,100,105}

\newcommand{\tierbadge}[2]{%
  \textcolor{#1}{\footnotesize\textsf{[#2]}}}
\newcommand{\Tzero}{\tierbadge{tierzero}{T0 Lean}}
\newcommand{\Tone}{\tierbadge{tierone}{T1 exact}}
\newcommand{\Ttwo}{\tierbadge{tiertwo}{T2 numeric}}
\newcommand{\Tthree}{\tierbadge{tierthree}{T3 analytic}}

% Graph identity of a statement, so a reader or an agent can run
% `workhouse why <id>` on anything asserted here.
\newcommand{\gid}[1]{\par\nobreak\smallskip\noindent
  {\footnotesize\textsf{Graph id:} \nolinkurl{#1}}\par\smallskip}
\newcommand{\gidc}[2]{\par\nobreak\smallskip\noindent
  {\footnotesize\textsf{Graph id:} \nolinkurl{#1}\quad
   \textsf{Reproduce:} \texttt{#2}}\par\smallskip}
\newcommand{\srcpath}[1]{\nolinkurl{#1}}

% An explicit, visible statement of what a theorem does NOT say.  Used
% wherever the corpus records a scope limit, which is most places.
\newcommand{\notasserted}[1]{\par\smallskip\noindent
  {\small\textbf{Does not assert.} #1}\par\smallskip}

% The obligation register's problem headers.
\newcommand{\obligation}[2]{\par\medskip\noindent
  \textbf{#1}\ \ \textit{#2}\par\smallskip}
% ---- end preamble.tex ----------------------------------------------

\pagestyle{fancy}\fancyhf{}
\fancyhead[L]{\small The WORKHOUSE program}
\fancyhead[R]{\small master edition, 12 September 2026}
\fancyfoot[C]{\thepage}

% ---- begin gen_macros.tex ------------------------------------------
% Generated by paper/master/generate_apparatus.py.  Do not edit.
% Every count the prose quotes is defined here, from the ledgers, so
% the paper cannot state a number the repository has moved past.
\newcommand{\LitPapers}{108}
\newcommand{\LitEdges}{152}
\newcommand{\LitVerified}{54}
\newcommand{\LitObtained}{34}
\newcommand{\RouteTotal}{111}
\newcommand{\RouteDone}{64}
\newcommand{\RouteLive}{20}
\newcommand{\RouteUntried}{12}
\newcommand{\RouteDead}{15}
\newcommand{\GapOpen}{20}
\newcommand{\GapDischarged}{5}
\newcommand{\ResTotal}{143}
\newcommand{\LeanTotal}{439}
\newcommand{\LeanSorry}{0}
\newcommand{\ContradictionTotal}{22}
\newcommand{\ContraFalsified}{4}
\newcommand{\ContraSuperseded}{3}
\newcommand{\ContraResolved}{15}
\newcommand{\ResConditional}{1}
\newcommand{\ResDisputed}{1}
\newcommand{\ResFalsified}{1}
\newcommand{\ResProven}{140}
% ---- end gen_macros.tex --------------------------------------------
% ---- begin gen_srcmacros.tex ---------------------------------------
% Generated by paper/master/generate_sources.py.  Do not edit.
% Counts for the source-provenance sections.
\newcommand{\CiteDocs}{40}
\newcommand{\DerivStatements}{332}
\newcommand{\DerivOpen}{63}
\newcommand{\LeanModules}{22}
\newcommand{\LibraryPapers}{260}
\newcommand{\LibraryVolOne}{99}
\newcommand{\LibraryVolTwo}{161}
\newcommand{\ExtractedPages}{777}
\newcommand{\ExtractedWorks}{28}
\newcommand{\ExtractedTopics}{7}
\newcommand{\SymbolCount}{28}
\newcommand{\AliasCount}{248}
\newcommand{\NoteArchives}{17}
\newcommand{\NoteReviews}{313}
\newcommand{\GraphEdges}{32503}
\newcommand{\GraphEdgeKinds}{0}
% ---- end gen_srcmacros.tex -----------------------------------------

\title{\bfseries The WORKHOUSE Program\\[4pt]
\large Exact \texorpdfstring{$\SU(N)$}{SU(N)} Flux-Band Spectra, the Fixed-Spacing Wilson
Construction,\\ and a Verified Register of What the Continuum Limit
Still Requires}
\author{Alexander Smith\\\small Independent Researcher}
\date{Master edition, 12 September 2026}

\begin{document}
\maketitle

% ---- begin 00_abstract.tex -----------------------------------------
\begin{abstract}
\noindent
This is the master reference edition of a four-year program on the
strong-coupling spectral structure of Hamiltonian $\SU(N)$ lattice gauge
theory and on the obligations that separate that structure from a
continuum Yang--Mills mass gap. It is organized so that its two halves can
be read, and used, independently.

Part~\ref{part:exact} is a closed body of exact mathematics. A
rank-one cellular subspace --- the kernel of the plaquette boundary map ---
carries the strong-coupling band, and its hopping amplitude is the
rational function
\[
  t_N = \frac{2N\,(N^2-4)}{(N^2-1)(2N^2-1)(4N^2-9)},
\]
positive for $N \ge 3$, vanishing identically at $N = 2$. The four
shared-link channel weights that produce it are order-two Weingarten
values rather than an isotropy assumption, so the second-order theory
rests on representation theory alone. At fourth order the effective
operator splits into Hodge terms and one geometric term coupling the
carrier to its complement, with the off-axis coefficient decided by
derivation across three independent implementations. The rank grading is
exact: every even order of the one-plaquette splitting carries the same
single power of $N$, with coefficients
$c_2 = -4$, $c_4 = -32$, $c_6 = -1748/3$, $c_8 = -123332/9$,
$c_{10} = -49593808/135$, and a four-channel cancellation that sums to
zero at every even order from four to ten; while the odd orders are
determinant families --- for even $N$ every odd order vanishes
identically, and for odd $N$ the first surviving odd order is $N-2$, with
a closed-form single vertex. At sixth order, a rational band shape is
proved to survive every local direct term. Of the statements in
Part~\ref{part:exact}, none is conditional on an open estimate.

Part~\ref{part:fixed} states the construction at fixed lattice spacing:
an actual infinite-volume Wilson transfer limit, a complete physical odd
band isomorphic to $\ell^2(\Z^3) \otimes \C^3$, and a literal
plaquette-source frame that is onto that entire band, on an explicit
interval of the coupling.

Part~\ref{part:continuum} is the part written for whoever works on this
next. It states six obstructions between the fixed-spacing construction
and the continuum limit --- five of them with explicit finite
counterexamples, four of them having killed arguments this program had
written down as complete --- of which the decisive one is that the
interval where the construction is proved and the trajectory along which
the continuum is taken are disjoint sets of couplings. It then gives the
refutation register: \ContraFalsified{} falsified claims with what
falsified each, and the mechanisms this program advanced and withdrew. It
closes with an obligation register of eight numbered problems and the
repository's full route record --- \RouteTotal{} routes, of which
\RouteDone{} are complete, \RouteLive{} live, \RouteUntried{} untried and
\RouteDead{} dead, each dead one carrying the reason it died.

No claim is made here about the existence of continuum Yang--Mills theory
or about a continuum mass gap. \S\ref{sec:breaks} is the argument that no
such claim currently follows, and \S\ref{sec:obligations} is the list of
what would have to change.

Every count quoted above is generated from the repository's ledgers
rather than typed, and every statement carries its
machine-verification tier and the command that re-establishes it.
Currently: \LeanTotal{} Lean~4 theorems with \LeanSorry{} occurrences of
\texttt{sorry}, \ResTotal{} registered analytic results, and
\LitVerified{} verified evidence edges to published work across
\LitObtained{} obtained papers.
\end{abstract}
% ---- end 00_abstract.tex -------------------------------------------

\tableofcontents
\bigskip

% ---- begin 01_introduction.tex -------------------------------------
\section{Introduction}
\label{sec:intro}

\subsection{What is actually here}

Strong-coupling expansions of Hamiltonian lattice gauge theory are old.
What is not old is the object this program organizes them around.

Write the cellular chain complex of the cubic lattice as
$C_3 \xrightarrow{\partial_3} C_2 \xrightarrow{\partial_2} C_1$, with
$C_2$ spanned by plaquettes. The subspace
\begin{equation}
  Z_2 = \ker \partial_2 \subset C_2
\end{equation}
--- the closed plaquette chains --- is the carrier of the strong-coupling
band. It is a homological object, fixed before any perturbation theory
is done, and the entire content of Parts~\ref{part:exact} and
\ref{part:fixed} is a description of what the Kogut--Susskind Hamiltonian
does to it, order by order and rank by rank.

Three things follow from taking that object seriously, and they are the
three claims of this document.

\emph{First: the coefficients are exact rational functions of the rank.}
Not asymptotic in $N$, not fitted, not computed rank by rank and
extrapolated. The second-order hopping amplitude is
\begin{equation}
  t_N = \frac{2N\,(N^2-4)}{(N^2-1)(2N^2-1)(4N^2-9)},
  \qquad t_2 = 0, \quad t_3 = \tfrac{5}{612},
  \label{eq:intro-tN}
\end{equation}
and the four channel weights behind it are order-two Weingarten values,
so \eqref{eq:intro-tN} imports nothing from any expansion. The factor
$N^2-4$ is not a coincidence of small rank: it is why $\SU(2)$ has no
band at this order at all. The same discipline runs through fourth
order, through the rank grading of \S\ref{sec:grading}, and into a
symbolic engine that works over $\Q(N)$ with $N$ an indeterminate rather
than a number.

\emph{Second: the rank grading has exact structure, and it comes from
cancellation.} For the $C$-parity splitting of the one-plaquette
diagonal, every even order carries the same single power of $N$ once the
coupling is written in the planar variable, and the leading coefficients
of the four contributing intermediate channels are in the ratio
$(-4, +2, +1, +1)$ and \emph{sum to zero}. The grading is a cancellation
theorem, not a power count. Complementarily, the odd orders vanish
identically for even $N$, and for odd $N$ the first surviving one is
order $N-2$, with a closed-form vertex. Neither statement is visible from
any single rank.

\emph{Third: the construction at fixed lattice spacing is complete, and
it does not reach the continuum.} \S\ref{sec:wilson} states an actual
infinite-volume Wilson transfer limit with a complete physical odd band
and a literal plaquette-source frame that is onto it. \S\ref{sec:breaks}
states why that is not a continuum mass gap, in six independent ways,
five of them with finite counterexamples. The most important is
structural rather than technical: the theorem holds on
$\abs{u} \le u_*/(10\,022\,400\,000 N)$ and the continuum trajectory runs
to $u \to \infty$, so the two are disjoint and no sharpening of the
constant changes that.

\subsection{Why the document is shaped this way}

Alex, the maintainer of this program, asked for a paper that would serve
as the base for further work on the mass gap rather than as a report of
work completed. That request has consequences for the form.

It means that failure is content. \S\ref{sec:breaks} and
\S\ref{sec:register} exist because a route this program has already
closed is more valuable to the next attempt than a route it has already
opened, and because at least four arguments were written down here as
complete and then found not to be. It means the route register in
\S\ref{sec:obligations} reproduces \RouteDead{} dead routes at length,
including their measured costs, rather than listing only the \RouteLive{}
live ones. And it means every count in the prose is a macro generated
from the repository's ledgers, so that a stale number is a build failure
rather than a quiet misstatement.

It also means the document distinguishes, everywhere and visibly, three
things that are routinely conflated: what a statement mathematically
establishes, what artefact backs it, and what a machine checks. A result
can be fully proved analytically and carry the weakest machine badge.
A Lean theorem can be checked by the kernel and prove something
trivial. This program has made both errors and the notation of
\S\ref{sec:howto} is the correction.

\subsection{Relation to the literature}

The external comparisons are recorded as an evidence map rather than a
citation list: \LitPapers{} indexed papers, \LitEdges{} recorded
relations, of which \LitVerified{} are verified against a named claim
here. A published paper is not authority in this program --- it is
unchecked until something checks it, exactly like an internal document.
What makes an external agreement evidence is that it was produced
without knowledge of this program, and \S\ref{sec:external} reports the
comparisons that hold and, separately, the two that contradict.

\subsection{What is claimed, and what is not}

Claimed: the contents of Part~\ref{part:exact}, unconditionally, at the
tiers recorded; and the contents of Part~\ref{part:fixed} on the stated
interval and under the stated hypotheses.

Not claimed: the existence of continuum four-dimensional Yang--Mills
theory; a continuum mass gap; Poincar\'e covariance; restoration of
rotational invariance from the hypercubic lattice; or the construction of
a continuum measure. None of these has an argument in this document, and
\S\ref{sec:breaks} explains for each of the six obstructions why the
present machinery does not supply one.
% ---- end 01_introduction.tex ---------------------------------------
% ---- begin 02_how_to_use.tex ---------------------------------------
\section{How to use this document}
\label{sec:howto}

This document has two audiences and it is worth saying which paragraph is
for whom.

A reader who wants the mathematics should read Part~\ref{part:exact},
which is self-contained and compiles on its own as a separate paper
(\srcpath{paper/master/workhouse_band_paper.tex}). Nothing in
Part~\ref{part:exact} depends on anything later, and nothing in it is
conditional on an open estimate.

A reader who intends to \emph{work} on the continuum problem should read
\S\ref{sec:breaks}, \S\ref{sec:register} and \S\ref{sec:obligations}
first, in that order, and only then come back for the mathematics. The
reason is economic. This program has spent four years discovering that
most of the obvious routes to the continuum limit are closed, and the
record of \emph{which} routes are closed and \emph{why} is worth more
than the record of what is open. Four sessions of one earlier effort were
spent on a single route that a static argument had already shown could
not decide the question it was aimed at; that experience is why the
route register in \S\ref{sec:obligations} exists and why dead routes are
reproduced in it at full length rather than deleted.

\subsection{What the badges mean}

Every statement in this document carries a machine-verification badge.

\begin{itemize}
  \item \Tzero{} --- Lean~4 compiles the statement from explicit
    definitions with no \texttt{sorry}, using only \texttt{propext},
    \texttt{Classical.choice} and \texttt{Quot.sound}. Nothing written
    in prose can weaken a T0 badge.
  \item \Tone{} --- the statement is re-derived symbolically, in exact
    arithmetic, from stated definitions.
  \item \Ttwo{} --- floating-point agreement within a tolerance that the
    check itself prints.
  \item \Tthree{} --- no single repository check covers the whole
    statement.
\end{itemize}

The badge is \emph{independent} of mathematical status. A theorem with a
complete analytic proof and no formalization carries \Tthree{}: the badge
says what the machine checks, not what the mathematics establishes. This
is the one convention a reader must internalize to read the document
correctly, because the opposite reading --- T3 as a synonym for
``unproved'' --- makes most of Part~\ref{part:continuum} look weaker than
it is, and the reading of T0 as a synonym for ``important'' makes it look
stronger. Both errors have been made in this program's own internal
documents and corrected.

Where a statement's scope is narrower than its headline suggests, a
\textbf{Does not assert} paragraph follows it. Those paragraphs are not
hedging. They are the part of each result that took longest to get right,
and several of them exist because an earlier version of the claim was
broader and was refuted.

\subsection{If you are the next agent on this program}

The following is an operating procedure, not advice. It is written for
the case --- now the normal case --- in which the next piece of work on
this program is done by an automated agent with a fresh context window.

\begin{enumerate}
  \item \textbf{Do not trust this document over the ledgers.} Every
    number quoted in the prose here is generated from
    \srcpath{ledger/*.yaml} by \srcpath{paper/master/generate\_apparatus.py}.
    If the ledger and this document disagree, the ledger is right and
    this document is stale; regenerate it. The counts in
    \S\ref{sec:obligations}, \S\ref{sec:external} and
    Appendix~\ref{app:verification} are macros, not typed literals,
    precisely so that this failure mode is visible rather than silent.

  \item \textbf{Establish the state before choosing a target.} Run
    \texttt{workhouse brief -{}-startup}, then \texttt{workhouse brief <gap> -{}-json}
    for the gap you mean to work on, and retain the snapshot. A saved
    graph in a fresh clone reports freshness \texttt{unknown} until local
    generation provenance exists; saved briefing mode never silently
    rebuilds it.

  \item \textbf{Read the dead routes for your gap before starting.}
    \texttt{workhouse why <gap>} prints them. \S\ref{sec:obligations}
    reproduces them here. Of \RouteTotal{} recorded routes across all
    gaps, \RouteDead{} are dead, and several died only after substantial
    expenditure with a measured cost recorded in the ledger. A route
    marked \texttt{cannot\_decide} is not merely unpromising: it is
    structurally incapable of resolving the claim it was aimed at, and
    re-running it will produce output that answers a different question.

  \item \textbf{Do not aim at the continuum limit directly.}
    \S\ref{sec:breaks} states six independent obstructions, five of them
    with explicit finite counterexamples. Each is a constraint on what a
    correct argument can look like, and an argument that does not
    explicitly clear all six is not a candidate. In particular, Break~5
    (\S\ref{sec:break5}) shows that the regime in which this program's
    strongest construction holds and the regime in which the continuum
    limit lives are disjoint sets of couplings. That is a statement about
    the shape of any possible proof, not a gap in the present one.

  \item \textbf{Pick a numbered Problem from \S\ref{sec:obligations}}, not
    a topic. Each Problem states its available hypotheses, what it
    unblocks downstream, and the specific next computation. Problems are
    ordered by the ledger's research judgement, which is not a claim that
    earlier ones are prerequisites for later ones.

  \item \textbf{Record the attempt, including its failure.} A route that
    dies is a result. Set its state in \srcpath{ledger/gaps.yaml}, name
    what closed it in \texttt{closed\_by}, and if the instrument could
    not have reached the claim, say so in \texttt{cannot\_decide}. The
    working agreement in \srcpath{CLAUDE.md} is the binding version of
    this instruction; this paragraph only explains why it is there.
\end{enumerate}

\subsection{What would falsify the contents of this document}

Three things, in decreasing order of how much they would cost.

First, an error in the Haar or Weingarten input to \S\ref{sec:second}
would propagate to every order and to both the fourth-order band and its
rank grading. That input is now representation theory rather than
assumption (\S\ref{sec:second}), and it is the load-bearing algebraic
fact of Part~\ref{part:exact}.

Second, a demonstration that the plaquette source family of
\S\ref{sec:wilson} is not total in the physical complement would empty
the fixed-spacing band theorem of physical content without making it
false --- this is exactly the failure mode that Break~4
(\S\ref{sec:break4}) exhibits in a three-dimensional model, and totality
is therefore asserted in that theorem as a proved conclusion rather than
assumed.

Third, and least costly, a published result reproducing any of the exact
coefficients of \S\ref{sec:grading} or \S\ref{sec:fourth} would remove
this program's claim to have derived them first without touching their
correctness. \S\ref{sec:external} records the searches already performed
against \LitPapers{} indexed papers.
% ---- end 02_how_to_use.tex -----------------------------------------
% ---- begin 03_sources.tex ------------------------------------------
\section{The source layers}
\label{sec:sources}

Every derivation in this program can be traced to bytes. This section
says how, because the tracing is not uniform: there are three source
systems, they certify different things, and reading a statement out of
the wrong one is the most common way to get this corpus wrong.

An agent that internalizes only one thing from this document should
internalize the table below.

\begin{center}
\begin{tabular}{@{}p{0.19\linewidth} p{0.34\linewidth} p{0.38\linewidth}@{}}
\toprule
Layer & What it holds & What it certifies \\
\midrule
\textbf{Theory graph} &
\DerivStatements{} derivation statements in \CiteDocs{} documents;
\ResTotal{} results; \LeanTotal{} theorems; \RouteTotal{} routes;
\GraphEdges{} edges in \GraphEdgeKinds{} relations &
That a statement was \emph{recorded} with a status, hypotheses, a
source locator, and a digest of the bytes it was read from. It does
not certify the mathematics; the tier does that separately. \\
\addlinespace
\textbf{Published work} &
\LitPapers{} papers in the evidence map, \LitVerified{} edges verified;
\LibraryPapers{} acquired papers in the reading library;
\ExtractedPages{} extracted pages across \ExtractedWorks{} works &
That an external result exists and, for a verified edge, that it was
read against a named claim here. Obtained is not read; read is not
checked. \\
\addlinespace
\textbf{Local archive} &
Four years of manuscripts, runs, notes and Lean, with
\AliasCount{} aliased documents, \NoteArchives{} note archives and
\NoteReviews{} recorded reviews &
That the bytes are preserved and pinned. Nothing about whether the
content is current, or correct. \\
\bottomrule
\end{tabular}
\end{center}

\subsection{Layer 1: the theory graph}

The graph is the index of record. Its node identifiers are prefixed,
and the prefix tells you what kind of object you are holding and
therefore what it can support.

\begin{description}[leftmargin=1.1em,style=nextline]
\item[\texttt{CITE:}] A source document, pinned by path and SHA-256.
  \CiteDocs{} of them. A \texttt{CITE:} record is bytes, not a claim.
\item[\texttt{DERIV:}] A statement located \emph{inside} a
  \texttt{CITE:} document, by heading and line range.
  \DerivStatements{} of them, every one with a locator. This is the
  finest granularity the corpus tracks, and it is the right unit to
  work at.
\item[\texttt{RESULT:}] A first-class analytic claim, independent of any
  one document. \ResTotal{} of them, each with status, evidence,
  hypotheses, scope, \texttt{depends\_on} and \texttt{bears\_on}.
\item[\texttt{LEAN:}] A registered theorem. \LeanTotal{} of them.
  \texttt{formalizes} links it to a catalogue record;
  \texttt{promotes} --- which only 49 carry --- means it proves an
  entire named check and therefore lifts that check's tier.
\item[\texttt{CHK:}] A registered invariant check, carrying its tier and
  the numbers it establishes. A \texttt{FINDING:} check is one that
  asserts a \emph{discrepancy} rather than an agreement.
\item[\texttt{RUN:}] A pinned computational run, with its own manifest.
\item[\texttt{ADR:}] A numbered decision record: what was decided, on
  what evidence, and what it retracts or amends.
\item[\texttt{CONST:}, \texttt{SYM:}] A named constant and its symbol
  registry entry --- canonical name, corpus spellings, code names, and
  exact values. Appendix~\ref{app:symbols}.
\item[\texttt{G}$n$, \texttt{C}$n$, \texttt{U}$n$, \texttt{R}$n$] Gaps,
  contradictions, unifying candidates and registry claims.
\end{description}

The edges matter as much as the nodes. \GraphEdges{} of them across
\GraphEdgeKinds{} relations, and three deserve separate mention because
they are routinely confused. \texttt{depends\_on} is a mathematical
input: if it falls, the dependent falls. \texttt{supported\_by} is
narrower --- a control that bears on part of a statement, and a result
with \texttt{supported\_by: []} has no registered machine support at
all, which is a fact worth checking before quoting a tier.
\texttt{bears\_on} is looser still: a topical connection, and it
establishes nothing.

\subsection{Layer 2: published work}

There are two published-work systems and they answer different
questions. Using the wrong one produces either a false novelty claim or
a false authority claim.

\emph{The evidence map} (\srcpath{literature/index.yaml}) is small,
curated and adversarial. \LitPapers{} papers, \LitEdges{} recorded
relations, \LitVerified{} of them verified against a named claim of
this program. An entry earns its place by naming a specific claim.
Query it with \texttt{workhouse lit -{}-for G19}. This is the only
layer from which an agreement or a contradiction may be quoted.

\emph{The reading library}
(\srcpath{literature/library/paper\_manifest.csv} and
\srcpath{volume2/}) is large and inclusive: \LibraryPapers{} acquired
papers --- \LibraryVolOne{} and \LibraryVolTwo{} --- each with tier,
category, journal reference, DOI, arXiv identifier, INSPIRE citation
count, local file, page count, SHA-256, abstract, and a curated
one-line \texttt{workhouse\_connection}. Appendix~\ref{app:library}
reproduces it. Its purpose is coverage: to answer ``has anyone done
this?'' before an agent spends a week. Its connection field is a
reading note and carries no evidential weight.

The gap between the two is the useful measurement. \LitObtained{}
papers have been read against a claim; \LibraryPapers{} have been
acquired. Most of the library has never been checked against anything
here, and a novelty claim asserted against the library is therefore
weaker than one asserted against the evidence map --- which is itself
weaker than priority, for the reasons in \S\ref{sec:external}.

\subsection{Layer 2b: the page-level extraction index}

This is the piece an agent is most likely to miss, and it is the one
that saves the most time.

\ExtractedPages{} pages of local PDFs have been extracted and indexed
at \srcpath{literature/yangmills/extraction/pages.jsonl}, covering
\ExtractedWorks{} works. Each record carries the paper identifier, the
PDF page number, the SHA-256 of the source \emph{and} of the extracted
text, the path of the reading copy, and the line positions of every hit
against \ExtractedTopics{} topic keys. Appendix~\ref{app:extraction}
tabulates it.

The workflow it supports:

\begin{enumerate}
  \item Search \srcpath{pages.jsonl} for a paper identifier and a topic
    key (\texttt{ground\_state}, \texttt{spectral\_gap},
    \texttt{curvature}, \texttt{reconstruction},
    \texttt{uniform\_limits}, \texttt{constructive},
    \texttt{gauge\_geometry}).
  \item Follow the page and line locator into the reading copy under
    \srcpath{literature/inbox/<ID>/extracted/}.
  \item Verify any equation you intend to use against the pinned PDF.
\end{enumerate}

Three cautions the index states about itself. A topic hit may be a
citation rather than a result, or an extraction error. Candidate
theorem labels are recorded on many pages, but a label can sit inside a
citation or carry an OCR fault --- the index is a search aid and
asserts nothing. And five pages extracted fewer than eighty characters
and are explicitly flagged as unreliable; their contents must not be
inferred from the text index. Exact equations from legacy scans must be
checked in the PDF.

\subsection{Layer 3: the local archive}

The archive is preserved rather than curated, and its governing
instruction is that nothing is erased. Failed attempts, superseded
drafts, duplicates and uncommitted work are all kept, because deleting
a refuted claim destroys the evidence that it was tried.

Consequently the field that matters when reading an archive document is
its \emph{standing}. \AliasCount{} documents carry a registered alias
and a standing, tabulated in Appendix~\ref{app:documents}. A superseded
document is not marked in its own text and reads exactly like a current
one; the register is the only thing that distinguishes them. Anything
under \srcpath{theory/superseded/} is retained for audit and must never
be read as current.

The maintainer's own research notes enter through a separate
discipline: \NoteArchives{} declared archives with \NoteReviews{}
recorded reviews, each carrying a content digest, a verdict from a
closed vocabulary, and a mandatory reason. An \texttt{import} is
byte-verified against its digest; an \texttt{extract} must name the
claims it entered through; a \texttt{set-aside} is a recorded judgement
and never a deletion. No verdict promotes anything past the default
machine tier.

\subsection{Finding things: the corpus joins on values, not concepts}

The single most important retrieval fact about this corpus:
\emph{its join keys are exact rationals and spellings, not concepts.}
No semantic query retrieves $109151/249696$, and the same quantity
appears in a manuscript, in code, and in a claim record under three
different names.

\texttt{workhouse search} exists for this. It matches by \emph{value}
rather than spelling --- searching $-10/96$ finds $-5/48$ --- across
the prose index, the code and certificate index, the claim catalogue,
and the curated alias table of Appendix~\ref{app:symbols}. It also
carries two warnings no text search can: names the register
\emph{forbids} because using them regenerates a dissolved
contradiction, and names coined inside this program that its own source
corpus never uses.

Two indexes sit underneath it, and the second exists because the first
could not see the sealed core: values such as $5/48$, $5/12$, $5/612$,
$11/306$ and $7/102$ can be absent from the prose index while their
computations are present in code.

\subsection{The queries worth knowing}

\begin{verbatim}
workhouse why <id>          everything recorded about one claim:
                            both sides of a dispute, every check and
                            its verdict, routes tried and which died
workhouse search '5/612'    search by exact VALUE, not by name
workhouse branches C2       both sides of a dispute, side by side
workhouse derive C2 G3      the evidence chain, registered edges only
workhouse lit --for G19     published work bearing on a claim
workhouse notes --queue     the next unreviewed notes, by signal
workhouse triage <path>     survey an unpinned collection
workhouse ask "<question>"  candidate finder over prose; ends in `why`
workhouse brief <gap> --json  target snapshot, with freshness
\end{verbatim}

\texttt{why} is the front door. It resolves a claim identifier, a
constant, a gap, an ADR, or a corpus file path, and it prints the
routes with their states --- which is how a dead route announces itself
before an agent re-runs it.

\subsection{What each layer cannot do}

\begin{itemize}
  \item The graph cannot tell you a statement is \emph{true}. It tells
    you what was recorded, by whom, from which bytes, under which
    hypotheses. Status, evidence and machine tier are three
    independent fields precisely so that this stays visible.
  \item The library cannot establish novelty. It establishes that a
    paper was acquired.
  \item The extraction index cannot establish what a paper says. It
    tells you which page to read.
  \item The archive cannot tell you what is current. Standing does.
  \item A digest cannot tell you an anchor resolves. Check both: a
    correct SHA-256 over a fabricated locator has occurred here.
\end{itemize}

\subsection{Operational rules live elsewhere, deliberately}

How to add an invariant, register a Lean theorem, record a route, order
the regeneration passes, and land work: those are specified in
\srcpath{CLAUDE.md} (the working agreement), \srcpath{AGENTS.md} (the
research posture) and \srcpath{CONTRIBUTING.md}. They are not restated
here.

That is a decision, not an omission. Those files are binding and they
change; a copy in a document that is rebuilt on its own schedule would
drift, and an agent following a stale copy of a rule is worse off than
one following no copy. This document carries the \emph{mathematical}
working knowledge --- which sources support what, how the derivations
connect, and where the traps are. For procedure, read the source of
record.

Two of its rules are worth naming here anyway, because they explain the
shape of everything above. A claim is established by its proof or
computation under explicit hypotheses, and its artefact must not claim
more than that argument establishes. And a disputed value closes by
derivation or not at all: both sides stay recorded, and no code path
may pick one, average them, or prefer whichever is an exact rational
because exactness reads as authority.
% ---- end 03_sources.tex --------------------------------------------

\part{Exact spectral structure at strong coupling}\label{part:exact}
% ---- begin 10_setting.tex ------------------------------------------
\section{Setting and conventions}
\label{sec:setting}

\subsection{The Hamiltonian}

Work on the cubic lattice $\Z^3$ with gauge group $G = \SU(N)$, one
group element $U_\ell$ per oriented link. The Kogut--Susskind
Hamiltonian is
\begin{equation}
  H \;=\; H_0 + V,
  \qquad
  H_0 = \tfrac12 \sum_{\ell} E_\ell^2,
  \qquad
  V = -\,u\, W,
  \label{eq:KS}
\end{equation}
where $E_\ell^2$ is the quadratic Casimir of the left-invariant electric
field on link $\ell$, and $W$ is the sum of plaquette characters in both
orientations,
\begin{equation}
  W = \sum_{p} \big(\chi_p + \bar\chi_p\big),
  \qquad
  \chi_p = \Tr \prod_{\ell \in \partial p} U_\ell .
\end{equation}

Strong-coupling perturbation theory is the expansion in $u$ around the
electric vacuum, organized as degenerate perturbation theory --- a des
Cloizeaux effective operator --- on the sector of interest.

\subsection{The coupling, and one erratum that matters}

The canonical coupling variable throughout is
\begin{equation}
  u \;=\; \frac{\beta_N}{2N},
  \label{eq:coupling}
\end{equation}
whose $\SU(3)$ specialization is $u = \beta_3/6$.

A warning, because it has corrupted transcriptions of this material
before. An archived line reading $Y = 2\beta/3 = 4u$ is a
definition-label erratum, not a change of variable. The coefficients
printed alongside it were \emph{already} in the canonical coordinate
\eqref{eq:coupling}. Applying a $4^r$ rescaling at order $r$ to
``correct'' them corrupts every order, and the repository carries an
explicit check that the printed towers reproduce their canonical values
verbatim.

For rank-uniform statements the natural variable is the planar one,
\begin{equation}
  \tau \;=\; \frac{2u}{N^2} \;=\; \frac{\beta_N}{N^3},
  \label{eq:tau}
\end{equation}
in which the grading theorems of \S\ref{sec:grading} take their simplest
form: every even order of the one-plaquette splitting then carries the
same single power of $N$.

\subsection{Two projections that must not be conflated}

Two distinct projections appear throughout, and several contradictions
in this program's own register were naming collisions between them.

The \emph{band} projection is onto the degenerate strong-coupling
subspace carrying the excitation of interest --- for the flux band, the
cellular subspace of \S\ref{sec:carrier}. Coefficients anchored to it
are band-kernel quantities.

The \emph{physical} projection is onto the gauge-invariant,
vacuum-subtracted spectrum at a definite momentum. Coefficients anchored
to it are physical quantities.

These are related by translation-local scalar shifts, and such a shift
changes neither the centered operator, nor its eigenvectors, nor the
bandwidth. It does change the scalar. Consequently the band-kernel
anchor and the physical $\Gamma$-point coefficient are two differently
anchored coordinates, not two competing estimates of one quantity, and
they are written $q^{(4)}_{\mathrm{band}}$ and $m^{(4)}_{\Gamma}$
respectively. Writing either as ``$m_4$'' regenerates a contradiction
that does not exist, and the repository's search tooling refuses the
name for that reason.

\subsection{Exactness}

Every coefficient displayed in Part~\ref{part:exact} is an exact
rational, or an exact rational function of $N$. Where a quantity is
known here only numerically it carries an explicit numerical marker in
the source registry and is never displayed as though exact. The single
most dangerous defect available in this material is a float that reads
as exact, and the boundary between the two is machine-guarded rather
than maintained by care.
% ---- end 10_setting.tex --------------------------------------------
% ---- begin 11_carrier.tex ------------------------------------------
\section{The homological carrier}
\label{sec:carrier}

\subsection{Where the flat band comes from}

Project the Kogut--Susskind Hamiltonian \eqref{eq:KS} onto the
fundamental, charge-odd, one-plaquette manifold. In three spatial
dimensions a fundamental plaquette excitation carries one of three face
orientations, and at second order neighbouring faces communicate
through a shared link. The projected problem is therefore a
three-orbital Bloch Hamiltonian, and it looks at first like a
tight-binding model.

It is not an engineered one. The signs of the matrix elements are fixed
by the orientation of the plaquette boundary and by charge conjugation,
and non-abelian group integration generates, exactly, the oriented
edge--face boundary Gram operator
\begin{equation}
  B(k)B(k)^{\dagger},
  \qquad B = \partial_2^{\dagger},
\end{equation}
with the coefficient of \S\ref{sec:second} in front of it. The
structure of the result is a clean separation:
\begin{equation}
  \boxed{\ \text{colour dynamics} \longrightarrow t_N,
  \qquad
  \text{cellular geometry} \longrightarrow BB^{\dagger}. \ }
  \label{eq:separation}
\end{equation}
Representation theory fixes the amplitude; the cubical chain complex
fixes the kernel and the spectrum. Neither borrows from the other.

\subsection{The projected spectrum}

\begin{theorem}[Second-order projected spectrum]
\label{thm:projspec}
The projected second-order spectrum consists of one
momentum-independent level together with a doubly degenerate branch
shifted by $t_N u^2 q(k)$, where
\begin{equation}
  q(k) \;=\; 4\sum_{j=1}^{3} \sin^2\!\big(k_j/2\big).
\end{equation}
\end{theorem}

The flat level is exact and it is exact for a reason: the flat carrier
is $\ker \partial_2$.

\begin{theorem}[The carrier and its multiplicity]
\label{thm:multiplicity}
Elementary cube boundaries are compact localized states of the carrier;
three wrapping sheets complete it on the torus. For $L \ge 3$ the exact
multiplicity is
\begin{equation}
  (L^3 - 1) + 3 \;=\; L^3 + 2 .
\end{equation}
\end{theorem}
\gid{RESULT:UNIVERSAL\_CELLULAR\_HODGE\_SPECTRUM}

The two contributions are structurally different and the distinction
matters later. The $L^3 - 1$ cube boundaries are local and finitely
supported. The three sheets are not: they wrap the torus, and they are
the entire zero-momentum content of the carrier. \S\ref{sec:wilson}
returns to the consequence --- that no finitely supported exactly-flat
local operator can see the $\Gamma$ point at all.

\subsection{Third order, and the coefficients that survive}

For $\SU(3)$, exact third-order perturbation theory stays inside the
same boundary-factorized algebra. It does not generate a new operator
structure; it renormalizes the coefficient in front of the same
$BB^{\dagger}$:
\begin{equation}
  H_{\mathrm{eff},-}
  \;=\; E_{\mathrm{flat}}(u)\, I
  \;+\; \left(\frac{5}{612}u^2 + \frac{1975}{124848}u^3\right) BB^{\dagger}
  \;+\; O(u^4),
  \label{eq:thirdorder}
\end{equation}
with
\begin{equation}
  E_{\mathrm{flat}}(u)
  \;=\; \frac{8}{3} + u + \frac{11}{306}u^2
        - \frac{109151}{249696}u^3 .
\end{equation}

That \eqref{eq:thirdorder} remains boundary-factorized through third
order is the non-trivial part. It is what makes the fourth order the
first place where the carrier can couple to its complement, and
\S\ref{sec:fourth} is about exactly that coupling.

\subsection{What is and is not novel here}

This program does not claim that incidence flatness, cellular compact
localized states, or homological counting of zero modes are new. All
three have a substantial prior literature --- line-graph constructions,
compact localized states, the necessity of non-contractible states on
periodic systems, rectangular factorizations and index arguments, and
closely related closed-surface structures in two-form spin liquids ---
and a recent independent construction treats face-graph and higher-cell
flat bands directly.

What is claimed is that the flatness here is \emph{dynamical}: it is
produced by exact strong-coupling matrix elements of non-abelian
$\SU(N)$ Hamiltonian lattice gauge theory, with the amplitude fixed by
group representation theory and the kernel fixed by the cubical chain
complex, as in \eqref{eq:separation}. The flat band is not put in; it
comes out.

\begin{scopenote}
Theorem~\ref{thm:projspec} and \eqref{eq:thirdorder} are finite-volume,
finite-order statements about a projected flux sector. They are not a
continuum glueball theorem and not a mass-gap theorem, and the source
manuscript says so in its own abstract.
\end{scopenote}
% ---- end 11_carrier.tex --------------------------------------------
% ---- begin 12_second_order.tex -------------------------------------
\section{Second order at every rank}
\label{sec:second}

This section is the algebraic base of the whole program. Everything in
it is proof-checked, and its conclusion --- that the band exists for
$N \ge 3$ and not for $N = 2$, with an exact amplitude --- follows from
representation theory with no input from any expansion.

\subsection{The shared-link weights are Weingarten values}

Two plaquettes contributing at second order either share a link or do
not. The amplitude for the shared-link case was for a long time taken
in this program with an isotropy premise: that the four contributing
representation channels carry weights in a fixed ratio. That premise
is not needed. It is a consequence of the order-two Weingarten values.

\begin{proposition}[The shared-link weights, from representation theory alone]
\label{prop:weingarten}
Let two plaquettes share exactly one link. Then:
\begin{enumerate}
  \item The six non-shared links collapse. Each plaquette contributes a
    product of three independent Haar links, and a product of
    independent Haar matrices is Haar. So the two-plaquette amplitude
    reduces to $\Tr(A U)\,\Tr(B U^{\pm 1})$ with $A$ and $B$ independent
    Haar, and integrating them out leaves a pure degree-$(2,2)$ moment of
    the shared link $U$.
  \item Two such moments settle both families. With
    $M_{\mathrm{direct}} = N^2$ and $M_{\mathrm{cross}} = N$, the like
    family splits as $(N+1)/(2N)$ and $(N-1)/(2N)$; and the singlet
    component of $U_{ij}\overline{U_{lk}}$ is
    $\delta_{il}\delta_{jk}/N$, of squared norm $1$, giving $1/N^2$.
\end{enumerate}
All four weights are exactly $d_R/N^2$, where $d_R$ is the dimension of
the channel.
\end{proposition}
\gidc{the shared-link weights are Weingarten, not an isotropy assumption}{workhouse verify --only 'the shared-link weights are Weingarten, not an isotropy assumption'}

\noindent\Tone{}. The relevant Weingarten pair is the inverse of the
$S_2$ Gram matrix and the index sums are explicit, so the check imports
nothing from the corpus. Registered as the discharge of the gap that
had carried this premise as the one unproved input of the second-order
chain; it was opened and closed on the same day, 28~August 2026, once
the question was asked in the right form.

\subsection{The four channels in closed form}

Write $C_F = (N^2-1)/(2N)$ for the fundamental Casimir. Each channel
$R$ contributes a resolvent weight $-(d_R/N^2)/(C_F + \tfrac12 \Delta C_2(R))$,
and all four evaluate in closed form.

\begin{theorem}[Per-channel resolvent weights]
\label{thm:channels}
\begin{align}
  w_{\mathbf 1}   &= -\frac{1/N^2}{C_F}
    &&= -\frac{2}{N(N^2-1)},
  \\[3pt]
  w_{\mathrm{Adj}} &= -\frac{(N^2-1)/N^2}{C_F + N/2}
    &&= -\frac{2(N^2-1)}{N(2N^2-1)},
  \\[3pt]
  w_{\mathrm{Sym}} &= -\frac{N(N+1)/2N^2}
                        {C_F + \frac{(N-1)(N+2)}{2N}}
    &&= -\frac{N+1}{(N-1)(2N+3)},
  \\[3pt]
  w_{\Lambda} &= -\frac{N(N-1)/2N^2}
                    {C_F + \frac{(N+1)(N-2)}{2N}}
    &&= -\frac{N-1}{(N+1)(2N-3)} .
\end{align}
\end{theorem}
\gid{LEAN:resolventWeight\_singlet, LEAN:resolventWeight\_adjoint, LEAN:resolventWeight\_sym, LEAN:resolventWeight\_antisym}

\noindent\Tzero{} --- all four, at
\srcpath{lean/Workhouse/Basic.lean:72,80,89,104}. Reproduce with
\texttt{make lean}.

\subsection{The two family sums}

The channels group into a \emph{mixed} family $\{\mathbf 1, \mathrm{Adj}\}$
and a \emph{like} family $\{\Lambda^2 F, \mathrm{Sym}^2 F\}$, and these
are the antiparallel and parallel channel sums.

\begin{theorem}[Family sums]
\label{thm:families}
\begin{equation}
  A_N \;=\; w_{\mathbf 1} + w_{\mathrm{Adj}}
    \;=\; -\,\frac{2N^3}{(N^2-1)(2N^2-1)},
  \qquad
  B_N \;=\; w_{\mathrm{Sym}} + w_{\Lambda}
    \;=\; -\,\frac{4N(N^2-2)}{(N^2-1)(4N^2-9)} .
\end{equation}
\end{theorem}
\gid{LEAN:mixed\_family\_sum, LEAN:like\_family\_sum}

\noindent\Tzero{}, at \srcpath{lean/Workhouse/Basic.lean:118,129}. The
mixed sum is obtained from the dimension and Casimir table rather than
by transcription from any source document, which is the point: the
family sums are re-derived, not quoted.

\subsection{The hopping amplitude}

\begin{theorem}[The all-rank second-order hopping amplitude]
\label{thm:tN}
\begin{equation}
  t_N \;=\; B_N - A_N
  \;=\; \frac{2N\,(N^2-4)}{(N^2-1)(2N^2-1)(4N^2-9)} .
  \label{eq:tN}
\end{equation}
It is strictly positive for every $N \ge 3$, and
\begin{equation}
  t_2 = 0, \qquad t_3 = \tfrac{5}{612}.
\end{equation}
\end{theorem}
\gidc{CONST:t\_N}{workhouse why t\_N}

\noindent\Tzero{} for the rank law and for both special values
(\srcpath{lean/Workhouse/Basic.lean}: \texttt{rank\_law\_numerator},
\texttt{hopping\_two}, \texttt{hopping\_three}). Clearing denominators,
the channel difference has numerator exactly $2N(N^2-4)$.

Two consequences deserve to be stated separately, because both are
structural rather than numerical.

\begin{corollary}[$\SU(2)$ has no band at this order]
\label{cor:su2}
The factor $N^2 - 4$ in \eqref{eq:tN} vanishes at $N = 2$. The exclusion
of $\SU(2)$ is therefore not a smallness but an identity: the parallel
and antiparallel channel sums coincide exactly at rank two.
\end{corollary}
\gid{LEAN:hopping\_two}

\begin{corollary}[The planar approach is from below]
\label{cor:deficit}
With $D = (N^2-1)(2N^2-1)(4N^2-9)$,
\begin{equation}
  D - 4N^3 \cdot 2N(N^2-4) \;=\; 2N^4 + 31N^2 - 9 \;>\; 0
  \qquad (N \ge 1),
\end{equation}
so $4N^3 t_N = 1 - (2N^4+31N^2-9)/D < 1$. Hence
$t_N < 1/(4N^3)$ for every $N$, and $t_N \sim 1/(4N^3)$: the
$N^{-3}$ cancellation is exact and what it leaves is a strictly positive
deficit.
\end{corollary}
\gid{LEAN:hopping\_deficit\_numerator}

\noindent\Tzero{}, at \srcpath{lean/Workhouse/Basic.lean:48}.
Corollary~\ref{cor:deficit} is the first appearance of the rank grading
that \S\ref{sec:grading} develops: the leading power of $N$ is fixed by
a cancellation, and the surviving term is a polynomial one can write
down.

\subsection{Independent corroboration of the channel weights}

The weights of Proposition~\ref{prop:weingarten} have an external
provenance, from work written without knowledge of this program.

An efficient-truncation study of $\SU(N_c)$ lattice gauge theory for
quantum simulation~\cite{CBB_2026} gives, in its Appendix~D, the exact
finite-rank plaquette matrix element
\begin{equation}
  \abs{M_\rho} = \sqrt{\frac{\dim\rho}{\dim A \,\dim R}} .
\end{equation}
Taking both factors fundamental gives $\abs{M_\rho}^2 = d_\rho/N^2$,
which is \emph{identically} the numerator of the weight formula used
here. The four shared-link channel weights, and with them $A_N$, $B_N$
and $t_N = B_N - A_N$, therefore rest on a matrix element an
independent group wrote down for an unrelated purpose. That truncation
is also the first published one retaining all four shared-link
irreducible representations.

The agreement is evidence rather than bookkeeping precisely because
neither side knew of the other. It does not make $t_N$ published: the
rank law \eqref{eq:tN}, the vanishing at $N = 2$, and the deficit of
Corollary~\ref{cor:deficit} are assembled here.

A caution attaches in the other direction. The leading-large-$N_c$
qutrit plaquette Hamiltonian~\cite{CB_2024} explicitly decouples the
charge-odd combination, so the sector in which $t_N$ lives is inert in
that model --- effectively $t_3 = 0$ at leading order in large $N_c$,
against the finite-$N$ channel-complete value $5/612$. That paper is
recorded in this program's index as not yet obtained, and the statement
here is taken from imported bridge documents rather than checked
against the paper itself.

\subsection{What this section does not establish}

\notasserted{Convergence of the strong-coupling series; anything about
the spectrum at finite $u$ beyond second order; and nothing about
$\SU(2)$ at higher orders, where Corollary~\ref{cor:su2} does not
apply. Theorem~\ref{thm:tN} is a statement about one coefficient of a
formal expansion, established exactly.}
% ---- end 12_second_order.tex ---------------------------------------
% ---- begin 13_fourth_order.tex -------------------------------------
\section{Fourth order, and how a disputed coefficient was decided}
\label{sec:fourth}

Fourth order is where the carrier first couples to its complement, and
it is where this program spent the largest single block of its effort.
The mathematics is one theorem and one number. The number was disputed
for months between two irreconcilable values, and the way it was settled
is, in this author's view, the more instructive half of the story.

\subsection{The structure}

The fourth-order effective operator splits into Hodge terms and one
geometric term coupling the carrier to its complement. Restricted to
range-one operator structures, the shape content collapses: each such
structure contributes a fixed shape coefficient, and every one of them
gives $B = D = 0$ identically. The tier collapse is simply what
range-one looks like.

\begin{theorem}[Fourth-order support, at $\SU(3)$]
\label{thm:tiercollapse}
The recorded fourth-order kernel has support
$\{I, U, S, S^2, R\}$, and the shape data closes in that span. Solved
over the whole Brillouin zone in exact rationals rather than fitted, the
shape table is determined by three signed amplitudes:
\begin{equation}
  C_{\mathrm{shp}} \;=\; -\tfrac{5}{96} - u - \tfrac12(\rho + \pi),
  \qquad
  u_2 = 2u, \qquad \nu = -\big(\tfrac{5}{48} + 4u\big),
  \label{eq:kernelform}
\end{equation}
where $u$ is the two-hop chain cumulant.
\end{theorem}
\gid{RESULT:TIER\_COLLAPSE\_ACTUAL\_H4\_SUPPORT, RESULT:TIER\_COLLAPSE\_IS\_R\_DEGREE}

\begin{scopenote}
``The recorded kernel'' means the $189$-record historical kernel, which
is $N = 3$. The operator splitting is an $\SU(3)$ statement; it is not
established at general rank, and framing it as an $\SU(N)$ identity
would be an overreach. The rank-uniform statements of this program are
the coefficient laws of \S\ref{sec:second} and \S\ref{sec:grading}, not
this splitting.
\end{scopenote}

Clearing the $1/q$ makes the shape ansatz a linear system over Laurent
polynomials in $e^{ik_j}$, which is why the whole zone can be solved
exactly instead of sampled.

There is a second, equivalent presentation. In the Hodge form the
two-hop weight is absorbed, $\tilde\pi = \pi + 2u$, and
\begin{equation}
  C_{\mathrm{shp}} \;=\; -\tfrac{5}{96} - \tfrac12\big(\rho + \tilde\pi\big),
  \label{eq:hodgeform}
\end{equation}
so that $u$ drops out of $C$ entirely and exactly two shared-link
amplitudes remain. Equations \eqref{eq:kernelform} and
\eqref{eq:hodgeform} are the same statement. They must not be
combined --- writing $u$ and $\tilde\pi$ in the same expression
double-counts the two-hop weight, and that error has been made.

\subsection{The dispute}

Two independently recorded kernels --- a historical pipeline from June
2026 and a later cold dump --- gave different values for the off-axis
shape coefficient $C_{\mathrm{shp}}$. Both were exact rationals. Neither
could be promoted, because the repository's governing rule is that a
dispute closes by derivation or not at all: code may not pick a side,
average them, or prefer whichever rational looks more authoritative.

Comparing the two kernels record by record localized the disagreement
rather than resolving it. Both have the same six weight orbits and the
same normalized shape table; both satisfy the same two relations above.
What differs is the skeleton scale and the \emph{signs} of two
amplitudes. One scale factor explains $144$ of the $189$ records to
$2\times10^{-12}$: the rotation bulk is structurally shared, and the
dispute is three amplitudes, not two kernels.

\subsection{How it was settled}

By building a third implementation that shares no primitive with either,
and by recovering the historical pipeline's own intermediate ledger.

The third implementation is a Wilson-loop calculus written from scratch:
the free Hamiltonian by the Fierz identity as a rewiring of index ports
--- a like pair swaps wires, an unlike pair is cut and crossed, which is
unitarity; Haar integrals by the $U(N)$ Weingarten function as the Gram
pseudoinverse on $S_n$, with $\det U = 1$ inserted on virtual slots for
the determinant families; and the resolvent by per-link irreducible
projectors, so that every component is an exact eigenvector of $H_0$
with a known energy. It reproduces the four second-order constants, the
two-hop weight with the correct sign on three chain types, and the
single-contact and shared-link dressings, in seconds --- and then gives
the corner cumulant to the digit.

Separately, a survey of the corpus located the historical pipeline's own
Stage-3I ledger: $4{,}221$ ordered words whose hash matches the input
the kernel copies quote. Its rooted assembly, ported verbatim, reproduces
the pinned $189$ records exactly. Grouping each word's rooted images by
the set of plaquettes it touches then yields the historical kernel's own
connected cumulants --- a reading the original pipeline never performed
on its own data.

\begin{theorem}[Adjudication of the off-axis coefficient]
\label{thm:cshp}
The historical kernel and the independent assembly agree on \emph{every}
cluster of the rotation record --- pair, fourteen chains, two fans, two
corners --- with identical rationals in both $C$-parity sectors, except
one. The exception is the adjacent-face cube completion. The historical
ledger holds $8$ of the $24$ insertion orderings, each with exactly the
weight the multi-loop model assigns it, namely
$(2,1,4,2,8,8,2,4)$ in units of $1/1536$, totalling $31/1536$; the
sixteen it lacks sum to $-25/512$. The opposite-face completion in the
same ledger has all $24$.

Consequently, in the kernel's basis,
\begin{equation}
  C_{\mathrm{shp}}
  \;=\; C_{\mathrm{historical}} + \frac{25}{1024}
  \;=\; -\,\frac{13035490122347}{550663802582400}
  \;=\; -0.0236723\ldots
  \label{eq:cshp}
\end{equation}
Substituted into the $24$ records and solved over the whole zone, this
gives $A = 5/48$ and $B = D = 0$ with no residual.
\end{theorem}
\gidc{C2, ADR:0024}{workhouse why C2}

\noindent\Tone{}, suite ``the third implementation and the historical
ledger'': nine exact checks, three of them registered findings.

The shortfall has a closed form at every rank. The adjacent-face
completion is
\begin{equation}
  c_4^{\,\mathrm{adj}}(N) = \frac{-106}{N(N^2-1)^3},
  \qquad
  c_4^{\,\mathrm{adj}}(3) = -\frac{53}{768},
\end{equation}
in ratio $53/80$ to the opposite-face completion $-160/(N(N^2-1)^3)$;
the $-106$ is the primitive $-88$ plus ten multi-loop orderings at
$-18$.

\subsection{Three things this episode establishes}

\emph{The disagreement was an enumeration shortfall, not a
disagreement about mathematics.} A pipeline dropped sixteen of
twenty-four orderings on one cluster and nowhere else. That is a
diagnosable defect, and it was diagnosed ordering by ordering rather
than adjudicated by preference.

\emph{A basis error can survive validation.} An earlier assembly of the
same amplitude was correct and was read in the conjugate basis, giving
a third value. Every $C$-odd component of the historical ledger is the
negative of the assembly's value in one traversal while every $C$-even
one is identical --- the signature of conjugating a single face. The
assembly was right; its basis was not. The decision record amends it
rather than withdrawing it.

\emph{The right answer arrived by a forbidden route first.} The value
\eqref{eq:cshp} is the one an all-rank continuation formula produces at
$N=3$ --- and that derivation was retracted, because the substitution it
used is not permitted. The retraction stands. That the forbidden route
happened to reach the number the legitimate assembly later confirmed is
not a vindication of the route, and the register keeps both facts.

\begin{scopenote}
Both recorded kernels remain in the registry with their verdicts: the
historical superseded, the cold dump falsified, beside the proven
assembled value. Nothing was deleted. One follow-up remains untried ---
the pair cluster from the third engine, pure-six family included ---
which can sharpen this result but not overturn it.
\end{scopenote}
% ---- end 13_fourth_order.tex ---------------------------------------
% ---- begin 14_rank_grading.tex -------------------------------------
\section{Rank grading: cancellation and determinant families}
\label{sec:grading}

The strong-coupling coefficients of \S\ref{sec:second} and
\S\ref{sec:fourth} are exact rational functions of $N$. That makes a
question available which is invisible at any single rank: how does each
order of the series scale with the rank, and why?

The answer has two halves. At even orders the scaling is fixed by an
exact cancellation among four representation channels. At odd orders
the coefficients vanish outright for even $N$, and for odd $N$ they
first appear at order $N-2$ with a closed-form vertex. Neither statement
can be seen from numerics at $N = 3$.

\subsection{The one-face grading law}

Let $S_m(N)$ be the $C$-parity splitting of the one-plaquette
strong-coupling diagonal at order $m$,
\begin{equation}
  S_m(N) \;=\; (\text{odd})_m - (\text{even})_m \;=\; -2\,K_m[0][1].
\end{equation}

\begin{theorem}[Single-face planar grading through tenth order]
\label{thm:grading}
At even orders $m = 2, 4, 6, 8, 10$ the exact rational function of $N$
satisfies
\begin{equation}
  S_m(N) \;=\; c_m\, N^{-(2m-1)}\big(1 + O(N^{-2})\big),
  \label{eq:grading}
\end{equation}
with
\begin{equation}
  c_2 = -4, \quad c_4 = -32, \quad c_6 = -\tfrac{1748}{3},
  \quad c_8 = -\tfrac{123332}{9}, \quad c_{10} = -\tfrac{49593808}{135}.
\end{equation}
Writing $u = N^2\tau/2$ as in \eqref{eq:tau}, the \emph{leading} part of
the order-$m$ contribution is $(c_m/2^m)\, N \tau^m$. Every even order
therefore carries the same single power of $N$, with
\begin{equation}
  b_2 = -1, \quad b_4 = -2, \quad b_6 = -\tfrac{437}{48},
  \quad b_8 = -\tfrac{30833}{576}, \quad b_{10} = -\tfrac{3099613}{8640},
  \qquad b_m = c_m/2^m .
\end{equation}
\end{theorem}
\gid{RESULT:SINGLE\_FACE\_PLANAR\_GRADING}

\noindent Recorded \texttt{proven} with \texttt{output-certified}
evidence. The exponent and the coefficients are exact limits of exact
rational functions of $N$, not numerical extrapolations. An independent
integer-rank engine reproduces $c_2, c_4, c_6, c_8$ by extrapolation to
within $1.4 \times 10^{-5}$ relative --- a cross-check by a different
method, not the source of the values.

The word \emph{leading} in Theorem~\ref{thm:grading} is doing work and
should not be dropped. Each $S_m$ is an exact rational function of $N$
with relative corrections in $N^{-2}$, so the reduction to $\tau$ is a
planar-limit statement, not an identity at finite rank. Explicitly, at
second order
\begin{equation}
  S_2(N) = -4N^{-3} - 4N^{-7},
  \qquad\text{so}\qquad
  \frac{S_2\,u^2}{N} = -\tau^2 - \frac{\tau^2}{N^4}.
\end{equation}
The single-face band per unit $N$ is therefore a function of $\tau$
alone \emph{as} $N \to \infty$ at fixed $\tau$, and not at finite $N$.
A statement of the second form would be false, and an earlier internal
summary of this result made it.

\subsection{The grading is a cancellation theorem}

Equation~\eqref{eq:grading} is a much smaller number than its
constituents, and the reason is exact.

\begin{theorem}[Four-channel cancellation]
\label{thm:cancellation}
Decomposing $S_m$ by the intermediate state of the perturbed vector
gives exactly four contributing channels at every even $m$ from $4$ to
$10$: the singlet $\big((),()\big)$, the adjoint
$\big((1),(1)\big)$, and the two-box states $\big((2),()\big)$ and
$\big((1,1),()\big)$. Each is of order $N^{-(m-1)}$, with leading
coefficients
\begin{equation}
  (-4,\ +2,\ +1,\ +1) \times 2^{\,m/2-2},
\end{equation}
which \emph{sum to zero}. At $m = 2$ only the singlet and the adjoint
contribute, at $-2$ and $+2$, and these also cancel.
\end{theorem}
\gid{RESULT:SINGLE\_FACE\_PLANAR\_GRADING}

So the individual channels are of order $N^{-(m-1)}$ and their leading
behaviour annihilates, leaving \eqref{eq:grading} at $N^{-(2m-1)}$ ---
two further orders down in the $N^{-2}$ expansion. At $m = 4$ this is
the statement that channels of order $N^{-3}$ conspire to produce a band
of order $N^{-7}$.

This is why the grading is stated as a theorem about a mechanism rather
than as an observed power. A power law read off from fitted data at
several ranks would carry no information about \emph{why} the leading
terms are absent, and would give no reason to expect
\eqref{eq:grading} to persist. The cancellation does.

\begin{scopenote}
Theorem~\ref{thm:grading} concerns the \emph{one-plaquette} diagonal
splitting at even orders $m \le 10$. It is not the cluster cumulant of
the multi-face conjecture, which is not discharged by it: the one-face
case at sixth order was already established separately, and the
shared-link pair route carries the multi-face statement. The general-$m$
law remains a conjecture. All statements are about coefficients at fixed
order, not about convergence.
\end{scopenote}

\subsection{Odd orders are determinant families}

\begin{theorem}[Centre parity]
\label{thm:centreparity}
For \emph{even} $N$, every odd-order matrix element of the des Cloizeaux
effective operator of the strong-coupling band vanishes identically ---
on every cluster, in both $C$-parity sectors, at every order.

For \emph{odd} $N$, an odd order $m$ vanishes unless $m \ge N-2$. At
$m = N-2$ the only surviving element is the one-plaquette vertex
\begin{equation}
  H_{N-2}(\bar F, F)
  \;=\; -\,\frac{\big(N/(N+1)\big)^{N-3}}{\big((N-3)!\big)^2},
  \label{eq:centrevertex}
\end{equation}
which is $-1$ at $N = 3$ and $-\tfrac{25}{144}$ at $N = 5$. Hops and
leakages vanish at that order.
\end{theorem}
\gid{RESULT:ODD\_ORDER\_CENTRE\_PARITY}

The mechanism is the centre. The $\SU(N)$ Haar integral of a word
vanishes unless every link's net flux is $0 \bmod N$; on a single
plaquette the four links are private, so the intermediate states are
characters with energy $2C_2(R)$, and the flux condition at odd order
can be met only once enough boxes have accumulated to close modulo $N$.
That threshold is $N-2$, and it is why the first odd order is a
function of the rank rather than a fixed number.

\begin{scopenote}
Statements about coefficients of the strong-coupling series at fixed
order, not about convergence. On a periodic lattice the link-set lemma
requires even extents. The even orders and the overlap theorem are not
addressed here.
\end{scopenote}

\subsection{Why this refuted a claim}

The register records a claim that centre-only circuits are dynamically
dark --- that configurations carrying only centre flux cannot contribute
to the dynamics. Theorem~\ref{thm:centreparity} and its fifth-order
companion refute it: the exact fifth-order determinant correction
$\delta c_{5,\mathrm{det}}$ is nonzero. The claim was falsified by an
exact computation rather than by an estimate, which is the only way a
claim of that shape can be settled.
% ---- end 14_rank_grading.tex ---------------------------------------
% ---- begin 15_symbolic.tex -----------------------------------------
\section{The symbolic-rank engine}
\label{sec:symbolic}

Every closed form in \S\ref{sec:second}--\S\ref{sec:grading} is a
rational function of $N$ computed as such, not interpolated through
ranks. This section explains how, because the method is the most
transferable thing in the program and because getting it wrong is the
default.

\subsection{The problem with sweeping ranks}

The natural procedure is to run an exact engine at $N = 3, 4, 5, \ldots$
and fit a rational function. That is what was done first. An assembled
cluster cumulant $\beta_N$ agreed with a candidate closed form
$P_{17}(N^2)\big/\big(N\,R_{20}(N^2)\big)$ at every rank from $4$ to
$70$: sixty-seven exact agreements between exact rationals.

Sixty-seven exact agreements do not prove the identity. Rational
interpolation is unique only once the degrees are bounded, and without
a degree bound the agreements constrain nothing about $N = 71$. The
review of that work made exactly this objection --- no ``for all $N$''
without a degree bound --- and it was correct.

A degree bound could have been argued: Weingarten denominators, energy
denominators, and a count of vertices together bound the degree, and
the sixty-seven agreements would then prove the identity by uniqueness.
That argument was not made. What was done instead makes it unnecessary,
and is better.

\subsection{Running the engine over \texorpdfstring{$\Q(N)$}{Q(N)}}

Inspect where the rank actually enters the exact engine. It enters in
exactly two places that are not rational arithmetic:
\begin{itemize}
  \item the charge predicate, $\text{charge} \equiv 0 \bmod N$;
  \item the flux-difference predicate, $\abs{\mathrm{diff}} \le N$.
\end{itemize}
Everywhere else the rank appears only through rational expressions in
$C_F = (N^2-1)/(2N)$, in $1/(2N)$, and in $N^{\text{closed}}$.

So make the two predicates overridable and promote everything else to
the rational function field:
\begin{itemize}
  \item scalars become elements of $\Q(N)$, carried as a numerator and
    a monic denominator in exact polynomial arithmetic;
  \item the Weingarten function becomes the symbolic inverse of the
    Gram matrix $N^{\text{cycles}}$ on $S_n$;
  \item the single-link spectra become the $\SU(N)$ quadratic Casimirs
    of the irreducible representations that link's flux content can
    carry.
\end{itemize}

A cumulant then emerges as \emph{one rational function of $N$},
computed rather than fitted. The per-rank engine is otherwise untouched
and its own test suite still holds --- which is the check that the
symbolic run is the same calculation and not a different one.

\begin{theorem}[Closed forms over $\Q(N)$]
\label{thm:symbolic}
The fourth-order cluster cumulants, the two-hop weight in both
$C$-parity sectors, and the assembled $\beta_N$ are exact elements of
$\Q(N)$ produced by the symbolic engine. The identity
$\beta_N = P_{17}(N^2)\big/\big(N\,R_{20}(N^2)\big)$ holds as an
identity of rational functions.
\end{theorem}
\gid{ADR:0029}

\noindent The final polynomial identity is promoted to \Tzero{}, and it
is worth saying exactly which identity that is. The corpus \emph{states}
$\beta_N = P_{17}(N^2)/(N R_{20}(N^2))$ and never derives it. What Lean
checks is that the closed forms the symbolic engine returns for the
three cumulants and the adjacent-face cube completion \emph{sum to} the
corpus's coefficient --- a polynomial identity between the engine's
output and the corpus's formula. It is not a Lean derivation of the
cumulants, and the Lean source says so in its own docstring. This
document repeats the qualification because an internal summary of the
same fact dropped it.

Every eigencomponent of every resolvent is verified to be an exact
eigenvector of every single-link $H_0$, symbolically --- a check that
has no per-rank analogue, because at a fixed rank a coincidence of
numbers is indistinguishable from an identity.

\subsection{Why this is the right shape of argument}

The methodological content generalizes beyond this program.

An identity between rational functions can be established in two ways:
verify it at enough points and bound the degrees, or compute in the
function field. The first requires an argument about the size of the
objects, which is a second piece of mathematics and is where the
difficulty hides. The second requires only that the algorithm's
dependence on the parameter be rational --- which is a property one can
\emph{inspect}, predicate by predicate, rather than estimate.

The title of the decision record is the summary: the degree bound is
proved by removing it.

\begin{scopenote}
The symbolic engine is valid at ranks where the generic partition
labelling applies. Exceptional low ranks, where representations
degenerate, are handled separately and are the reason $\SU(2)$ and in
places $\SU(4)$ carry their own entries in the register.

This degree bound must not be confused with the \emph{other} degree
bound in this program's history --- the proposed kinematic mechanism
for the fourth-order tier collapse, which was retracted
(\S\ref{sec:fourth}, \S\ref{sec:register}). That one was withdrawn
because it was wrong. This one was removed because it became
unnecessary.
\end{scopenote}
% ---- end 15_symbolic.tex -------------------------------------------
% ---- begin 15_sixth_order.tex --------------------------------------
\section{Sixth order: a shape that cannot be removed}
\label{sec:sixth}

The fourth-order shape data closes in a five-element span. The question
this section answers is whether the sixth order stays inside it. It does
not, and the proof is a divisibility argument that needs no knowledge of
the sixth-order dynamics at all.

\subsection{Conventions}

In the Bloch variables $z_j$ write
\begin{equation}
  a_j = 2 - z_j - z_j^{-1},
  \qquad q = \sum_j a_j,
  \qquad e_2 = \sum_{i<j} a_i a_j,
  \qquad e_3 = a_1 a_2 a_3 ,
\end{equation}
and let $U = \psi\psi^{*}$ with $\psi^{*}\psi = q$, $P_c = U/q$,
$Q_c = I - P_c$, $S = L_{\downarrow} - 4I$, and $R$ the cross-plane part
of $S$.

Two warnings, both of which have caused errors here. The invariants
$e_2, e_3$ are Bloch invariants, \emph{not} representation Casimirs;
and $U$ is the up-Laplacian, not its normalized projector, so its kernel
is the complementary plane rather than the carrier line. Cleared
identities hold at $q = 0$; normalized carrier formulas require
$q > 0$.

\subsection{The complete folded operator}

Let $P_E$ be the isolated electric model-space projector,
$D = Q_E/(E_0 - H_0)$, and $V_0$ the perturbation. If
$P_E V_0 P_E = a P_E$, set $V = V_0 - aI$ \emph{on the whole Hilbert
space}: shifting the complementary block as well is essential. With
$\chi_0 = P_E$, intermediate normalization gives the exact recursions
\begin{equation}
  K_n = P_E V \chi_{n-1},
  \qquad
  \chi_n = D V \chi_{n-1} - D \sum_{j=1}^{n-1} \chi_{n-j} K_j ,
\end{equation}
and with $M = P_E + \chi^{*}\chi$ the canonical Hermitian effective
operator is $M^{1/2} K M^{-1/2}$. Writing
\begin{equation}
  A_m = P_E V (DV)^{m-1} P_E,
  \qquad
  C_2 = P_E V D^3 V P_E,
\end{equation}
and $B_m$ for the sum of $A_m$ with exactly one of its $m-1$ slots
replaced by $D^2$:

\begin{theorem}[The eighteen-word sixth-order fold]
\label{thm:folds}
\begin{equation}
  H_6 = A_6 - \tfrac12\{A_4,B_2\} - \tfrac12\{A_3,B_3\}
        - \tfrac12\{A_2,B_4\} + \tfrac12\{A_2^2,C_2\}
        + \tfrac38\{A_2,B_2^2\} + \tfrac14 B_2 A_2 B_2 .
  \label{eq:F6}
\end{equation}
Expanded, this is exactly eighteen words in $P_E, V, D$. The same
recursion reproduces the established $H_4 = A_4 - \tfrac12\{A_2,B_2\}$.
\end{theorem}
\gid{RESULT:G9\_HERMITIAN\_FOLDED\_SUPPORT}

The $\{A_3,B_3\}$ term is present and it matters: dropping odd electric
histories gives an incomplete formula. No odd order is dropped anywhere
in the recursion. Independent rational matrix checks verify both the
invariant-subspace equation and the normalization through degree six,
with a rank-two $P_E$ and non-commuting effective coefficients.

\notasserted{Equation~\eqref{eq:F6} is the universal folded operator
support in this convention. It does not evaluate Wilson plaquette matrix
elements, establish a specific support, or carry out connected
multi-plaquette rooted subtraction. Those require the actual electric
and Haar implementation on each support.}

\subsection{Exact word symbols}

In a polynomial carrier gauge, exact symmetrization gives each word its
cleared symbol in $\Q[q, e_2, e_3]$. The table below is computed, not
inferred, and all forty words of length zero through three agree
coefficient by coefficient with independently implemented spatial
Laurent operators.

\begin{center}
\begin{tabular}{@{}ll@{\qquad}ll@{}}
\toprule
Word & $\sigma(W)$ & Word & $\sigma(W)$ \\
\midrule
$I$        & $q$          & $RR$        & $q e_2 + 3 e_3$ \\
$U$        & $q^2$        & $RUR$       & $4 e_2^2$ \\
$S$        & $-4q$        & $RSR$       & $(q-4)(q e_2 + 3e_3) - 4e_2^2$ \\
$S^2$      & $16q$        & $RRR$       & $-2e_2^2 - 2q e_3$ \\
$R$        & $-2e_2$      & $URR,\,RRU$ & $q^2 e_2 + 3 q e_3$ \\
$UR,\,RU$  & $-2q e_2$    &             & \\
\bottomrule
\end{tabular}
\end{center}

A word census must retain coefficients and combine them before testing
shape membership: $\sigma(RUR + 2\,RRR) = -4q e_3$, although both
individual words contain $e_2^2$. Support is not unique modulo operator
identities, and the invariant object being tested is the complete
carrier polynomial. Missing direct input is represented as unknown, not
as zero.

\subsection{The band coefficient at sixth order}

At fixed nonzero momentum, with scalar anchors removed,
\begin{equation}
  H_{\mathrm{eff}}(u) = u^2 t_3\, q\, Q_c + u^3 H_3 + u^4 H_4
    + u^5 H_5 + u^6 H_6 + O(u^7).
\end{equation}
The recorded $\SU(3)$ third-order factorization gives
$Q_c H_3 P_c = 0$, so the off-diagonal block first appears at order
$u^2$ in the scaled operator, from $H_4$. Expanding the Schur complement
for the scalar carrier branch,
\begin{equation}
  \lambda_6 = \frac{\sigma(H_6)}{q}
    - \frac{q\,\sigma(H_4^2) - \sigma(H_4)^2}{t_3\, q^3}.
  \label{eq:S6}
\end{equation}
This is not an assumption that $H_5$ vanishes. $H_5$ first couples to
$H_4$ at order $u^7$ in $H_{\mathrm{eff}}$, as does the $uH_3$
correction to the complementary inverse, and unknown off-diagonal $H_6$
contributes later still. The decoupling of $H_3$ is an explicit recorded
input rather than an assumption, and any scalar vacuum subtraction is
part of $H_6$ and changes nothing.

Evaluating all twenty-five ordered pairs from the actual $H_4$ support
$\{I,U,S,S^2,R\}$, including their projected subtraction, gives zero for
twenty-four of them. Only $RR$ survives, in agreement with direct
composition of the full assembled spatial kernel. With
\begin{equation}
  \kappa = \frac{4C_{\mathrm{shp}}^2}{t_3}
  = \frac{169924002729806205028788409}{619343593697933825385600000} > 0,
  \qquad t_3 = \tfrac{5}{612},
\end{equation}
the known dynamical mixing support is
$-(\kappa/q)\,RR + (\kappa/q^2)\,RUR$, with normalized expectation
\begin{equation}
  \lambda_6^{\mathrm{mix}}
  = -\kappa\left(\frac{e_2}{q} + \frac{3e_3}{q^2}
      - \frac{4e_2^2}{q^3}\right).
\end{equation}

\notasserted{The $RUR$ here arises from inserting $P_c = U/q$ in the
projected subtraction. It is not evidence that a direct six-insertion
electric history carries the local word $RUR$, and the ratio $1:3:-4$
belongs to this mixing term rather than to the full coefficient.}

\subsection{Non-cancellation}

\begin{theorem}[An extra rational shape survives every local direct term]
\label{thm:noncancellation}
Assume the order-six electric-shell effective operator is a
finite-range, translation-covariant stencil with Laurent-polynomial
entries in this basis, and write $F = \sigma(H_6)$. Then \eqref{eq:S6}
gives
\begin{equation}
  q^3 \lambda_6 = q^2 F - \kappa\big(q^2 e_2 + 3q e_3 - 4e_2^2\big),
  \label{eq:N6}
\end{equation}
and $q$ does not divide the right-hand side, for \emph{any} Laurent
polynomial $F$.
\end{theorem}

\begin{proof}
Modulo $q$, the right-hand side of \eqref{eq:N6} is $4\kappa e_2^2$.
Evaluate at
\begin{equation}
  z_1 = z_2 = -1, \qquad z_3 = 5 + 2\sqrt{6},
\end{equation}
all nonzero, so the point lies in the Laurent ring. Since
$z_3^{-1} = 5 - 2\sqrt6$, we get $a = (4, 4, -8)$, hence $q = 0$ and
$e_2 = -48$. The value there is $4\kappa \cdot 48^2 = 9216\,\kappa \neq 0$,
and a Laurent polynomial divisible by $q$ would vanish at that point.
\end{proof}
\gid{RESULT:G9\_LOCAL\_SIXTH\_NONCANCELLATION}

\noindent\Tone{}. Consequently $\lambda_6$ cannot lie in the
fourth-order shape span $\{1, q, e_2, 4e_2/q, e_3/q\}$, and no local
direct $RUR$ term removes the extra rational shape: a direct $h\,RUR$
contributes $4h e_2^2/q$, which cannot cancel an $e_2^2/q^3$ component.
For symmetric $F$ the remainder modulo $q^2$ is exactly
$4\kappa e_2^2 - 3\kappa q e_3$, so both $e_3/q^2$ and $e_2^2/q^3$ are
present. The $e_2/q$ term alone \emph{can} be cancelled, by
$F = \kappa e_2$ --- for instance a local $H_6 = -(\kappa/2)R$ --- so
the $1:3:-4$ ratio is not a constraint on all sixth-order shapes.

Note what the proof does not use: any knowledge of $F$. The extra shape
is forced by the fourth-order data alone, through the Schur complement,
whatever the sixth-order dynamics turns out to be.

\subsection{A withdrawn derivation}

An earlier version of this work enumerated closed rooted walks --- $90$
four-step unit-edge walks, of which $66$ retrace immediately and $24$
are planar squares, together with $192$ non-retracing six-step walks
using all three axes --- and read a sixth-order amplitude off that
census.

The derivation is withdrawn. Edge walks are not ordered plaquette
insertions with electric states, Haar weights, projectors and folds, and
no map between them was supplied; the census therefore determines no net
coefficient of any of the candidate words. The old amplitude formula was
assigned rather than derived. The old carrier function used
$\sigma(R) = 0$ and $\sigma(U) = q$, where the actual unnormalized
symbols are $-2e_2$ and $q^2$; it returned a nonzero expression for a
defect it separately reported as zero. The implementation now computes
the symbols and tests the defect, and the old interface raises an
explicit missing-dynamics error rather than returning a number.

What is withdrawn is a claimed derivation, not the possibility of a
physical amplitude of that shape. Theorem~\ref{thm:noncancellation} is
independent of the census entirely. Computing the direct term on a
connected multi-plane support is Problem~\ref{prob:g9}, and it is the
cheapest decisive computation currently open in this program.
% ---- end 15_sixth_order.tex ----------------------------------------
% ---- begin 16_wick_shell.tex ---------------------------------------
\section{A rank-uniform local shell construction}
\label{sec:wick}

This section is a separate result in its own regime, included because it
is exact, rank-uniform, and independent of everything above --- not
because it continues it. The object is the formal local Wilson
oscillator expansion around a single plaquette, and the point is that
its spectrum admits a finite exact recursion at every perturbative
order with coefficients in $\Q[N, N^{-1}]$, needing no Gram-matrix
inversion.

Conflating this construction with the flux band of
\S\ref{sec:carrier}--\S\ref{sec:sixth}, or with the Wilson transfer
operator of \S\ref{sec:wilson}, would be an error. They are four
different objects --- the physical $\SU(N)$ transfer band, the $\SU(3)$
internal Hamiltonian carrier, the conditional block, and this local
oscillator --- and no cross-regime continuation between them is
asserted anywhere in this document.

\subsection{The exact algebra}

Let $X$ be a traceless Hermitian $N \times N$ matrix with density
proportional to $e^{-\Tr(X^2)}$, write $P_k = \Tr(X^k)$ with $P_0 = N$
and $P_1 = 0$, and use a traceless Hermitian basis normalized by
$\Tr(T_a T_b) = \delta_{ab}$. The covariance is
\begin{equation}
  \mathbb{E}\big[X_{ab} X_{cd}\big]
  = \tfrac12\left(\delta_{ad}\delta_{bc}
    - \tfrac{1}{N}\delta_{ab}\delta_{cd}\right).
\end{equation}
With $D = \tfrac12\Delta$ and $E$ the polynomial Euler operator, the
completeness identity for this basis gives the exact relations
\begin{align}
  D P_k &= \frac{k}{2}\sum_{l=0}^{k-2} P_l P_{k-2-l}
           - \frac{k(k-1)}{2N} P_{k-2},
  \\
  D(fg) &= (Df)g + f(Dg) + \nabla f \cdot \nabla g,
  \\
  \nabla P_k \cdot \nabla P_l
        &= k\,l\left(P_{k+l-2} - \frac{1}{N} P_{k-1}P_{l-1}\right).
\end{align}
These are exact at every integer rank, including ranks at which some
traces become dependent --- which is the property that makes the
construction rank-uniform rather than rank-by-rank. For homogeneous $F$
of positive degree $d$, Gaussian integration by parts gives
$\mathbb{E}_\mu[F] = \mathbb{E}_\mu[DF]/d$, and with
$\mathbb{E}_\mu[1] = 1$ and vanishing odd moments this is a finite exact
moment recursion.

\subsection{Shell projectors without a Gram inverse}

\begin{proposition}[Terminating Wick conjugation]
\label{prop:wick}
Put $L = E - D$. Since $[E, D] = -2D$, the terminating exponential
$W = \exp(-D/2)$ satisfies
\begin{equation}
  L\,W = W\,E .
\end{equation}
Hence, with $H_d$ the extraction of the homogeneous degree-$d$ part,
\begin{equation}
  \Pi_d = \exp(-D/2)\, H_d \,\exp(D/2)
\end{equation}
is the exact oscillator-shell projector, and on each finite polynomial
space $R_s = \sum_{d \neq s} \Pi_d/(s-d)$ is the exact shell resolvent.
\end{proposition}

The content of Proposition~\ref{prop:wick} is that $L$ and $E$ are
conjugate by a \emph{terminating} exponential, so the shell
decomposition of the interacting operator is the shell decomposition of
the Euler operator, transported. No Gram matrix is formed and none is
inverted, which is precisely the step that makes rank-uniformity
delicate elsewhere in this program.

\subsection{The coefficients}

\begin{theorem}[All-rank local coefficients]
\label{thm:wickcoeffs}
The terminating conjugation gives exact shell projectors and a unique
all-order simple-branch formal recursion in $\Q[N, N^{-1}]$, with state
degree at most $s + 4r$. The two formal local gaps have exact all-rank
coefficients through $\beta^{-2}$, including $c_3^{-}$ and
$c_4^{\pm}$. The five corrections $c_0, \ldots, c_4$ are strictly
negative at every allowed rank, by positive shifted numerator
polynomials.
\end{theorem}
\gid{RESULT:LOCAL\_CLASS\_WICK\_RECURRENCE, RESULT:LOCAL\_CLASS\_WICK\_COEFFICIENTS, RESULT:LOCAL\_CLASS\_WICK\_SIGNS}

The construction applies to the vacuum, the first charge-even shell for
$N \ge 2$, and the first charge-odd shell for $N \ge 3$. It recovers the
previously archived $c_0$, $c_1$, $c_2$ formulas in both sectors and the
archived all-rank $c_3^{+}$, and additionally derives $c_3^{-}$ and both
all-rank $c_4$ formulas. The certificate carries full polynomial
eigen-equation residuals rather than only agreement with the archived
constants --- a distinction worth stating, because agreement with a
previous value is a weaker check than satisfying the equation.

The recovery is by an independent route: the archive engine uses
full-GUE cut-and-join and then subtracts the trace, while this one
applies $D$ directly in traceless coordinates.

\subsection{A companion algebraic identity}

\begin{theorem}[Shifted-power carrier symbols]
\label{thm:rsmr}
For the cubic Laurent operators with $S = L_{\downarrow} - 4I$,
\begin{equation}
  \sigma(R S^m R) = (q-4)^m\big(q e_2 + 3 e_3\big) - 4\,\Pi_m(q)\, e_2^2
\end{equation}
for every natural $m$, with $\Pi$ defined polynomially also at $q = 0$.
\end{theorem}
\gid{RESULT:RSM\_R\_CARRIER\_SYMBOL\_MASTER\_THEOREM}

\noindent This closes the word table of \S\ref{sec:sixth} at all powers,
and is formalized in Lean as an all-power operator statement.

\begin{scopenote}
Theorem~\ref{thm:wickcoeffs} concerns the \emph{formal} local class
oscillator only. No error bound for the compact-group eigenvalues is
claimed, and local spectral coefficients do not by themselves identify
the interacting Wilson spectrum, an interacting-volume estimate, or a
continuum source measure. Theorem~\ref{thm:rsmr} is an algebra identity
and carries no sixth-order dynamical population claim.
\end{scopenote}
% ---- end 16_wick_shell.tex -----------------------------------------
% ---- begin 17_external.tex -----------------------------------------
\section{External comparison}
\label{sec:external}

Published work enters this program as evidence rather than as authority.
A paper is unchecked until something checks it, exactly like an internal
document; the reason an external agreement counts for more is not that
it is published but that it was produced without knowledge of this
program, so agreement is not bookkeeping.

The map holds \LitPapers{} indexed papers and \LitEdges{} recorded
relations to named claims here. \LitVerified{} of those edges are
verified --- the paper was obtained, read, and compared against the
claim --- across \LitObtained{} distinct papers. The remainder are
recorded targets for papers not yet obtained, which is itself useful: a
paper heavily cited within this web and still unobtained is
automatically the next acquisition target.

An entry earns its place by naming a specific claim and saying what
relationship it has to it. \texttt{workhouse lit --for G19} answers the
query in the other direction, and validation rejects an entry whose
target does not resolve.

\subsection{A 1989 table, and rationals computed cold}

The sharpest external check available to this program is a
Hamiltonian strong-coupling series for $\SU(3)$ glueball masses
published in 1989~\cite{HAMER_1989}, whose Table~1 is pinned here by
digest. Under the normalization bridge $m_n = 2^{\,n-1} a_n$, its
coefficients are compared against the $\Gamma$-point coefficients this
program derives as exact rationals.

In the $C$-odd ($1^{+-}$) channel --- the sector of the carrier, whose
series is \eqref{eq:thirdorder} --- the agreement runs through third
order:

\begin{center}
\begin{tabular}{@{}l l l l@{}}
\toprule
Order & Published & This program & Gap (bound) \\
\midrule
$n=0$ & $\tfrac{16}{3}/2$ & $\tfrac83 = E_{\mathrm{flat}}(0)$ & exact \\
$n=1$ & $a_1 = +1$ & the $u$ coefficient $+1$ & exact \\
$n=2$ & $2a_2 = 0.0359477124184$ & $\tfrac{11}{306}$
      & $9.94\times10^{-14}$ ($2\times10^{-13}$) \\
$n=3$ & $4a_3 = -0.437135556836$ & $-\tfrac{109151}{249696}$
      & $1.11\times10^{-12}$ ($3\times10^{-12}$) \\
\bottomrule
\end{tabular}
\end{center}

\noindent In the $C$-even ($0^{++}$) channel, $2a_2 = -0.21274509804$
against $-\tfrac{217}{1020}$ with gap $7.84\times10^{-13}$, and
$4a_3 = -0.1039004229144$ against $-\tfrac{54049}{520200}$ with gap
$1.36\times10^{-13}$. The sign structure matches as well: $a_1 = -1$ in
the scalar channel, where the carrier's is $+1$.

Each bound above is the printed half-ulp of the published table scaled
with margin, so these are agreements to the precision at which the 1989
table was printed --- not approximate corroborations. The rationals
$11/306$ and $109151/249696$ are the third-order coefficients of
$E_{\mathrm{flat}}$ in \eqref{eq:thirdorder}, computed here from Haar
integration and the cellular complex with no knowledge of the table.

At fourth order, $8a_4 = -0.7751458630184$ against
$m^{(4)}_{\Gamma} = -0.7751458630189173$, gap $5.17\times10^{-13}$
against a bound of $5.3\times10^{-13}$.

\begin{scopenote}
These are \Ttwo{} checks: float agreement within a stated tolerance, to
the precision of a printed table. They corroborate the coefficients;
they do not prove them, and the exact statements of
\S\ref{sec:second}--\S\ref{sec:grading} do not rest on them. Note also
that the pinned revision of the source compares only $2a_2$ and $4a_3$
in its own text --- the fourth-order agreement is registered here and
not leaned on there.
\end{scopenote}

\subsection{The two contradicting edges}

Two edges in the map are recorded as \emph{contradicts}, both from a
1976 strong-coupling series, and both against this program: one against
the third-order $C$-odd coefficient, one against the $\Gamma$-point
fourth-order value. They originate in a disagreement the 1989 paper
itself reports at orders $x^3$ and $x^4$ with the series it supersedes.

Both are recorded as not yet obtained. Until that primary source is
read, this program has two unresolved external contradictions against
coefficients it otherwise regards as settled, and says so rather than
resolving them by the age of the sources. Obtaining that paper is
listed with the acquisition targets below.

\subsection{How far a novelty claim here can be trusted}

Not as far as the size of the index suggests, and the limitation is
worth stating precisely rather than as a disclaimer.

Only \LitVerified{} of \LitEdges{} recorded edges are verified. The
unverified remainder is not a random sample: it includes several of the
\emph{nearest} neighbours to this program's own results. The main
Nuclear Physics B paper of the 1981 strong-coupling glueball series is
recorded as not obtained --- only its errata were --- and the 1974 and
1979 large-$N$ papers that bear most directly on the rank grading of
\S\ref{sec:grading} are likewise unobtained, with the ledger noting for
one of them that its definitions were used from memory and only their
sign consequences checked.

Consequently: where this document says a result is novel, that means
\emph{no verified entry in this index contains it}, which is a weaker
statement than priority. Two of the strongest results are affected. The
rank-grading cancellation of \S\ref{sec:grading} sits closest to
large-$N$ counting arguments in papers this program has not read; and
the fixed-spacing transfer construction of \S\ref{sec:wilson} sits
alongside the classical positivity and self-adjoint transfer-matrix
results of Osterwalder and Seiler, which the index does register.

Acquiring the unobtained near neighbours is cheap relative to
everything else in \S\ref{sec:obligations}, and it is the one item on
that list that could \emph{subtract} from the program rather than add
to it. That is a reason to do it, not a reason to defer it.

% ---- begin gen_literature.tex --------------------------------------
% Generated by paper/master/generate_apparatus.py.  Do not edit.
\begin{longtable}{@{}l l l l l@{}}
\caption{Verified external-evidence edges. Each row is a published result that was obtained, read, and checked against a named claim of this program. Relation is the ledger's own classification; \emph{contradicts} and \emph{confusable} are listed first because they are the rows that cost the program something.}\label{tab:literature}\\
\toprule
Paper & Ref. & Year & Bears on & Relation \\
\midrule
\endfirsthead
\toprule Paper & Ref. & Year & Bears on & Relation \\ \midrule \endhead
\bottomrule
\endfoot
\texttt{\small HSB\_\allowbreak{}2000} & \cite{HSB_2000} & 2000 & \texttt{HAMER\_\allowbreak{}A4\_\allowbreak{}NUM} & confusable with \\
\texttt{\small HAMER\_\allowbreak{}1989} & \cite{HAMER_1989} & 1989 & \texttt{BAND\_\allowbreak{}EVEN\_\allowbreak{}BOTTOM} & corroborates \\
\texttt{\small HAMER\_\allowbreak{}1989} & \cite{HAMER_1989} & 1989 & \texttt{C1} & corroborates \\
\texttt{\small LLL\_\allowbreak{}2006} & \cite{LLL_2006} & 2006 & \texttt{C1} & corroborates \\
\texttt{\small HAMER\_\allowbreak{}1989} & \cite{HAMER_1989} & 1989 & \texttt{D\_\allowbreak{}3} & corroborates \\
\texttt{\small SCHIERHOLZ\_\allowbreak{}1988} & \cite{SCHIERHOLZ_1988} & 1988 & \texttt{G18} & corroborates \\
\texttt{\small HAMER\_\allowbreak{}1989} & \cite{HAMER_1989} & 1989 & \texttt{M3\_\allowbreak{}EVEN\_\allowbreak{}K0} & corroborates \\
\texttt{\small KPS\_\allowbreak{}1981} & \cite{KPS_1981} & 1981 & \texttt{R14} & corroborates \\
\texttt{\small CBB\_\allowbreak{}2026} & \cite{CBB_2026} & 2026 & \texttt{R2} & corroborates \\
\texttt{\small KPS\_\allowbreak{}1981} & \cite{KPS_1981} & 1981 & \texttt{G7} & supplies a value \\
\texttt{\small HAMER\_\allowbreak{}1989} & \cite{HAMER_1989} & 1989 & \texttt{HAMER\_\allowbreak{}A4\_\allowbreak{}NUM} & supplies a value \\
\texttt{\small MUNSTER\_\allowbreak{}1985\_\allowbreak{}TM} & \cite{MUNSTER_1985_TM} & 1985 & \texttt{C2} & supplies a comparison \\
\texttt{\small CAO\_\allowbreak{}2023} & \cite{CAO_2023} & 2023 & \texttt{C7} & supplies a comparison \\
\texttt{\small LEMOINE\_\allowbreak{}2026} & \cite{LEMOINE_2026} & 2026 & \texttt{C7} & supplies a comparison \\
\texttt{\small HAZRA\_\allowbreak{}2026} & \cite{HAZRA_2026} & 2026 & \texttt{G14} & supplies a comparison \\
\texttt{\small AT\_\allowbreak{}2020\_\allowbreak{}SU3} & \cite{AT_2020_SU3} & 2020 & \texttt{G18} & supplies a comparison \\
\texttt{\small BESIII\_\allowbreak{}2026} & \cite{BESIII_2026} & 2026 & \texttt{G18} & supplies a comparison \\
\texttt{\small CHEN\_\allowbreak{}2006} & \cite{CHEN_2006} & 2006 & \texttt{G18} & supplies a comparison \\
\texttt{\small LLL\_\allowbreak{}2006} & \cite{LLL_2006} & 2006 & \texttt{G18} & supplies a comparison \\
\texttt{\small MORNINGSTAR\_\allowbreak{}2025} & \cite{MORNINGSTAR_2025} & 2025 & \texttt{G18} & supplies a comparison \\
\texttt{\small MP\_\allowbreak{}1999} & \cite{MP_1999} & 1999 & \texttt{G18} & supplies a comparison \\
\texttt{\small AT\_\allowbreak{}2021\_\allowbreak{}SUN} & \cite{AT_2021_SUN} & 2021 & \texttt{G19} & supplies a comparison \\
\texttt{\small BB\_\allowbreak{}1983} & \cite{BB_1983} & 1983 & \texttt{G19} & supplies a comparison \\
\texttt{\small FLPS\_\allowbreak{}2026} & \cite{FLPS_2026} & 2026 & \texttt{G19} & supplies a comparison \\
\texttt{\small URZUA\_\allowbreak{}2019} & \cite{URZUA_2019} & 2019 & \texttt{G19} & supplies a comparison \\
\texttt{\small BORGA\_\allowbreak{}2024} & \cite{BORGA_2024} & 2024 & \texttt{G5} & supplies a comparison \\
\texttt{\small CAO\_\allowbreak{}2023} & \cite{CAO_2023} & 2023 & \texttt{G5} & supplies a comparison \\
\texttt{\small OBZ\_\allowbreak{}1985} & \cite{OBZ_1985} & 1985 & \texttt{G5} & supplies a comparison \\
\texttt{\small AT\_\allowbreak{}2021\_\allowbreak{}SUN} & \cite{AT_2021_SUN} & 2021 & \texttt{G6} & supplies a comparison \\
\texttt{\small CM\_\allowbreak{}2003} & \cite{CM_2003} & 2003 & \texttt{G6} & supplies a comparison \\
\texttt{\small BALAJI\_\allowbreak{}2026} & \cite{BALAJI_2026} & 2026 & \texttt{R2} & supplies a comparison \\
\texttt{\small CKS\_\allowbreak{}2021} & \cite{CKS_2021} & 2021 & \texttt{U1} & supplies a comparison \\
\texttt{\small DRS\_\allowbreak{}2021} & \cite{DRS_2021} & 2021 & \texttt{U1} & supplies a comparison \\
\texttt{\small KRS\_\allowbreak{}2023} & \cite{KRS_2023} & 2023 & \texttt{U1} & supplies a comparison \\
\texttt{\small SZH\_\allowbreak{}1997} & \cite{SZH_1997} & 1997 & \texttt{U1} & supplies a comparison \\
\texttt{\small CS\_\allowbreak{}2006} & \cite{CS_2006} & 2006 & \texttt{C7} & supplies a method \\
\texttt{\small JW\_\allowbreak{}2006} & \cite{JW_2006} & 2006 & \texttt{G17} & supplies a method \\
\texttt{\small UELTSCHI\_\allowbreak{}2004\_\allowbreak{}CLUSTER} & \cite{UELTSCHI_2004_CLUSTER} & 2004 & \texttt{G17} & supplies a method \\
\texttt{\small YAROTSKY\_\allowbreak{}2004\_\allowbreak{}GS} & \cite{YAROTSKY_2004_GS} & 2004 & \texttt{G17} & supplies a method \\
\texttt{\small BB\_\allowbreak{}1983} & \cite{BB_1983} & 1983 & \texttt{G18} & supplies a method \\
\texttt{\small JW\_\allowbreak{}2006} & \cite{JW_2006} & 2006 & \texttt{G18} & supplies a method \\
\texttt{\small NSY\_\allowbreak{}2019} & \cite{NSY_2019} & 2019 & \texttt{G18} & supplies a method \\
\texttt{\small YAROTSKY\_\allowbreak{}2004\_\allowbreak{}QP} & \cite{YAROTSKY_2004_QP} & 2004 & \texttt{G18} & supplies a method \\
\texttt{\small JW\_\allowbreak{}2006} & \cite{JW_2006} & 2006 & \texttt{G19} & supplies a method \\
\texttt{\small KLEIN\_\allowbreak{}ROSENBERGER\_\allowbreak{}2020\_\allowbreak{}WKB} & \cite{KLEIN_ROSENBERGER_2020_WKB} & 2020 & \texttt{G19} & supplies a method \\
\texttt{\small SCHIERHOLZ\_\allowbreak{}1988} & \cite{SCHIERHOLZ_1988} & 1988 & \texttt{G19} & supplies a method \\
\texttt{\small UELTSCHI\_\allowbreak{}2004\_\allowbreak{}CLUSTER} & \cite{UELTSCHI_2004_CLUSTER} & 2004 & \texttt{G19} & supplies a method \\
\texttt{\small YAROTSKY\_\allowbreak{}2004\_\allowbreak{}GS} & \cite{YAROTSKY_2004_GS} & 2004 & \texttt{G19} & supplies a method \\
\texttt{\small YAROTSKY\_\allowbreak{}2004\_\allowbreak{}QP} & \cite{YAROTSKY_2004_QP} & 2004 & \texttt{G19} & supplies a method \\
\texttt{\small JW\_\allowbreak{}2006} & \cite{JW_2006} & 2006 & \texttt{G20} & supplies a method \\
\texttt{\small JW\_\allowbreak{}2006} & \cite{JW_2006} & 2006 & \texttt{G22} & supplies a method \\
\texttt{\small JW\_\allowbreak{}2006} & \cite{JW_2006} & 2006 & \texttt{G23} & supplies a method \\
\texttt{\small MUNSTER\_\allowbreak{}1985\_\allowbreak{}TM} & \cite{MUNSTER_1985_TM} & 1985 & \texttt{G3} & supplies a method \\
\texttt{\small KRS\_\allowbreak{}2023} & \cite{KRS_2023} & 2023 & \texttt{R2} & supplies a method \\
\end{longtable}

\nocite{*}

% ---- end gen_literature.tex ----------------------------------------
% ---- end 17_external.tex -------------------------------------------

\part{The construction at fixed lattice spacing}\label{part:fixed}
% ---- begin 20_wilson.tex -------------------------------------------
\section{The construction at fixed lattice spacing}
\label{sec:wilson}

This is the strongest constructive statement in the program. It is also
the one most often misread, so it is stated here with its hypotheses,
its interval, and an explicit list of what it does not assert.

Throughout this section $\varepsilon$ denotes the temporal mesh. It is
\emph{not} the planar variable $\tau$ of \eqref{eq:tau}; the two are
unrelated and the symbols are kept distinct deliberately.

\subsection{The theorem}

\begin{theorem}[Infinite-volume Wilson band and literal source frame]
\label{thm:wilsonband}
Under the hypotheses listed below, the Perron-normalized Wilson transfer
operator has a common strong limit along open-box and centered-periodic
exhaustions, in the transported product representation. The limit is a
positive contraction fixing the vacuum; it differs from the free product
in norm by at most $1/998$, and its norm on the non-vacuum subspace is
at most $\tfrac45 + \tfrac1{998}$.

On the coupling interval
\begin{equation}
  \abs{u} \;\le\; \frac{u_*}{10\,022\,400\,000\,N},
  \label{eq:wilson-interval}
\end{equation}
the complete infinite-volume physical odd band is isomorphic to
$\ell^2(\Z^3) \otimes \C^3$. The projected literal plaquette-source
synthesis is \emph{onto} that entire band, with Gram bounds
\begin{equation}
  \tfrac{9}{16}\,I \;<\; \mathcal{G} \;\le\; \tfrac{81}{64}\,I .
\end{equation}
The Euclidean Osterwalder--Schrader band and the quantum GNS band
coincide, intertwining the physical transfer and the sources.
\end{theorem}
\gid{RESULT:WILSON\_INFINITE\_PHYSICAL\_BAND}

\noindent Recorded \texttt{proven} with \texttt{analytic} evidence;
\Tthree{} for whole-statement machine certification, with its component
estimates checked across several suites.

\subsection{Hypotheses}

\begin{enumerate}
  \item All calibrated kinetic, positive Wilson-kernel, incidence and
    common-disk hypotheses of the weighted-activity and cardinality-chart
    theorems.
  \item Fixed $\SU(N)$ with $N \ge 3$, and fixed \emph{spatial} lattice
    spacing. Compatible induced open boxes, or centered periodic
    exhaustions, carrying the same local data.
  \item Temporal mesh
    $0 < \varepsilon \le \varepsilon_0 \le \log(5/4)/(4\gamma)$, with
    $m = \lceil \log(5/4)/(\gamma\varepsilon) \rceil$, and the full
    neutral physical odd free-window premise.
  \item $\gamma = C_F/2 = (N^2-1)/(4N)$, the \emph{exact} calibrated
    fundamental-link energy. This is a hypothesis about an identity, not
    a gap lower bound one is free to weaken: substituting a smaller
    number invalidates the conclusion rather than weakening it.
  \item Literal odd sources
    $\mathcal{O}_p = (\chi_p - \bar\chi_p)/\sqrt2$, with free
    orthonormal plaquette vectors of norm at most $\sqrt2\,N$.
  \item The calibrated representation-uniform kinetic-window mesh
    premise, whose $N$-dependent admissible temporal threshold is the
    previously established one. Every estimate is uniform in volume and
    in meshes satisfying these premises, and the common coupling
    interval depends on fixed $N$ and is nonzero.
\end{enumerate}

\noindent The interval \eqref{eq:wilson-interval} unpacks as
$u_0 = u_*/1252800000$ and $u_1 = u_0/(8N)$, with the accompanying
constants $\eta_{\mathrm{act}} = 1/2500$, $\epsilon = 1/998$,
$g_* = 1024/15625$ and $q_0 = \tfrac45 + \tfrac1{998} < 1$. The bound
$\eta_{\mathrm{act}} \le 1/2500$ on the integrated generator norm at
weight $1 + \log 2$ is a \emph{hypothesis} of the source
specialization, not a conclusion of the theorem; it is listed here to
prevent the reading in which the theorem is taken to establish it.

\subsection{Why totality is the substance}

The clause that carries the physical content of
Theorem~\ref{thm:wilsonband} is that the source synthesis is \emph{onto}
the band, not merely bounded into it.

Break~\ref{sec:break4} is the reason. A local probe family can decay at
one rate while the true spectrum contains a much slower state the probe
cannot see; the three-dimensional model of Proposition~\ref{prop:break4}
exhibits a factor of twenty-one between the observed rate and the actual
gap. A construction that produced a band and a decay estimate, without
totality, would prove nothing about the spectrum.

Totality here is a conclusion. The tagged source norm controls the
entire synthesis operator; close projections and a Neumann inverse
establish surjectivity in infinite dimension; and reflection-positive
history completion contains the whole band by selective source
cyclicity. Fine time steps are the positive $m$-th root of the block.

\subsection{What Theorem~\ref{thm:wilsonband} does not assert}

\notasserted{Global equality of the Osterwalder--Schrader and full
quantum high-energy spaces. Global transfer injectivity. Any spatial
continuum limit. Extension beyond $\SU(N)$, $N \ge 3$. The positive
energy gap is in electric-time units, and its conversion to physical
units is not part of the statement.}

Most importantly, and separately from the list above: the interval
\eqref{eq:wilson-interval} is a neighbourhood of $u = 0$, and the
continuum limit is taken along a trajectory on which $u \to \infty$.
Theorem~\ref{thm:wilsonband} therefore supplies no control whatever on
the continuum limit. This is Break~\ref{sec:break5}, it is structural,
and no improvement of the constant $10\,022\,400\,000$ addresses it.
A version of exactly that inference was made in this program's own
drafts and retracted; see \S\ref{sec:register}.

\subsection{The relation to the Hamiltonian carrier}

Theorem~\ref{thm:wilsonband} is a statement about the Wilson transfer
operator. The Hamiltonian carrier of Part~\ref{part:exact} is a
different object, spectrally identified with the band only on momentum
patches away from zero.

The obstruction at zero momentum is proved, not merely encountered, and
it is a statement about \emph{chains}:
\begin{enumerate}
  \item every vector in the image of the cube boundary map has exactly
    zero zone-average component;
  \item the zero-momentum carrier content is therefore carried entirely
    by the three wrapping sheets of \S\ref{sec:carrier};
  \item a sheet's normalized overlap with any fixed local observable is
    exactly $1/L$;
  \item compact two-cycles on open windows are exactly cube boundaries
    --- checked directly for small windows, and Alexander duality in
    general.
\end{enumerate}
The repository's own prose summarizes (i)--(iv) as ``every finitely
supported exactly-flat operator is $\Gamma$-blind''. That phrasing is a
convenient shorthand and this document does not adopt it: the four
statements are about chain vectors, and read literally as a claim about
arbitrary operators the shorthand is false. What follows from (i)--(iv)
is the thing actually needed --- that an interpolator for the
zero-momentum protected level must leave the exactly-flat local algebra,
coupling either through in-sector directions that are flat-degenerate at
$\Gamma$ anyway, or through out-of-sector multi-plaquette content.

\subsection{A prediction, sharper than the measurement it explains}

The published few-percent overlap of a local operator with the lightest
state is not a constant: it degrades toward the continuum. For a local
lattice source whose lowest-dimension continuum operator has dimension
$d$, the overlap law is
\begin{equation}
  Z_0(a) \;\sim\; \frac{\abs{f_d}^2\, a^{2d-3}}{2M},
  \label{eq:overlaplaw}
\end{equation}
which at $d = 4$ gives $a^5$ and so reproduces the published
measurement --- the agreement that makes \eqref{eq:overlaplaw} worth
recording at all.

But $a^5$ is the wrong exponent for \emph{this} carrier. The C-odd
carrier's leading continuum operator is $d^{abc}F^aF^bF^c$, of dimension
six, so the applicable exponent is $a^{2\cdot 6-3} = a^9$. The
extrapolation is brutal: $Z_0 \approx 1.2\times10^{-6}$ at
$a\sqrt\sigma = 0.10$ and $2.4\times10^{-9}$ at $0.05$. Still a power,
so formally survivable, and a serious practical obstacle to measuring
this state with a local operator.

The same law is the argument for smearing, quantitatively: at fixed
physical radius $R$, $Z_0 = \abs{f_R}^2 R^{2d-3}/(2M)$ is
$a$-independent. A smeared basis is structural here, not a convenience.

\begin{scopenote}
Equation~\eqref{eq:overlaplaw} is asserted, not proved here. What is
checked is the dimension arithmetic and the agreement of its $d = 4$
case with the published $a^5$. An earlier version of this program's
notes quoted $a^5$ for the carrier itself; that is corrected above.
\end{scopenote}

\begin{scopenote}
The out-of-sector dressing, the transfer-matrix identification at zero
momentum, and everything non-perturbative remain open. A particle
statement lives at fixed nonzero momentum, where no collapse occurs; the
$L^{-2}$ statement in the literature is a minimal-grid-momentum
statement governing the approach to $\Gamma$, not a fixed-momentum one.
\end{scopenote}
% ---- end 20_wilson.tex ---------------------------------------------

\part{The continuum limit: what is required, and what is ruled out}\label{part:continuum}
% ---- begin 30_continuum.tex ----------------------------------------
\section{The continuum problem}
\label{sec:continuum}

Part~\ref{part:fixed} constructs an infinite-volume theory with a
positive gap at fixed lattice spacing. This part is about the distance
between that and a continuum quantum Yang--Mills theory with a mass gap,
and it is written to be useful to someone trying to close that distance.

\subsection{The target, and what has no argument here}

The statement in question is the existence of a non-trivial quantum
field theory on $\R^{1,3}$ for $G = \SU(N)$, $N \ge 2$, satisfying the
Wightman axioms, with
\begin{equation}
  \spec(H)\big|_{\Omega^{\perp}} \subset [m, \infty),
  \qquad m > 0 .
\end{equation}

Four of its ingredients have no argument anywhere in this program and
are stated here so that no reader has to infer their absence.

\begin{enumerate}
  \item \textbf{Poincar\'e covariance.} No construction of a unitary
    representation of the covering group of $\mathcal{P}_+^{\uparrow}$ is
    given.
  \item \textbf{Restoration of rotational invariance.} The hypercubic
    lattice breaks it; nothing here restores it in the limit.
  \item \textbf{The continuum measure.} It is never constructed.
  \item \textbf{Reconstruction.} Even granting a limiting kernel, the
    passage to a Wightman theory is not carried out.
\end{enumerate}

What \emph{is} carried out is a set of estimates on the lattice side,
and this part states them, states why the six obstructions of
\S\ref{sec:breaks} block their assembly, and states in
\S\ref{sec:obligations} what would have to be proved.

\subsection{The three routes that are live}

The estimates in this program group into three approaches, which are
alternative rather than sequential. The ledger records them as
independent routes and explicitly notes that no chain of necessary
conditions runs between them.

\emph{The conditional-score route.} Control the conditional law of a
link given its complement, in the small-coupling limit, by dominating
the score of the conditional measure. The available inputs are strong:
exact heat-bath conditional laws, synchronized radial dilation following
the conditional minimum curve, the antipodal magnetic normal geometry,
exponential suppression outside the complete minimizing set, and a
uniform Agmon action gap on the complement. The obligation is
Problem~\ref{prob:m10}, and what stands between the inputs and the
conclusion is polynomial amplitude control and the tube moments, not a
missing structural idea.

\emph{The transport route.} Compare the fine and coarse theories by
transporting the vacuum and the sources along a common energy domain
against a physical clock. Here the recorded failure is instructive: an
earlier closure used a vacuum-only skew generator, and the source
generator transports the ground state rather than annihilating it. The
corrected obligation is Problem~\ref{prob:r10}, and its companion,
Problem~\ref{prob:grid}, is the requirement that the comparison's
constants not degrade with the number of interacting blocks.

\emph{The moving-time route.} Extract a limiting gap from approximate
sources and one growing observation time per cutoff. The implication
from its hypotheses to a gap is proved, and sharp: pairwise inequalities
give the optimal normalized rate, the horizon penalty and a witnessing
time, and one-atom examples show the rate cannot be improved. Its
hypotheses are the difficulty --- see Break~\ref{sec:break3}, which shows
the theorem returns exactly zero when the background error rate fails to
exceed twice the probe tilt rate, an inequality no calculation certifies
for four-dimensional Yang--Mills.

\subsection{Why none of the three currently closes}

Each route reaches a point where a summability or uniformity statement
is required, and at each such point the natural estimate is one power
short.

That is not a coincidence of technique. Break~\ref{sec:break1} exhibits
the arithmetic: on an asymptotically free trajectory the running
coupling gives increments of size $O(g_k^2) = O(1/k)$, the normalization
of the scale derivative carries no small factor, and $\sum 1/k$
diverges. Break~\ref{sec:break2} exhibits the same shortfall in the
polymer variable, where a contracting activity and a scale-independent
physical mass cannot be read off one schedule. Problem~\ref{prob:grid}
is the same shortfall again in the block count.

Underneath all three sits Break~\ref{sec:break5}: the estimates are
proved near $u = 0$ and the continuum lives at $u \to \infty$. A
summable increment would connect scales \emph{within} a regime; it would
still not move the construction between regimes. That is why
Problem~\ref{prob:trajectory} is marked load-bearing and the others are
not.
% ---- end 30_continuum.tex ------------------------------------------
% ---- begin 34_breaks.tex -------------------------------------------
\section{Six obstructions, with counterexamples}
\label{sec:breaks}

This section states the obstructions that any argument from the
fixed-spacing construction of \S\ref{sec:wilson} to a continuum mass gap
must clear. They were produced by a deliberate adversarial audit of this
program's own drafts, and four of them killed arguments that had been
written down here as complete.

They are stated as mathematics, not as caveats. Five carry explicit
finite counterexamples, and a finite counterexample to a proof strategy
is a permanent fact about that strategy. A reader who wants to know
where the remaining difficulty of the problem actually sits should read
this section and \S\ref{sec:obligations} and can safely skip everything
else in Part~\ref{part:continuum}.

The counterexamples are not rhetorical. Those of
Break~\ref{sec:break4} are registered checks in this repository's own
test suite --- \srcpath{tests/test\_os\_kernel\_gap.py}, six checks,
passing --- and the two collapse conditions of Break~\ref{sec:break3}
are checked in \srcpath{tests/test\_moving\_time\_gap.py}. An
obstruction that a repository asserts against itself, and then runs in
continuous integration, is harder to forget than one recorded in prose.

\subsection{Break 1: an \texorpdfstring{$O(g_k^2)$}{O(g\_k\textasciicircum 2)} increment bound does not give convergence}
\label{sec:break1}

A renormalization-group argument for the continuum limit produces a
sequence of effective actions $S_k^{\mathrm{eff}}$ and observable
expectations $m_k(\mathcal{O})$ indexed by scale, and controls the
one-step increment by the running coupling:
\begin{equation}
  \abs{m_{k+1}(\mathcal{O}) - m_k(\mathcal{O})}
  \;\le\; C_{\mathcal{O}}\, g_k^2
  \;\le\; \frac{C_{\mathcal{O}}'}{k + k_0},
  \label{eq:break1-bound}
\end{equation}
the second inequality being asymptotic freedom, $g_k^2 = \kappa/(k+k_0)$.
The step that must not be taken is to conclude from
\eqref{eq:break1-bound} that $(m_k)$ is Cauchy. Telescoping gives only
\begin{equation}
  \abs{m_n(\mathcal{O}) - m_m(\mathcal{O})}
  \;\le\; C_{\mathcal{O}}' \sum_{k=m}^{n-1} \frac{1}{k+k_0}
  \;\asymp\; C_{\mathcal{O}}' \log\!\left(\frac{n}{m}\right),
\end{equation}
which does not tend to zero as $m, n \to \infty$.

\begin{proposition}[Counterexample]
\label{prop:break1}
Let $x_k = \sin\!\big(\log(k+1)\big)$. Then $x_k \in [-1,1]$ and, by the
mean value theorem, $\abs{x_{k+1} - x_k} \le 1/(k+1) \to 0$, yet the set
of subsequential limits of $(x_k)$ is the whole interval $[-1,1]$.
\end{proposition}

So the increment bound \eqref{eq:break1-bound} is consistent with the
observable expectations oscillating forever across a fixed interval. What
is needed is one more power: $\abs{m_{k+1} - m_k} = O(g_k^4) = O(k^{-2})$.

The obstruction to obtaining it is not technical. On the asymptotically
free trajectory $u_k = (\alpha + \beta k)^{-1}$,
\begin{equation}
  \Delta u_k = -\frac{\beta}{(\alpha+\beta k)\big(\alpha + \beta(k+1)\big)},
  \qquad
  \frac{\abs{\Delta u_k}}{4 u_k u_{k+1}} = \frac{\beta}{4} = O(1),
\end{equation}
and differentiating an expectation along the trajectory,
\begin{equation}
  \frac{d}{ds}\,\mu_s(\mathcal{O}_s)
  = \mu_s(\partial_s \mathcal{O}_s)
  - \operatorname{Cov}_{\mu_s}\!\big(\mathcal{O}_s, \partial_s S_s\big),
  \label{eq:break1-cov}
\end{equation}
the normalization carries no small factor. The natural size of the
variation is therefore $O(g_k^2)$, and a bounded-covariance estimate
cannot improve it. Converting $1/k$ into $1/k^2$ requires an exact
cancellation of the leading partition-function covariance in
\eqref{eq:break1-cov}. No such cancellation is proved here.

\begin{scopenote}
This obstruction applies to the multiscale route recorded as\\
\texttt{\small DERIV:YANGMILLS\_CONTINUUM\_\allowbreak{}BALABAN\_\allowbreak{}MULTISCALE\_\allowbreak{}PROOF:\allowbreak{}THEOREM\_7\_1},\\
whose Cauchy input is exactly what \eqref{eq:break1-bound} fails to
supply. It does not bear on any statement in
Part~\ref{part:exact} or \ref{part:fixed}, which are at fixed scale.
\end{scopenote}

\subsection{Break 2: the polymer activity and the mass schedule are algebraically incompatible}
\label{sec:break2}

A multiscale polymer cluster expansion carries an activity parameter
$\kappa_k$ that must contract for the expansion to converge, say
\begin{equation}
  \kappa_{k+1} = \tfrac{3}{4}\,\kappa_k,
  \qquad \kappa_k = \kappa_0 \left(\tfrac{3}{4}\right)^{k} \longrightarrow 0,
\end{equation}
while a physical mass is extracted at each scale by
\begin{equation}
  \frac{\kappa_k - \log C_4}{a_k} = c_0 \Lambda > 0
  \quad\text{for all } k \ge 0,
  \qquad a_k = a_0 L^{-k},\ L = 3,
  \label{eq:break2-mass}
\end{equation}
with $C_4 > 1$ the polymer coordination constant.

\begin{proposition}
\label{prop:break2}
The two displays above are incompatible. Since $\kappa_k \to 0$ and
$\log C_4 > 0$, there is a finite
\[
  k_* = \left\lceil \frac{\log\!\big(\kappa_0/\log C_4\big)}{\log(4/3)} \right\rceil
\]
with $\kappa_k - \log C_4 < 0$ for every $k > k_*$; beyond that scale the
cluster sum $\sum_n e^{-(\kappa_k - \log C_4)n}$ diverges and the
expansion does not exist. Setting $C_4 = 1$, which discards polymer
entropy, does not repair it: then
$\kappa_k / a_k = (\kappa_0/a_0)(3L/4)^k = (\kappa_0/a_0)(9/4)^k \to \infty$,
so \eqref{eq:break2-mass} cannot equal a finite constant either.
\end{proposition}

The content of Proposition~\ref{prop:break2} is that a contracting
activity and a scale-independent physical mass cannot both be read off
the same schedule. Any repair must decouple them --- the mass must come
from a quantity that is not the activity divided by the spacing.

\subsection{Break 3: the moving-time gap can be exactly zero}
\label{sec:break3}

The moving-time theorem of \S\ref{sec:continuum} extracts a limiting
gap from approximate sources and one growing observation time per
cutoff. Under the error bound
$K_n(t) \le A a_n^{-p} e^{-mt} + B a_n^{r}$ and probe tilt
$\abs{\alpha_n} \ge c_\alpha a_n^{s}$, it gives
\begin{equation}
  M_* = \frac{m\,(r - 2s)}{p + r}
  \qquad\text{at}\qquad
  t_n = \frac{p+r}{m}\,\log\frac{1}{a_n}.
  \label{eq:break3}
\end{equation}
Two ways this returns nothing.

\emph{Zero margin.} If $r \le 2s$ then $M_* \le 0$: the theorem is
sharp and it yields exactly zero. Here $r$ is the convergence rate of the
coarse-graining background and $s$ the rate at which probes must shrink
to stay in the quadratic basin of the Wilson action. No calculation in
this program, or to this author's knowledge anywhere, certifies
$r > 2s$ for four-dimensional Yang--Mills. The inequality $r > 2s$ is
therefore a hypothesis with no current route, not a technical condition.

\emph{Finite horizon.} If the observation time is capped by the box,
$t_n \le c_{\max}\log(1/a_n)$, the optimum becomes
\begin{equation}
  M_{\mathrm{opt}} = \min\!\left(m - \frac{p+2s}{c_{\mathrm{opt}}},\
  \frac{r-2s}{c_{\mathrm{opt}}}\right),
  \qquad
  c_{\mathrm{opt}} = \min\!\left(c_{\max},\ \frac{p+r}{m}\right),
\end{equation}
and $c_{\max} \le (p+2s)/m$ forces $M_{\mathrm{opt}} \le 0$. Both
collapses are registered checks rather than remarks.

\subsection{Break 4: a local probe can be blind to the lightest state}
\label{sec:break4}

Suppose one measures a plaquette correlator and observes
\begin{equation}
  \ip{\mathcal{O}_{\mathrm{plaq}}}{e^{-tH}\mathcal{O}_{\mathrm{plaq}}}
  \le C e^{-M_{\mathrm{obs}} t},
  \qquad M_{\mathrm{obs}} = \log 4 \approx 1.386 .
\end{equation}
This does not bound the spectrum.

\begin{proposition}[Dark sector]
\label{prop:break4}
On $\C^3$ let
\[
  T = \operatorname{diag}\!\left(1,\ \tfrac{15}{16},\ \tfrac14\right),
  \qquad \Omega = e_1,\qquad v_{\mathrm{obs}} = e_3 .
\]
Then $\ip{v_{\mathrm{obs}}}{T^t v_{\mathrm{obs}}} = 4^{-t} = e^{-t\log 4}$,
while the true first excited energy is
$-\log(15/16) \approx 0.0645$, smaller by a factor above $21$. The probe
has zero overlap with the state realizing the gap.
\end{proposition}

Consequently a decay rate for local plaquettes is a statement about the
gap only after the plaquette family is proved \emph{total} in
$\Omega^{\perp}$. This is why totality is a conclusion, not a hypothesis,
in the fixed-spacing band theorem (\S\ref{sec:wilson}); and it is why a
continuum argument may not inherit totality from the fixed-spacing
statement without proving it again in the limit.

A second, subtler version of the same failure concerns positivity rather
than overlap.

\begin{proposition}[Osterwalder--Schrader null states]
\label{prop:break4b}
Let $(x,y) \in \{-1,1\}^2$ be uniform and let time reflection act by
$\Theta(x,y) = (y,x)$. The centered observable $F(x,y) = y$ has Euclidean
variance $\mathbb{E}[F^2] = 1 > 0$ but physical norm
$\norm{F}_{\mathrm{OS}}^2 = \ip{\Theta F}{F} = \mathbb{E}[xy] = 0$.
\end{proposition}

So positive Euclidean variance does not certify that a probe exists in
the physical Hilbert space at all. An algorithm that selects probes by
Euclidean variance can select states that vanish after reconstruction.

A third version concerns the generator rather than the states: with
$H_n = \operatorname{diag}(0,n)$ one has
$e^{-tH_n} \to \operatorname{diag}(1,0)$ for every $t>0$, so the limiting
semigroup is not strongly continuous at $t=0$ and no self-adjoint
generator exists on the second coordinate. Convergence of Gram matrices
alone therefore does not produce a limiting Hamiltonian.

\subsection{Break 5: the proved regime and the continuum regime are disjoint}
\label{sec:break5}

This is the structural obstruction, and the one that most changes what a
correct argument must look like.

The fixed-spacing band theorem (\S\ref{sec:wilson}) holds on
\begin{equation}
  \abs{u} \;\le\; \frac{u_*}{10\,022\,400\,000\,N},
  \label{eq:break5-u}
\end{equation}
and the spatial-closure estimates hold for
$0 \le \lambda \le \lambda_c$ with
$\lambda_c \in \left(\tfrac{1}{73}, \tfrac{1}{72}\right)$. Both are
neighbourhoods of zero in the strong-coupling variable.

The continuum limit lives elsewhere. With
$g^{-2}(a) = \beta_0 \log(1/a\Lambda)$ and $u = 2/\big(g^2(a)N\big)$,
\begin{equation}
  u(a) \longrightarrow \infty,
  \qquad
  \lambda(a) \sim g^{-4}(a) \longrightarrow \infty
  \qquad (a \to 0).
\end{equation}
So the interval on which the theorem is proved and the trajectory along
which the continuum is taken have empty intersection:
\begin{equation}
  \left[0, \lambda_c\right] \,\cap\, \left[\lambda(a), \infty\right) = \emptyset
  \quad\text{for all sufficiently small } a .
\end{equation}

\begin{scopenote}
The fixed-spacing results establish a strictly positive gap at fixed
lattice spacing in a neighbourhood of strong coupling. They supply no
control on the weak-coupling trajectory. An inference from
\eqref{eq:break5-u} to a continuum mass gap is a regime error, and this
program has made and retracted a version of it.
\end{scopenote}

What Break~5 demands of a future argument is a \emph{controlled
trajectory}: an estimate that survives continuation in the coupling, not
a better constant at $u \approx 0$. Improving the constant
$10\,022\,400\,000$ in \eqref{eq:break5-u} by any finite factor does not
move the boundary of the difficulty. This is why Problem~\ref{prob:trajectory}
in \S\ref{sec:obligations}, and not a sharper fixed-spacing estimate, is
the load-bearing open problem.

\subsection{Break 6: diffusion time is not Euclidean time}
\label{sec:break6}

A recurring route derives a spectral gap from curvature: positive
Bakry--\'Emery--Ricci curvature on the configuration manifold gives a
logarithmic Sobolev inequality, which gives a spectral gap
$\lambda_{\mathrm{diff}} > 0$ for the Langevin generator. The step that
fails is the identification of $\lambda_{\mathrm{diff}}$ with a physical
mass.

Langevin time is a fictitious parabolic parameter; Euclidean time is a
spacetime coordinate of the reconstructed theory. In free models the two
are related by $\Delta_{\mathrm{phys}} = \sqrt{\lambda_{\mathrm{diff}}}$,
but that relation is a consequence of Gaussian structure and not a
general isometry. For non-abelian gauge theory no such correspondence is
available here.

The companion obstruction concerns the coarse-graining itself. A scale
map built from a gauge-covariant Markov kernel that preserves both
reflection positivity and locality is not available: any such kernel
either forces the holonomies to commute --- reducing the theory to an
abelian one --- or destroys reflection positivity, so that the
reconstructed inner product is indefinite. This is why the scale route
actually used in \S\ref{sec:continuum} is a \emph{deterministic}
pushforward rather than a Markov kernel, and why the two-spin
counterexample $\mathbb{E}[s_- s_+] = -1$ is recorded as a proposition
rather than as an anecdote.

\begin{scopenote}
Break~6 does not refute the curvature computations themselves. The Haar
Ricci constant and the Wilson Hessian bound are correct statements about
the configuration manifold. What is refuted is their use as a route to
$\Delta_{\mathrm{phys}}$, and the curvature results are consequently
recorded in this program as inputs to coercivity estimates rather than
as a mass-gap mechanism.
\end{scopenote}

\subsection{What survives}

\begin{enumerate}
  \item The exact strong-coupling spectral theory of
    Part~\ref{part:exact}, in its entirety. None of the six obstructions
    touches a fixed-order statement.
  \item The fixed-spacing infinite-volume Wilson band and its literal
    source frame (\S\ref{sec:wilson}), on the interval
    \eqref{eq:break5-u}, with the scope stated there.
  \item The moving-time spectral exclusion as an abstract interface: the
    implication from its hypotheses to a limiting gap is proved and
    sharp. Break~3 is a statement about whether its hypotheses hold for
    Yang--Mills, not about the implication.
\end{enumerate}

What does not survive is any inference from the above to a continuum
Yang--Mills mass gap. The program states that plainly, and
\S\ref{sec:obligations} is the list of what would have to change.
% ---- end 34_breaks.tex ---------------------------------------------
% ---- begin 35_register.tex -----------------------------------------
\section{The refutation register}
\label{sec:register}

A refuted claim is evidence that it was tried. This section reproduces
the register of every claim this program made, or inherited from its own
sources, and then lost --- together with the exact thing that killed it.

The register carries \ContradictionTotal{} entries in three states.
\ContraFalsified{} are \emph{falsified}: a counterexample or an exact
computation contradicted them. \ContraSuperseded{} are
\emph{superseded}: replaced by a sharper statement. \ContraResolved{} are
\emph{resolved}: the apparent conflict dissolved, in most cases because
the two sides turned out to be differently anchored quantities rather
than competing values for the same quantity.

Three features of how this register is maintained matter to a reader
deciding how much to trust the rest of the document.

First, a dispute closes by derivation or it does not close. Both values
stay recorded side by side, and no code path is permitted to pick one,
average them, or prefer whichever is an exact rational because exactness
reads as authority. The fourth-order off-axis coefficient closed that way:
every cluster of the rotation record reproduces the historical pipeline's
own number, and the single cluster that does not is explained ordering by
ordering (\S\ref{sec:fourth}). Its two recorded values remain in the
registry with their verdicts attached.

Second, a resolved entry is not a win. The largest class here is
resolutions in which a contradiction turned out to be a naming
collision --- two differently anchored coordinates both called $m_4$, for
instance. Those cost real time and produced no mathematics, and the
register's response was to forbid the ambiguous name rather than to
record a victory.

Third, this program has retracted a mechanism of its own. The
\emph{tier collapse} --- the observation that two of the four
fourth-order shape coefficients vanish at every solved rank --- was
proposed here as a kinematic degree bound with a falsifiable
prediction, and the mechanism was withdrawn a fortnight later when one
projection turned out to be uncounted (\S\ref{sec:fourth}). The
withdrawal is a numbered decision record rather than a silent edit,
because the next attempt at that mechanism needs to know which step
failed. The reformulation survives; the mechanism and its prediction do
not.

That is a different object from the degree bound of
\S\ref{sec:symbolic}, which was never retracted --- it was made
unnecessary. The two are easy to confuse and this document keeps them
apart deliberately.

\medskip
% ---- begin gen_verification_register.tex ---------------------------
% Generated by paper/master/generate_apparatus.py.  Do not edit.
\subsection*{Falsified (4)}
\addcontentsline{toc}{subsection}{Falsified}
\noindent Falsified --- killed by a counterexample or an exact computation. These are the four claims this program asserted, or inherited, and then lost.\par\smallskip
\begin{itemize}[leftmargin=1.4em]
  \item \texttt{C6} \textbf{F-2 mobility rule.} The promotion r\_\allowbreak{}physical = w\_\allowbreak{}min - 2 is falsified by the pentagonal isotropic cap hop at r=4 with w\_\allowbreak{}min-2=5. The scoped bound r >= w\_\allowbreak{}min - 2 survives, with survival gates.
  \item \texttt{C7} \textbf{Stranded-flux zero backend.} Claimed Haar zero for all 20 pentagonal histories refuted by exact Weingarten/network contractions (8/8).
  \item \texttt{C18} \textbf{Center-only circuits are dynamically dark.} Refuted by the exact fifth-order determinant correction delta c\_\allowbreak{}5,det != 0
  \item \texttt{C19} \textbf{Global fixed-window firewall.} Disproved by the Bernoulli no-go; replaced by the rooted conditional calculus
\end{itemize}

\subsection*{Superseded (3)}
\addcontentsline{toc}{subsection}{Superseded}
\noindent Superseded --- replaced by a sharper statement. The original is not wrong so much as no longer the best available.\par\smallskip
\begin{itemize}[leftmargin=1.4em]
  \item \texttt{C5} \textbf{String-tension sign at O(u\textasciicircum{}3).} -61/408 (the earlier +61/408 was an odd-order convention error)
  \item \texttt{C11} \textbf{Compact-gap constants.} Corrections: c\_\allowbreak{}0 = -3(N\textasciicircum{}2-3)/(16N); radial moment difference 3N\textasciicircum{}2+1; dilation eps\textasciicircum{}4 = N/(2*beta);
  \item \texttt{C12} \textbf{Curvature misreads.} Two-stage. The kappa values are RADIAL DIRECTIONAL curvatures, not Hessian entries at Gamma; a genuine Hessian needs beta = 2*alpha, violated by the historical kernel.
\end{itemize}

\subsection*{Resolved (15)}
\addcontentsline{toc}{subsection}{Resolved}
\noindent Resolved --- the collision was dissolved. In most of these the two sides turned out to be differently anchored quantities rather than competing values.\par\smallskip
\begin{itemize}[leftmargin=1.4em]
  \item \texttt{C1} \textbf{Fourth-order rest scalar --- dissolved, an anchoring distinction.} NOT a contradiction. q\_\allowbreak{}band\textasciicircum{}(4) and m\_\allowbreak{}Gamma\textasciicircum{}(4) are not two competing estimates of the same scalar coordinate: the first is a band-kernel anchor, the second a vacuum-subtracted physical Gamma-point coefficient. They are differently anchored coordinates related by a translation-local scalar shift, and such a shift cannot change the centered operator, its eigenvectors, the SOS factorization, the mobility coefficients, or the bandwidth.
  \item \texttt{C2} \textbf{Fourth-order off-axis coefficient C\_\allowbreak{}shp --- resolved by derivation 2026-09-04.} C\_\allowbreak{}shp = C\_\allowbreak{}SHP\_\allowbreak{}HISTORICAL + 25/1024 = -13035490122347/550663802582400 = -0.0236723, the cluster-assembled coefficient read in the kernel's own basis (ADR 0024). Both recorded kernels are wrong on the rotation amplitude: the historical by the sixteen orderings of the adjacent-face cube completion its own Stage-3I ledger lacks (8 of 24, each with the right weight; the shortfall is exactly 25/512 on rho), the cold dump on u, pi and rho alike.
  \item \texttt{C3} \textbf{Status of the fourth-order "theorem".} Newest status controls. The \S{}5.2 values are sealed; everything else is C1/C2.
  \item \texttt{C4} \textbf{Coupling conventions.} The archived Y = 2*beta/3 = 4u line is a definition/label error, not a rescaling rule. Printed coefficients were already in u = beta\_\allowbreak{}lat/6. Never apply 4**n factors to them.
  \item \texttt{C8} \textbf{Pentagonal degenerate space.} E\_\allowbreak{}cap = 10/3 != E\_\allowbreak{}side = 8/3; choose the physical degenerate eigenspace first
  \item \texttt{C9} \textbf{SU(4) exceptional correction.} Only the flat-branch eigenvalue identity holds, not scalarity of the matrix
  \item \texttt{C10} \textbf{Determinant sectors shift only q\_\allowbreak{}N.} True for N>=4, false at SU(3): Delta\_\allowbreak{}beta\_\allowbreak{}3 = -25/64 and Delta\_\allowbreak{}q\_\allowbreak{}3 = -16863189551/76406976000.
  \item \texttt{C13} \textbf{C-even second-order hopping.} -481/612 superseded by -11/306 (the vacuum-mediated route was omitted)
  \item \texttt{C14} \textbf{Betti "+2".} Reinterpreted as b\_\allowbreak{}2 - b\_\allowbreak{}3; the open box has no extra states
  \item \texttt{C15} \textbf{Tetrahedral local coefficient.} Re-derived 2026-08-21 (tetrahedral Haar-resolvent suite): the asserted -8/(N(N\textasciicircum{}2-1)) is exact at symbolic N. The derivation reconstructs the primitive law's three stated ingredients --- 1/N per merge from the fundamental-pair Haar moment (independent of shared-path length, derived not assumed), the certified electric convention E(L) = L*C\_\allowbreak{}F/2, and exhaustive temporal orderings --- and reproduces every certified or printed instance first (sealed-core cube -5/48 with the (-8,-8,-4) temporal classes, prism 64 with its cap-sector 24 dissolved as a sector\,\ldots
  \item \texttt{C16} \textbf{SU(5)/SU(6) fourth-order completeness.} The payloads were shipped after all, and are now machine-checked (restored-payloads suite, 2026-08-21): the SU(6) determinant certificate re-derives Delta q\_\allowbreak{}6 = 6/343 exactly, momentum-independent with A, B, and the bandwidth unchanged, and the SU(5) stage-1 summary pins the 895,524-pair scan with zero determinant sectors. Two upstream word manifests remain absent (their SHAs are recorded in the certificates), so the evidence level is output-certified, not cold-reproduced.
  \item \texttt{C17} \textbf{Fixed-rank holdout overclaim.} Corrected: q\_\allowbreak{}N matches stored values for N=7..18, but full kernels only at N=7,8; the A/B sample ledger remains open.
  \item \texttt{C20} \textbf{RUN15 internal display artifact.} Cosmetic; the gate value is the intended exact one
  \item \texttt{C21} \textbf{Monte Carlo gate integrity.} next14.json's "23/23 gates" is downgraded: one gate is a literal truth value, the JSON is not source-hash bound (non-RFC NaN tokens), and no raw ensemble exists. Numbers stand as structured numerical evidence; the certificate claim does not.
  \item \texttt{C22} \textbf{Gate-85 equality.} Not a numerical contradiction but a status correction: the final diagonal shift was chosen to produce exactly that equality, so gate 85 certifies internal bookkeeping, not independent agreement.
\end{itemize}


% status histogram: {'proven': 140, 'disputed': 1, 'falsified': 1, 'conditional': 1}
% evidence histogram: {'analytic': 141, 'record-backed': 1, 'numerical': 1}
% ---- end gen_verification_register.tex -----------------------------

\subsection{Claims withdrawn by this program}

Distinct from the register above, the following are claims this
repository itself advanced and then withdrew. They are listed separately
because their failure mode is different: nobody contradicted them, the
argument was found not to close.

\begin{itemize}
  \item \textbf{The tier collapse as a kinematic degree bound.}
    Proposed with a falsifiable prediction --- that the next shape tier
    would first appear at sixth order --- and retracted eleven days
    later. The count bounded the vertices of $H_4$; the carrier
    projection $d^{\dagger}(\cdot)d$ contributes one further power of
    the Bloch variable that the argument did not carry. Recorded as
    \srcpath{docs/decisions/0005-retracting-the-degree-bound.md}. The
    reformulation stands, the mechanism does not, and the collapse is
    now established by the range-one argument of \S\ref{sec:fourth}
    instead.
  \item \textbf{An unrestricted exponential-source radius.} Inferred from
    independent bounded sources; the inference does not hold, and the
    reformulation is Problem~\ref{prob:g17}.
  \item \textbf{Closure of the source-energy transport budget.} Claimed
    complete, rejected on review: a vacuum-only skew generator does not
    control a source that the generator transports rather than
    annihilates. The corrected partial results are retained and the
    obligation is Problem~\ref{prob:r10}.
\end{itemize}

\noindent The falsified entries above --- the mobility rule, the
stranded-flux backend, the dark centre-only circuits and the global
fixed-window firewall --- are a different category and are not repeated
here: each was killed by a counterexample or an exact computation
rather than by an argument failing to close.

Each of these was, at some point, written down in this program's own
documents as established. That is the reason the machine-verification
tier of \S\ref{sec:howto} is computed rather than assigned, and the
reason the present document quotes its counts from the ledgers by macro
rather than by hand.
% ---- end 35_register.tex -------------------------------------------
% ---- begin 36_obligations.tex --------------------------------------
\section{The obligation register}
\label{sec:obligations}

This is the operative section of the document. It states, as numbered
problems, what would have to be proved; and it then reproduces the
repository's full route register, including every route that is dead and
the reason it died.

The numbered problems come first because they are curated: they are the
routes this program judges to be live and decisive. The register that
follows is not curated --- it is the whole record, \RouteTotal{} routes
across \GapOpen{} open and \GapDischarged{} discharged gaps, of which
\RouteDone{} are complete, \RouteLive{} live, \RouteUntried{} untried and
\RouteDead{} dead.

\subsection{The problems}

Ordering expresses research judgement. It is \emph{not} a claim that an
earlier problem is a prerequisite for a later one; the ledger records
explicit inputs separately, and each problem below states its own.

\begin{problem}[Conditional-score domination, uniformly in the coupling]
\label{prob:m10}
For the synchronized transport of the four-face $\SU(2)$ conditional
model, establish
\begin{equation}
  K_g(w) \;\le\; C_0 + C_1 W_g(w)
  \qquad \nu_g\text{-almost everywhere},
  \label{eq:prob-m10}
\end{equation}
with $C_0, C_1$ independent of $g$ as $g \to 0$, retaining the
chart-complement, rare-fiber and antipodal contributions.

\emph{Available.} The actual source moment; the interval-obstruction
repair; conditional-score continuity; synchronized conditional dilation;
the antipodal Hessian spectrum, gauge-normal coercivity, angular
reduction and gauge-score invariance --- all proved. Exponentially small
complement suppression outside a neighbourhood of the complete minimizing
set, and a uniform Agmon action gap on the complement, are proved.

\emph{Missing.} Polynomial relative-amplitude control; the tube moments;
the actual normal/angular comparison.

\emph{Consequence.} With the established source-potential moment,
\eqref{eq:prob-m10} makes the uniform weighted score bound applicable.
It does \emph{not} supply source-energy transport or interacting-grid
uniformity; those are Problems~\ref{prob:r10} and \ref{prob:grid}.

\emph{Next computation.} For the synchronized field, obtain the actual
conditional tube moments and the differentiated amplitude bound; then use
the established antipodal magnetic normal and angular geometry to prove
the true-ground score estimate, controlling the outside deviation against
the same reference. Compare the resulting complete variance with
$1 + g^{-2}\mathbb{E}(V \mid w)$, or prove a sufficient source-weighted
pairing directly.

\emph{Dead here.} Inferring uniform domination from endpoint requirements
on arbitrary cutoffs; differentiating a global fixed-rank physical
Coulomb source projection across flat holonomies; bounding the global
true-ground Fisher matrix by $C_0 + C_1\sqrt{u}\,v(U)$.
\end{problem}

\begin{problem}[Source--vacuum transport energy jets]
\label{prob:r10}
Bound the complete compact source and vacuum transport on a common energy
domain against a physical clock. A vacuum-only skew generator is
insufficient: the source generator transports the ground state, it does
not annihilate it.

\emph{Available.} The ground jets; the ground-frame generator; the
score--Poisson identity; the conditional derivative; $H_4$ derived from
$H_0$; the projection-jet recurrence --- all proved.

\emph{Next computation.} Decompose the complete generator into vacuum and
conditional-source parts; retain commutator and energy-domain terms
through two derivatives; test the proposed powers of $g$ on the full
energy graph rather than only on ground vectors.

\emph{Recorded failure.} An earlier closure of this problem was claimed
and rejected on exactly the ground that the vacuum-only generator does
not see the source. The rejection is preserved; see \S\ref{sec:register}.
\end{problem}

\begin{problem}[Interacting-grid comparison with uniform constants]
\label{prob:grid}
On the actual coupled Wilson grid, with the true vacuum, the full fast
space, source metrics, interfaces and a common physical clock, prove a
comparison whose constants are uniform in the number of interacting
blocks.

\emph{Available.} The literal coarse form; covariant block sources; the
form Schur scale comparison --- all proved.

\emph{Next computation.} Two genuinely coupled adjacent blocks. Compute
the full conditional and source form including the interface and the high
retained directions, and isolate a local estimate whose constants sum
without a block-count factor.

\emph{Why uniformity is the whole problem.} Fixed-block estimates are
available in quantity. A constant that degrades with block count
telescopes to nothing, which is the same arithmetic failure as
Break~\ref{sec:break1} in a different variable.
\end{problem}

\begin{problem}[The signed conditional defect in the actual vacuum]
\label{prob:sc17}
On compatible finite periodic $\SU(2)$ blocks with the actual positive
Wilson vacuum, bound the full signed conditional reference defect,
retaining all localization and domain terms.

\emph{Available.} The first and third residual estimates, proved.

\emph{Next computation.} On the smallest coupled block, derive
outside-pressure and projector derivatives alongside the magnetic and
metric terms; check the Gaussian cancellation exactly; then bound the
remaining full expression with its cutoff commutators.
\end{problem}

\begin{problem}[Summable normalized RG and source increments]
\label{prob:cauchy}
For the actual normalized RG map and the renormalized observables after
marginal matching, with all measures and source spaces identified, prove
a cross-scale increment bound with scale-uniform remainder constants that
is \emph{summable}.

\emph{Available.} Contractive drift on a common space, proved.

\emph{This is Break~\ref{sec:break1}.} The required improvement is from
$O(g_k^2)$ to $O(g_k^4)$, and \eqref{eq:break1-cov} shows the
normalization supplies no small factor, so the improvement must come from
an exact cancellation of the leading partition-function covariance.

\emph{Next computation.} Work out one genuine normalized blocking step
with its observable transport; isolate the marginal contribution; and
calculate the surviving common-norm increment before claiming a
power-improved bound.
\end{problem}

\begin{problem}[A controlled trajectory out of the proved interval]
\label{prob:trajectory}
Connect the fixed-spacing construction of \S\ref{sec:wilson}, proved on
$\abs{u} \le u_*/(10\,022\,400\,000\,N)$, to the weak-coupling
trajectory along which the continuum limit is taken.

\emph{This is Break~\ref{sec:break5}, and it is load-bearing.} The two
regimes are disjoint. Borel--Pad\'e consistency of the strong-coupling
series is not such a trajectory: agreement of a series with a resummation
is not continuation of an estimate.

\emph{What a solution looks like.} An estimate whose constants survive
continuation in the coupling, or a second construction valid at large $u$
together with a proof that the two constructions describe the same
theory on an overlapping window. The second route has a named
prerequisite already discharged in part --- the dimension-five operator
irrelevance lemma --- and the coupling window between the two
constructions is the object to close.

\emph{Not a solution.} Any improvement of the finite constant
$10\,022\,400\,000$.
\end{problem}

\begin{problem}[A valid source-family formulation]
\label{prob:g17}
Identify the exact source operator, normalization and norm consumed by
the selected spectral comparison, and derive a local, connected or
normalized replacement \emph{together with} a proof that the replacement
preserves that comparison and a nonzero physical source weight.

\emph{Refuted approach.} Inferring an unrestricted exponential-source
radius from independent bounded sources. Linear susceptibility bounds do
not bound the centered exponential source.

\emph{Available.} The raw $\alpha$-dependent partition bound for bounded
sources, established. Scalar Kotecky--Preiss inequalities and raw moment
estimates remain useful ingredients but construct neither the tilted
activities nor a gap.

\emph{Next computation.} Compute the candidate source norm and its
spectral normalization in the explicit independent product law; compare
with $\big(M(2\alpha)^n - M(\alpha)^{2n}\big)/n$; then state the precise
interacting inequality that remains.
\end{problem}

\begin{problem}[One complete sixth-order cluster]
\label{prob:g9}
At fixed $\SU(3)$, on one smallest admissible connected plaquette
support, evaluate the eighteen electric folds with exact Haar and Gram
quotients, physical projections, electric denominators and rooted
subtraction.

\emph{Available.} The Hodge--Feshbach splitting; the fourth-order support
theorem; the Hermitian folded support; local sixth-order
non-cancellation; the one-face sixth order --- all proved
(\S\ref{sec:sixth}).

\emph{Consequence.} The extra rational band-shape component is already
forced under the local-shell and third-order hypotheses. Computing the
direct term determines its independent polynomial shapes and
coefficients, and decides whether the local structure survives. It
changes neither the fourth-order support theorem nor the global
sixth-order census.

\emph{Next computation.} Evaluate the eighteen folds on a connected
multi-plane support with all Haar weights and rooted subtractions. Reduce
the direct term separately; then combine it with the known mixing
numerator at common denominator, retaining both the direct polynomial
shape coefficients and the invariant remainder.

\emph{Unlike the others, this one is finite and priced.} It is the
cheapest decisive computation on this list.
\end{problem}

\subsection{The full route register}

What follows is generated from \srcpath{ledger/gaps.yaml}. Dead routes are
included deliberately and at length. A route marked \emph{cannot decide}
is structurally unable to resolve the named claim, which is a stronger
statement than having failed to.

\medskip
% ---- begin gen_obligations.tex -------------------------------------
% Generated by paper/master/generate_apparatus.py.  Do not edit.
% One subsection per gap that is not discharged, then the discharged
% ones in brief.  Dead routes are reproduced in full: they are the
% most expensive information in the ledger.

\subsection*{G1 --- Ship or restore machine-readable payloads that are only record-backed}
\addcontentsline{toc}{subsection}{G1 --- Ship or restore machine-readable payloads that are only record-backed}
\noindent\emph{Tier 0 (days); partial.}\par
\smallskip\noindent SU(3) 189-record kernel with reference SHA; Q\_\allowbreak{}32 and P\_\allowbreak{}402 ledgers; SU(5)/SU(6) exceptional certificates;\par
\smallskip\noindent\textbf{Ledger status.} largely discharged 2026-08-21 --- three of five families were shipped all along and are now machine-checked (restored-payloads suite): the 189-record kernel with both reference SHAs, the Q\_\allowbreak{}32/P\_\allowbreak{}402 Newton ledgers, and the SU(5)/SU(6) certificates. Still absent: the pentagonal dual-cold bundle (manifest row A60 records the SHA any restoration must hash to) and the tetrahedral certificate, for which no reference SHA exists anywhere --- that one was G5's to close, not this gap's, and G5 closed it by re-derivation on 2026-08-21 (the artifact itself remains absent). A\,\ldots\par

\subsection*{G25 --- The charge-even sector is registered and unchecked}
\addcontentsline{toc}{subsection}{G25 --- The charge-even sector is registered and unchecked}
\noindent\emph{Tier 1 (weeks); partial.}\par
\smallskip\noindent What remains. The suite reaches orders 2 and 3, which is where the corpus's C-even ledger stops. Fourth order is untouched on this side: the pre-registered O(u\textasciicircum{}4) criterion in the corpus is a statement about C-odd flatness, and no C-even fourth-order coefficient is registered here at all, so the sector's fourth order is not T3 -- it is absent.\par
\smallskip\noindent\textbf{Ledger status.} opened and largely discharged 2026-08-28. The C-even sector was registered and inert: LEAK\_\allowbreak{}2, LEAK\_\allowbreak{}2\_\allowbreak{}EVEN and LEAK\_\allowbreak{}3\_\allowbreak{}EVEN were read by no check body, band\_\allowbreak{}assembly() was dead code, and the only C-even geometry the repository checked was two hand-written matrices at Gamma and R. The `charge-even band, exactly` suite (15 checks, all T1) closes the second- and third-order half. Both entries of constants.DECLARED\_\allowbreak{}COINCIDENCES are now checked rather than only declared: -11/306 is ell\_\allowbreak{}N\,\ldots\par
\smallskip\noindent\textbf{Depends on.} \texttt{G3}\par
\smallskip\noindent\textbf{Live routes.}
\begin{itemize}[leftmargin=1.4em]
  \item \textbf{fourth-order domino, the U4 falsifier.} PROGRESS 2026-09-12, part (2) DONE: conn\_\allowbreak{}vac is derived and no longer T3. The linked two-face vacuum weight is now computed exactly by two routes that agree (swap-odd suite, new T1 check). At SU(3) the third engine's loop calculus gives the two-face vacuum sector 0, 0, -3/2, 9/16, -54321/837760 against twice the one face, so the linked cumulant is 0, 0, 0, 0, -327/83776: the corpus's omega4, identical for the coplanar and the perpendicular pair, with the linked O(u\textasciicircum{}2) and O(u\textasciicircum{}3) weights exactly zero as the v10a.7 gates reported in floats.
\end{itemize}

\subsection*{G4 --- Pentagonal O(u\textasciicircum{}4) closure in-corpus}
\addcontentsline{toc}{subsection}{G4 --- Pentagonal O(u\textasciicircum{}4) closure in-corpus}
\noindent\emph{Tier 1 (weeks); partial.}\par
\smallskip\noindent Fierz closure of the 20-history cut spaces, Gram quotient, physical resolvents, summed compression.\par
\smallskip\noindent\textbf{Ledger status.} partial --- the raw history Gram/Haar frontier is closed (8/8 gates, cold rerun in settlement/); the Fierz closure, physical quotient, and summed resolvent compression remain (see progress.remaining).\par

\subsection*{G6 --- B\_\allowbreak{}N fixed-rank holdout ledger for N = 7..18}
\addcontentsline{toc}{subsection}{G6 --- B\_\allowbreak{}N fixed-rank holdout ledger for N = 7..18}
\noindent\emph{Tier 1 (weeks); open.}\par

\subsection*{G7 --- Native torelon reruns of sigma\_\allowbreak{}5 and sigma\_\allowbreak{}6}
\addcontentsline{toc}{subsection}{G7 --- Native torelon reruns of sigma\_\allowbreak{}5 and sigma\_\allowbreak{}6}
\noindent\emph{Tier 1 (weeks); open.}\par
\smallskip\noindent Replace the KPS historical targets.\par

\subsection*{G8 --- Cold one-shot regeneration of the all-rank fourth-order bundle}
\addcontentsline{toc}{subsection}{G8 --- Cold one-shot regeneration of the all-rank fourth-order bundle}
\noindent\emph{Tier 1 (weeks); open.}\par
\smallskip\noindent Promotes output-certified to cold-certified.\par

\subsection*{G10 --- Full fifth-order pentagonal coefficient}
\addcontentsline{toc}{subsection}{G10 --- Full fifth-order pentagonal coefficient}
\noindent\emph{Tier 2 (months); open.}\par
\smallskip\noindent Support-resolved folds for the 572 return histories before rooted subtraction.\par

\subsection*{G11 --- Interval rigor and near-Gamma touching gates}
\addcontentsline{toc}{subsection}{G11 --- Interval rigor and near-Gamma touching gates}
\noindent\emph{Tier 2 (months); partial.}\par
\smallskip\noindent Outward-rounded pi. Without this, a fixed-momentum coefficient theorem must NOT be promoted to a uniformly isolated near-Gamma band theorem: for |k| <\textasciitilde{} u the O(u\textasciicircum{}2) splitting \textasciitilde{} u\textasciicircum{}2|k|\textasciicircum{}2 meets the O(u\textasciicircum{}4) scale.\par
\smallskip\noindent\textbf{Ledger status.} partially discharged\par

\subsection*{G12 --- Hyperhoneycomb O(u) candidate}
\addcontentsline{toc}{subsection}{G12 --- Hyperhoneycomb O(u) candidate}
\noindent\emph{Tier 2 (months); open.}\par
\smallskip\noindent Tests the framework beyond cubic cells.\par

\subsection*{G9 --- m\_\allowbreak{}6: global sixth-order geometry census}
\addcontentsline{toc}{subsection}{G9 --- m\_\allowbreak{}6: global sixth-order geometry census}
\noindent\emph{Tier 2 (months); open.}\par
\smallskip\noindent Contraction with external-memory sharding; folded weights and carrier census pre-cleared.\par
\smallskip\noindent\textbf{Depends on.} \texttt{G3}\par
\smallskip\noindent\textbf{Live routes.}
\begin{itemize}[leftmargin=1.4em]
  \item \textbf{Assemble one complete sixth-order cluster and test its extra carrier-shape component.} The September 11 continuation derives all 18 universal Hermitian electric folds, evaluates the complete one-face SU(3) rooted coefficient -2055143/35123200, and reduces all 25 H4 mixing pairs. Under the recorded third-order factorization and local-shell hypotheses, an additional rational shape cannot cancel against any local direct H6: q\textasciicircum{}3 lambda6 modulo q\textasciicircum{}2 is 4 kappa e2\textasciicircum{}2-3 kappa q e3 for symmetric polynomial direct F. The remaining task is the direct connected multi-plane H6 evaluation, including its Haar/Gram quotient, electric\,\ldots
\end{itemize}

\subsection*{G13 --- Classification theorem for shortest physical temporal histories}
\addcontentsline{toc}{subsection}{G13 --- Classification theorem for shortest physical temporal histories}
\noindent\emph{Tier 3 (unbounded); open.}\par
\smallskip\noindent Determine r\_\allowbreak{}escape from the energy-decorated incidence/Feshbach graph without exponential enumeration. Would subsume tetra/prism/cube/pentagonal/ hyperhoneycomb as one theorem, and includes proving partition-independence of the linked/folded assembly.\par

\subsection*{G14 --- Fourth-order support mechanism and its higher-order extension}
\addcontentsline{toc}{subsection}{G14 --- Fourth-order support mechanism and its higher-order extension}
\noindent\emph{Tier 3 (unbounded); partial.}\par
\smallskip\noindent Derive the actual supported operator words from complete physical histories, with the Haar quotient, electric resolvents and folded assembly retained. Use the established fourth-order support theorem as the baseline for a bounded sixth-order calculation. The Hodge splitting alone does not assert an all-length word classification or a dynamical sixth-order coefficient.\par
\smallskip\noindent\textbf{Ledger status.} Fourth-order tier collapse is established by RESULT:\allowbreak{}TIER\_\allowbreak{}COLLAPSE\_\allowbreak{}ACTUAL\_\allowbreak{}H4\_\allowbreak{}SUPPORT: the recorded kernel has support \{I,U,S,S\textasciicircum{}2,R\}, with cleared carrier span \{q,q\textasciicircum{}2,e\_\allowbreak{}2\}. The broader R-degree inference is falsified by sigma(UR)=-2q e\_\allowbreak{}2. The open successors concern history-to-word support, higher multiplicities and the separate proper-return identification in other geometries; the September channel computations are established inputs.\par
\smallskip\noindent\textbf{Live routes.}
\begin{itemize}[leftmargin=1.4em]
  \item \textbf{prove channel-wise universality from the Haar contraction -- that a channel's value depends on which links are doubled and in which irreps, not on where on the cluster they sit -- for u, the fan and the single-contact dressing.} Opened 2026-09-04 (ADR 0029). With every channel an identity in Q(N), the mechanism of the tier collapse is a statement about the channels alone: B is unpopulated by the carrier projection, D = 0 is u2 = 2u, and u2 = 2u is the L chain equalling the straight chain channel by channel. What is left is a proof of that equality from the contraction rather than a computation of it.
\end{itemize}
\smallskip\noindent\textbf{Untried routes.}
\begin{itemize}[leftmargin=1.4em]
  \item \textbf{derive the actual fourth-order Hodge support from the histories, including the absence of UR, RU and repeated R insertions.} Reviewed 2026-09-10. The narrower actual support \{I,U,S,S\textasciicircum{}2,R\} already proves the recorded kernel's B=D=0. Explain this support from the complete electric/Haar and folded histories, including the absence of UR, RU and repeated R insertions;
  \item \textbf{test the Hodge-summand mechanism in a second geometry by computing the tetrahedral Q-projected proper-return contributions, the first clause of U7's falsifier.} 
  \item \textbf{audit the premise that D - (1/2)\{H2, V2\} is the Hermitian fourth-order effective Hamiltonian on a cluster with PVP = 0 (the pair element at N >= 4), against a direct Bloch or Schrieffer-Wolff expansion on the two-face cluster at N = 5.} Taken as given by every lens of the 2026-09-04 adversarial verification of ADR 0027; unrefuted, unaudited.
  \item \textbf{audit the rank-generic resolvent at N >= 5 -- link\_\allowbreak{}spectrum roots, the dropping of the E0 component, vacuum intermediates -- against the SU(3) Casimir table route continued to N = 4, 5 by hand.} Taken as given by every lens of the 2026-09-04 adversarial verification of ADR 0027; unrefuted, unaudited. Partly covered later the same day (ADR 0029): over Q(N) every eigencomponent of every resolvent is verified to be an exact eigenvector of every single-link H0 with the SU(N) Casimir energies (21,556 components), so the spectra are no longer read off characteristic-polynomial roots at all;
  \item \textbf{the completeness of the eleven-cumulant decomposition of beta\_\allowbreak{}N at general N rests on the charge-counting argument alone (a connected four-insertion history from P to Q is P Qbar X Xbar, the pair, or the cube); enumerate the fourth-order support at N = 5 from the engine and compare with the 189 Hodge-form keys.} 
  \item \textbf{the N = 4 row of the beta record rests on the determinant trick (pseudoinverse Weingarten at n > N plus the epsilon insertion) in the three-face cumulants; audit that path at N = 4 against an independent epsilon-network reduction, and make the (4,0)-family pair element feasible the same way.} Sharpened 2026-09-04 (ADR 0029): the closed forms derived over Q(N) are regular at N = 4 and agree with the N = 4 row of every per-rank record, so whatever the determinant trick contributes there is what the generic rational function gives. The symbolic run also names the N = 4 degeneracy: the state P\textasciicircum{}3 with every link of the plaquette in Lambda\textasciicircum{}3 F, which is Fbar at N = 4, has energy E0 exactly there (the one resolvent denominator with an integer root >= 4). The audit of the path itself is still open.
\end{itemize}
\smallskip\noindent\textbf{Dead routes, and why they died.}
\begin{itemize}[leftmargin=1.4em]
  \item \textbf{derive why the fourth-order dynamics is R-linear -- why no four-insertion history inserts the cross-plane operator twice -- which is what the tier collapse now reduces to.} The proposed reduction is false: UR contains one R but its symbol -2q e\_\allowbreak{}2 populates the B tier. The actual fourth-order collapse instead follows from support \{I,U,S,S\textasciicircum{}2,R\}. The step name is retained to preserve its existing route id;
\end{itemize}
\smallskip\noindent\textbf{Completed routes (2).} read the mechanism off u(N) -- decompose the two-hop cumulant by intermediate irrep and match its denominator factors to resolvents; the same channel decomposition for the corner and fan dressings of the rotation orbit.\par

\subsection*{G15 --- Symbolic-N Gram transcript}
\addcontentsline{toc}{subsection}{G15 --- Symbolic-N Gram transcript}
\noindent\emph{Tier 3 (unbounded); open.}\par
\smallskip\noindent Promotes c\_\allowbreak{}1, c\_\allowbreak{}2 to unrestricted fixed-N theorems; one more order of the equivariant harmonic-well lemma.\par

\subsection*{G16 --- Rank-uniform control in tau = beta/N\textasciicircum{}3}
\addcontentsline{toc}{subsection}{G16 --- Rank-uniform control in tau = beta/N\textasciicircum{}3}
\noindent\emph{Tier 3 (unbounded); open.}\par
\smallskip\noindent Neither fixed-regime series supplies an overlap theorem. PROGRESS 2026-09-11 (ADR 0046, runs/planar\_\allowbreak{}band\_\allowbreak{}2026-09-11, suite "the planar limit of the fourth-order band (G16)", five T1 checks): the strong-coupling side is now exact through fourth order. Every fourth-order cumulant of the beta\_\allowbreak{}N assembly is O(N\textasciicircum{}-7) in both sectors while its resolvent channels are O(N\textasciicircum{}-3), the N\textasciicircum{}-3 and N\textasciicircum{}-5 channel totals vanishing identically in all sixteen cluster/sector pairs;\par

\subsection*{G17 --- Source-family stability and a valid source-radius formulation}
\addcontentsline{toc}{subsection}{G17 --- Source-family stability and a valid source-radius formulation}
\noindent\emph{Tier 3 (unbounded); open; \textbf{load-bearing}.}\par
\smallskip\noindent Identify the exact source operator, normalization and norm consumed by the selected spectral comparison. Derive a local, connected or normalized replacement only with a proof that it preserves that comparison and nonzero physical source weight. Linear susceptibility bounds do not bound the centered exponential source.\par
\smallskip\noindent\textbf{Ledger status.} The raw alpha-dependent partition bound is established by DERIV:\allowbreak{}TRACK\_\allowbreak{}A\_\allowbreak{}SUPPORTED:\allowbreak{}G17\_\allowbreak{}P with K\_\allowbreak{}alpha=exp(|alpha|/2). The unrestricted centered exponential-source radius is not a consequence of clustering: DERIV:\allowbreak{}TRACK\_\allowbreak{}A\_\allowbreak{}SUPPORTED:\allowbreak{}G17\_\allowbreak{}I gives divergence even for independent bounded centered sources. A valid source-family formulation and its actual interacting estimate remain open. The fixed-spacing infinite-volume Wilson band and literal-source\,\ldots\par
\smallskip\noindent\textbf{Untried routes.}
\begin{itemize}[leftmargin=1.4em]
  \item \textbf{Reformulate the source family while preserving its downstream spectral comparison.} Choose a precisely stated local, connected or normalized source formulation and prove the required relationship to the actual spectral comparison, including nonzero physical weight. Test it in the independent product model before seeking an interacting estimate. This is a candidate route with a source-identification obligation, not a claim that the original unrestricted exponential radius can be recovered from covariance summability.
\end{itemize}
\smallskip\noindent\textbf{Dead routes, and why they died.}
\begin{itemize}[leftmargin=1.4em]
  \item \textbf{Infer an unrestricted exponential-source radius from independent bounded sources.} DERIV:\allowbreak{}TRACK\_\allowbreak{}A\_\allowbreak{}SUPPORTED:\allowbreak{}G17\_\allowbreak{}I constructs an independent bounded centered source whose exponential variance divided by footprint cardinality tends to infinity. This refutes the stated inference, not every localized or normalized source formulation.
\end{itemize}
\smallskip\noindent\textbf{Completed routes (1).} Establish the raw alpha-dependent partition bound for bounded sources.\par

\subsection*{G19 --- Continuum limit}
\addcontentsline{toc}{subsection}{G19 --- Continuum limit}
\noindent\emph{Tier 3 (unbounded); open; \textbf{load-bearing}.}\par
\smallskip\noindent The actual infinite-volume Wilson construction and complete physical source band are now established at fixed spatial spacing on a common small-u interval. A controlled scale trajectory connecting it to large u, renormalized continuum correlations and a positive physical mass gap remain open. Borel-Pade consistency is not that trajectory.\par
\smallskip\noindent\textbf{Ledger status.} Reconciled 2026-09-09 after reviewing 43 distinct G19 document versions. Nonlinear finite-cell shells and scores, complete Gaussian fast/source bounds, fixed-spacing matching, and conditional continuum/OS transport are retained. The newly cited Balaban continuum proof stops at (7.6): its O(1/k) expectation increments do not imply the Cauchy conclusion in Theorem 7.1, which invokes a nonexistent O(1/k\textasciicircum{}2) estimate. Its (6.10) weight kappa\_\allowbreak{}k=(3/4)\textasciicircum{}k kappa\_\allowbreak{}0 also fails the positive entropy margin required in (8.14).\par
\smallskip\noindent\textbf{Live routes.}
\begin{itemize}[leftmargin=1.4em]
  \item \textbf{Prove the actual conditional-score domination M10.} The source-potential moment M8-M9 is established. The open target is the actual small-g conditional score K\_\allowbreak{}g<=C0+C1 g\textasciicircum{}-2 E(V|w), including rare fibers, with uniform nonnegative C0,C1. This sufficient route implies the weighted score/Hardy estimate M11-M15;
  \item \textbf{Control complete source-vacuum transport energy jets R10.} Raw H\_\allowbreak{}g derivatives and vacuum-only ground jets are established. Prove the complete source/vacuum generator bounds R10, ||A\textasciicircum{}(r)(g)||\_\allowbreak{}(q\_\allowbreak{}g to q\_\allowbreak{}g)<=d\_\allowbreak{}r g\textasciicircum{}(-r-1) for r=0,1,2, or use the source's direct mixed-form R11 alternative. These are alternative sufficient entrances to the complete residual budget, not cumulative prerequisites.
  \item \textbf{Control the interacting quantum comparison in redundant compact coarse variables.} Use the actual smooth matrix observation with its exact coarse gauge action and boundary source lift. Prove interacting true-vacuum marginal and full complementary energy estimates, including nonflat localization, noncommuting compact slow dynamics, the complete direct/exchange/metric/source correction and high retained tails. The ambient redundant projector is not a physical quantum Schur projection.
  \item \textbf{Bound the nonlinear excess above the exact quadratic memory on complete physical source windows.} Combine the established dynamic cubic energy and exact full fast Green identity with quartic magnetic, electric metric, Haar, moving-source, baseline cross and fast-form variation terms. Prove bounds for the actual spatial coefficients and the complete interacting remainder with retained harmonic localization, actual compact dynamics or explicit cancellations. Full source algebra, high retained directions and a summable physical-clock trajectory remain open.
  \item \textbf{Transport the endpoint and localized true-ground estimates through actual interacting scale blocks.} The additive complete endpoint and selected static-score comparisons are established inputs with different scopes. Prove uniform actual interacting fast/source or memory estimates and compare the resulting endpoint law with a consistent coarse quantum theory in one physical clock. For a static Schur route, retain its high-source coupling rather than inferring a full gap from the chosen chart.
  \item \textbf{Connect the actual Wilson fixed-scale construction to a controlled continuum trajectory.} Actual fixed-spacing infrared transfer and sources remain the controlled endpoint. The ultraviolet side now has a general fixed-cell physical first-cluster theorem and all-size planar and three-dimensional harmonic interface bounds, with planar low-mode/source frames and retained periodic harmonic directions. Gaussian Schur elimination retains memory and induced mass, giving a quantitative sorted-frequency approximation without small interface coupling.
  \item \textbf{Control the true interacting Wilson score and forms on the relevant energy space.} The exact SU(2) central true-ground Fisher matrix grows at least linearly in u, so the proposed global C0+C1 sqrt(u)v(U) estimate is no longer a live target. Prove an actual full-Q energy weight and a resolvent-weighted or coarse-gradient-energy score estimate on the relevant source window. A good-region small ratio plus controlled bad-region gradient-energy leakage yields a conditional small window Schur loss.
  \item \textbf{prove actual interacting W6 and coarse/source matching through spatial refinement.} The exact selected inverse-pairing and compatible source transport theorems remain available. R4a supplies a flat conditioned Gaussian reference budget; R12a identifies its zero complete pressure defect. \emph{Cannot decide:} \texttt{G19}.
  \item \textbf{supply the dimension-five irrelevance lemma, then close the coupling window between the two constructions.} The two rigorous constructions in this repository approach the continuum from opposite ends of the coupling axis, and the target is now the interval between them. From the strong-coupling side SC17 and BA20-BA25 give a volume-uniform physical gap for the actual vacuum, and the thermodynamic continuation gives an infinite-volume vacuum with a closed generator and spectrum in \{0\} union [4epsilon/3, infinity), for k/epsilon at most lambda\_\allowbreak{}c, the root of 6561 l\textasciicircum{}4 + 1458 l\textasciicircum{}3 - 81 l\textasciicircum{}2 - 72 l + 1 in (1/73, 1/72). From the\,\ldots
  \item \textbf{Extend the explicit square current to the complete finite-coupling residual.} Prove the complete transported finite-g current remainder against the actual interacting energy, with reference source synthesis and the required uniform parameter bounds. The no-diagonal criterion then supplies the selected-inverse estimate. Compatible source domains and metric transport remain explicit, and coupled-volume extension is a further step.
  \item \textbf{Transfer a physical continuum gap from one growing-time estimate per cutoff and approximate sources.} The original spectral-measure theorem and sharp rate budget are proved. The compatible-kernel successor now constructs the limiting physical semigroup and total centered defining space from finite Gram limits and fixed-time compatibility; positive-time divisibility supplies zero-time continuity, so a separate vague-limit or frame estimate is then unnecessary.
\end{itemize}
\smallskip\noindent\textbf{Untried routes.}
\begin{itemize}[leftmargin=1.4em]
  \item \textbf{Establish the actual interacting-grid fast and source comparison with uniform constants.} Construct the true interacting-vacuum fast floor and complete local centered/source comparison uniformly in block count, retaining interfaces, Gauss constraints, nonreducing source directions and the normalized coarse operator. A fixed-square bound, additive tensorization or harmonic floor alone does not prove this target. The chosen physical scale trajectory and summable errors must then satisfy the closed-form iteration theorem's actual hypotheses.
  \item \textbf{Bound the complete signed SC17 conditional defect in the actual vacuum.} Instantiate the actual uniform smallness input for the Hessian defect D\_\allowbreak{}B=Hess(V\_\allowbreak{}B+(H\_\allowbreak{}out Omega)/Omega)-2epsilon A\textasciicircum{}2, with pressure, metric, projector, moving-source and cutoff/product-domain terms assigned once. Preserve signed cancellations before bounding the complete expression. The conditional Riccati and block-geometry consequences are available, and the thermodynamic/physical-time constructions already stand at fixed spacing;
  \item \textbf{Prove normalized RG and source increments with a summable common-space remainder.} RESULT:\allowbreak{}COMMON\_\allowbreak{}SPACE\_\allowbreak{}CONTRACTIVE\_\allowbreak{}DRIFT supplies the proved abstract convergence criterion; its actual Wilson map and source hypotheses remain open. Construct the actual RG transformation and matched normalized observable maps on common spaces, subtract the running marginal terms, and bound the remaining expectation/source increments by a summable sequence.
\end{itemize}
\smallskip\noindent\textbf{Dead routes, and why they died.}
\begin{itemize}[leftmargin=1.4em]
  \item \textbf{Infer actual uniform M10 from the endpoint requirements on arbitrary Q8 cutoffs.} A permitted smooth radial cutoff freezes a fixed small regular Q while moving its actual conditional center. Full conditional concentration and a finite-flow speed contradiction rule out M10 for that choice. This is an existence counterexample within Q8, not a proof against every admissible transport or the synchronized successor.
  \item \textbf{Differentiate a global fixed-rank physical Coulomb source projection across flat holonomies.} An explicit SU2 stabilizer change gives a norm-one projection jump. The directions are gauge off identity and physical harmonic at identity, so this does not refute the uniform fast bound. Keep the smooth redundant actual observation and impose Gauss exactly instead of differentiating the changing Coulomb quotient. \emph{Cannot decide:} \texttt{G19}, \texttt{G23}.
  \item \textbf{Bound the global true-ground Fisher matrix by C0+C1 sqrt(u)v(U).} The exact actual SU(2) bouquet gives weighted Fisher at least (64/3)u-O(sqrt(u)) at U=-I, where v(U)=4. Smoothness makes the violation persist on open neighborhoods, so uniform almost-everywhere variants fail. The generic intrinsic-score and full-gap theorem remains valid under its explicit hypotheses, and other weighted or energy-restricted score estimates remain available. \emph{Cannot decide:} \texttt{G19}, \texttt{G23}.
  \item \textbf{Infer a scale contraction from the forward averaging derivative and a uniform raw fiber diffusion gap.} The preserved EX-014 Section 1 STEP 6 reverses the pullback derivative: a linear arithmetic-average test requires reverse gradient constant m, not 1/m. Correct quotient metrics remove the artificial gain. The actual Wilson horizontal score also has an order-g geometric term in a physical two-holonomy example; \emph{Cannot decide:} \texttt{G19}, \texttt{G23}.
  \item \textbf{Griffiths-type monotonicity of the electric-flux sector free energies in beta.} proposed 2026-09-03 (Alex): a correlation inequality that holds at every beta and makes the electric-flux free energies F\_\allowbreak{}e(beta) monotone in beta, sector by sector -- the sectors are Z\_\allowbreak{}N-valued, so the non-abelian obstruction to Griffiths' inequalities for Wilson loops need not apply -- would carry the gap from strong to weak coupling the way Griffiths carries Ising order across phases. Closed the same day by a sign, before the inequality itself was attempted: the Ising argument needs the monotone quantity to be INCREASING in beta, and the certified series says\,\ldots \emph{Cannot decide:} \texttt{G19}.
  \item \textbf{compare the whole compact Wilson fast form from above by a Gaussian form after finite-rank retention.} W1-W3 use actual compact SU2 spectral growth and finite-codimension min-max to refute this implementation at every fixed positive coupling. Compact eigenvalues grow quadratically in the irrep label while the Gaussian radial reference grows linearly. This does not refute a compatible compact reference or selected-inverse comparison. \emph{Cannot decide:} \texttt{G19}.
  \item \textbf{infer a regulator-uniform scalar single-site angle margin from the full square-root Gaussian precision.} For M\_\allowbreak{}m=sqrt(L+m\textasciicircum{}2) on a periodic scalar lattice, (L+m\textasciicircum{}2)1=m\textasciicircum{}2 1 gives M\_\allowbreak{}m 1 = m 1 exactly, so the row sums are m; every off-diagonal entry is negative by the heat representation, hence sum\_\allowbreak{}(f!=e)|M\_\allowbreak{}m(e,f)| equals M\_\allowbreak{}m(e,e)-m and the conditional maximal-correlation row is exactly kappa\_\allowbreak{}m = 1 - m/M\_\allowbreak{}m(e,e). Since M\_\allowbreak{}m(e,e) tends to a finite positive constant, kappa\_\allowbreak{}m tends to one as m vanishes: numerically .6154 at m=1,\,\ldots \emph{Cannot decide:} \texttt{G19}.
\end{itemize}
\smallskip\noindent\textbf{Completed routes (28).} Derive the actual conditional continuity score and locate the constrained Agmon centers; Synchronize compact radial dilation with the actual conditional minimum curve; Resolve the antipodal conditional magnetic normal geometry; Control time-integrated conditional cubic energy with retained-mean localization; Identify the actual nonreducing Gaussian fast inverse with its energy prior; Resolve the unlocalized local harmonic cubic exchange target; Keep the exact quadratic endpoint and its complete low-frequency source normalization; Control the entire conditioned quantum Gaussian fiber at each fixed block scale; Determine the first actual Wilson ground and chosen-chart source correction; Extend the actual path-source fast bound over every flat holonomy; Realize the uniform transverse harmonic source as an actual local covariant path observation; Compare the complete actual endpoint spectrum with all retained source directions included; Establish actual local true-ground score loss on complete additive selected source charts; Identify the true-vacuum literal coarse form and additive physical fast space; Lift the full harmonic boundary inequality to the entire Gaussian quantum fast form; Establish the actual two-face vertical, coarse and full-vacuum complementary forms; Retain planar and periodic three-dimensional harmonic interfaces while isolating box-scale fast modes; Derive an energy comparison that retains arbitrary cross coupling and the induced norm; construct actual discrete-time Wilson blocks and cancel disconnected vacuum factors at a fixed observable; control the complete marked shell and onto literal source frame at fixed spatial spacing; sum temporal Wilson matching and transport the carrier with its specified observable overlap; assemble the full spatial Schur excess, second Wilson and source jets and accumulated physical-clock budget; replace the failed global fast-form comparison by an admissible selected-resolvent Schur identity; reconcile the broader G19 corpus and derive the flat reference and Cauchy repair inputs; audit the GPU and Lean closure claims against their actual operators and hypotheses; Construct the complete first residual current for the fixed twelve-edge square; Derive the compatible-kernel reconstruction and moving-time full-space implication; Optimize finite decay channels and time-amplified errors in the physical clock.\par

\subsection*{G20 --- SAFE-region curvature certificate}
\addcontentsline{toc}{subsection}{G20 --- SAFE-region curvature certificate}
\noindent\emph{Tier 3 (unbounded); open.}\par
\smallskip\noindent The notes program's quantitative core: a certified lower bound on the projected physical Hessian H\_\allowbreak{}phys over a SAFE ball around the identity (claimed constants kappa* = 0.25, delta = 0.006, headline min eigenvalue 0.248), feeding curvature -> LSI -> gap. As reviewed, the headline is UNREPRODUCED: the sole source table is contradicted by the archive's own scans and independent recomputation (single-link Haar profile is flat at 1/4), the operator's code has never been executed, and no lattice size is stated anywhere. What discharging requires: run the imported\,\ldots\par

\subsection*{G21 --- Davies/Combes-Thomas O(m) decay for the massive Maxwell 1-form kernel}
\addcontentsline{toc}{subsection}{G21 --- Davies/Combes-Thomas O(m) decay for the massive Maxwell 1-form kernel}
\noindent\emph{Tier 3 (unbounded); partial.}\par
\smallskip\noindent The strongest result in the notes archive, and the first to survive a full-proof audit: |G(b,b')| <= (2/m\textasciicircum{}2) exp(-eta dist) for G = (m\textasciicircum{}2 + alpha d1* d1)\textasciicircum{}-1 on lattice 1-forms of a finite periodic lattice, with eta = arcosh(1 + m\textasciicircum{}2/(2 alpha C)) = 2 arsinh(m/(2 sqrt(alpha C))) -- an O(m) exponent where naive resolvent bounds give O(m\textasciicircum{}2). RE-ANCHORED after the audit: the canonical constant C is the WEIGHT-PROFILE row-sum C\_\allowbreak{}bdy(M\_\allowbreak{}1) of Def A.9.4 (Appendix H, Prop H.4.3); the proved\,\ldots\par
\smallskip\noindent\textbf{Ledger status.} Appendix H audited in full 2026-08-22, proposition by proposition: the proof is sound end to end, and the no-factor-2 perturbation bound (Lemma H.3.3) is correctly derived and numerically confirmed with near-tight margin (<= 0.97). Remaining open: the standing T2 sweep at several (m, alpha), the H.3.3 mechanism check, the arsinh-form float guard, and the deferred upstream Prop A.9.5.\par

\subsection*{G22 --- Coercivity via Cartan alignment}
\addcontentsline{toc}{subsection}{G22 --- Coercivity via Cartan alignment}
\noindent\emph{Tier 3 (unbounded); open.}\par
\smallskip\noindent The notes program's crux conjecture: on the rough set (average plaquette energy >= eps), configurations far from Cartan alignment have force norm bounded below, volume-uniformly. The original 6-constraints-vs-3-DOF counting argument is refuted (stationarity at a link is 3 equations, and three orthogonal antipodal pairs sum to zero with no common Cartan direction -- a registered FINDING check), so whatever truth the conjecture has must come from the gauge-coupled structure across links. Falsifier, required by this register's own discipline: exhibit a configuration with B\_\allowbreak{}avg >=\,\ldots\par
\smallskip\noindent\textbf{Ledger status.} The stated finite matrix model has Simon's transverse confinement bound H\_\allowbreak{}m>=T/2+g sqrt(m-1) sum|A\_\allowbreak{}i|, finite-dimensional compact resolvent and a raw ground-energy floor. These results are retained with their exact scalar split. WILSON\_\allowbreak{}G19\_\allowbreak{}REVIEW identifies the remaining derivation: realize this comparison for the actual gauge-coupled physical operator and prove volume-uniform constants after vacuum subtraction.\par
\smallskip\noindent\textbf{Live routes.}
\begin{itemize}[leftmargin=1.4em]
  \item \textbf{identify the transverse quantum form and uniform constants for the actual gauge-coupled physical operator.} Identify the actual physical transverse form and retained kinetic allocation. The cited finite-model proof uses finite-dimensional Rellich compactness and bounds E0 before subtraction; it supplies neither this operator identification nor uniform vacuum-complement coercivity for the coupled lattice.
\end{itemize}
\smallskip\noindent\textbf{Completed routes (1).} derive transverse confinement along classical flat directions in the stated finite matrix model.\par

\subsection*{G23 --- The OS bridge with uniform constants}
\addcontentsline{toc}{subsection}{G23 --- The OS bridge with uniform constants}
\noindent\emph{Tier 3 (unbounded); open.}\par
\smallskip\noindent Historical notes-program spine as a conditional chain: local curvature -> LSI -> diffusion gap -> exponential clustering -> OS reconstruction -> Hamiltonian mass gap, with two named holes that no archive document closes. (1) Uniformity: every implication holds at fixed cutoff and volume once a curvature floor rho\_\allowbreak{}0 is granted; nothing supplies rho\_\allowbreak{}0 uniform in a and volume.\par
\smallskip\noindent\textbf{Ledger status.} The actual fixed-scale Wilson transfer and complete physical odd source band remain established. The exact history theorem gives an isometry and transfer intertwiner for a true reflection/time-compatible pushforward. Gaussian observability now proves that its range is a history-frequency subspace, not an equal-time coordinate count; an arbitrarily small observed residue can make the whole physical two-mode SU(2) history space visible.\par
\smallskip\noindent\textbf{Live routes.}
\begin{itemize}[leftmargin=1.4em]
  \item \textbf{Transport the nonlinear excess through a compatible actual generated-history hierarchy.} Use the actual full fast inverse and bounded-mean energy input in the complete nonlinear generated coefficient, retaining the harmonic witness and spatial cancellations. Establish uniform interacting remainder and complete source/history matching with a common physical clock. The chosen local source chart need not generate the whole matrix observation sigma algebra;
  \item \textbf{Compare the actual generated history dynamics with a controlled physical low-energy model.} The actual history map, jointly invariant algebra, generated measure, memory and physical clock remain explicit. Literal true-vacuum source ranges now have complete low-window frames and fast floors in additive Wilson models, including common Gauss, and in the regulated full Gaussian boundary model. These equal-time ranges need not be the reducing OS-history range.
  \item \textbf{control the actual physical trial residual after local vacuum cancellation.} WR14-WR16 construct a gauge-invariant trial with a nonnegative shared-link residual, but its global sum is not a centered gap budget. Connected residual and complete boundary-source control are still needed. Independent copies refute a naive extensive residual loss.
  \item \textbf{prove physical conditional assembly and improved coercivity through spatial refinement.} Fixed-spacing substeps are complete: BA20--BA26 Cauchy-Schwarz true-vacuum bound eliminates the multiplicity penalty on single-link partitions (kappa <= 10/229, gap > 2.8679 epsilon). SC17 bootstrap closes the connected spatial estimate for lambda <= 1/73, establishing weighted Hessian row <= 2/7 at weight 1025/1024, projection angle row kappa <= 10658/24923, C\_\allowbreak{}AT <= 24923/14265 < 7/4, and physical gap > 4 epsilon / 3. The rational budgets are machine-certified by the supplied Lean lemmas; \emph{Cannot decide:} \texttt{G19}, \texttt{G23}.
\end{itemize}
\smallskip\noindent\textbf{Untried routes.}
\begin{itemize}[leftmargin=1.4em]
  \item \textbf{establish tunable physical-time localization on a complete observable class.} Prove the localization and bad-event bounds for every radius with controlled prefactors and complete physical observables. A diffusion gap alone does not identify Euclidean-time spectral decay. A source with an unobserved slow mode falsifies an incomplete-source argument.
\end{itemize}
\smallskip\noindent\textbf{Dead routes, and why they died.}
\begin{itemize}[leftmargin=1.4em]
  \item \textbf{transfer an isolated physical rotor gap to a fixed finite unreduced boundary frame.} WR20 embeds the entire single-link electric Dirichlet form using the true patch ground and its Haar marginal. The unreduced boundary gap is at most 3epsilon and the number of modes below a fixed target grows as epsilon\textasciicircum{}(-3/2). This refutes the fixed finite frame, not a physical gap or assembly with full compact boundary retention. \emph{Cannot decide:} \texttt{G19}, \texttt{G23}.
  \item \textbf{replace physical spatial control by globally positive classical Wilson orbit curvature at weak coupling.} The exact regular SU2 plaquette curvature is 8/3-beta/2 at theta=2pi/3. Horizontal geodesics extend the failure to the stated finite cubic lattices. The classical weight is not the physical ground and Wilson flow has an expansion witness 8/7. \emph{Cannot decide:} \texttt{G19}, \texttt{G23}.
  \item \textbf{obtain a positive physical bound by discarding WR25 local remainders before vacuum subtraction.} VA1-VA8 give an actual connected periodic SU2 Wilson test at L>=3, k/epsilon>=32. The discarded vacuum remainder is at least 15delta p0 |E| with p0>0 independent of volume. A physical mean-zero local phase has expectation of A-F at most delta(4-15p0 |E|), negative at large volume. \emph{Cannot decide:} \texttt{G19}, \texttt{G23}.
\end{itemize}
\smallskip\noindent\textbf{Completed routes (15).} Identify the actual nonreducing Gaussian fast inverse with its energy prior; Identify exact Gaussian endpoint memory and its first chosen-source nonlinear coefficient; Retain exact endpoint memory while comparing the entire additive physical low source window; Determine what the exact harmonic history block observes before assigning an eliminated space; repair negative Wilson curvature and derive the complete compact rotor gap; retain the full compact boundary and prove the actual Wilson-link fast barrier; identify the finite ground transform and retain the full trial residual gap budget; derive the physical spectral-weight bound and optimized localization exponent; align conditional projections with the true Wilson vacuum and establish their local floor; optimize the actual conditional block gap and derive its connected angle and pressure identities; prove actual true-vacuum block assembly in the explicit electric strong-coupling window; derive an all-coupling background magnetic comparison and charged physical-time covariance bound; prove SC17 spatial decay and actual conditional assembly through its explicit strong-coupling threshold; construct the SC17 thermodynamic ground law and its closed physical form with nonzero plaquette overlap; identify actual SC17 thermodynamic physical-time correlations with the closed ground-law semigroup.\par

\subsection*{Discharged gaps}
\addcontentsline{toc}{subsection}{Discharged gaps}
\noindent Finished work stays recorded rather than re-queued. Each entry names what closed it.\par\smallskip
\begin{itemize}[leftmargin=1.4em]
  \item \texttt{G18} \textbf{The spectral bridge at fixed lattice spacing.} Fixed-spacing closure only; continuum source transport remains G19.
  \item \texttt{G2} \textbf{Regenerate all displayed tables in canonical u.} discharged 2026-08-21 --- the conversion statement is now arithmetic (coupling-erratum suite): the printed towers reproduce 4*Delta(3u/2) verbatim in u = beta/6 and would be off by exactly 4**r under the archived Y reading.
  \item \texttt{G24} \textbf{Derive the shared-link isotropy premise, the one unproved input of the second-order chain.} discharged 2026-08-28, the day it was opened. The premise is not an assumption: it follows from the order-2 Weingarten values.
  \item \texttt{G3} \textbf{Fourth-order adjudication --- decide C\_\allowbreak{}shp where it actually lives.} DISCHARGED 2026-09-04 (ADR 0024). C\_\allowbreak{}shp lives in the rotation amplitude rho, and rho is now decided: the historical pipeline's own Stage-3I ledger, the 2026-09-02 cluster assembly and a third engine agree on every cluster of the rotation record but the adjacent-face cube completion, where the historical ledger holds 8 of 24 orderings;
  \item \texttt{G5} \textbf{Tetrahedral local Haar-resolvent coefficient.} discharged 2026-08-21 --- the r = 2 instance of the primitive completion law is re-derived exactly at symbolic N (tetrahedral Haar-resolvent suite, workhouse.cellular): c\_\allowbreak{}2 = -8/(N(N\textasciicircum{}2-1)), from two temporal histories through the single 4-link intermediate loop, under the same convention that reproduces the sealed-core cube row, the prism 64/24 sector pair, and the pentagonal 120-history count. The falsification test the mobility theorem proposed did NOT kill the coefficient: the r = w\_\allowbreak{}min - 2 = 2 allowance carries nonzero primitive Haar-resolvent\,\ldots
\end{itemize}

% ---- end gen_obligations.tex ---------------------------------------
% ---- end 36_obligations.tex ----------------------------------------
% ---- begin 37_attack.tex -------------------------------------------
\section{Working the problems}
\label{sec:attack}

\S\ref{sec:obligations} states what must be proved.
This section is about how to start, and it is written for an agent with
a fresh context window and a week.

\subsection{Choosing}

Three properties make a target worth taking, and they are independent.

\emph{Decisive} --- solving it resolves a contradiction, discharges a
gap, or unblocks another problem. \emph{Ready} --- no unresolved input
is recorded. \emph{Bounded} --- someone has measured, or can cheaply
estimate, what it costs.

Most of the problems in \S\ref{sec:obligations} are decisive and ready.
Almost none are bounded, and that is the honest difficulty: they are
open analytic estimates, not computations waiting for a machine. The
two exceptions are worth naming, because an agent looking for
tractable decisive work should start at one of them.

\begin{itemize}
  \item \textbf{Problem~\ref{prob:g9}}, one complete sixth-order
    cluster. Finite, priced, and it determines an object that is
    already proved to exist (\S\ref{sec:sixth}): the extra rational
    band shape survives every local direct term, so computing the
    direct term settles its coefficients rather than gambling on
    whether there is anything there.
  \item \textbf{Acquiring the unobtained near neighbours}
    (\S\ref{sec:external}). Days, not weeks. It is the only item that
    can \emph{subtract} from the program, which is exactly why it
    should not keep being deferred in favour of items that can only
    add.
\end{itemize}

\subsection{What the published methods actually supply}

Several of the continuum problems have a named external method
attached. The attachment is real but narrower than it looks in a
citation list, and the narrowing is the useful part.

\begin{description}[leftmargin=1.1em,style=nextline]
\item[Quantum-perturbation Riesz projection, and ground-state methods.]
  Verified against the fixed-spacing bridge and against the
  source-family problem. This is the machinery behind the discharged
  fixed-spacing construction; it is method, not result, and it is
  already spent --- it got the program to \S\ref{sec:wilson} and does
  not by itself go further.
\item[Cluster expansion methods.] Verified against the source-family
  problem and the continuum gap. Break~\ref{sec:break2} is the
  cautionary result here: the expansion's activity and a
  scale-independent physical mass cannot be read off one schedule.
  Scalar Kotecky--Preiss inequalities and raw moment estimates remain
  useful ingredients, but they construct neither the tilted activities
  nor a gap.
\item[WKB and tunnelling for operators on a vector bundle.] Supplies
  the local weighted comparison with actual Dirichlet eigenfunctions
  used in the conditional-transport route. It does \emph{not} supply
  the all-fiber domination of Problem~\ref{prob:m10}, nor the
  differentiated coupling-score asymptotics. The scalar one-well
  application is what has been used; the vector case is what is
  needed.
\item[Renormalization-group multiscale methods.] Supply the shape of
  Problem~\ref{prob:cauchy} and nothing about its missing power.
  Break~\ref{sec:break1} localizes the difficulty to an exact
  cancellation of the leading partition-function covariance.
\item[Osterwalder--Seiler positivity and transfer matrices.] Registered
  in the evidence map, and the classical background against which the
  fixed-spacing construction of \S\ref{sec:wilson} should be read.
  Novelty claims in that section must clear it.
\item[Stochastic-analysis approaches at strong coupling.] Present in
  the reading library, with mass-gap and exponential-clustering results
  at explicit small-coupling bounds. These live in the \emph{same}
  regime as \S\ref{sec:wilson} --- neighbourhoods of zero coupling ---
  and therefore also do not cross Break~\ref{sec:break5}. They are
  comparison, not a route.
\end{description}

The pattern is worth stating plainly: every external method that has
been verified against this program supplies control \emph{at fixed
scale or small coupling}. None of them supplies continuation in the
coupling, which is the one thing Problem~\ref{prob:trajectory} needs.
An agent that finds a method claiming otherwise should treat that as
the most interesting thing it found that week.

\subsection{A procedure for one problem}

\begin{enumerate}
  \item \textbf{Resolve the target to graph identifiers.} Run
    \texttt{workhouse why} on the gap and on each named input. Retain
    the briefing. You now have the routes, their states, and the
    \texttt{cannot\_decide} annotations.
  \item \textbf{Read the dead routes first} (\S\ref{sec:obligations}).
    If your idea matches one, read why it died before deciding it
    died for a reason that no longer applies. Sometimes that is true;
    the register records the reason so the judgement can be made
    rather than guessed.
  \item \textbf{Locate the available lemmas as \texttt{DERIV}
    statements}, not as prose. Appendix~\ref{app:concordance} gives
    each one a document, a digest and a line range. Read the
    hypotheses; the corpus records them separately from the statement
    because they are what gets dropped.
  \item \textbf{Check the literature both ways.} The evidence map for
    what has been verified against this claim; the reading library
    (Appendix~\ref{app:library}) and the extraction index
    (Appendix~\ref{app:extraction}) for whether the estimate you are
    about to attempt already exists somewhere.
  \item \textbf{Do the smallest instance first.} Every decisive-next-test
    in \S\ref{sec:obligations} is phrased this way --- two coupled
    blocks, one cluster, one blocking step --- because the program's
    own experience is that the general case hides which term is the
    obstruction and the smallest instance exposes it.
  \item \textbf{Record the attempt, including its failure.} A dead
    route is a result. The procedure is in the working agreement.
\end{enumerate}

\subsection{Failure modes with a track record}

Each of these has cost this program real time. They are listed in
descending order of how much.

\begin{description}[leftmargin=1.1em,style=nextline]
\item[Running an instrument that cannot decide the question.]
  Four sessions went into a sealed sweep that a static argument had
  already shown emits only a scalar at one momentum, a quantity
  structurally incapable of constraining the coefficient it was aimed
  at. The register now carries \texttt{cannot\_decide} for exactly this
  reason. Check it before running anything expensive.
\item[Interpolating without a degree bound.] Sixty-seven exact
  agreements between exact rationals do not determine a rational
  function. The fix that worked was not to prove the bound but to
  remove the need for it by computing in the function field
  (\S\ref{sec:symbolic}).
\item[Counting vertices and forgetting the projection.] The retracted
  tier-collapse mechanism (\S\ref{sec:fourth}) bounded the degree of
  the operator and omitted the power contributed by the carrier
  projection itself. Any degree argument in this setting must count the
  projection.
\item[A correct calculation in the wrong basis.] An assembled amplitude
  was right and was read in the conjugate basis, producing a third
  value in a two-value dispute. The signature was visible in the data:
  every $C$-odd component negated, every $C$-even one identical. When a
  new value appears in an existing dispute, test for conjugation before
  believing it.
\item[Quoting a tier from a record instead of a run.] A registered
  axiom set is evidence; it is not a build. If a tier is load-bearing
  for what you are about to claim, re-establish it --- the commands
  are in Appendix~\ref{app:repro} and take seconds.
\item[Treating a summary as a source.] Measured in this program at
  thirteen out of thirteen (Appendix~\ref{app:provenance}). Cite the
  primary, and check that the anchor resolves rather than only that the
  digest matches.
\end{description}

\subsection{Where new mathematics is cheap}

The mass-gap problems are hard by construction. The structural
questions below are not, and several are decisive for the theory even
though none of them is on the continuum path. An agent looking to
produce new results rather than to close the Clay problem should read
this list first.

\begin{description}[leftmargin=1.1em,style=nextline]
\item[The general grading law.] \S\ref{sec:grading} establishes
  $S_m(N) = c_m N^{-(2m-1)}(1+O(N^{-2}))$ for even $m \le 10$, by a
  four-channel cancellation with leading coefficients in the ratio
  $(-4,+2,+1,+1)$ summing to zero. The law for general $m$ is a
  conjecture. The cancellation is a statement about four specific
  intermediate representations, which is the kind of statement a
  character-theoretic argument settles at all orders at once. This is
  the most likely place in the program for a clean general theorem.
\item[The multi-face cumulant.] The one-face law is proved; the cluster
  cumulant version is not discharged, and the shared-link pair route
  carries it. Distinct from the item above and often conflated with it.
\item[The fifth-order pentagonal coefficient.] Registered open, finite,
  and it interacts with the centre-parity theorem of
  \S\ref{sec:grading}: the odd orders are determinant families, and the
  fifth-order determinant correction is what refuted the
  centre-only-circuits-are-dark claim. Computing the full coefficient
  is bounded work with a known consumer.
\item[The rank holdout ledger.] Fixed-rank values for a mid-range band
  of $N$. Pure computation, and it is the natural holdout against every
  closed form the symbolic engine has produced --- exactly the kind of
  check that would have caught an interpolation error before the
  function-field method removed the risk.
\item[Native torelon reruns.] Two string-tension coefficients, by a
  method the program already has.
\item[The hyperhoneycomb candidate.] An alternative tessellation with a
  first-order term. The reading library has a substantial
  Hamiltonian-and-quantum-simulation category, including improved
  honeycomb and hyperhoneycomb Hamiltonians, so the external groundwork
  exists.
\item[Classification of shortest physical temporal histories.] Open and
  unbounded, but it would consolidate the mobility taxonomy --- which
  is where the falsified $F{-}2$ rule lived, and where a replacement is
  still missing.
\item[The symbolic Gram transcript.] Would make the function-field
  method of \S\ref{sec:symbolic} auditable end to end rather than
  trusted at its output.
\end{description}

Three of these --- the general grading law, the fifth-order
coefficient, and the multi-face cumulant --- are in the same
neighbourhood as results this program has already proved, with the same
tools. That is the definition of cheap here, and it is where the next
novel theorem most plausibly comes from.
% ---- end 37_attack.tex ---------------------------------------------

\appendix
\part*{Appendices}
\addcontentsline{toc}{part}{Appendices}
% ---- begin A0_reproducibility.tex ----------------------------------
\section{Reproducibility}
\label{app:repro}

Every claim in this document that carries a \Tzero{}, \Tone{} or
\Ttwo{} badge can be re-established from a clean checkout in seconds to
minutes. This appendix says how, and --- more usefully --- says what each
kind of re-establishment does and does not prove.

\subsection{From a clean checkout}

\begin{verbatim}
git clone https://github.com/ats314/WORKHOUSE
cd WORKHOUSE
uv sync --all-extras --frozen
uv run --no-sync workhouse --help
\end{verbatim}

Then:

\begin{verbatim}
make verify        # every registered check: T1 and T2, a few seconds
make lean          # T0: proof-check the Lean core (requires elan)
make status        # the contradiction and gap registers
workhouse verify --tier 1              # only the exact re-derivations
workhouse verify -v                    # every tolerance, printed
workhouse verify --only 'h_4^side'     # one claim, with its numbers
\end{verbatim}

A single claim is the useful unit. \texttt{workhouse verify -{}-only} takes a
claim name and prints the numbers it establishes, which is what makes a
disagreement arguable rather than merely reported. A certification a
reader cannot cheaply reproduce is a claim of authority, not evidence,
and the design of this repository takes that seriously enough to attach
the reproducing command to every row of \srcpath{CERTIFIED.md}.

\subsection{Querying a claim}

\begin{verbatim}
workhouse why C2                   # everything recorded about one id
workhouse why t_N                  # a constant: value, formalizations, sources
workhouse branches C2              # both sides of a dispute, side by side
workhouse derive C2 G3 --out f.md  # evidence chain, registered edges only
workhouse search '5/612'           # search by exact VALUE, not by name
workhouse lit --for G19            # published work bearing on a claim
workhouse brief G19 --json         # target snapshot with freshness
\end{verbatim}

\texttt{search} matches by value rather than spelling, which matters
because the join keys of this corpus are exact rationals and no
natural-language query retrieves $109151/249696$. It also carries two
warnings a text search cannot: names the register forbids because they
regenerate a dissolved contradiction, and names coined in this program
that its own source corpus never uses.

\subsection{What each tier actually establishes}

\begin{description}[leftmargin=1.2em,style=nextline]
\item[\Tzero{} --- Lean 4, no \texttt{sorry}.]
  The statement is checked by the Lean kernel from explicit definitions,
  with axioms restricted to \texttt{propext}, \texttt{Classical.choice}
  and \texttt{Quot.sound}. The axiom set is exported and validated rather
  than asserted: \srcpath{scripts/export\_lean\_dependencies.py} builds
  strictly and records elaborated constants and transitive axioms into
  \srcpath{ledger/lean\_dependencies.json}, and the loader checks these
  against declaration records and source fingerprints. Identifiers
  mentioned in comments are not proof dependencies.

  What a T0 badge does \emph{not} establish is importance. Several
  registered theorems formalize arithmetic evaluations of named
  rationals. The registry records for each theorem whether it proves an
  entire named check (\texttt{promotes}) or one ingredient of a larger
  statement (\texttt{lean\_support}); supporting a statement never
  silently promotes the whole statement's tier.

\item[\Tone{} --- exact re-derivation.]
  The statement is recomputed symbolically in exact arithmetic from
  stated definitions. Corpus rationals are \texttt{sympy.Rational};
  values the corpus records only as floats carry a \texttt{\_NUM} suffix
  and are Python floats. A float that reads as exact is the most
  dangerous defect available in this repository, and the boundary is
  guarded by a dedicated test.

\item[\Ttwo{} --- numerical agreement.]
  Float agreement within a tolerance the check itself prints. A T2 pass
  is that and nothing more. Tolerances are never widened to make a
  disagreement disappear; a disagreement is registered as an explicit
  finding that asserts the discrepancy.

\item[\Tthree{} --- no whole-statement machine certification.]
  The default. An analytic result may be fully proved and still be T3:
  the badge reports machine coverage, not mathematical status. Status,
  evidence and tier are three independent fields in
  \srcpath{ledger/results.yaml} for exactly this reason.
\end{description}

\subsection{The one inference this apparatus does not license}

A finite check certifies the statement it checks. It does not certify a
general theorem of which that statement is an instance, and the converse
also holds: a complete analytic proof is not weakened by absent Lean
coverage.

This cuts both ways in practice, and both directions have caused errors
here. A suite of exact checks passing at ranks $N = 3, 4, 5, 6$ is not a
statement about all $N$ --- which is why the symbolic-$N$ engine of
\S\ref{sec:symbolic} exists, and why its degree bound was retracted when
one projection turned out to be uncounted. Conversely, a separate Lean
library in this program's archive comprises seventy-one modules with no
\texttt{sorry} and was found on audit to prove mostly tautologies; a
theorem count is not a measure of coverage, and no count in this
document should be read as one.

\subsection{Regenerating this document}

\begin{verbatim}
python paper/master/generate_apparatus.py
tectonic paper/master/workhouse_master.tex        # the monograph
tectonic paper/master/workhouse_band_paper.tex    # Part I, standalone
\end{verbatim}

The generator reads \srcpath{literature/index.yaml} and
\srcpath{ledger/\{gaps,results,theorems,contradictions\}.yaml} and writes
the bibliography, the external-evidence table, the route register, the
results table and the count macros. Every number the prose quotes is one
of those macros. If a count here disagrees with the repository, the
generator has not been re-run; the repository is right.
% ---- end A0_reproducibility.tex ------------------------------------
% ---- begin A1_verification.tex -------------------------------------
\section{Verification index}
\label{app:verification}

Table~\ref{tab:results} lists every registered analytic result with its
recorded status and evidence level. The two columns are independent
fields and are meant to be read against each other: \texttt{proven} with
\texttt{record-backed} evidence means the argument exists and the machine
artefact does not.

Of \ResTotal{} registered results, \ResProven{} are recorded
\texttt{proven}, \ResConditional{} \texttt{conditional},
\ResDisputed{} \texttt{disputed} and \ResFalsified{} \texttt{falsified}.
The Lean tree carries \LeanTotal{} registered theorems with
\LeanSorry{} occurrences of \texttt{sorry}.

A caution on both counts. \ResProven{} out of \ResTotal{} is a high
proportion because a claim is not registered here until its argument is
written down; the register measures what has been recorded, not what
fraction of the problem is solved. Likewise \LeanTotal{} theorems is a
count of declarations, not of mathematical content --- see
\S\ref{app:repro}, which explains why a theorem count is not a coverage
measure and cites a case in this program's own archive where seventy-one
\texttt{sorry}-free modules turned out to prove tautologies.

% ---- begin gen_verification.tex ------------------------------------
% Generated by paper/master/generate_apparatus.py.  Do not edit.
\begin{longtable}{@{}p{0.44\linewidth} l l@{}}
\caption{Analytic results by recorded status and evidence level. \emph{Status} is what the argument establishes; \emph{evidence} is what artefact backs it. The two are deliberately independent: a result can be \texttt{proven} in status and \texttt{record-backed} in evidence, meaning the proof exists and the machine artefact does not.}\label{tab:results}\\
\toprule
Result & Status & Evidence \\
\midrule
\endfirsthead
\toprule Result & Status & Evidence \\ \midrule \endhead
\bottomrule
\endfoot
\texttt{\scriptsize ANISOTROPY\_\allowbreak{}THREE\_\allowbreak{}COORDINATE\_\allowbreak{}SHARP\_\allowbreak{}BOUND} & proven & analytic \\
\texttt{\scriptsize ARCHIVE\_\allowbreak{}CENTER\_\allowbreak{}DEFECT\_\allowbreak{}OBSTRUCTIONS} & proven & analytic \\
\texttt{\scriptsize ARCHIVE\_\allowbreak{}CONCENTRATION\_\allowbreak{}LOCALIZATION} & proven & analytic \\
\texttt{\scriptsize ARCHIVE\_\allowbreak{}FINITE\_\allowbreak{}RANGE\_\allowbreak{}INVERSE} & proven & analytic \\
\texttt{\scriptsize ARCHIVE\_\allowbreak{}GLOBAL\_\allowbreak{}PAIRING\_\allowbreak{}OBSTRUCTION} & proven & analytic \\
\texttt{\scriptsize ARCHIVE\_\allowbreak{}HORIZONTAL\_\allowbreak{}RESTRICTION\_\allowbreak{}BOUNDARY} & proven & analytic \\
\texttt{\scriptsize ARCHIVE\_\allowbreak{}LYAPUNOV\_\allowbreak{}PATCHING} & proven & analytic \\
\texttt{\scriptsize ARCHIVE\_\allowbreak{}REFLECTION\_\allowbreak{}PERMANENCE} & proven & analytic \\
\texttt{\scriptsize ARCHIVE\_\allowbreak{}SIGNED\_\allowbreak{}SU2\_\allowbreak{}DIFFUSION} & proven & analytic \\
\texttt{\scriptsize ARCHIVE\_\allowbreak{}SMOOTH\_\allowbreak{}PROFILES} & proven & analytic \\
\texttt{\scriptsize ARCHIVE\_\allowbreak{}SOURCE\_\allowbreak{}TILT\_\allowbreak{}AND\_\allowbreak{}ROOTED\_\allowbreak{}BOUNDS} & proven & analytic \\
\texttt{\scriptsize ASSEMBLED\_\allowbreak{}Q4\_\allowbreak{}COERCIVITY} & proven & analytic \\
\texttt{\scriptsize BALABAN\_\allowbreak{}MULTISCALE\_\allowbreak{}CAUCHY\_\allowbreak{}SUMMABILITY} & proven & analytic \\
\texttt{\scriptsize COMMON\_\allowbreak{}SPACE\_\allowbreak{}CONTRACTIVE\_\allowbreak{}DRIFT} & proven & analytic \\
\texttt{\scriptsize COMPACT\_\allowbreak{}ROTOR\_\allowbreak{}FAST\_\allowbreak{}GAP} & proven & analytic \\
\texttt{\scriptsize CONDITIONAL\_\allowbreak{}GRADIENT\_\allowbreak{}REPAIR} & proven & analytic \\
\texttt{\scriptsize CONDITIONAL\_\allowbreak{}QUANTUM\_\allowbreak{}PATH\_\allowbreak{}COVARIANCE} & proven & analytic \\
\texttt{\scriptsize CREATOR\_\allowbreak{}PARENT\_\allowbreak{}GAP} & proven & analytic \\
\texttt{\scriptsize CREATOR\_\allowbreak{}VELOCITY\_\allowbreak{}INVERSION} & proven & analytic \\
\texttt{\scriptsize DIMENSION\_\allowbreak{}FIVE\_\allowbreak{}OPERATOR\_\allowbreak{}IRRELEVANCE} & proven & analytic \\
\texttt{\scriptsize DYNAMIC\_\allowbreak{}CONDITIONAL\_\allowbreak{}CUBIC\_\allowbreak{}ENERGY} & proven & analytic \\
\texttt{\scriptsize EARLIER\_\allowbreak{}B6\_\allowbreak{}RETAINED\_\allowbreak{}KRYLOV\_\allowbreak{}CAMPAIGN} & conditional & numerical \\
\texttt{\scriptsize EARLIER\_\allowbreak{}CONDITIONAL\_\allowbreak{}SPECTRAL\_\allowbreak{}FLOOR} & proven & analytic \\
\texttt{\scriptsize EARLIER\_\allowbreak{}CORNER\_\allowbreak{}BLOCKING\_\allowbreak{}REFLECTION\_\allowbreak{}DEFECT} & proven & analytic \\
\texttt{\scriptsize EARLIER\_\allowbreak{}CUBIC\_\allowbreak{}QUOTIENT\_\allowbreak{}REGULARITY} & proven & analytic \\
\texttt{\scriptsize EARLIER\_\allowbreak{}D4\_\allowbreak{}DISJOINT\_\allowbreak{}STAPLE\_\allowbreak{}COORDINATES} & proven & analytic \\
\texttt{\scriptsize EARLIER\_\allowbreak{}DECORATED\_\allowbreak{}HISTORY\_\allowbreak{}MERGING} & proven & analytic \\
\texttt{\scriptsize EARLIER\_\allowbreak{}DETERMINANT\_\allowbreak{}PARITY\_\allowbreak{}REPAIR} & proven & analytic \\
\texttt{\scriptsize EARLIER\_\allowbreak{}EQUAL\_\allowbreak{}RAY\_\allowbreak{}IDENTIFICATION} & proven & analytic \\
\texttt{\scriptsize EARLIER\_\allowbreak{}FINITE\_\allowbreak{}ORDER\_\allowbreak{}SUPPORT} & proven & analytic \\
\texttt{\scriptsize EARLIER\_\allowbreak{}GRAM\_\allowbreak{}NULL\_\allowbreak{}OPERATOR\_\allowbreak{}QUOTIENT} & proven & analytic \\
\texttt{\scriptsize EARLIER\_\allowbreak{}ISOLATED\_\allowbreak{}POSITIVE\_\allowbreak{}MATRIX\_\allowbreak{}ATOM} & proven & analytic \\
\texttt{\scriptsize EARLIER\_\allowbreak{}JOINT\_\allowbreak{}RANK\_\allowbreak{}VOLUME\_\allowbreak{}SCALING} & proven & analytic \\
\texttt{\scriptsize EARLIER\_\allowbreak{}LOCALIZED\_\allowbreak{}FRAME\_\allowbreak{}OBSTRUCTION} & proven & analytic \\
\texttt{\scriptsize EARLIER\_\allowbreak{}MOVING\_\allowbreak{}ADJOINT\_\allowbreak{}FORCE\_\allowbreak{}DERIVATIVE} & proven & analytic \\
\texttt{\scriptsize EARLIER\_\allowbreak{}NESTED\_\allowbreak{}QUOTIENT\_\allowbreak{}REDUCTION} & proven & analytic \\
\texttt{\scriptsize EARLIER\_\allowbreak{}POSITIVE\_\allowbreak{}SECTOR\_\allowbreak{}PHASE\_\allowbreak{}ISOLATION} & proven & analytic \\
\texttt{\scriptsize EARLIER\_\allowbreak{}REFLECTION\_\allowbreak{}ADAPTED\_\allowbreak{}BLOCKING} & proven & analytic \\
\texttt{\scriptsize EARLIER\_\allowbreak{}SHELL\_\allowbreak{}SYMMETRY\_\allowbreak{}RITZ\_\allowbreak{}CONTROL} & proven & analytic \\
\texttt{\scriptsize EARLIER\_\allowbreak{}SPHERE\_\allowbreak{}EQUAL\_\allowbreak{}COUPLING\_\allowbreak{}COVARIANCE} & proven & analytic \\
\texttt{\scriptsize EARLIER\_\allowbreak{}SU3\_\allowbreak{}POSITIVITY\_\allowbreak{}BOUNDARIES} & proven & analytic \\
\texttt{\scriptsize EARLIER\_\allowbreak{}VSU\_\allowbreak{}BOUNDED\_\allowbreak{}DOMAIN\_\allowbreak{}WELLPOSEDNESS} & proven & analytic \\
\texttt{\scriptsize EXPONENTIAL\_\allowbreak{}SOURCE\_\allowbreak{}TOTALITY} & proven & analytic \\
\texttt{\scriptsize FESHBACH\_\allowbreak{}RESOLVENT\_\allowbreak{}COMPARISON} & proven & analytic \\
\texttt{\scriptsize FORM\_\allowbreak{}SCHUR\_\allowbreak{}SCALE\_\allowbreak{}COMPARISON} & proven & analytic \\
\texttt{\scriptsize FULL\_\allowbreak{}GRAPH\_\allowbreak{}SOURCE\_\allowbreak{}CONGRUENCE} & proven & analytic \\
\texttt{\scriptsize G17\_\allowbreak{}CORRECTED\_\allowbreak{}KP\_\allowbreak{}POLYMER\_\allowbreak{}STABILITY} & proven & analytic \\
\texttt{\scriptsize G9\_\allowbreak{}COMBINED\_\allowbreak{}MIXING\_\allowbreak{}SUPPORT} & proven & analytic \\
\texttt{\scriptsize G9\_\allowbreak{}HERMITIAN\_\allowbreak{}FOLDED\_\allowbreak{}SUPPORT} & proven & analytic \\
\texttt{\scriptsize G9\_\allowbreak{}LOCAL\_\allowbreak{}SIXTH\_\allowbreak{}NONCANCELLATION} & proven & analytic \\
\texttt{\scriptsize G9\_\allowbreak{}ONE\_\allowbreak{}FACE\_\allowbreak{}SIXTH} & proven & analytic \\
\texttt{\scriptsize G9\_\allowbreak{}SIXTH\_\allowbreak{}ORDER\_\allowbreak{}RUR\_\allowbreak{}ACTIVATION} & disputed & record-backed \\
\texttt{\scriptsize GAUSSIAN\_\allowbreak{}FULL\_\allowbreak{}FAST\_\allowbreak{}GREEN} & proven & analytic \\
\texttt{\scriptsize GAUSSIAN\_\allowbreak{}OS\_\allowbreak{}OBSERVABILITY} & proven & analytic \\
\texttt{\scriptsize GAUSSIAN\_\allowbreak{}PATH\_\allowbreak{}ENDPOINT\_\allowbreak{}BASELINE} & proven & analytic \\
\texttt{\scriptsize GAUSSIAN\_\allowbreak{}QUANTUM\_\allowbreak{}FAST\_\allowbreak{}SOURCES} & proven & analytic \\
\texttt{\scriptsize GAUSSIAN\_\allowbreak{}SCHUR\_\allowbreak{}MEMORY} & proven & analytic \\
\texttt{\scriptsize GROUND\_\allowbreak{}MARGINAL\_\allowbreak{}SCHUR\_\allowbreak{}SCORE} & proven & analytic \\
\texttt{\scriptsize HODGE\_\allowbreak{}FESHBACH\_\allowbreak{}SPLITTING} & proven & analytic \\
\texttt{\scriptsize HODGE\_\allowbreak{}INDUCED\_\allowbreak{}EIGHTH\_\allowbreak{}ORDER\_\allowbreak{}MOMENT} & proven & analytic \\
\texttt{\scriptsize HODGE\_\allowbreak{}LEAKAGE\_\allowbreak{}ZERO\_\allowbreak{}EXTENSION} & proven & analytic \\
\texttt{\scriptsize HODGE\_\allowbreak{}RANK\_\allowbreak{}ONE\_\allowbreak{}SECULAR\_\allowbreak{}CUBIC} & proven & analytic \\
\texttt{\scriptsize INTRINSIC\_\allowbreak{}SCORE\_\allowbreak{}CURRENT} & proven & analytic \\
\texttt{\scriptsize LITERAL\_\allowbreak{}ENDPOINT\_\allowbreak{}SPECTRAL\_\allowbreak{}COMPARISON} & proven & analytic \\
\texttt{\scriptsize LOCAL\_\allowbreak{}CLASS\_\allowbreak{}WICK\_\allowbreak{}COEFFICIENTS} & proven & analytic \\
\texttt{\scriptsize LOCAL\_\allowbreak{}CLASS\_\allowbreak{}WICK\_\allowbreak{}RECURRENCE} & proven & analytic \\
\texttt{\scriptsize LOCAL\_\allowbreak{}CLASS\_\allowbreak{}WICK\_\allowbreak{}SIGNS} & proven & analytic \\
\texttt{\scriptsize LOCAL\_\allowbreak{}GRADIENT\_\allowbreak{}EXCITATION\_\allowbreak{}SUPPORT} & proven & analytic \\
\texttt{\scriptsize MESOSCOPIC\_\allowbreak{}NORMALIZED\_\allowbreak{}SOURCE} & proven & analytic \\
\texttt{\scriptsize MOVING\_\allowbreak{}TIME\_\allowbreak{}MULTICHANNEL\_\allowbreak{}RATE} & proven & analytic \\
\texttt{\scriptsize MOVING\_\allowbreak{}TIME\_\allowbreak{}MULTICHANNEL\_\allowbreak{}SHARPNESS} & proven & analytic \\
\texttt{\scriptsize MOVING\_\allowbreak{}TIME\_\allowbreak{}SPECTRAL\_\allowbreak{}GAP} & proven & analytic \\
\texttt{\scriptsize ODD\_\allowbreak{}ORDER\_\allowbreak{}CENTRE\_\allowbreak{}PARITY} & proven & analytic \\
\texttt{\scriptsize ORDERED\_\allowbreak{}CONTOUR\_\allowbreak{}ACTIVITIES} & proven & analytic \\
\texttt{\scriptsize OS\_\allowbreak{}HISTORY\_\allowbreak{}BLOCK\_\allowbreak{}INTERTWINER} & proven & analytic \\
\texttt{\scriptsize OS\_\allowbreak{}KERNEL\_\allowbreak{}MOVING\_\allowbreak{}TIME\_\allowbreak{}GAP} & proven & analytic \\
\texttt{\scriptsize PREDICTABLE\_\allowbreak{}SOURCE\_\allowbreak{}DRIFT} & proven & analytic \\
\texttt{\scriptsize PROJECTIVE\_\allowbreak{}SOURCE\_\allowbreak{}MARTINGALE} & proven & analytic \\
\texttt{\scriptsize REVERSE\_\allowbreak{}MARTINGALE\_\allowbreak{}BLOCKING} & proven & analytic \\
\texttt{\scriptsize RICCATI\_\allowbreak{}RESIDUAL\_\allowbreak{}CERTIFICATE} & proven & analytic \\
\texttt{\scriptsize ROUGH\_\allowbreak{}GAUGE\_\allowbreak{}COERCIVITY\_\allowbreak{}LOWER\_\allowbreak{}BOUND} & proven & analytic \\
\texttt{\scriptsize RSM\_\allowbreak{}R\_\allowbreak{}CARRIER\_\allowbreak{}SYMBOL\_\allowbreak{}MASTER\_\allowbreak{}THEOREM} & proven & analytic \\
\texttt{\scriptsize SELECTED\_\allowbreak{}INVERSE\_\allowbreak{}WITHOUT\_\allowbreak{}DIAGONAL} & proven & analytic \\
\texttt{\scriptsize SHARP\_\allowbreak{}CUTOFF\_\allowbreak{}SOURCE\_\allowbreak{}GAP\_\allowbreak{}BUDGET} & proven & analytic \\
\texttt{\scriptsize SOURCE\_\allowbreak{}COMPLEMENT\_\allowbreak{}COMPLETE\_\allowbreak{}WINDOW} & proven & analytic \\
\texttt{\scriptsize SOURCE\_\allowbreak{}CONGRUENCE\_\allowbreak{}METRIC\_\allowbreak{}TRANSPORT} & proven & analytic \\
\texttt{\scriptsize SOURCE\_\allowbreak{}CURRENT\_\allowbreak{}BOUNDED\_\allowbreak{}W6} & proven & analytic \\
\texttt{\scriptsize SOURCE\_\allowbreak{}CURRENT\_\allowbreak{}VARIATIONAL\_\allowbreak{}ENERGY} & proven & analytic \\
\texttt{\scriptsize SQUARE\_\allowbreak{}FIRST\_\allowbreak{}CURRENT\_\allowbreak{}BUDGET} & proven & analytic \\
\texttt{\scriptsize SQUARE\_\allowbreak{}FIRST\_\allowbreak{}RESIDUAL\_\allowbreak{}CURRENT} & proven & analytic \\
\texttt{\scriptsize SQUARE\_\allowbreak{}SCALAR\_\allowbreak{}SOURCE\_\allowbreak{}ODD\_\allowbreak{}SCHUR\_\allowbreak{}VANISHING} & proven & analytic \\
\texttt{\scriptsize TETRAHEDRAL\_\allowbreak{}FLUX\_\allowbreak{}PROJECTION\_\allowbreak{}BOUNDARY} & proven & analytic \\
\texttt{\scriptsize TETRAHEDRAL\_\allowbreak{}S4\_\allowbreak{}COMMUTANT\_\allowbreak{}RESOLUTION} & proven & analytic \\
\texttt{\scriptsize TIER\_\allowbreak{}COLLAPSE\_\allowbreak{}ACTUAL\_\allowbreak{}H4\_\allowbreak{}SUPPORT} & proven & analytic \\
\texttt{\scriptsize TIER\_\allowbreak{}COLLAPSE\_\allowbreak{}IS\_\allowbreak{}R\_\allowbreak{}DEGREE} & falsified & analytic \\
\texttt{\scriptsize TRUE\_\allowbreak{}GROUND\_\allowbreak{}CENTER\_\allowbreak{}SCORE\_\allowbreak{}OBSTRUCTION} & proven & analytic \\
\texttt{\scriptsize UNIFORM\_\allowbreak{}MARKED\_\allowbreak{}CLUSTERING} & proven & analytic \\
\texttt{\scriptsize UNIVERSAL\_\allowbreak{}CELLULAR\_\allowbreak{}HODGE\_\allowbreak{}SPECTRUM} & proven & analytic \\
\texttt{\scriptsize W6\_\allowbreak{}ACTUAL\_\allowbreak{}CENTERED\_\allowbreak{}COMPLEMENT\_\allowbreak{}SUPPRESSION} & proven & analytic \\
\texttt{\scriptsize W6\_\allowbreak{}ANGULAR\_\allowbreak{}REFERENCE\_\allowbreak{}SCORE\_\allowbreak{}BUDGET} & proven & analytic \\
\texttt{\scriptsize W6\_\allowbreak{}R10\_\allowbreak{}CONDITIONAL\_\allowbreak{}DERIVATIVE} & proven & analytic \\
\texttt{\scriptsize W6\_\allowbreak{}R10\_\allowbreak{}H4\_\allowbreak{}FROM\_\allowbreak{}H0} & proven & analytic \\
\texttt{\scriptsize W6\_\allowbreak{}R10\_\allowbreak{}PROJECTION\_\allowbreak{}JET\_\allowbreak{}RECURRENCE} & proven & analytic \\
\texttt{\scriptsize W6\_\allowbreak{}UNIFORM\_\allowbreak{}AGMON\_\allowbreak{}COMPLEMENT\_\allowbreak{}GAP} & proven & analytic \\
\texttt{\scriptsize WILSON\_\allowbreak{}ACTIVITY\_\allowbreak{}EXTRACTION} & proven & analytic \\
\texttt{\scriptsize WILSON\_\allowbreak{}ACTUAL\_\allowbreak{}BLOCK\_\allowbreak{}FAST\_\allowbreak{}COMPLEMENT} & proven & analytic \\
\texttt{\scriptsize WILSON\_\allowbreak{}BLOCK\_\allowbreak{}SCORE\_\allowbreak{}OBSTRUCTION} & proven & analytic \\
\texttt{\scriptsize WILSON\_\allowbreak{}CARDINALITY\_\allowbreak{}CHART} & proven & analytic \\
\texttt{\scriptsize WILSON\_\allowbreak{}COEFFICIENT\_\allowbreak{}LOCALITY} & proven & analytic \\
\texttt{\scriptsize WILSON\_\allowbreak{}COMMON\_\allowbreak{}GAUSS\_\allowbreak{}LITERAL\_\allowbreak{}FAST\_\allowbreak{}COMPLEMENT} & proven & analytic \\
\texttt{\scriptsize WILSON\_\allowbreak{}COVARIANT\_\allowbreak{}BLOCK\_\allowbreak{}SOURCES} & proven & analytic \\
\texttt{\scriptsize WILSON\_\allowbreak{}CREATOR\_\allowbreak{}LIMIT} & proven & analytic \\
\texttt{\scriptsize WILSON\_\allowbreak{}CUBIC\_\allowbreak{}GROUND\_\allowbreak{}TRANSFER} & proven & analytic \\
\texttt{\scriptsize WILSON\_\allowbreak{}ENDPOINT\_\allowbreak{}COMPLETE\_\allowbreak{}CLUSTER} & proven & analytic \\
\texttt{\scriptsize WILSON\_\allowbreak{}ENDPOINT\_\allowbreak{}EQUATION} & proven & analytic \\
\texttt{\scriptsize WILSON\_\allowbreak{}FINITE\_\allowbreak{}CELL\_\allowbreak{}PHYSICAL\_\allowbreak{}GAP} & proven & analytic \\
\texttt{\scriptsize WILSON\_\allowbreak{}FIXED\_\allowbreak{}ORDER\_\allowbreak{}CHART} & proven & analytic \\
\texttt{\scriptsize WILSON\_\allowbreak{}FLAT\_\allowbreak{}BACKGROUND\_\allowbreak{}FAST\_\allowbreak{}SOURCES} & proven & analytic \\
\texttt{\scriptsize WILSON\_\allowbreak{}GLOBAL\_\allowbreak{}VERTICAL\_\allowbreak{}BARRIER} & proven & analytic \\
\texttt{\scriptsize WILSON\_\allowbreak{}GROUND\_\allowbreak{}BUNDLE\_\allowbreak{}RELATIVE\_\allowbreak{}FORM} & proven & analytic \\
\texttt{\scriptsize WILSON\_\allowbreak{}HARMONIC\_\allowbreak{}BOUNDARY} & proven & analytic \\
\texttt{\scriptsize WILSON\_\allowbreak{}HARMONIC\_\allowbreak{}CUBIC\_\allowbreak{}OBSTRUCTION} & proven & analytic \\
\texttt{\scriptsize WILSON\_\allowbreak{}INFINITE\_\allowbreak{}PHYSICAL\_\allowbreak{}BAND} & proven & analytic \\
\texttt{\scriptsize WILSON\_\allowbreak{}KINETIC\_\allowbreak{}WINDOW} & proven & analytic \\
\texttt{\scriptsize WILSON\_\allowbreak{}LITERAL\_\allowbreak{}COARSE\_\allowbreak{}FORM} & proven & analytic \\
\texttt{\scriptsize WILSON\_\allowbreak{}LITERAL\_\allowbreak{}FAST\_\allowbreak{}COMPLEMENT} & proven & analytic \\
\texttt{\scriptsize WILSON\_\allowbreak{}LOCALIZED\_\allowbreak{}TRUE\_\allowbreak{}GROUND\_\allowbreak{}SCORE} & proven & analytic \\
\texttt{\scriptsize WILSON\_\allowbreak{}PARENT\_\allowbreak{}GNS} & proven & analytic \\
\texttt{\scriptsize WILSON\_\allowbreak{}PHYSICAL\_\allowbreak{}SOURCE\_\allowbreak{}RANK\_\allowbreak{}REPAIR} & proven & analytic \\
\texttt{\scriptsize WILSON\_\allowbreak{}ROOTED\_\allowbreak{}CONTRACTION} & proven & analytic \\
\texttt{\scriptsize WILSON\_\allowbreak{}SAME\_\allowbreak{}WEIGHT\_\allowbreak{}OBSTRUCTION} & proven & analytic \\
\texttt{\scriptsize WILSON\_\allowbreak{}SECOND\_\allowbreak{}ORDER\_\allowbreak{}CHART} & proven & analytic \\
\texttt{\scriptsize WILSON\_\allowbreak{}STRIP\_\allowbreak{}BO\_\allowbreak{}FIRST\_\allowbreak{}TERM} & proven & analytic \\
\texttt{\scriptsize WILSON\_\allowbreak{}SYMMETRIC\_\allowbreak{}CREATORS} & proven & analytic \\
\texttt{\scriptsize WILSON\_\allowbreak{}THREE\_\allowbreak{}DIMENSIONAL\_\allowbreak{}HARMONIC\_\allowbreak{}BOUNDARY} & proven & analytic \\
\texttt{\scriptsize WILSON\_\allowbreak{}TRANSFER\_\allowbreak{}RESOLVENT} & proven & analytic \\
\texttt{\scriptsize WILSON\_\allowbreak{}TWO\_\allowbreak{}SQUARE\_\allowbreak{}PHYSICAL\_\allowbreak{}SHELLS} & proven & analytic \\
\texttt{\scriptsize WILSON\_\allowbreak{}TWO\_\allowbreak{}STRIP\_\allowbreak{}PHYSICAL\_\allowbreak{}SPLITTING} & proven & analytic \\
\texttt{\scriptsize WILSON\_\allowbreak{}UNIFORM\_\allowbreak{}FINITE\_\allowbreak{}SHELL} & proven & analytic \\
\texttt{\scriptsize WILSON\_\allowbreak{}VACUUM\_\allowbreak{}COMPRESSION} & proven & analytic \\
\texttt{\scriptsize WILSON\_\allowbreak{}VACUUM\_\allowbreak{}SPECTRAL\_\allowbreak{}FLOW} & proven & analytic \\
\texttt{\scriptsize WILSON\_\allowbreak{}VERTICAL\_\allowbreak{}FAST\_\allowbreak{}ENERGY} & proven & analytic \\
\texttt{\scriptsize WILSON\_\allowbreak{}WEIGHTED\_\allowbreak{}ACTIVITIES} & proven & analytic \\
\end{longtable}
% ---- end gen_verification.tex --------------------------------------
% ---- end A1_verification.tex ---------------------------------------
% ---- begin A2_provenance.tex ---------------------------------------
\section{Provenance}
\label{app:provenance}

The verification repository is one layer over a larger research archive.
Both are needed to read a result's history, and they answer different
questions: the repository says what is currently checked, the archive
says what was actually done and when.

\subsection{The two layers}

The repository (\srcpath{https://github.com/ats314/WORKHOUSE}) holds the
ledgers, the invariant suites, the Lean tree, the generated views
(\srcpath{FRONTIER.md}, \srcpath{CERTIFIED.md}, \srcpath{index/}) and the
manuscripts. Its generated views are maintained by the verifier commands
and are never hand-edited: a test fails if either has drifted, because a
generated file that has drifted still reads as current.

The archive holds the primary sources: manuscripts, proof collections,
appendices, simulation campaigns, a separate Lean library, reference
papers and the raw computational runs. Its organizing rule is that
nothing is erased --- failed attempts, duplicates, superseded drafts and
uncommitted work are all preserved, because deleting a refuted claim
destroys the evidence that it was tried.

\subsection{Reading a source document}

Documents under \srcpath{theory/}, \srcpath{corpus-import/} and
\srcpath{settlement/} state what someone believed at the time. They are
immutable evidence and are never edited to make a check pass; when a
check disagrees with a document, the disagreement is registered as a
finding. \srcpath{theory/SHA256SUMS} pins the contents.
\srcpath{theory/superseded/} holds documents kept for the audit trail,
and none of them should be read as current.

This matters more than it usually would, because parts of the corpus
were written with AI assistance and some of it is wrong. The repository
has already caught a reversed tensor-product identity, a one-ulp
transcription error, a tolerance quoted tighter than its own data, and a
mechanism proposed here and then retracted. Confidence of phrasing
carries no evidential weight in either direction.

\subsection{Synthesis documents are not evidence}

A specific hazard, recorded because it has cost this program time. Several
large synthesis documents in the archive were produced by automated
agents summarizing primary sources. Audit found the primary sources
honest about their own scope and the synthesis layer systematically
stripping their caveats --- upgrading conditional statements to proven,
attaching Lean support that does not cover the statement it is attached
to, and citing a formalization library whose theorems are tautologies.

The rule this produced: cite the primary source, never the synthesis of
it, and verify that a claimed anchor resolves rather than only that its
hash matches. A correct SHA-256 over a fabricated locator is a failure
mode that has actually occurred here.

\subsection{A measurement of that hazard}

The present edition was assembled with a deliberate test of the effect
above, and the result is worth recording because it is quantitative.

Thirteen result clusters of this corpus were summarized independently,
each by an agent instructed to extract exact statements with their
recorded status, hypotheses and verification tier, and explicitly told
not to upgrade anything the ledger records as conditional. Each summary
was then audited by a second agent instructed to refute it against the
primary files.

All thirteen summaries were found to overstate in at least one respect.
The failure modes recurred:

\begin{itemize}
  \item a status upgraded past what the ledger records;
  \item a hypothesis present in the ledger and absent from the summary
    --- most often a periodicity condition, a rank restriction, or a
    regime bound;
  \item a verification tier asserted with no check behind it, or with a
    check that tests a tautology;
  \item novelty asserted against an index the same summary elsewhere
    notes is incomplete;
  \item a cited anchor that does not resolve to the named object, with
    the surrounding hash correct.
\end{itemize}

Five of those findings changed statements in this document, and each is
named at the point of correction rather than silently applied:
\S\ref{sec:grading} on the planar reduction, \S\ref{sec:fourth} on the
rank scope of the operator splitting, \S\ref{sec:symbolic} on what the
formal proof of the coefficient identity actually checks,
\S\ref{sec:wilson} on chains versus operators, and \S\ref{sec:wilson}
again on the overlap exponent.

The lesson for the next agent is not that summarization is useless ---
the extractions found real material that a direct reading missed,
including the strongest external corroboration in
\S\ref{sec:external}. It is that a summary of this corpus should be
treated as a \emph{lead}, never as a citation, and that the audit step
is not optional overhead. It changed one statement in every three that
mattered.
% ---- end A2_provenance.tex -----------------------------------------
% ---- begin B0_concordance.tex --------------------------------------
\section{Derivation concordance}
\label{app:concordance}

Every registered derivation statement, grouped by the document it lives
in, with the digest of the bytes it was read from and the heading and
line range that locate it inside that document.

\DerivStatements{} statements across \CiteDocs{} documents. The
\emph{Formal} column is \texttt{full} when a Lean theorem covers the
whole statement, \texttt{partial} when a theorem is registered as
supporting one ingredient of it, and \texttt{---} when neither. Partial
support never promotes the whole statement's tier, which is why the
two are separate columns of the underlying registry rather than one.

A caution on using this table. The status recorded here describes the
\emph{source argument} as the corpus assessed it, not an independent
re-derivation. Empty formal coverage means no complete registered
formalization exists in this inventory --- it is not a judgement that
the analytic proof is invalid, and reading it as one would invert the
meaning of the default tier (\S\ref{sec:howto}).

% ---- begin gen_concordance.tex -------------------------------------
% Generated by paper/master/generate_sources.py.  Do not edit.
\noindent The concordance below covers all 332 registered derivation statements across 40 source documents. Each document carries the SHA-256 of the bytes its statements were read from; if the file has changed, the digest has changed and the statements must be re-located rather than assumed.\par\medskip
\subsection*{CITE:\allowbreak{}ARCHIVE\_\allowbreak{}FUNCTIONAL}
\addcontentsline{toc}{subsection}{ARCHIVE\_\allowbreak{}FUNCTIONAL}
\noindent\srcpath{docs/derivations/archive-functional-inequalities.md}\\
{\footnotesize SHA-256 \texttt{048fbf01cdc55ee4}\ldots}\par\smallskip
\begin{longtable}{@{}p{0.30\linewidth} p{0.30\linewidth} l l@{}}
\toprule Statement & Locator & Status & Formal \\ \midrule
\endfirsthead
\toprule Statement & Locator & Status & Formal \\ \midrule \endhead
\bottomrule\endfoot
{\scriptsize\ttfamily F1} & {\scriptsize \#\# F1. Trace defect, faithfulness and the vacuum set (lines 13-22)} & {\scriptsize proven} & {\scriptsize ---} \\
{\scriptsize\ttfamily F2} & {\scriptsize \#\# F2. Global gradient domination without a quotient singularity (lines 23-37)} & {\scriptsize proven} & {\scriptsize ---} \\
{\scriptsize\ttfamily F3} & {\scriptsize \#\# F3. Plaquette lifts preserve link derivatives (lines 38-53)} & {\scriptsize proven} & {\scriptsize ---} \\
{\scriptsize\ttfamily F4} & {\scriptsize \#\# F4. Aggregate defect comparisons (lines 54-64)} & {\scriptsize proven} & {\scriptsize ---} \\
{\scriptsize\ttfamily F5} & {\scriptsize \#\# F5. Exact drift identity and the squared-profile budget (lines 65-82)} & {\scriptsize proven} & {\scriptsize ---} \\
{\scriptsize\ttfamily F6} & {\scriptsize \#\# F6. General profiles with vanishing derivative at zero (lines 83-96)} & {\scriptsize proven} & {\scriptsize ---} \\
{\scriptsize\ttfamily F7} & {\scriptsize \#\# F7. Linear pairing implies a genuinely uniform small-set drift (lines 97-114)} & {\scriptsize proven} & {\scriptsize ---} \\
{\scriptsize\ttfamily F8} & {\scriptsize \#\# F8. Lyapunov energy conversion as a completed square (lines 115-129)} & {\scriptsize proven} & {\scriptsize ---} \\
{\scriptsize\ttfamily F9} & {\scriptsize \#\# F9. Local-to-global Poincare patching (lines 130-144)} & {\scriptsize proven} & {\scriptsize ---} \\
{\scriptsize\ttfamily F10} & {\scriptsize \#\# F10. Averaging gives a volume-scale gradient bound (lines 145-156)} & {\scriptsize proven} & {\scriptsize ---} \\
{\scriptsize\ttfamily F11} & {\scriptsize \#\# F11. Herbst concentration and the mean threshold (lines 157-172)} & {\scriptsize proven} & {\scriptsize ---} \\
{\scriptsize\ttfamily F12} & {\scriptsize \#\# F12. Covariance localization with an explicit error budget (lines 173-190)} & {\scriptsize proven} & {\scriptsize ---} \\
{\scriptsize\ttfamily F13} & {\scriptsize \#\# F13. Absorbing a volume tail into a distance exponent (lines 191-203)} & {\scriptsize proven} & {\scriptsize ---} \\
{\scriptsize\ttfamily F14} & {\scriptsize \#\# F14. Averaged smallness does not imply plaquettewise smallness (lines 204-213)} & {\scriptsize proven} & {\scriptsize ---} \\
{\scriptsize\ttfamily F15} & {\scriptsize \#\# F15. Center configurations refute the proposed Laplacian identity (lines 214-229)} & {\scriptsize proven} & {\scriptsize ---} \\
{\scriptsize\ttfamily F16} & {\scriptsize \#\# F16. Extensive center defects obstruct uniform pairing coercivity (lines 230-270)} & {\scriptsize proven} & {\scriptsize ---} \\
{\scriptsize\ttfamily F17} & {\scriptsize \#\# F17. The exact SU(2) diffusion budget retains favorable signs (lines 271-311)} & {\scriptsize proven} & {\scriptsize ---} \\
\end{longtable}

\subsection*{CITE:\allowbreak{}ARCHIVE\_\allowbreak{}SOURCE\_\allowbreak{}TRANSPORT}
\addcontentsline{toc}{subsection}{ARCHIVE\_\allowbreak{}SOURCE\_\allowbreak{}TRANSPORT}
\noindent\srcpath{docs/derivations/archive-source-transport.md}\\
{\footnotesize SHA-256 \texttt{4a9d953d888e1b0e}\ldots}\par\smallskip
\begin{longtable}{@{}p{0.30\linewidth} p{0.30\linewidth} l l@{}}
\toprule Statement & Locator & Status & Formal \\ \midrule
\endfirsthead
\toprule Statement & Locator & Status & Formal \\ \midrule \endhead
\bottomrule\endfoot
{\scriptsize\ttfamily S1} & {\scriptsize \#\# S1. Inhomogeneous source tilting is a partition ratio (lines 10-19)} & {\scriptsize proven} & {\scriptsize ---} \\
{\scriptsize\ttfamily S2} & {\scriptsize \#\# S2. Source response is covariance and log partition is convex (lines 20-31)} & {\scriptsize proven} & {\scriptsize ---} \\
{\scriptsize\ttfamily S3} & {\scriptsize \#\# S3. Pointwise envelopes do not order normalized tilts (lines 32-44)} & {\scriptsize proven} & {\scriptsize ---} \\
{\scriptsize\ttfamily S4} & {\scriptsize \#\# S4. Bounded log likelihood ratios transfer tilted estimates (lines 45-57)} & {\scriptsize proven} & {\scriptsize ---} \\
{\scriptsize\ttfamily S5} & {\scriptsize \#\# S5. Free energy implies conditional rooted-capacity summability (lines 58-78)} & {\scriptsize proven} & {\scriptsize ---} \\
{\scriptsize\ttfamily S6} & {\scriptsize \#\# S6. A conservative rooted-animal count (lines 79-94)} & {\scriptsize proven} & {\scriptsize ---} \\
{\scriptsize\ttfamily S7} & {\scriptsize \#\# S7. Reflection-positive deterministic pushforward (lines 95-107)} & {\scriptsize proven} & {\scriptsize ---} \\
{\scriptsize\ttfamily S8} & {\scriptsize \#\# S8. Projective and weak limits preserve the stated OS forms (lines 108-122)} & {\scriptsize proven} & {\scriptsize ---} \\
{\scriptsize\ttfamily S9} & {\scriptsize \#\# S9. Reflection equivariance alone has a counterexample (lines 123-134)} & {\scriptsize proven} & {\scriptsize ---} \\
{\scriptsize\ttfamily S10} & {\scriptsize \#\# S10. Block Schur bound with row and column control (lines 135-147)} & {\scriptsize proven} & {\scriptsize ---} \\
{\scriptsize\ttfamily S11} & {\scriptsize \#\# S11. Finite-range inverse transport with an explicit exponent (lines 148-162)} & {\scriptsize proven} & {\scriptsize ---} \\
{\scriptsize\ttfamily S12} & {\scriptsize \#\# S12. Massive Maxwell row sums from incidence counting (lines 163-178)} & {\scriptsize proven} & {\scriptsize ---} \\
{\scriptsize\ttfamily S13} & {\scriptsize \#\# S13. Cochain horizontals reduce the massive Maxwell operator (lines 179-186)} & {\scriptsize proven} & {\scriptsize ---} \\
{\scriptsize\ttfamily S14} & {\scriptsize \#\# S14. Restriction need not preserve an ambient local kernel (lines 187-201)} & {\scriptsize proven} & {\scriptsize ---} \\
\end{longtable}

\subsection*{CITE:\allowbreak{}EARLIER\_\allowbreak{}PROOF\_\allowbreak{}RECOVERY}
\addcontentsline{toc}{subsection}{EARLIER\_\allowbreak{}PROOF\_\allowbreak{}RECOVERY}
\noindent\srcpath{docs/derivations/earlier-proof-recovery.md}\\
{\footnotesize SHA-256 \texttt{4b5b067ee7697cf5}\ldots}\par\smallskip
\begin{longtable}{@{}p{0.30\linewidth} p{0.30\linewidth} l l@{}}
\toprule Statement & Locator & Status & Formal \\ \midrule
\endfirsthead
\toprule Statement & Locator & Status & Formal \\ \midrule \endhead
\bottomrule\endfoot
{\scriptsize\ttfamily CUBIC\_\allowbreak{}QUOTIENT\_\allowbreak{}REGULARITY} & {\scriptsize \#\# E1. Cubic quotient regularity trichotomy} & {\scriptsize proven} & {\scriptsize ---} \\
{\scriptsize\ttfamily EQUAL\_\allowbreak{}RAY\_\allowbreak{}IDENTIFICATION} & {\scriptsize \#\# E2. Equal-ray recovery and an independent holdout} & {\scriptsize proven} & {\scriptsize ---} \\
{\scriptsize\ttfamily LOCALIZED\_\allowbreak{}FRAME\_\allowbreak{}OBSTRUCTION} & {\scriptsize \#\# E3. No uniformly localized finite translation frame} & {\scriptsize proven} & {\scriptsize ---} \\
{\scriptsize\ttfamily JOINT\_\allowbreak{}RANK\_\allowbreak{}VOLUME\_\allowbreak{}SCALING} & {\scriptsize \#\# E4. Sharp joint rank-volume scaling} & {\scriptsize proven} & {\scriptsize ---} \\
{\scriptsize\ttfamily FINITE\_\allowbreak{}ORDER\_\allowbreak{}SUPPORT} & {\scriptsize \#\# E5a. Bounded support at a fixed order} & {\scriptsize proven} & {\scriptsize ---} \\
{\scriptsize\ttfamily GRAM\_\allowbreak{}NULL\_\allowbreak{}OPERATOR\_\allowbreak{}QUOTIENT} & {\scriptsize \#\# E5b. Gram-null decoupling} & {\scriptsize proven} & {\scriptsize ---} \\
{\scriptsize\ttfamily DECORATED\_\allowbreak{}HISTORY\_\allowbreak{}MERGING} & {\scriptsize \#\# E5c. Exact merging of decorated histories} & {\scriptsize proven} & {\scriptsize ---} \\
{\scriptsize\ttfamily NESTED\_\allowbreak{}QUOTIENT\_\allowbreak{}REDUCTION} & {\scriptsize \#\# E5d. Finite-order nested-quotient spectral reduction} & {\scriptsize proven} & {\scriptsize ---} \\
{\scriptsize\ttfamily SHELL\_\allowbreak{}SYMMETRY\_\allowbreak{}RITZ\_\allowbreak{}CONTROL} & {\scriptsize \#\# E6a. Symmetry and residual control for a reduced shell} & {\scriptsize proven} & {\scriptsize ---} \\
{\scriptsize\ttfamily B6\_\allowbreak{}RETAINED\_\allowbreak{}KRYLOV\_\allowbreak{}CAMPAIGN} & {\scriptsize \#\# E6b. Retained B6 numerical campaign} & {\scriptsize conditional} & {\scriptsize ---} \\
{\scriptsize\ttfamily ISOLATED\_\allowbreak{}POSITIVE\_\allowbreak{}MATRIX\_\allowbreak{}ATOM} & {\scriptsize \#\# E7. Isolated positive matrix atom persists} & {\scriptsize proven} & {\scriptsize ---} \\
{\scriptsize\ttfamily CONDITIONAL\_\allowbreak{}SPECTRAL\_\allowbreak{}FLOOR} & {\scriptsize \#\# E8. Conditional spectral floor and defect contraction} & {\scriptsize proven} & {\scriptsize ---} \\
{\scriptsize\ttfamily D4\_\allowbreak{}DISJOINT\_\allowbreak{}STAPLE\_\allowbreak{}COORDINATES} & {\scriptsize \#\# E9. Six disjoint staple coordinates in four dimensions} & {\scriptsize proven} & {\scriptsize ---} \\
{\scriptsize\ttfamily POSITIVE\_\allowbreak{}SECTOR\_\allowbreak{}PHASE\_\allowbreak{}ISOLATION} & {\scriptsize \#\# E10. Positive-sector phase isolation} & {\scriptsize proven} & {\scriptsize ---} \\
{\scriptsize\ttfamily VSU\_\allowbreak{}BOUNDED\_\allowbreak{}DOMAIN\_\allowbreak{}WELLPOSEDNESS} & {\scriptsize \#\# E11. Convex VSU action on a bounded domain} & {\scriptsize proven} & {\scriptsize ---} \\
{\scriptsize\ttfamily DETERMINANT\_\allowbreak{}PARITY\_\allowbreak{}REPAIR} & {\scriptsize \#\# E12. Determinant reduction with the missing parity restored} & {\scriptsize proven} & {\scriptsize ---} \\
{\scriptsize\ttfamily MOVING\_\allowbreak{}ADJOINT\_\allowbreak{}FORCE\_\allowbreak{}DERIVATIVE} & {\scriptsize \#\# E13. A moving adjoint force has two derivative terms} & {\scriptsize proven} & {\scriptsize ---} \\
{\scriptsize\ttfamily SPHERE\_\allowbreak{}EQUAL\_\allowbreak{}COUPLING\_\allowbreak{}COVARIANCE} & {\scriptsize \#\# E14. Equal-coupling SU(3) faces on a sphere have nonnegative covariance} & {\scriptsize proven} & {\scriptsize ---} \\
{\scriptsize\ttfamily SU3\_\allowbreak{}POSITIVITY\_\allowbreak{}BOUNDARIES} & {\scriptsize \#\# E15. Exact SU(3) boundaries of the covariance argument} & {\scriptsize proven} & {\scriptsize ---} \\
{\scriptsize\ttfamily REFLECTION\_\allowbreak{}ADAPTED\_\allowbreak{}BLOCKING} & {\scriptsize \#\# E16. Reflection-adapted deterministic blocking preserves positivity} & {\scriptsize proven} & {\scriptsize ---} \\
{\scriptsize\ttfamily CORNER\_\allowbreak{}BLOCKING\_\allowbreak{}REFLECTION\_\allowbreak{}DEFECT} & {\scriptsize \#\# E17. Literal corner paths fail the required reflection identity} & {\scriptsize proven} & {\scriptsize ---} \\
{\scriptsize\ttfamily DIPOLAR\_\allowbreak{}REMAINDER} & {\scriptsize \#\# E18. Explicit remaining application obligations} & {\scriptsize open} & {\scriptsize ---} \\
{\scriptsize\ttfamily PHYSICAL\_\allowbreak{}CARRIER\_\allowbreak{}REALIZATION} & {\scriptsize \#\# E18. Explicit remaining application obligations} & {\scriptsize open} & {\scriptsize ---} \\
{\scriptsize\ttfamily STAPLE\_\allowbreak{}TRANSVERSALITY\_\allowbreak{}REPAIR} & {\scriptsize \#\# E18. Explicit remaining application obligations} & {\scriptsize open} & {\scriptsize ---} \\
{\scriptsize\ttfamily VSU\_\allowbreak{}WHOLE\_\allowbreak{}SPACE} & {\scriptsize \#\# E18. Explicit remaining application obligations} & {\scriptsize open} & {\scriptsize ---} \\
{\scriptsize\ttfamily OMITTED\_\allowbreak{}DETERMINANT\_\allowbreak{}PARITY} & {\scriptsize \#\# E12. Determinant reduction with the missing parity restored} & {\scriptsize falsified} & {\scriptsize ---} \\
{\scriptsize\ttfamily OMITTED\_\allowbreak{}OWN\_\allowbreak{}FORCE\_\allowbreak{}DERIVATIVE} & {\scriptsize \#\# E13. A moving adjoint force has two derivative terms} & {\scriptsize falsified} & {\scriptsize ---} \\
\end{longtable}

\subsection*{CITE:\allowbreak{}HODGE\_\allowbreak{}POLYMER\_\allowbreak{}SUPPORTED}
\addcontentsline{toc}{subsection}{HODGE\_\allowbreak{}POLYMER\_\allowbreak{}SUPPORTED}
\noindent\srcpath{docs/derivations/hodge-polymer-supported-statements.md}\\
{\footnotesize SHA-256 \texttt{e0221c892ff185e7}\ldots}\par\smallskip
\begin{longtable}{@{}p{0.30\linewidth} p{0.30\linewidth} l l@{}}
\toprule Statement & Locator & Status & Formal \\ \midrule
\endfirsthead
\toprule Statement & Locator & Status & Formal \\ \midrule \endhead
\bottomrule\endfoot
{\scriptsize\ttfamily H1} & {\scriptsize \#\# H1: complementary kernels from Hodge relations (from line 17)} & {\scriptsize proven} & {\scriptsize full} \\
{\scriptsize\ttfamily H2} & {\scriptsize \#\# H2: carrier action and absence of Hodge coupling (from line 33)} & {\scriptsize proven} & {\scriptsize full} \\
{\scriptsize\ttfamily H3} & {\scriptsize \#\# H3: rank-one normalization, excitation and factorization (from line 45)} & {\scriptsize proven} & {\scriptsize full} \\
{\scriptsize\ttfamily H4} & {\scriptsize \#\# H4: finite word patterns and supplied polynomial identities (from line 71)} & {\scriptsize proven} & {\scriptsize full} \\
{\scriptsize\ttfamily P1} & {\scriptsize \#\# P1: geometric-series majorant (from line 93)} & {\scriptsize proven} & {\scriptsize full} \\
{\scriptsize\ttfamily P2} & {\scriptsize \#\# P2: parameter bounds and the separate tree budget (from line 107)} & {\scriptsize proven} & {\scriptsize full} \\
\end{longtable}

\subsection*{CITE:\allowbreak{}LOCAL\_\allowbreak{}CLASS\_\allowbreak{}WICK\_\allowbreak{}SPECTRUM}
\addcontentsline{toc}{subsection}{LOCAL\_\allowbreak{}CLASS\_\allowbreak{}WICK\_\allowbreak{}SPECTRUM}
\noindent\srcpath{docs/derivations/local-class-wick-spectrum.md}\\
{\footnotesize SHA-256 \texttt{5cde9ee4fe88bd47}\ldots}\par\smallskip
\begin{longtable}{@{}p{0.30\linewidth} p{0.30\linewidth} l l@{}}
\toprule Statement & Locator & Status & Formal \\ \midrule
\endfirsthead
\toprule Statement & Locator & Status & Formal \\ \midrule \endhead
\bottomrule\endfoot
{\scriptsize\ttfamily RECURRENCE} & {\scriptsize \#\# 3. All-order constructive theorem} & {\scriptsize proven} & {\scriptsize ---} \\
{\scriptsize\ttfamily COEFFICIENTS} & {\scriptsize \#\# 4. Recovered coefficients and the next odd term} & {\scriptsize proven} & {\scriptsize ---} \\
{\scriptsize\ttfamily SIGNS} & {\scriptsize \#\# 4.2 Sign corollary} & {\scriptsize proven} & {\scriptsize ---} \\
\end{longtable}

\subsection*{CITE:\allowbreak{}MOVING\_\allowbreak{}TIME\_\allowbreak{}SPECTRAL\_\allowbreak{}GAP}
\addcontentsline{toc}{subsection}{MOVING\_\allowbreak{}TIME\_\allowbreak{}SPECTRAL\_\allowbreak{}GAP}
\noindent\srcpath{docs/derivations/moving-time-spectral-gap.md}\\
{\footnotesize SHA-256 \texttt{d0a20c25e023bff7}\ldots}\par\smallskip
\begin{longtable}{@{}p{0.30\linewidth} p{0.30\linewidth} l l@{}}
\toprule Statement & Locator & Status & Formal \\ \midrule
\endfirsthead
\toprule Statement & Locator & Status & Formal \\ \midrule \endhead
\bottomrule\endfoot
{\scriptsize\ttfamily PROJECTION} & {\scriptsize \#\# 1. Premises and the finite-cutoff projection estimate (MT1) (lines 17-43)} & {\scriptsize proven} & {\scriptsize ---} \\
{\scriptsize\ttfamily VAGUE\_\allowbreak{}LIMIT} & {\scriptsize \#\# 2. Moving-time exclusion and totality (MT2) (lines 44-91)} & {\scriptsize proven} & {\scriptsize ---} \\
{\scriptsize\ttfamily SHRINKING\_\allowbreak{}TILTS} & {\scriptsize \#\# 3. Shrinking exponential sources without a Taylor rate penalty (MT3) (lines 92-126)} & {\scriptsize proven} & {\scriptsize ---} \\
{\scriptsize\ttfamily OPTIMAL\_\allowbreak{}BUDGET} & {\scriptsize \#\# 4. Optimal power-law budget and finite observation horizon (MT4) (lines 127-179)} & {\scriptsize proven} & {\scriptsize ---} \\
{\scriptsize\ttfamily SHARPNESS} & {\scriptsize \#\# 5. Sharpness and falsifiers (MT5) (lines 180-234)} & {\scriptsize proven} & {\scriptsize ---} \\
\end{longtable}

\subsection*{CITE:\allowbreak{}OS\_\allowbreak{}KERNEL\_\allowbreak{}GAP}
\addcontentsline{toc}{subsection}{OS\_\allowbreak{}KERNEL\_\allowbreak{}GAP}
\noindent\srcpath{docs/derivations/os-kernel-moving-time-gap.md}\\
{\footnotesize SHA-256 \texttt{d7088d36f67a1f8c}\ldots}\par\smallskip
\begin{longtable}{@{}p{0.30\linewidth} p{0.30\linewidth} l l@{}}
\toprule Statement & Locator & Status & Formal \\ \midrule
\endfirsthead
\toprule Statement & Locator & Status & Formal \\ \midrule \endhead
\bottomrule\endfoot
{\scriptsize\ttfamily KERNEL} & {\scriptsize \#\# K1. Construct the physical semigroup from compatible limiting kernels (lines 10-61)} & {\scriptsize proven} & {\scriptsize ---} \\
{\scriptsize\ttfamily LATE\_\allowbreak{}TIME} & {\scriptsize \#\# K2. A single late-time bound controls every limiting fixed time (lines 62-110)} & {\scriptsize proven} & {\scriptsize ---} \\
{\scriptsize\ttfamily WEIGHT} & {\scriptsize \#\# K3. Nontrivial finite-energy weight from one separated-time lower bound (lines 111-132)} & {\scriptsize proven} & {\scriptsize ---} \\
{\scriptsize\ttfamily LIMITATIONS} & {\scriptsize \#\# K4. Falsifiers and actual Wilson application check (lines 133-175)} & {\scriptsize proven} & {\scriptsize ---} \\
{\scriptsize\ttfamily POSITIVE\_\allowbreak{}HISTORY} & {\scriptsize \#\# K5. Positive-time history separation supplies zero-time continuity (lines 176-221)} & {\scriptsize proven} & {\scriptsize ---} \\
\end{longtable}

\subsection*{CITE:\allowbreak{}PBH\_\allowbreak{}WILSON\_\allowbreak{}TEST}
\addcontentsline{toc}{subsection}{PBH\_\allowbreak{}WILSON\_\allowbreak{}TEST}
\noindent\srcpath{docs/derivations/yangmills-pbh-wilson-test.md}\\
{\footnotesize SHA-256 \texttt{3122ef912820b5a0}\ldots}\par\smallskip
\begin{longtable}{@{}p{0.30\linewidth} p{0.30\linewidth} l l@{}}
\toprule Statement & Locator & Status & Formal \\ \midrule
\endfirsthead
\toprule Statement & Locator & Status & Formal \\ \midrule \endhead
\bottomrule\endfoot
{\scriptsize\ttfamily PBH1\_\allowbreak{}PBH3} & {\scriptsize \#\# PBH-1. The actual measure and horizontal form (lines 8-57)} & {\scriptsize proven} & {\scriptsize ---} \\
{\scriptsize\ttfamily PBH4} & {\scriptsize \#\# PBH-2. Explicit regular-point counterexample (lines 58-113)} & {\scriptsize proven} & {\scriptsize ---} \\
{\scriptsize\ttfamily PBH5} & {\scriptsize \#\# PBH-3. The obstruction is not confined to a one-plaquette lattice (lines 114-163)} & {\scriptsize proven} & {\scriptsize ---} \\
{\scriptsize\ttfamily PBH6} & {\scriptsize \#\# PBH-4. Classical Wilson weight is not the quantum ground-state weight (lines 164-198)} & {\scriptsize proven} & {\scriptsize ---} \\
{\scriptsize\ttfamily FLOW\_\allowbreak{}COUNTEREXAMPLE} & {\scriptsize \#\# PBH-5. The stated gradient-flow shortcut also fails an exact test (lines 199-230)} & {\scriptsize proven} & {\scriptsize ---} \\
\end{longtable}

\subsection*{CITE:\allowbreak{}REVIEW\_\allowbreak{}SPECTRAL\_\allowbreak{}BUDGETS}
\addcontentsline{toc}{subsection}{REVIEW\_\allowbreak{}SPECTRAL\_\allowbreak{}BUDGETS}
\noindent\srcpath{docs/derivations/review-derived-spectral-budgets.md}\\
{\footnotesize SHA-256 \texttt{96037da9bfc02943}\ldots}\par\smallskip
\begin{longtable}{@{}p{0.30\linewidth} p{0.30\linewidth} l l@{}}
\toprule Statement & Locator & Status & Formal \\ \midrule
\endfirsthead
\toprule Statement & Locator & Status & Formal \\ \midrule \endhead
\bottomrule\endfoot
{\scriptsize\ttfamily Q1} & {\scriptsize \#\# Q1. The assembled fourth-order polynomial is coercive (lines 11-42)} & {\scriptsize proven} & {\scriptsize ---} \\
{\scriptsize\ttfamily M1} & {\scriptsize \#\# M1. Several decay channels and time-amplified remainders (lines 43-89)} & {\scriptsize proven} & {\scriptsize ---} \\
{\scriptsize\ttfamily M2} & {\scriptsize \#\# M2. Closed optimal rate and a witnessing observation time (lines 90-125)} & {\scriptsize proven} & {\scriptsize ---} \\
{\scriptsize\ttfamily M3} & {\scriptsize \#\# M3. Sharpness, failure cases and physical scope (lines 126-171)} & {\scriptsize proven} & {\scriptsize ---} \\
\end{longtable}

\subsection*{CITE:\allowbreak{}SEPTEMBER10\_\allowbreak{}UPSTREAM\_\allowbreak{}SCALARS}
\addcontentsline{toc}{subsection}{SEPTEMBER10\_\allowbreak{}UPSTREAM\_\allowbreak{}SCALARS}
\noindent\srcpath{docs/derivations/september10-upstream-scalar-inventory.md}\\
{\footnotesize SHA-256 \texttt{75a6b78ebaed289f}\ldots}\par\smallskip
\begin{longtable}{@{}p{0.30\linewidth} p{0.30\linewidth} l l@{}}
\toprule Statement & Locator & Status & Formal \\ \midrule
\endfirsthead
\toprule Statement & Locator & Status & Formal \\ \midrule \endhead
\bottomrule\endfoot
{\scriptsize\ttfamily KP\_\allowbreak{}SCALARS} & {\scriptsize \#\# U1 Tilted KP scalar ingredients} & {\scriptsize proven} & {\scriptsize partial} \\
{\scriptsize\ttfamily COERCIVITY\_\allowbreak{}SCALARS} & {\scriptsize \#\# U2 Coercivity scalar ingredients} & {\scriptsize proven} & {\scriptsize partial} \\
\end{longtable}

\subsection*{CITE:\allowbreak{}SEPT\_\allowbreak{}DOC\_\allowbreak{}YANGMILLS\_\allowbreak{}FORMAL\_\allowbreak{}BRIDGES}
\addcontentsline{toc}{subsection}{SEPT\_\allowbreak{}DOC\_\allowbreak{}YANGMILLS\_\allowbreak{}FORMAL\_\allowbreak{}BRIDGES}
\noindent\srcpath{docs/derivations/yangmills-formal-bridges.md}\\
{\footnotesize SHA-256 \texttt{bb3245e82e5dc857}\ldots}\par\smallskip
\begin{longtable}{@{}p{0.30\linewidth} p{0.30\linewidth} l l@{}}
\toprule Statement & Locator & Status & Formal \\ \midrule
\endfirsthead
\toprule Statement & Locator & Status & Formal \\ \midrule \endhead
\bottomrule\endfoot
{\scriptsize\ttfamily F1} & {\scriptsize \#\# Exact finite ground-state identification (lines 14-78)} & {\scriptsize proven} & {\scriptsize full} \\
{\scriptsize\ttfamily F2} & {\scriptsize \#\# Exact finite ground-state identification (lines 14-78)} & {\scriptsize proven} & {\scriptsize full} \\
{\scriptsize\ttfamily F3} & {\scriptsize \#\# Exact finite ground-state identification (lines 14-78)} & {\scriptsize proven} & {\scriptsize full} \\
{\scriptsize\ttfamily F4} & {\scriptsize \#\# A trial residual needs a floor and a mean (lines 79-101)} & {\scriptsize proven} & {\scriptsize full} \\
{\scriptsize\ttfamily F5} & {\scriptsize \#\# The pure quartic Gaussian certificate (lines 102-137)} & {\scriptsize proven} & {\scriptsize full} \\
{\scriptsize\ttfamily F6} & {\scriptsize \#\# The pure quartic Gaussian certificate (lines 102-137)} & {\scriptsize proven} & {\scriptsize full} \\
{\scriptsize\ttfamily F7} & {\scriptsize \#\# Algebra for Simon's transverse mechanism (lines 138-169)} & {\scriptsize proven} & {\scriptsize full} \\
{\scriptsize\ttfamily SLICE\_\allowbreak{}RADIAL\_\allowbreak{}ALGEBRA} & {\scriptsize \#\# Algebra for Simon's transverse mechanism (lines 138-169)} & {\scriptsize proven} & {\scriptsize ---} \\
\end{longtable}

\subsection*{CITE:\allowbreak{}SEPT\_\allowbreak{}DOC\_\allowbreak{}YANGMILLS\_\allowbreak{}RESOLVENT\_\allowbreak{}LOCALIZATION\_\allowbreak{}AND\_\allowbreak{}DIRICHLET\_\allowbreak{}GAP}
\addcontentsline{toc}{subsection}{SEPT\_\allowbreak{}DOC\_\allowbreak{}YANGMILLS\_\allowbreak{}RESOLVENT\_\allowbreak{}LOCALIZATION\_\allowbreak{}AND\_\allowbreak{}DIRICHLET\_\allowbreak{}GAP}
\noindent\srcpath{docs/derivations/yangmills-resolvent-localization-and-dirichlet-gap.md}\\
{\footnotesize SHA-256 \texttt{876b76c1aee3eb8e}\ldots}\par\smallskip
\begin{longtable}{@{}p{0.30\linewidth} p{0.30\linewidth} l l@{}}
\toprule Statement & Locator & Status & Formal \\ \midrule
\endfirsthead
\toprule Statement & Locator & Status & Formal \\ \midrule \endhead
\bottomrule\endfoot
{\scriptsize\ttfamily SUBMITTED\_\allowbreak{}W6\_\allowbreak{}CLOSURE} & {\scriptsize \#\# 1. Summary of Results (lines 20-39)} & {\scriptsize disputed} & {\scriptsize ---} \\
{\scriptsize\ttfamily RESOLVENT\_\allowbreak{}IDENTITY} & {\scriptsize \#\#\# 2.1 The Resolvent Identity (lines 42-54)} & {\scriptsize conditional} & {\scriptsize partial} \\
{\scriptsize\ttfamily LOCAL\_\allowbreak{}KERNEL\_\allowbreak{}IMPLICATION} & {\scriptsize \#\#\# 2.3 Volume-Uniformity of the Pairing (lines 62-80)} & {\scriptsize conditional} & {\scriptsize partial} \\
{\scriptsize\ttfamily GROUND\_\allowbreak{}TRANSFORM} & {\scriptsize \#\#\# 3.2 True Ground Conjugation and Exact Cancellation (lines 104-123)} & {\scriptsize proven} & {\scriptsize ---} \\
{\scriptsize\ttfamily SUBMITTED\_\allowbreak{}CLOSURES} & {\scriptsize \#\# 4. Status and Graph Integration (lines 131-135)} & {\scriptsize superseded} & {\scriptsize ---} \\
\end{longtable}

\subsection*{CITE:\allowbreak{}SOURCE\_\allowbreak{}CURRENT\_\allowbreak{}SPECTRAL\_\allowbreak{}BRIDGES}
\addcontentsline{toc}{subsection}{SOURCE\_\allowbreak{}CURRENT\_\allowbreak{}SPECTRAL\_\allowbreak{}BRIDGES}
\noindent\srcpath{docs/derivations/source-currents-and-spectral-totality.md}\\
{\footnotesize SHA-256 \texttt{860c81a3fb665d6b}\ldots}\par\smallskip
\begin{longtable}{@{}p{0.30\linewidth} p{0.30\linewidth} l l@{}}
\toprule Statement & Locator & Status & Formal \\ \midrule
\endfirsthead
\toprule Statement & Locator & Status & Formal \\ \midrule \endhead
\bottomrule\endfoot
{\scriptsize\ttfamily SCB0A} & {\scriptsize \#\# SCB0A. Degenerate energy variational certificate (lines 21-34)} & {\scriptsize proven} & {\scriptsize partial} \\
{\scriptsize\ttfamily SCB0B} & {\scriptsize \#\# SCB0B. Bounded-inverse W6 composition (lines 35-59)} & {\scriptsize proven} & {\scriptsize partial} \\
{\scriptsize\ttfamily SCB1} & {\scriptsize \#\# SCB1. The complete first current in the actual square convention (lines 60-114)} & {\scriptsize proven} & {\scriptsize ---} \\
{\scriptsize\ttfamily SCB2} & {\scriptsize \#\# SCB2. Exact current moments and an all-source constant (lines 115-172)} & {\scriptsize proven} & {\scriptsize partial} \\
{\scriptsize\ttfamily SCB3} & {\scriptsize \#\# SCB3. A selected-inverse implication without a separate diagonal bound (lines 173-197)} & {\scriptsize proven} & {\scriptsize partial} \\
{\scriptsize\ttfamily SCB4} & {\scriptsize \#\# SCB4. Intrinsic score currents and exact centering corrections (lines 198-257)} & {\scriptsize proven} & {\scriptsize ---} \\
{\scriptsize\ttfamily SCB5} & {\scriptsize \#\# SCB5. Tiny exponential sources imply totality by differentiation (lines 258-298)} & {\scriptsize proven} & {\scriptsize partial} \\
{\scriptsize\ttfamily SCB6} & {\scriptsize \#\# SCB6. Uniform normalized mesoscopic source amplitudes (lines 299-316)} & {\scriptsize proven} & {\scriptsize ---} \\
\end{longtable}

\subsection*{CITE:\allowbreak{}THEORY\_\allowbreak{}GEOMETRY\_\allowbreak{}RICCATI\_\allowbreak{}BRIDGES}
\addcontentsline{toc}{subsection}{THEORY\_\allowbreak{}GEOMETRY\_\allowbreak{}RICCATI\_\allowbreak{}BRIDGES}
\noindent\srcpath{docs/derivations/theory-geometry-and-riccati-bridges.md}\\
{\footnotesize SHA-256 \texttt{7216dd8c2f263782}\ldots}\par\smallskip
\begin{longtable}{@{}p{0.30\linewidth} p{0.30\linewidth} l l@{}}
\toprule Statement & Locator & Status & Formal \\ \midrule
\endfirsthead
\toprule Statement & Locator & Status & Formal \\ \midrule \endhead
\bottomrule\endfoot
{\scriptsize\ttfamily TG1} & {\scriptsize \#\# TG1: The sharp three-coordinate anisotropy constant (lines 14-60)} & {\scriptsize proven} & {\scriptsize partial} \\
{\scriptsize\ttfamily TG2} & {\scriptsize \#\# TG2: Zero extension and integrated energy bounds (lines 61-93)} & {\scriptsize proven} & {\scriptsize ---} \\
{\scriptsize\ttfamily TG3} & {\scriptsize \#\# TG3: One secular cubic for the recorded Hodge support (lines 94-136)} & {\scriptsize proven} & {\scriptsize partial} \\
{\scriptsize\ttfamily TG4} & {\scriptsize \#\# TG4: The induced eighth-order moment in the same truncation (lines 137-166)} & {\scriptsize proven} & {\scriptsize partial} \\
{\scriptsize\ttfamily TG5} & {\scriptsize \#\# TG5: Complete source motion is a Schur congruence (lines 167-222)} & {\scriptsize proven} & {\scriptsize partial} \\
{\scriptsize\ttfamily TG6} & {\scriptsize \#\# TG6: The metric and finite source resolvent must travel with the form (lines 223-274)} & {\scriptsize proven} & {\scriptsize partial} \\
{\scriptsize\ttfamily TG7} & {\scriptsize \#\# TG7: Scalar square-source parity eliminates formal odd Schur jets (lines 275-313)} & {\scriptsize proven} & {\scriptsize ---} \\
{\scriptsize\ttfamily TG8} & {\scriptsize \#\# TG8: A residual certificate for the abstract Riccati fixed point (lines 314-350)} & {\scriptsize proven} & {\scriptsize partial} \\
\end{longtable}

\subsection*{CITE:\allowbreak{}TRACK\_\allowbreak{}A\_\allowbreak{}SUPPORTED}
\addcontentsline{toc}{subsection}{TRACK\_\allowbreak{}A\_\allowbreak{}SUPPORTED}
\noindent\srcpath{docs/derivations/track-a-supported-statements.md}\\
{\footnotesize SHA-256 \texttt{26efd60d345eb4c9}\ldots}\par\smallskip
\begin{longtable}{@{}p{0.30\linewidth} p{0.30\linewidth} l l@{}}
\toprule Statement & Locator & Status & Formal \\ \midrule
\endfirsthead
\toprule Statement & Locator & Status & Formal \\ \midrule \endhead
\bottomrule\endfoot
{\scriptsize\ttfamily W6\_\allowbreak{}V} & {\scriptsize \#\# W6-V: actual variational residual identity and dual energy (from line 18)} & {\scriptsize proven} & {\scriptsize full} \\
{\scriptsize\ttfamily W6\_\allowbreak{}F} & {\scriptsize \#\# W6-F: bounded-operator factorization controls the residual functional (from line 43)} & {\scriptsize proven} & {\scriptsize full} \\
{\scriptsize\ttfamily W6\_\allowbreak{}E} & {\scriptsize \#\# W6-E: sufficient selected-inverse error certificate (from line 61)} & {\scriptsize proven} & {\scriptsize full} \\
{\scriptsize\ttfamily W6\_\allowbreak{}C} & {\scriptsize \#\# W6-C: exact Gaussian coefficient substitutions (from line 76)} & {\scriptsize proven} & {\scriptsize full} \\
{\scriptsize\ttfamily SC17\_\allowbreak{}R} & {\scriptsize \#\# SC17-R: Riccati roots and the continuous barrier (from line 92)} & {\scriptsize proven} & {\scriptsize full} \\
{\scriptsize\ttfamily SC17\_\allowbreak{}B} & {\scriptsize \#\# SC17-B: constructed noncommutative Banach-algebra fixed point (from line 112)} & {\scriptsize proven} & {\scriptsize full} \\
{\scriptsize\ttfamily SC17\_\allowbreak{}P} & {\scriptsize \#\# SC17-P: exact interval arithmetic under the proposed inputs (from line 132)} & {\scriptsize proven} & {\scriptsize full} \\
{\scriptsize\ttfamily G17\_\allowbreak{}P} & {\scriptsize \#\# G17-P: corrected raw partition estimate for actual probability measures (from line 151)} & {\scriptsize proven} & {\scriptsize full} \\
{\scriptsize\ttfamily G17\_\allowbreak{}V} & {\scriptsize \#\# G17-V: centered L2 source, variance, and bounded footprints (from line 168)} & {\scriptsize proven} & {\scriptsize full} \\
{\scriptsize\ttfamily G17\_\allowbreak{}O} & {\scriptsize \#\# G17-O: strict obstruction to a constant equal to one (from line 185)} & {\scriptsize proven} & {\scriptsize full} \\
{\scriptsize\ttfamily G17\_\allowbreak{}I} & {\scriptsize \#\# G17-I: independent product sources have unbounded radius (from line 206)} & {\scriptsize proven} & {\scriptsize full} \\
\end{longtable}

\subsection*{CITE:\allowbreak{}UNIVERSAL\_\allowbreak{}CELLULAR\_\allowbreak{}HODGE\_\allowbreak{}TETRAHEDRAL}
\addcontentsline{toc}{subsection}{UNIVERSAL\_\allowbreak{}CELLULAR\_\allowbreak{}HODGE\_\allowbreak{}TETRAHEDRAL}
\noindent\srcpath{docs/derivations/universal-cellular-hodge-tetrahedral.md}\\
{\footnotesize SHA-256 \texttt{427974f8eda3ae77}\ldots}\par\smallskip
\begin{longtable}{@{}p{0.30\linewidth} p{0.30\linewidth} l l@{}}
\toprule Statement & Locator & Status & Formal \\ \midrule
\endfirsthead
\toprule Statement & Locator & Status & Formal \\ \midrule \endhead
\bottomrule\endfoot
{\scriptsize\ttfamily SPECTRUM\_\allowbreak{}EIGENVALUE} & {\scriptsize \#\#\# 2.2 Universal spectral identities (lines 45-58)} & {\scriptsize proven} & {\scriptsize partial} \\
{\scriptsize\ttfamily UP\_\allowbreak{}HARMONIC\_\allowbreak{}EXCURSION} & {\scriptsize \#\#\# 2.3 Universal Up-Harmonicity of Excursions (lines 59-67)} & {\scriptsize proven} & {\scriptsize full} \\
{\scriptsize\ttfamily TOTAL\_\allowbreak{}LAPLACIAN\_\allowbreak{}DIAGONAL} & {\scriptsize \#\#\# 2.4 Total Laplacian Diagonal \& Regularity Criterion (lines 68-83)} & {\scriptsize proven} & {\scriptsize ---} \\
{\scriptsize\ttfamily TETRAHEDRAL\_\allowbreak{}S4\_\allowbreak{}COMMUTANT} & {\scriptsize \#\#\# 3.1 Tetrahedral Hodge duality and 3.2 The S4 commutant (lines 86-119)} & {\scriptsize proven} & {\scriptsize partial} \\
{\scriptsize\ttfamily PROPER\_\allowbreak{}RETURN\_\allowbreak{}Q\_\allowbreak{}ANNIHILATION} & {\scriptsize \#\#\# 3.4 Open physical-history identification (lines 133-156)} & {\scriptsize open} & {\scriptsize partial} \\
{\scriptsize\ttfamily RETAINED\_\allowbreak{}VECTOR\_\allowbreak{}RANK\_\allowbreak{}ONE} & {\scriptsize \#\#\# 3.3 Feshbach Intermediate Annihilation of Retained Vectors (lines 120-132)} & {\scriptsize proven} & {\scriptsize full} \\
{\scriptsize\ttfamily RSMR\_\allowbreak{}MASTER\_\allowbreak{}THEOREM} & {\scriptsize \#\#\# 4.3 The carrier symbol (lines 196-221)} & {\scriptsize proven} & {\scriptsize partial} \\
{\scriptsize\ttfamily SHIFTED\_\allowbreak{}OPERATOR\_\allowbreak{}POWERS} & {\scriptsize \#\#\# 4.2 Polynomial recurrence and operator powers (lines 172-195)} & {\scriptsize proven} & {\scriptsize full} \\
\end{longtable}

\subsection*{CITE:\allowbreak{}W6\_\allowbreak{}ANTIPODAL\_\allowbreak{}MAGNETIC\_\allowbreak{}GEOMETRY}
\addcontentsline{toc}{subsection}{W6\_\allowbreak{}ANTIPODAL\_\allowbreak{}MAGNETIC\_\allowbreak{}GEOMETRY}
\noindent\srcpath{docs/derivations/w6-antipodal-magnetic-geometry.md}\\
{\footnotesize SHA-256 \texttt{5b3e0b2312f11ee8}\ldots}\par\smallskip
\begin{longtable}{@{}p{0.30\linewidth} p{0.30\linewidth} l l@{}}
\toprule Statement & Locator & Status & Formal \\ \midrule
\endfirsthead
\toprule Statement & Locator & Status & Formal \\ \midrule \endhead
\bottomrule\endfoot
{\scriptsize\ttfamily HESSIAN\_\allowbreak{}SPECTRUM} & {\scriptsize A1-A6; exact fixed-Q chart, quadratic form and full spectrum} & {\scriptsize proven} & {\scriptsize ---} \\
{\scriptsize\ttfamily GAUGE\_\allowbreak{}NORMAL\_\allowbreak{}COERCIVITY} & {\scriptsize A7-A11; complete minimum sphere, tangent kernel and original electric metric} & {\scriptsize proven} & {\scriptsize ---} \\
{\scriptsize\ttfamily ANGULAR\_\allowbreak{}REDUCTION} & {\scriptsize A12-A13; exact old-orbit restriction and compact normal-bundle implicit function} & {\scriptsize proven} & {\scriptsize ---} \\
{\scriptsize\ttfamily GAUGE\_\allowbreak{}SCORE\_\allowbreak{}INVARIANCE} & {\scriptsize A14-A15; actual ground and Haar half-density symmetry} & {\scriptsize proven} & {\scriptsize ---} \\
\end{longtable}

\subsection*{CITE:\allowbreak{}W6\_\allowbreak{}CONDITIONAL\_\allowbreak{}SCORE\_\allowbreak{}TAIL\_\allowbreak{}CONTROL}
\addcontentsline{toc}{subsection}{W6\_\allowbreak{}CONDITIONAL\_\allowbreak{}SCORE\_\allowbreak{}TAIL\_\allowbreak{}CONTROL}
\noindent\srcpath{docs/derivations/w6-conditional-score-tail-control.md}\\
{\footnotesize SHA-256 \texttt{d04ae01d29d33cc5}\ldots}\par\smallskip
\begin{longtable}{@{}p{0.30\linewidth} p{0.30\linewidth} l l@{}}
\toprule Statement & Locator & Status & Formal \\ \midrule
\endfirsthead
\toprule Statement & Locator & Status & Formal \\ \midrule \endhead
\bottomrule\endfoot
{\scriptsize\ttfamily INTERVAL\_\allowbreak{}OBSTRUCTION\_\allowbreak{}REPAIR} & {\scriptsize \#\# 1. What the actual compact endpoints do prove (lines 67-284)} & {\scriptsize proven} & {\scriptsize ---} \\
{\scriptsize\ttfamily SCORE\_\allowbreak{}DOMINATION\_\allowbreak{}M10} & {\scriptsize \#\#\# M.4. The precise remaining conditional-score inequality (lines 473-495)} & {\scriptsize open} & {\scriptsize ---} \\
{\scriptsize\ttfamily ACTUAL\_\allowbreak{}SOURCE\_\allowbreak{}MOMENT} & {\scriptsize \#\# 6. Uniform actual source potential moment and sharper score criterion (lines 325-586)} & {\scriptsize proven} & {\scriptsize ---} \\
{\scriptsize\ttfamily SCORE\_\allowbreak{}WEIGHT\_\allowbreak{}FROM\_\allowbreak{}M10} & {\scriptsize \#\#\# M.4. The precise remaining conditional-score inequality and M.5 (lines 497-563)} & {\scriptsize conditional} & {\scriptsize ---} \\
\end{longtable}

\subsection*{CITE:\allowbreak{}W6\_\allowbreak{}CONDITIONAL\_\allowbreak{}TRANSPORT\_\allowbreak{}OBSTRUCTION}
\addcontentsline{toc}{subsection}{W6\_\allowbreak{}CONDITIONAL\_\allowbreak{}TRANSPORT\_\allowbreak{}OBSTRUCTION}
\noindent\srcpath{docs/derivations/w6-conditional-transport-obstruction.md}\\
{\footnotesize SHA-256 \texttt{33f5cefc2105ec3e}\ldots}\par\smallskip
\begin{longtable}{@{}p{0.30\linewidth} p{0.30\linewidth} l l@{}}
\toprule Statement & Locator & Status & Formal \\ \midrule
\endfirsthead
\toprule Statement & Locator & Status & Formal \\ \midrule \endhead
\bottomrule\endfoot
{\scriptsize\ttfamily SCORE\_\allowbreak{}CONTINUITY} & {\scriptsize \#\# S1 Exact conditional score (lines 15-72)} & {\scriptsize proven} & {\scriptsize ---} \\
{\scriptsize\ttfamily AGMON\_\allowbreak{}GEOMETRY} & {\scriptsize \#\# S2 Exact constrained geometry (lines 73-140)} & {\scriptsize proven} & {\scriptsize ---} \\
{\scriptsize\ttfamily CONDITIONAL\_\allowbreak{}CONCENTRATION} & {\scriptsize \#\# S3 Actual conditional concentration (lines 141-213)} & {\scriptsize proven} & {\scriptsize ---} \\
{\scriptsize\ttfamily ARBITRARY\_\allowbreak{}Q8\_\allowbreak{}OBSTRUCTION} & {\scriptsize \#\# S4 A transport obstruction (lines 214-284)} & {\scriptsize proven} & {\scriptsize ---} \\
{\scriptsize\ttfamily SYNCHRONIZED\_\allowbreak{}TANGENCY} & {\scriptsize \#\# S5 Synchronized tangency (lines 285-316)} & {\scriptsize proven} & {\scriptsize ---} \\
{\scriptsize\ttfamily LOCAL\_\allowbreak{}GLOBAL\_\allowbreak{}SCORE\_\allowbreak{}CRITERION} & {\scriptsize \#\# S6 Remaining all-fiber estimate (lines 317-375)} & {\scriptsize conditional} & {\scriptsize ---} \\
{\scriptsize\ttfamily M10\_\allowbreak{}SYNCHRONIZED\_\allowbreak{}SCORE} & {\scriptsize \#\# S6 Remaining all-fiber estimate (lines 317-375), explicitly open synchronized M10 target} & {\scriptsize open} & {\scriptsize ---} \\
\end{longtable}

\subsection*{CITE:\allowbreak{}W6\_\allowbreak{}GROUND\_\allowbreak{}JETS\_\allowbreak{}TRANSPORT\_\allowbreak{}BUDGET}
\addcontentsline{toc}{subsection}{W6\_\allowbreak{}GROUND\_\allowbreak{}JETS\_\allowbreak{}TRANSPORT\_\allowbreak{}BUDGET}
\noindent\srcpath{docs/derivations/w6-ground-jets-and-transport-budget.md}\\
{\footnotesize SHA-256 \texttt{a1638b09ba941b0c}\ldots}\par\smallskip
\begin{longtable}{@{}p{0.30\linewidth} p{0.30\linewidth} l l@{}}
\toprule Statement & Locator & Status & Formal \\ \midrule
\endfirsthead
\toprule Statement & Locator & Status & Formal \\ \midrule \endhead
\bottomrule\endfoot
{\scriptsize\ttfamily GROUND\_\allowbreak{}JETS} & {\scriptsize \#\# 2. Raw electric and magnetic derivatives already have the required scaling (lines 51-198)} & {\scriptsize proven} & {\scriptsize ---} \\
{\scriptsize\ttfamily SOURCE\_\allowbreak{}ENERGY\_\allowbreak{}JETS\_\allowbreak{}R10} & {\scriptsize \#\# 4. An explicit bridge from the actual transport generator to M1,M2,M3 (lines 199-218); section 5 (lines 271-309)} & {\scriptsize open} & {\scriptsize ---} \\
{\scriptsize\ttfamily COMPLETE\_\allowbreak{}BUDGET} & {\scriptsize \#\# 4. An explicit bridge from the actual transport generator to M1,M2,M3 (lines 199-367)} & {\scriptsize conditional} & {\scriptsize ---} \\
{\scriptsize\ttfamily INTERACTING\_\allowbreak{}GRID\_\allowbreak{}COMPARISON} & {\scriptsize \#\# 6. Source energies, subdivision, interacting grids, and the physical clock (lines 351-366)} & {\scriptsize open} & {\scriptsize ---} \\
\end{longtable}

\subsection*{CITE:\allowbreak{}W6\_\allowbreak{}SOURCE\_\allowbreak{}ENERGY\_\allowbreak{}JETS\_\allowbreak{}R10}
\addcontentsline{toc}{subsection}{W6\_\allowbreak{}SOURCE\_\allowbreak{}ENERGY\_\allowbreak{}JETS\_\allowbreak{}R10}
\noindent\srcpath{docs/derivations/w6-source-energy-jets-r10.md}\\
{\footnotesize SHA-256 \texttt{7cb2ad710a35e1dd}\ldots}\par\smallskip
\begin{longtable}{@{}p{0.30\linewidth} p{0.30\linewidth} l l@{}}
\toprule Statement & Locator & Status & Formal \\ \midrule
\endfirsthead
\toprule Statement & Locator & Status & Formal \\ \midrule \endhead
\bottomrule\endfoot
{\scriptsize\ttfamily H0\_\allowbreak{}ENERGY\_\allowbreak{}CONTRACTION} & {\scriptsize \#\# 6. Scaling defects and unresolved estimates (C6)} & {\scriptsize open} & {\scriptsize ---} \\
{\scriptsize\ttfamily H1\_\allowbreak{}FIBERWISE\_\allowbreak{}SCORE\_\allowbreak{}VARIANCE} & {\scriptsize \#\# 6. Scaling defects and unresolved estimates (C6)} & {\scriptsize open} & {\scriptsize ---} \\
{\scriptsize\ttfamily VACUUM\_\allowbreak{}CROSS\_\allowbreak{}DERIVATIVES} & {\scriptsize \#\# 6. Scaling defects and unresolved estimates (C6)} & {\scriptsize open} & {\scriptsize ---} \\
{\scriptsize\ttfamily KATO\_\allowbreak{}DERIVATIVES} & {\scriptsize \#\# 6. Scaling defects and unresolved estimates (C6)} & {\scriptsize open} & {\scriptsize ---} \\
{\scriptsize\ttfamily CONDITIONAL\_\allowbreak{}DERIVATIVE} & {\scriptsize \#\# 2. Correct conditional differentiation (C1)} & {\scriptsize proven} & {\scriptsize ---} \\
{\scriptsize\ttfamily HORIZONTAL\_\allowbreak{}SCORE\_\allowbreak{}SUFFICIENT} & {\scriptsize \#\# 3. A sufficient horizontal-score estimate for H0 (C3)} & {\scriptsize proven} & {\scriptsize ---} \\
{\scriptsize\ttfamily H4\_\allowbreak{}FROM\_\allowbreak{}H0} & {\scriptsize \#\# 4. H4 follows from H0 and the established ground jet (C4)} & {\scriptsize proven} & {\scriptsize ---} \\
{\scriptsize\ttfamily PROJECTION\_\allowbreak{}JET\_\allowbreak{}RECURRENCE} & {\scriptsize \#\# 5. Complete projection-jet recurrence (C5)} & {\scriptsize proven} & {\scriptsize ---} \\
{\scriptsize\ttfamily SCALING\_\allowbreak{}OBSTRUCTIONS} & {\scriptsize \#\# 6. Scaling defects and unresolved estimates (C6)} & {\scriptsize proven} & {\scriptsize ---} \\
\end{longtable}

\subsection*{CITE:\allowbreak{}W6\_\allowbreak{}SOURCE\_\allowbreak{}GENERATOR\_\allowbreak{}SCORE\_\allowbreak{}FRAME}
\addcontentsline{toc}{subsection}{W6\_\allowbreak{}SOURCE\_\allowbreak{}GENERATOR\_\allowbreak{}SCORE\_\allowbreak{}FRAME}
\noindent\srcpath{docs/derivations/w6-source-generator-score-frame.md}\\
{\footnotesize SHA-256 \texttt{20b070cb78311187}\ldots}\par\smallskip
\begin{longtable}{@{}p{0.30\linewidth} p{0.30\linewidth} l l@{}}
\toprule Statement & Locator & Status & Formal \\ \midrule
\endfirsthead
\toprule Statement & Locator & Status & Formal \\ \midrule \endhead
\bottomrule\endfoot
{\scriptsize\ttfamily GROUND\_\allowbreak{}FRAME\_\allowbreak{}GENERATOR} & {\scriptsize \#\# 2. The ground-state frame (SF1-SF2) and \#\# 3. The complete R9 generator in the ground-state frame (SF3-SF5)} & {\scriptsize proven} & {\scriptsize ---} \\
{\scriptsize\ttfamily SCORE\_\allowbreak{}POISSON} & {\scriptsize \#\# 4. The score Poisson equation (SF6-SF7)} & {\scriptsize proven} & {\scriptsize ---} \\
{\scriptsize\ttfamily FIBER\_\allowbreak{}ENERGY\_\allowbreak{}IDENTITY} & {\scriptsize \#\# 4. The score Poisson equation (SF6-SF7), conditional divergence identity and fiber identity} & {\scriptsize proven} & {\scriptsize ---} \\
{\scriptsize\ttfamily R10\_\allowbreak{}ORDER\_\allowbreak{}ZERO\_\allowbreak{}REDUCTION} & {\scriptsize \#\# 5. Exact reduction of the order-zero R10 bound (SF8-SF9)} & {\scriptsize conditional} & {\scriptsize ---} \\
{\scriptsize\ttfamily CONDITIONAL\_\allowbreak{}ENERGY\_\allowbreak{}H0\_\allowbreak{}H4} & {\scriptsize \#\# 5. Exact reduction of the order-zero R10 bound (SF8-SF9), conditional-energy hypotheses} & {\scriptsize open} & {\scriptsize ---} \\
\end{longtable}

\subsection*{CITE:\allowbreak{}W6\_\allowbreak{}SYNCHRONIZED\_\allowbreak{}M10\_\allowbreak{}DOMINATION}
\addcontentsline{toc}{subsection}{W6\_\allowbreak{}SYNCHRONIZED\_\allowbreak{}M10\_\allowbreak{}DOMINATION}
\noindent\srcpath{docs/derivations/w6-synchronized-m10-domination.md}\\
{\footnotesize SHA-256 \texttt{e81031e9f7d13467}\ldots}\par\smallskip
\begin{longtable}{@{}p{0.30\linewidth} p{0.30\linewidth} l l@{}}
\toprule Statement & Locator & Status & Formal \\ \midrule
\endfirsthead
\toprule Statement & Locator & Status & Formal \\ \midrule \endhead
\bottomrule\endfoot
{\scriptsize\ttfamily TANGENCY\_\allowbreak{}DRIFT\_\allowbreak{}CANCELLATION} & {\scriptsize \#\# 2. Synchronized tangency (lines 34-51)} & {\scriptsize proven} & {\scriptsize ---} \\
{\scriptsize\ttfamily TUBE\_\allowbreak{}VARIANCE\_\allowbreak{}DOMINATION} & {\scriptsize \#\# 3. Small-angle tube implication (lines 52-84)} & {\scriptsize proven} & {\scriptsize ---} \\
{\scriptsize\ttfamily ANTIPODAL\_\allowbreak{}SCORE\_\allowbreak{}REGULARITY} & {\scriptsize \#\# 4. Antipodal geometry and a uniform angular reference calculation (lines 85-143)} & {\scriptsize proven} & {\scriptsize ---} \\
{\scriptsize\ttfamily AMPLITUDE\_\allowbreak{}GRADIENT\_\allowbreak{}SCALING} & {\scriptsize \#\# 5. Correct parameter-derivative criterion (lines 144-171)} & {\scriptsize open} & {\scriptsize ---} \\
{\scriptsize\ttfamily OUTSIDE\_\allowbreak{}COMPLEMENT\_\allowbreak{}SUPPRESSION} & {\scriptsize \#\# 6. Correct complement argument and compact action gap (lines 172-261)} & {\scriptsize proven} & {\scriptsize ---} \\
{\scriptsize\ttfamily FULL\_\allowbreak{}SCORE\_\allowbreak{}DOMINATION\_\allowbreak{}M10} & {\scriptsize \#\# 7. Assembly criterion and actual status (lines 262-286)} & {\scriptsize open} & {\scriptsize ---} \\
{\scriptsize\ttfamily LOG\_\allowbreak{}PARAMETER\_\allowbreak{}CRITERION} & {\scriptsize \#\# 5. Correct parameter-derivative criterion (lines 144-171)} & {\scriptsize proven} & {\scriptsize ---} \\
{\scriptsize\ttfamily COMPACT\_\allowbreak{}ACTION\_\allowbreak{}GAP} & {\scriptsize \#\# 6. Correct complement argument and compact action gap (lines 172-261)} & {\scriptsize proven} & {\scriptsize ---} \\
{\scriptsize\ttfamily NORMALIZED\_\allowbreak{}COMPLEMENT\_\allowbreak{}CRITERION} & {\scriptsize \#\# 6. Correct complement argument and compact action gap (lines 172-261)} & {\scriptsize proven} & {\scriptsize ---} \\
{\scriptsize\ttfamily ANGULAR\_\allowbreak{}REFERENCE\_\allowbreak{}BUDGET} & {\scriptsize \#\# 4. Antipodal geometry and a uniform angular reference calculation (lines 85-143)} & {\scriptsize proven} & {\scriptsize ---} \\
{\scriptsize\ttfamily ASSEMBLY\_\allowbreak{}CRITERION} & {\scriptsize \#\# 7. Assembly criterion and actual status (lines 262-286)} & {\scriptsize proven} & {\scriptsize ---} \\
{\scriptsize\ttfamily ACTUAL\_\allowbreak{}WEIGHTED\_\allowbreak{}JET\_\allowbreak{}AND\_\allowbreak{}LOWER\_\allowbreak{}BOUND} & {\scriptsize \#\# 6. Correct complement argument and compact action gap (lines 172-261)} & {\scriptsize proven} & {\scriptsize ---} \\
\end{longtable}

\subsection*{CITE:\allowbreak{}WILSON\_\allowbreak{}BACKGROUND}
\addcontentsline{toc}{subsection}{WILSON\_\allowbreak{}BACKGROUND}
\noindent\srcpath{docs/derivations/wilson-background-covariance-continuation.md}\\
{\footnotesize SHA-256 \texttt{05ee05dc5acf99fa}\ldots}\par\smallskip
\begin{longtable}{@{}p{0.30\linewidth} p{0.30\linewidth} l l@{}}
\toprule Statement & Locator & Status & Formal \\ \midrule
\endfirsthead
\toprule Statement & Locator & Status & Formal \\ \midrule \endhead
\bottomrule\endfoot
{\scriptsize\ttfamily BF1\_\allowbreak{}BF3} & {\scriptsize \#\# BF1. The energy-dependent relative-form statement (lines 48-102)} & {\scriptsize conditional} & {\scriptsize partial} \\
{\scriptsize\ttfamily BF4\_\allowbreak{}BF5} & {\scriptsize \#\# BF4. An actual Wilson small-field estimate, uniform in volume (lines 103-151)} & {\scriptsize proven} & {\scriptsize ---} \\
{\scriptsize\ttfamily BF6} & {\scriptsize \#\# BF6. Centering at an actual constrained minimizer (lines 152-188)} & {\scriptsize proven} & {\scriptsize ---} \\
{\scriptsize\ttfamily SC1\_\allowbreak{}SC4} & {\scriptsize \#\# 1. Setup and exact equations (lines 282-343)} & {\scriptsize proven} & {\scriptsize ---} \\
{\scriptsize\ttfamily SC5\_\allowbreak{}SC9} & {\scriptsize \#\# 2. A gauge-orbit gap that does not assume the physical gap (lines 344-415)} & {\scriptsize proven} & {\scriptsize ---} \\
{\scriptsize\ttfamily SC10\_\allowbreak{}SC15} & {\scriptsize \#\# 3. Uniform pointwise first and mixed derivatives at every coupling (lines 416-468)} & {\scriptsize proven} & {\scriptsize ---} \\
{\scriptsize\ttfamily SC16} & {\scriptsize \#\# 4. Direct finite-time form, without an implicit long-time interchange (lines 469-493)} & {\scriptsize proven} & {\scriptsize ---} \\
{\scriptsize\ttfamily SC17} & {\scriptsize \#\# 5. The neutral channel and the exact remaining spatial bound (lines 494-534)} & {\scriptsize conditional} & {\scriptsize ---} \\
{\scriptsize\ttfamily BC1} & {\scriptsize \#\# BC1. Why a fast inverse alone cannot finish the transport (lines 535-544)} & {\scriptsize proven} & {\scriptsize ---} \\
{\scriptsize\ttfamily BC2\_\allowbreak{}AUDIT} & {\scriptsize \#\# BC2. The additional multiscale document is evidence, not a completed input (lines 545-569)} & {\scriptsize proven} & {\scriptsize ---} \\
\end{longtable}

\subsection*{CITE:\allowbreak{}WILSON\_\allowbreak{}MARKED}
\addcontentsline{toc}{subsection}{WILSON\_\allowbreak{}MARKED}
\noindent\srcpath{docs/derivations/wilson-marked-transfer.md}\\
{\footnotesize SHA-256 \texttt{e7096d99ced64935}\ldots}\par\smallskip
\begin{longtable}{@{}p{0.30\linewidth} p{0.30\linewidth} l l@{}}
\toprule Statement & Locator & Status & Formal \\ \midrule
\endfirsthead
\toprule Statement & Locator & Status & Formal \\ \midrule \endhead
\bottomrule\endfoot
{\scriptsize\ttfamily WT1} & {\scriptsize \#\# WT-1. Exact physical-time blocks (lines 14-49)} & {\scriptsize proven} & {\scriptsize ---} \\
{\scriptsize\ttfamily WT2} & {\scriptsize \#\# WT-2. An all-orders activated-block bound (lines 50-88)} & {\scriptsize proven} & {\scriptsize ---} \\
{\scriptsize\ttfamily WT3} & {\scriptsize \#\# WT-3. Configurations and exact vacuum cancellation (lines 89-121)} & {\scriptsize proven} & {\scriptsize ---} \\
{\scriptsize\ttfamily WT4} & {\scriptsize \#\# WT-4. A uniform convergent vacuum polymer expansion (lines 122-171)} & {\scriptsize proven} & {\scriptsize ---} \\
{\scriptsize\ttfamily WT5} & {\scriptsize \#\# WT-5. Uniform control of free propagation and marked contour denominators (lines 172-208)} & {\scriptsize proven} & {\scriptsize ---} \\
{\scriptsize\ttfamily WT6} & {\scriptsize \#\# WT-6. The remaining complete-shell estimate and conditional transport (lines 209-248)} & {\scriptsize open} & {\scriptsize ---} \\
\end{longtable}

\subsection*{CITE:\allowbreak{}WILSON\_\allowbreak{}SC17}
\addcontentsline{toc}{subsection}{WILSON\_\allowbreak{}SC17}
\noindent\srcpath{docs/derivations/wilson-sc17-spatial-closure.md}\\
{\footnotesize SHA-256 \texttt{d02e906ba0d52d20}\ldots}\par\smallskip
\begin{longtable}{@{}p{0.30\linewidth} p{0.30\linewidth} l l@{}}
\toprule Statement & Locator & Status & Formal \\ \midrule
\endfirsthead
\toprule Statement & Locator & Status & Formal \\ \midrule \endhead
\bottomrule\endfoot
{\scriptsize\ttfamily REANCHORING} & {\scriptsize \#\# Endpoint reanchoring removes the neutral obstruction for norm bounds (lines 54-65)} & {\scriptsize proven} & {\scriptsize ---} \\
{\scriptsize\ttfamily CHARGED\_\allowbreak{}EVOLUTION} & {\scriptsize \#\# Why the charged contraction survives the time-dependent drift (lines 66-105)} & {\scriptsize proven} & {\scriptsize ---} \\
{\scriptsize\ttfamily H1\_\allowbreak{}H2} & {\scriptsize \#\# Actual weighted Volterra inequality (lines 106-131)} & {\scriptsize proven} & {\scriptsize ---} \\
{\scriptsize\ttfamily SMALL\_\allowbreak{}WINDOW} & {\scriptsize \#\# Explicit new regime (lines 132-166)} & {\scriptsize proven} & {\scriptsize ---} \\
{\scriptsize\ttfamily BARE\_\allowbreak{}MAJORANT\_\allowbreak{}WALL} & {\scriptsize \#\# Exact wall for this majorant (lines 167-189)} & {\scriptsize proven} & {\scriptsize ---} \\
{\scriptsize\ttfamily CURVATURE} & {\scriptsize \#\# Actual-vacuum curvature corollary (lines 190-212)} & {\scriptsize proven} & {\scriptsize ---} \\
{\scriptsize\ttfamily ANGLES} & {\scriptsize \#\# Sharper projection angles via covariance and the derived conditional gaps (lines 213-256)} & {\scriptsize proven} & {\scriptsize ---} \\
{\scriptsize\ttfamily OPTIMIZED\_\allowbreak{}INTERVAL} & {\scriptsize \#\# Maximized bare-Hessian bootstrap: a larger interval (lines 257-298)} & {\scriptsize proven} & {\scriptsize ---} \\
{\scriptsize\ttfamily RATIONAL\_\allowbreak{}WINDOW} & {\scriptsize \#\#\# Explicit rational larger-window corollary (lines 299-320)} & {\scriptsize proven} & {\scriptsize ---} \\
{\scriptsize\ttfamily GAUSSIAN\_\allowbreak{}PRECISION} & {\scriptsize \#\# 1. The actual flat quadratic ground precision (lines 323-358)} & {\scriptsize proven} & {\scriptsize ---} \\
{\scriptsize\ttfamily FRACTIONAL\_\allowbreak{}KERNEL} & {\scriptsize \#\# 2. Exact lattice fractional kernel and a sufficient weighted norm (lines 359-395)} & {\scriptsize proven} & {\scriptsize ---} \\
{\scriptsize\ttfamily GAUSSIAN\_\allowbreak{}ANGLE} & {\scriptsize \#\# 3. An exact scalar Gaussian angle test (lines 396-422)} & {\scriptsize proven} & {\scriptsize ---} \\
{\scriptsize\ttfamily FAST\_\allowbreak{}PRECISION\_\allowbreak{}CUSP} & {\scriptsize \#\# 4. Exact actual averaged-path fast conditional precision, ell=2 (lines 423-466)} & {\scriptsize proven} & {\scriptsize ---} \\
{\scriptsize\ttfamily R4A} & {\scriptsize \#\# 6. The existing dynamic covariance supplies the flat reference budget (lines 489-534)} & {\scriptsize proven} & {\scriptsize ---} \\
{\scriptsize\ttfamily R1\_\allowbreak{}R3} & {\scriptsize \#\# 1. Exact Euclidean equation and its projected cross terms (lines 537-573)} & {\scriptsize proven} & {\scriptsize ---} \\
{\scriptsize\ttfamily R4\_\allowbreak{}R9} & {\scriptsize \#\# 2. A rigorous conditional small-defect criterion (lines 574-648)} & {\scriptsize conditional} & {\scriptsize partial} \\
{\scriptsize\ttfamily R10\_\allowbreak{}R11} & {\scriptsize \#\# 3. An explicit true-angle consequence if the full Hessian defect is known (lines 649-697)} & {\scriptsize conditional} & {\scriptsize ---} \\
{\scriptsize\ttfamily R12\_\allowbreak{}R14} & {\scriptsize \#\# 4. The actual quantum pressure and cutoff defects are part of the input (lines 698-756)} & {\scriptsize conditional} & {\scriptsize ---} \\
{\scriptsize\ttfamily ACTUAL\_\allowbreak{}REFERENCE\_\allowbreak{}DEFECT\_\allowbreak{}BOUND} & {\scriptsize \#\# 4. The actual quantum pressure and cutoff defects are part of the input (lines 698-756), with the R5-R7 budget in section 2} & {\scriptsize open} & {\scriptsize ---} \\
{\scriptsize\ttfamily R15} & {\scriptsize \#\# 5. The averaged-path cusp is an analytic compression, not an eigenvalue crossing (lines 757-780)} & {\scriptsize proven} & {\scriptsize ---} \\
{\scriptsize\ttfamily PROPOSED\_\allowbreak{}QUANTUM\_\allowbreak{}PARAMETER\_\allowbreak{}WINDOW} & {\scriptsize \#\# 6. Formal repair of the spatial closure for lambda >= 1/73 via the quantum reference defect (lines 781-821)} & {\scriptsize conditional} & {\scriptsize partial} \\
\end{longtable}

\subsection*{CITE:\allowbreak{}WILSON\_\allowbreak{}SC17\_\allowbreak{}LIMIT}
\addcontentsline{toc}{subsection}{WILSON\_\allowbreak{}SC17\_\allowbreak{}LIMIT}
\noindent\srcpath{docs/derivations/wilson-sc17-thermodynamic-limit.md}\\
{\footnotesize SHA-256 \texttt{6d72876a35a86676}\ldots}\par\smallskip
\begin{longtable}{@{}p{0.30\linewidth} p{0.30\linewidth} l l@{}}
\toprule Statement & Locator & Status & Formal \\ \midrule
\endfirsthead
\toprule Statement & Locator & Status & Formal \\ \midrule \endhead
\bottomrule\endfoot
{\scriptsize\ttfamily T1} & {\scriptsize \#\# T1. Uniform setting including an actual cut interpolation (lines 28-95)} & {\scriptsize proven} & {\scriptsize ---} \\
{\scriptsize\ttfamily T2\_\allowbreak{}T3} & {\scriptsize \#\# T2. Weighted gradient contraction from the actual charged equations (lines 96-131)} & {\scriptsize proven} & {\scriptsize ---} \\
{\scriptsize\ttfamily T4\_\allowbreak{}T10} & {\scriptsize \#\# T3. The normalized actual vacuum response has no boundary-area loss (lines 132-195)} & {\scriptsize proven} & {\scriptsize ---} \\
{\scriptsize\ttfamily T11} & {\scriptsize \#\# T4. Actual exponential covariance localization (lines 196-210)} & {\scriptsize proven} & {\scriptsize ---} \\
{\scriptsize\ttfamily T12} & {\scriptsize \#\# T5. Fixed-lambda thermodynamic Cauchy convergence (lines 211-264)} & {\scriptsize proven} & {\scriptsize partial} \\
{\scriptsize\ttfamily T13} & {\scriptsize \#\# T7. Endpoint continuity also gives a thermodynamic limit at lambda\_\allowbreak{}c (lines 280-311)} & {\scriptsize proven} & {\scriptsize ---} \\
{\scriptsize\ttfamily SCORE\_\allowbreak{}LIMIT} & {\scriptsize \#\# Uniform score limits discharge the form construction's local input (lines 312-337)} & {\scriptsize proven} & {\scriptsize partial} \\
{\scriptsize\ttfamily IF3} & {\scriptsize \#\# 2. Exact infinite-volume integration by parts (lines 371-392)} & {\scriptsize proven} & {\scriptsize partial} \\
{\scriptsize\ttfamily IF4\_\allowbreak{}CLOSABLE} & {\scriptsize \#\# 3. Closability and the canonical closed form (lines 393-436)} & {\scriptsize proven} & {\scriptsize partial} \\
{\scriptsize\ttfamily IF4\_\allowbreak{}IF7} & {\scriptsize \#\# 3. Closability and the canonical closed form (lines 393-436)} & {\scriptsize proven} & {\scriptsize partial} \\
{\scriptsize\ttfamily PHYSICAL\_\allowbreak{}RESTRICTION} & {\scriptsize \#\# 4. Gauge invariance and the physical closed subspace (lines 437-460)} & {\scriptsize proven} & {\scriptsize ---} \\
{\scriptsize\ttfamily IF8\_\allowbreak{}LIMIT} & {\scriptsize \#\# 5. Passing the full and physical Poincare floor (lines 461-479)} & {\scriptsize proven} & {\scriptsize partial} \\
{\scriptsize\ttfamily IF8\_\allowbreak{}IF9} & {\scriptsize \#\# 5. Passing the full and physical Poincare floor (lines 461-479)} & {\scriptsize proven} & {\scriptsize partial} \\
{\scriptsize\ttfamily IF10\_\allowbreak{}IF13} & {\scriptsize \#\# 6. Actual nonzero local observable overlap with a bounded energy band (lines 480-528)} & {\scriptsize proven} & {\scriptsize partial} \\
\end{longtable}

\subsection*{CITE:\allowbreak{}WILSON\_\allowbreak{}SC17\_\allowbreak{}TIME}
\addcontentsline{toc}{subsection}{WILSON\_\allowbreak{}SC17\_\allowbreak{}TIME}
\noindent\srcpath{docs/derivations/wilson-sc17-physical-time-limit.md}\\
{\footnotesize SHA-256 \texttt{f842a6b9f826a997}\ldots}\par\smallskip
\begin{longtable}{@{}p{0.30\linewidth} p{0.30\linewidth} l l@{}}
\toprule Statement & Locator & Status & Formal \\ \midrule
\endfirsthead
\toprule Statement & Locator & Status & Formal \\ \midrule \endhead
\bottomrule\endfoot
{\scriptsize\ttfamily P4\_\allowbreak{}P6} & {\scriptsize \#\# P2. A limiting influence kernel and weighted Lipschitz estimate (lines 75-133)} & {\scriptsize proven} & {\scriptsize ---} \\
{\scriptsize\ttfamily P7\_\allowbreak{}P13} & {\scriptsize \#\# P3. A global extension with exact sphere invariance (lines 134-220)} & {\scriptsize proven} & {\scriptsize ---} \\
{\scriptsize\ttfamily P14} & {\scriptsize \#\# P4. Uniform Hilbert convergence of the actual finite drifts (lines 221-258)} & {\scriptsize proven} & {\scriptsize ---} \\
{\scriptsize\ttfamily P15} & {\scriptsize \#\# P5. Strong finite-time process comparison (lines 259-293)} & {\scriptsize proven} & {\scriptsize ---} \\
{\scriptsize\ttfamily P16} & {\scriptsize \#\# P6. Stationary path laws and actual multitime vacuum correlations (lines 294-332)} & {\scriptsize proven} & {\scriptsize ---} \\
{\scriptsize\ttfamily P17\_\allowbreak{}P18} & {\scriptsize \#\# P7. Identification with the already closed gradient form (lines 333-440)} & {\scriptsize proven} & {\scriptsize ---} \\
{\scriptsize\ttfamily PHYSICAL\_\allowbreak{}TIME\_\allowbreak{}GAP} & {\scriptsize \#\# P8. Physical restriction and the fixed-spacing time limit (lines 441-470)} & {\scriptsize proven} & {\scriptsize ---} \\
{\scriptsize\ttfamily P19\_\allowbreak{}P20} & {\scriptsize \#\# P9. Endpoint extension by a summable nonexponential weight (lines 471-525)} & {\scriptsize proven} & {\scriptsize ---} \\
\end{longtable}

\subsection*{CITE:\allowbreak{}WILSON\_\allowbreak{}SELECTED}
\addcontentsline{toc}{subsection}{WILSON\_\allowbreak{}SELECTED}
\noindent\srcpath{docs/derivations/wilson-selected-inverse-wall.md}\\
{\footnotesize SHA-256 \texttt{c0473aca76790ef7}\ldots}\par\smallskip
\begin{longtable}{@{}p{0.30\linewidth} p{0.30\linewidth} l l@{}}
\toprule Statement & Locator & Status & Formal \\ \midrule
\endfirsthead
\toprule Statement & Locator & Status & Formal \\ \midrule \endhead
\bottomrule\endfoot
{\scriptsize\ttfamily W1\_\allowbreak{}W3\_\allowbreak{}OBSTRUCTION} & {\scriptsize \#\# 2. First attempted step: global relative Gaussian fast-form control (lines 33-100)} & {\scriptsize proven} & {\scriptsize ---} \\
{\scriptsize\ttfamily W4\_\allowbreak{}W5} & {\scriptsize \#\# 3. Repair: only compare inverses on the graph force (lines 101-157)} & {\scriptsize conditional} & {\scriptsize partial} \\
{\scriptsize\ttfamily SIGNED\_\allowbreak{}PAIRING\_\allowbreak{}OBSTRUCTIONS} & {\scriptsize \#\# 4. Attempt to establish the repaired estimate for Wilson (lines 158-220)} & {\scriptsize proven} & {\scriptsize ---} \\
{\scriptsize\ttfamily W6} & {\scriptsize \#\# 5. The exact point at which this argument cannot proceed (lines 221-264)} & {\scriptsize open} & {\scriptsize partial} \\
\end{longtable}

\subsection*{CITE:\allowbreak{}WILSON\_\allowbreak{}SHELL}
\addcontentsline{toc}{subsection}{WILSON\_\allowbreak{}SHELL}
\noindent\srcpath{docs/derivations/wilson-marked-shell-transport.md}\\
{\footnotesize SHA-256 \texttt{a887ed9de154cec0}\ldots}\par\smallskip
\begin{longtable}{@{}p{0.30\linewidth} p{0.30\linewidth} l l@{}}
\toprule Statement & Locator & Status & Formal \\ \midrule
\endfirsthead
\toprule Statement & Locator & Status & Formal \\ \midrule \endhead
\bottomrule\endfoot
{\scriptsize\ttfamily S2\_\allowbreak{}S7} & {\scriptsize \#\# 2. The expansion with a fixed observable (lines 76-169)} & {\scriptsize proven} & {\scriptsize ---} \\
{\scriptsize\ttfamily S8\_\allowbreak{}S9} & {\scriptsize \#\# 3. Complex anchored operator bound (lines 170-218)} & {\scriptsize proven} & {\scriptsize ---} \\
{\scriptsize\ttfamily S10\_\allowbreak{}S15} & {\scriptsize \#\# 4. Actual complete shell and a holomorphic literal-source frame (lines 219-311)} & {\scriptsize proven} & {\scriptsize ---} \\
{\scriptsize\ttfamily S16\_\allowbreak{}S17} & {\scriptsize \#\# 5. From finite-order locality to a uniform weighted Taylor bound (lines 312-390)} & {\scriptsize proven} & {\scriptsize ---} \\
{\scriptsize\ttfamily S18\_\allowbreak{}S20} & {\scriptsize \#\# 6. Sum the matching and transport the carrier (lines 391-461)} & {\scriptsize conditional} & {\scriptsize ---} \\
{\scriptsize\ttfamily S21\_\allowbreak{}S22} & {\scriptsize \#\# 7. Observable overlap, with the quantifiers retained (lines 462-493)} & {\scriptsize proven} & {\scriptsize ---} \\
\end{longtable}

\subsection*{CITE:\allowbreak{}WILSON\_\allowbreak{}SPATIAL}
\addcontentsline{toc}{subsection}{WILSON\_\allowbreak{}SPATIAL}
\noindent\srcpath{docs/derivations/wilson-spatial-schur-excess.md}\\
{\footnotesize SHA-256 \texttt{5ddb7f812a9e6f53}\ldots}\par\smallskip
\begin{longtable}{@{}p{0.30\linewidth} p{0.30\linewidth} l l@{}}
\toprule Statement & Locator & Status & Formal \\ \midrule
\endfirsthead
\toprule Statement & Locator & Status & Formal \\ \midrule \endhead
\bottomrule\endfoot
{\scriptsize\ttfamily SP1\_\allowbreak{}SP6} & {\scriptsize \#\# 2. Exact excess about the full reference graph (lines 42-114)} & {\scriptsize conditional} & {\scriptsize partial} \\
{\scriptsize\ttfamily SP7\_\allowbreak{}SP10} & {\scriptsize \#\# 3. The complete second-order coefficient and a uniform remainder (lines 115-182)} & {\scriptsize conditional} & {\scriptsize partial} \\
{\scriptsize\ttfamily SP11\_\allowbreak{}SP16} & {\scriptsize \#\# 4. What W\_\allowbreak{}1 and W\_\allowbreak{}2 mean for the actual Wilson operator (lines 183-306)} & {\scriptsize proven} & {\scriptsize ---} \\
{\scriptsize\ttfamily SP17\_\allowbreak{}SP19} & {\scriptsize \#\# 5. Second-order literal-source normalization and harmonic control (lines 307-411)} & {\scriptsize proven} & {\scriptsize ---} \\
{\scriptsize\ttfamily SP20\_\allowbreak{}SP24} & {\scriptsize \#\# 6. Iterating energies and observable amplitudes in the physical clock (lines 412-477)} & {\scriptsize conditional} & {\scriptsize ---} \\
{\scriptsize\ttfamily SP25} & {\scriptsize \#\# 7. Why the complete second-order subtraction matters (lines 478-519)} & {\scriptsize proven} & {\scriptsize ---} \\
\end{longtable}

\subsection*{CITE:\allowbreak{}WILSON\_\allowbreak{}TRUE\_\allowbreak{}BLOCK}
\addcontentsline{toc}{subsection}{WILSON\_\allowbreak{}TRUE\_\allowbreak{}BLOCK}
\noindent\srcpath{docs/derivations/wilson-true-vacuum-block-estimates.md}\\
{\footnotesize SHA-256 \texttt{23aebb949cd2b6da}\ldots}\par\smallskip
\begin{longtable}{@{}p{0.30\linewidth} p{0.30\linewidth} l l@{}}
\toprule Statement & Locator & Status & Formal \\ \midrule
\endfirsthead
\toprule Statement & Locator & Status & Formal \\ \midrule \endhead
\bottomrule\endfoot
{\scriptsize\ttfamily BA1} & {\scriptsize \#\# BA1. A specified block partition (lines 10-28)} & {\scriptsize proven} & {\scriptsize ---} \\
{\scriptsize\ttfamily BA2\_\allowbreak{}BA5} & {\scriptsize \#\# BA2. Optimizing the true conditional gap (lines 29-94)} & {\scriptsize proven} & {\scriptsize partial} \\
{\scriptsize\ttfamily BA6\_\allowbreak{}BA8} & {\scriptsize \#\# BA3. Reduce the actual projection angle to a mixed ground-amplitude ratio (lines 95-141)} & {\scriptsize proven} & {\scriptsize partial} \\
{\scriptsize\ttfamily BA9\_\allowbreak{}BA10} & {\scriptsize \#\# BA4. What a frozen-block spectral computation would additionally require (lines 142-168)} & {\scriptsize conditional} & {\scriptsize ---} \\
{\scriptsize\ttfamily BA11\_\allowbreak{}BA12} & {\scriptsize \#\# BA5. Exact finite-time derivation of the connected quantity (lines 169-201)} & {\scriptsize proven} & {\scriptsize ---} \\
{\scriptsize\ttfamily BA13\_\allowbreak{}BA16} & {\scriptsize \#\# BA6. The first connected coefficients for the actual coupled lattice (lines 202-269)} & {\scriptsize proven} & {\scriptsize ---} \\
{\scriptsize\ttfamily BA20\_\allowbreak{}BA26} & {\scriptsize \#\# BA7. Actual angle closure in an explicit strong-coupling window (lines 270-405)} & {\scriptsize proven} & {\scriptsize ---} \\
\end{longtable}

\subsection*{CITE:\allowbreak{}WILSON\_\allowbreak{}TRUE\_\allowbreak{}BLOCK\_\allowbreak{}CHECK}
\addcontentsline{toc}{subsection}{WILSON\_\allowbreak{}TRUE\_\allowbreak{}BLOCK\_\allowbreak{}CHECK}
\noindent\srcpath{docs/derivations/wilson-true-blocks-independent-check.md}\\
{\footnotesize SHA-256 \texttt{6984f428f801c70e}\ldots}\par\smallskip
\begin{longtable}{@{}p{0.30\linewidth} p{0.30\linewidth} l l@{}}
\toprule Statement & Locator & Status & Formal \\ \midrule
\endfirsthead
\toprule Statement & Locator & Status & Formal \\ \midrule \endhead
\bottomrule\endfoot
{\scriptsize\ttfamily ENVELOPE\_\allowbreak{}AND\_\allowbreak{}CROSS\_\allowbreak{}RATIO} & {\scriptsize \#\# CB1. What was re-derived independently (lines 17-56)} & {\scriptsize proven} & {\scriptsize ---} \\
{\scriptsize\ttfamily CB1} & {\scriptsize \#\# CB2. The single-link partition is stronger, and closes VA12 directly (lines 57-75)} & {\scriptsize proven} & {\scriptsize ---} \\
{\scriptsize\ttfamily FREE\_\allowbreak{}BUDGET\_\allowbreak{}LOOSENESS} & {\scriptsize \#\# CB3. The angle budget is not graded in lambda (lines 76-87)} & {\scriptsize proven} & {\scriptsize ---} \\
{\scriptsize\ttfamily CB2} & {\scriptsize \#\# CB4. The BA26 wall is more robust than its own derivation (lines 88-109)} & {\scriptsize proven} & {\scriptsize ---} \\
{\scriptsize\ttfamily NONCHAINING} & {\scriptsize \#\# CB5. The two boxed constants must not be chained (lines 110-121)} & {\scriptsize proven} & {\scriptsize ---} \\
\end{longtable}

\subsection*{CITE:\allowbreak{}WILSON\_\allowbreak{}VACUUM\_\allowbreak{}ASSEMBLY}
\addcontentsline{toc}{subsection}{WILSON\_\allowbreak{}VACUUM\_\allowbreak{}ASSEMBLY}
\noindent\srcpath{docs/derivations/wilson-vacuum-aligned-assembly.md}\\
{\footnotesize SHA-256 \texttt{6c21bda892b88872}\ldots}\par\smallskip
\begin{longtable}{@{}p{0.30\linewidth} p{0.30\linewidth} l l@{}}
\toprule Statement & Locator & Status & Formal \\ \midrule
\endfirsthead
\toprule Statement & Locator & Status & Formal \\ \midrule \endhead
\bottomrule\endfoot
{\scriptsize\ttfamily VA1\_\allowbreak{}VA4} & {\scriptsize \#\# VA1. A local finite-energy bound for the actual quantum vacuum (lines 17-74)} & {\scriptsize proven} & {\scriptsize partial} \\
{\scriptsize\ttfamily VA5\_\allowbreak{}VA7} & {\scriptsize \#\# VA2. The WR25 remainder is extensive on actual connected lattices (lines 75-130)} & {\scriptsize proven} & {\scriptsize partial} \\
{\scriptsize\ttfamily VA8} & {\scriptsize \#\# VA3. WR26 is false on the physical vacuum complement at large volume (lines 131-164)} & {\scriptsize proven} & {\scriptsize ---} \\
{\scriptsize\ttfamily VA9\_\allowbreak{}VA11} & {\scriptsize \#\# VA4. Exact cancellation before comparison (lines 165-210)} & {\scriptsize proven} & {\scriptsize partial} \\
{\scriptsize\ttfamily VA12\_\allowbreak{}VA13} & {\scriptsize \#\# VA5. Projection geometry, with the Gauss law retained (lines 211-251)} & {\scriptsize conditional} & {\scriptsize partial} \\
{\scriptsize\ttfamily VA14\_\allowbreak{}VA19} & {\scriptsize \#\# VA6. The next estimate and the attempts to establish it (lines 252-415)} & {\scriptsize conditional} & {\scriptsize partial} \\
\end{longtable}

\subsection*{CITE:\allowbreak{}WILSON\_\allowbreak{}WEIGHTED\_\allowbreak{}REPAIR}
\addcontentsline{toc}{subsection}{WILSON\_\allowbreak{}WEIGHTED\_\allowbreak{}REPAIR}
\noindent\srcpath{docs/derivations/wilson-weighted-repair-and-rotor-gap.md}\\
{\footnotesize SHA-256 \texttt{dac4e8fd2d13f3e7}\ldots}\par\smallskip
\begin{longtable}{@{}p{0.30\linewidth} p{0.30\linewidth} l l@{}}
\toprule Statement & Locator & Status & Formal \\ \midrule
\endfirsthead
\toprule Statement & Locator & Status & Formal \\ \midrule \endhead
\bottomrule\endfoot
{\scriptsize\ttfamily WR1\_\allowbreak{}WR3} & {\scriptsize \#\# WR1. Absorb negative curvature in the derivative energy (lines 15-66)} & {\scriptsize proven} & {\scriptsize ---} \\
{\scriptsize\ttfamily WR4\_\allowbreak{}WR6} & {\scriptsize \#\# WR2. A smooth non-Gaussian trial for the physical rotor (lines 67-99)} & {\scriptsize proven} & {\scriptsize ---} \\
{\scriptsize\ttfamily WR7\_\allowbreak{}WR12} & {\scriptsize \#\# WR3. A physical gap on the complete rotor space, at every coupling (lines 100-161)} & {\scriptsize proven} & {\scriptsize ---} \\
{\scriptsize\ttfamily ROTOR\_\allowbreak{}VOLUME} & {\scriptsize \#\# WR4. Volume passage that is already justified (lines 162-175)} & {\scriptsize proven} & {\scriptsize ---} \\
{\scriptsize\ttfamily WR13\_\allowbreak{}WR16} & {\scriptsize \#\# WR5. Actual shared-link residual, and a positive interacting repair (lines 176-228)} & {\scriptsize proven} & {\scriptsize ---} \\
{\scriptsize\ttfamily WR17\_\allowbreak{}WR21} & {\scriptsize \#\# WR6. The precise remaining derivation (lines 229-308)} & {\scriptsize conditional} & {\scriptsize ---} \\
{\scriptsize\ttfamily WR22\_\allowbreak{}WR24} & {\scriptsize \#\# WR7. Retain the full boundary: an actual-link fast barrier (lines 309-356)} & {\scriptsize proven} & {\scriptsize ---} \\
{\scriptsize\ttfamily WR25\_\allowbreak{}WR28} & {\scriptsize \#\# WR8. Where the repaired local bounds still require a new proof (lines 357-443)} & {\scriptsize proven} & {\scriptsize ---} \\
\end{longtable}

\subsection*{CITE:\allowbreak{}YM\_\allowbreak{}BALABAN\_\allowbreak{}MULTISCALE}
\addcontentsline{toc}{subsection}{YM\_\allowbreak{}BALABAN\_\allowbreak{}MULTISCALE}
\noindent\srcpath{docs/derivations/yangmills-continuum-balaban-multiscale-proof.md}\\
{\footnotesize SHA-256 \texttt{55296902489318e2}\ldots}\par\smallskip
\begin{longtable}{@{}p{0.30\linewidth} p{0.30\linewidth} l l@{}}
\toprule Statement & Locator & Status & Formal \\ \midrule
\endfirsthead
\toprule Statement & Locator & Status & Formal \\ \midrule \endhead
\bottomrule\endfoot
{\scriptsize\ttfamily LEMMA\_\allowbreak{}3\_\allowbreak{}2\_\allowbreak{}ADJOINT} & {\scriptsize \#\#\# 3.1 Geometric Covariant Averaging Operator \$Q\_\allowbreak{}k\$ (lines 84-110)} & {\scriptsize proven} & {\scriptsize ---} \\
{\scriptsize\ttfamily PROPOSITION\_\allowbreak{}3\_\allowbreak{}4} & {\scriptsize \#\#\# 3.2 The Background Field Variational Problem (lines 111-139)} & {\scriptsize disputed} & {\scriptsize ---} \\
{\scriptsize\ttfamily THEOREM\_\allowbreak{}4\_\allowbreak{}2} & {\scriptsize \#\#\# 4.2 Super-Exponential Suppression of Large Fields (lines 157-180)} & {\scriptsize conditional} & {\scriptsize ---} \\
{\scriptsize\ttfamily GAUGE\_\allowbreak{}HESSIAN} & {\scriptsize \#\#\# 5.1 Fluctuation Gauge Fixing (lines 183-197)} & {\scriptsize disputed} & {\scriptsize ---} \\
{\scriptsize\ttfamily GHOST\_\allowbreak{}5\_\allowbreak{}8} & {\scriptsize \#\#\# 5.2 Faddeev--Popov Ghost Determinant (lines 198-208)} & {\scriptsize conditional} & {\scriptsize ---} \\
{\scriptsize\ttfamily THEOREM\_\allowbreak{}5\_\allowbreak{}1} & {\scriptsize \#\#\# 5.3 Combes--Thomas Resolvent Localization for Continuous \$\textbackslash{}mathcal\{H\}\_\allowbreak{}\{B\_\allowbreak{}k\}\$ (lines 209-248)} & {\scriptsize disputed} & {\scriptsize ---} \\
{\scriptsize\ttfamily BANACH\_\allowbreak{}ALGEBRA} & {\scriptsize \#\#\# 6.2 The Polymer Banach Algebra \$\textbackslash{}mathcal\{E\}\_\allowbreak{}k\$ (lines 260-268)} & {\scriptsize proven} & {\scriptsize ---} \\
{\scriptsize\ttfamily THEOREM\_\allowbreak{}6\_\allowbreak{}2} & {\scriptsize \#\#\# 6.3 The Mayer Tree-Graph Expansion \& Cluster Contraction (lines 269-295)} & {\scriptsize conditional} & {\scriptsize ---} \\
{\scriptsize\ttfamily THEOREM\_\allowbreak{}6\_\allowbreak{}3} & {\scriptsize \#\#\# 6.4 Asymptotic Freedom and Infinite Scale Stability (lines 296-314)} & {\scriptsize conditional} & {\scriptsize ---} \\
{\scriptsize\ttfamily THEOREM\_\allowbreak{}7\_\allowbreak{}1} & {\scriptsize \#\#\# 7.1 Cauchy Sequence of Effective Actions (lines 317-340)} & {\scriptsize open} & {\scriptsize ---} \\
{\scriptsize\ttfamily THEOREM\_\allowbreak{}7\_\allowbreak{}2} & {\scriptsize \#\#\# 7.2 Construction of the Limiting Continuum Measure \$d\textbackslash{}mu\_\allowbreak{}\{\textbackslash{}text\{YM\}\}\$ (lines 341-363)} & {\scriptsize conditional} & {\scriptsize ---} \\
{\scriptsize\ttfamily OS0} & {\scriptsize \#\#\# 8.1 OS0: Analyticity and Temperedness (lines 368-374)} & {\scriptsize conditional} & {\scriptsize ---} \\
{\scriptsize\ttfamily OS1} & {\scriptsize \#\#\# 8.2 OS1: Euclidean Covariance (lines 375-382)} & {\scriptsize open} & {\scriptsize ---} \\
{\scriptsize\ttfamily THEOREM\_\allowbreak{}8\_\allowbreak{}1} & {\scriptsize \#\#\# 8.3 OS2: Reflection Positivity (lines 383-408)} & {\scriptsize conditional} & {\scriptsize ---} \\
{\scriptsize\ttfamily OS3} & {\scriptsize \#\#\# 8.4 OS3: Permutation Symmetry (lines 409-411)} & {\scriptsize conditional} & {\scriptsize ---} \\
{\scriptsize\ttfamily THEOREM\_\allowbreak{}8\_\allowbreak{}2} & {\scriptsize \#\#\# 8.5 OS4: Exponential Clustering and Mass Gap (lines 412-436)} & {\scriptsize open} & {\scriptsize ---} \\
{\scriptsize\ttfamily OS\_\allowbreak{}RECONSTRUCTION} & {\scriptsize \#\#\# 9.1 The Physical Hilbert Space and Hamiltonian (lines 439-453)} & {\scriptsize conditional} & {\scriptsize ---} \\
{\scriptsize\ttfamily THEOREM\_\allowbreak{}9\_\allowbreak{}1} & {\scriptsize \#\#\# 9.2 The Physical Spectral Mass Gap (lines 454-478)} & {\scriptsize conditional} & {\scriptsize ---} \\
{\scriptsize\ttfamily SCATTERING} & {\scriptsize \#\#\# 9.3 Relativistic Wightman Fields and Non-Trivial Scattering (\$S \textbackslash{}ne I\$) (lines 479-501)} & {\scriptsize conditional} & {\scriptsize ---} \\
\end{longtable}

\subsection*{CITE:\allowbreak{}YM\_\allowbreak{}FLAT\_\allowbreak{}DIRECTIONS}
\addcontentsline{toc}{subsection}{YM\_\allowbreak{}FLAT\_\allowbreak{}DIRECTIONS}
\noindent\srcpath{docs/derivations/yangmills-simon-flat-directions.md}\\
{\footnotesize SHA-256 \texttt{850d8bf0ea4b6192}\ldots}\par\smallskip
\begin{longtable}{@{}p{0.30\linewidth} p{0.30\linewidth} l l@{}}
\toprule Statement & Locator & Status & Formal \\ \midrule
\endfirsthead
\toprule Statement & Locator & Status & Formal \\ \midrule \endhead
\bottomrule\endfoot
{\scriptsize\ttfamily OSCILLATOR} & {\scriptsize \#\# 1. The scalar oscillator, including its constant (lines 102-138)} & {\scriptsize proven} & {\scriptsize ---} \\
{\scriptsize\ttfamily TRANSVERSE} & {\scriptsize \#\# 2. SU(2) has two transverse oscillator modes per slice (lines 139-167)} & {\scriptsize proven} & {\scriptsize ---} \\
{\scriptsize\ttfamily SLICE\_\allowbreak{}AGGREGATION} & {\scriptsize \#\# 3. Allocate the pair potential without losing longitudinal control (lines 168-212)} & {\scriptsize proven} & {\scriptsize ---} \\
{\scriptsize\ttfamily COMPACT\_\allowbreak{}RESOLVENT} & {\scriptsize \#\# 4. Compact resolvent follows from the retained gradient and radial tail (lines 213-233)} & {\scriptsize proven} & {\scriptsize ---} \\
{\scriptsize\ttfamily GROUND\_\allowbreak{}FLOOR} & {\scriptsize \#\# 5. An explicit ground-energy lower bound (lines 234-275)} & {\scriptsize proven} & {\scriptsize ---} \\
{\scriptsize\ttfamily JACOBIAN\_\allowbreak{}VALLEY} & {\scriptsize \#\# 6. What Simon's Jacobian condition measures (lines 276-324)} & {\scriptsize proven} & {\scriptsize ---} \\
{\scriptsize\ttfamily SCALING} & {\scriptsize \#\# 7. Energy scale, vacuum subtraction and the research gaps (lines 325-368)} & {\scriptsize proven} & {\scriptsize ---} \\
\end{longtable}

\subsection*{CITE:\allowbreak{}YM\_\allowbreak{}GPU\_\allowbreak{}AUDIT}
\addcontentsline{toc}{subsection}{YM\_\allowbreak{}GPU\_\allowbreak{}AUDIT}
\noindent\srcpath{docs/derivations/yangmills-gpu-resolution-audit.md}\\
{\footnotesize SHA-256 \texttt{45a72f1f90d9da8f}\ldots}\par\smallskip
\begin{longtable}{@{}p{0.30\linewidth} p{0.30\linewidth} l l@{}}
\toprule Statement & Locator & Status & Formal \\ \midrule
\endfirsthead
\toprule Statement & Locator & Status & Formal \\ \midrule \endhead
\bottomrule\endfoot
{\scriptsize\ttfamily GA1\_\allowbreak{}GA4} & {\scriptsize \#\# GA1. Phase 1 is a massive scalar model with an exact uniform theorem (lines 24-76)} & {\scriptsize proven} & {\scriptsize ---} \\
{\scriptsize\ttfamily GA5\_\allowbreak{}GA7} & {\scriptsize \#\# GA2. Phase 2's alleged domino is an exact tensor sum (lines 77-129)} & {\scriptsize proven} & {\scriptsize ---} \\
{\scriptsize\ttfamily RADIAL\_\allowbreak{}DISCRETIZATION} & {\scriptsize \#\# GA3. Phase 3 compares two radial discretizations, not a transfer matrix (lines 130-161)} & {\scriptsize proven} & {\scriptsize ---} \\
{\scriptsize\ttfamily GA8\_\allowbreak{}GA9} & {\scriptsize \#\# GA4. The additional Phase 4 has a valid conditional scalar limit (lines 162-209)} & {\scriptsize conditional} & {\scriptsize ---} \\
{\scriptsize\ttfamily GA10} & {\scriptsize \#\# GA5. Phase 5 keeps an admissible cube basis but changes its magnetic operator (lines 210-249)} & {\scriptsize proven} & {\scriptsize partial} \\
{\scriptsize\ttfamily GA11} & {\scriptsize \#\# GA6. What the supplied Lean additions actually formalize (lines 250-286)} & {\scriptsize proven} & {\scriptsize ---} \\
{\scriptsize\ttfamily GA12} & {\scriptsize \#\# GA7. The incoming smearing script measures spatial power, not spectral overlap (lines 287-327)} & {\scriptsize proven} & {\scriptsize ---} \\
{\scriptsize\ttfamily GA13} & {\scriptsize \#\# GA8. The later G17 script inserts both its cluster coefficients and its mass (lines 328-363)} & {\scriptsize proven} & {\scriptsize ---} \\
{\scriptsize\ttfamily GA14\_\allowbreak{}GA19} & {\scriptsize \#\# GA9. Phase 8 supplies an assumed block kernel, not conditional projections (lines 364-500)} & {\scriptsize proven} & {\scriptsize ---} \\
{\scriptsize\ttfamily GA20\_\allowbreak{}GA21} & {\scriptsize \#\# GA10. Phase 9 verifies a Gram identity, an inserted exponential and a Haar moment (lines 501-563)} & {\scriptsize proven} & {\scriptsize partial} \\
\end{longtable}

\subsection*{CITE:\allowbreak{}YM\_\allowbreak{}GROUND\_\allowbreak{}STATE}
\addcontentsline{toc}{subsection}{YM\_\allowbreak{}GROUND\_\allowbreak{}STATE}
\noindent\srcpath{docs/derivations/yangmills-weighted-curvature.md}\\
{\footnotesize SHA-256 \texttt{75c853725c68db30}\ldots}\par\smallskip
\begin{longtable}{@{}p{0.30\linewidth} p{0.30\linewidth} l l@{}}
\toprule Statement & Locator & Status & Formal \\ \midrule
\endfirsthead
\toprule Statement & Locator & Status & Formal \\ \midrule \endhead
\bottomrule\endfoot
{\scriptsize\ttfamily GST1} & {\scriptsize \#\# GST-1: exact conjugation with a trial-state residual (lines 32-70)} & {\scriptsize proven} & {\scriptsize partial} \\
{\scriptsize\ttfamily GST2} & {\scriptsize \#\# GST-2: weighted energy and the needed domain statement (lines 71-91)} & {\scriptsize proven} & {\scriptsize partial} \\
{\scriptsize\ttfamily GST3} & {\scriptsize \#\# GST-3: curvature normalization and a spectral argument without a discrete excited state (lines 92-130)} & {\scriptsize conditional} & {\scriptsize partial} \\
{\scriptsize\ttfamily GST4} & {\scriptsize \#\# GST-4: the finite matrix identity and its exact hypotheses (lines 131-146)} & {\scriptsize proven} & {\scriptsize full} \\
{\scriptsize\ttfamily GST5} & {\scriptsize \#\# GST-5: oscillator calibration (lines 147-166)} & {\scriptsize proven} & {\scriptsize ---} \\
{\scriptsize\ttfamily GST6} & {\scriptsize \#\# GST-6: the residual-mean certificate (lines 167-190)} & {\scriptsize conditional} & {\scriptsize ---} \\
{\scriptsize\ttfamily GST7} & {\scriptsize \#\# GST-7: a complete algebraic certificate for the pure quartic oscillator (lines 191-244)} & {\scriptsize proven} & {\scriptsize ---} \\
{\scriptsize\ttfamily GST8} & {\scriptsize \#\# GST-8: falsifiers and the Yang--Mills boundary (lines 245-266)} & {\scriptsize proven} & {\scriptsize ---} \\
{\scriptsize\ttfamily GST9} & {\scriptsize \#\# GST-9: an exact volume-accumulation falsifier (lines 267-300)} & {\scriptsize proven} & {\scriptsize ---} \\
\end{longtable}

\subsection*{CITE:\allowbreak{}YM\_\allowbreak{}RECONSTRUCTION}
\addcontentsline{toc}{subsection}{YM\_\allowbreak{}RECONSTRUCTION}
\noindent\srcpath{docs/derivations/yangmills-reconstruction.md}\\
{\footnotesize SHA-256 \texttt{a180dddf14179838}\ldots}\par\smallskip
\begin{longtable}{@{}p{0.30\linewidth} p{0.30\linewidth} l l@{}}
\toprule Statement & Locator & Status & Formal \\ \midrule
\endfirsthead
\toprule Statement & Locator & Status & Formal \\ \midrule \endhead
\bottomrule\endfoot
{\scriptsize\ttfamily R1} & {\scriptsize \#\# 1. Source statements and the operator to be bounded (lines 5-27)} & {\scriptsize proven} & {\scriptsize partial} \\
{\scriptsize\ttfamily R2} & {\scriptsize \#\# 2. A finite reflection-positive model one can inspect exactly (lines 28-38)} & {\scriptsize proven} & {\scriptsize full} \\
{\scriptsize\ttfamily R3} & {\scriptsize \#\# 3. Observable completeness is a testable interface (lines 39-51)} & {\scriptsize proven} & {\scriptsize full} \\
{\scriptsize\ttfamily R4\_\allowbreak{}R5} & {\scriptsize \#\# 4. What a localization plateau actually proves (lines 52-89)} & {\scriptsize proven} & {\scriptsize partial} \\
{\scriptsize\ttfamily R6\_\allowbreak{}R7} & {\scriptsize \#\# 5. A precise sufficient repair, with its remaining hypothesis exposed (lines 90-111)} & {\scriptsize conditional} & {\scriptsize partial} \\
\end{longtable}

% ---- end gen_concordance.tex ---------------------------------------

\section{Statements that do not close}
\label{app:open}

\DerivOpen{} of the \DerivStatements{} registered statements are
recorded as open, conditional, or disputed. They are extracted here
with their hypotheses and the corpus's own sentence on what remains,
because this is the shortest complete list of places where an agent can
do mathematics that the program has already established is missing.

Read alongside \S\ref{sec:obligations}, which curates the subset the
program judges decisive, and \S\ref{sec:attack}, which is about how to
start on one.

% ---- begin gen_open_statements.tex ---------------------------------
% Generated by paper/master/generate_sources.py.  Do not edit.
\noindent 63 of the registered derivation statements are not recorded as proven. Each is listed with the document it lives in, its locator, and the corpus's own sentence on what remains. These are the statements an agent can act on directly.\par\medskip
\begingroup\sloppy
\noindent\textbf{\texttt{\small DERIV:\allowbreak{}EARLIER\_\allowbreak{}PROOF\_\allowbreak{}RECOVERY:\allowbreak{}DIPOLAR\_\allowbreak{}REMAINDER}} \hfill {\footnotesize open}\par
{\small Establish the archived O(|x|\textasciicircum{}-4) remainder of the dipolar lattice kernel with the required Fourier-cutoff and derivative bounds.}\par
{\footnotesize \srcpath{docs/derivations/earlier-proof-recovery.md} --- \#\# E18. Explicit remaining application obligations}\par
{\footnotesize \textit{Hypotheses.} The exact source assumptions and required estimates described in E18.}\par
{\footnotesize \textit{Remaining.} Explicit pending application or extension; neither the exact finite controls nor the narrower recovered theorem discharge it.}\par
\smallskip
\noindent\textbf{\texttt{\small DERIV:\allowbreak{}EARLIER\_\allowbreak{}PROOF\_\allowbreak{}RECOVERY:\allowbreak{}PHYSICAL\_\allowbreak{}CARRIER\_\allowbreak{}REALIZATION}} \hfill {\footnotesize open}\par
{\small Construct actual cutoff sources, residue floor, collapsing isolated intervals, true-sector gap, physical clock and limiting identifications needed to apply the positive matrix atom theorem.}\par
{\footnotesize \srcpath{docs/derivations/earlier-proof-recovery.md} --- \#\# E18. Explicit remaining application obligations}\par
{\footnotesize \textit{Hypotheses.} The exact source assumptions and required estimates described in E18.}\par
{\footnotesize \textit{Remaining.} Explicit pending application or extension; neither the exact finite controls nor the narrower recovered theorem discharge it.}\par
\smallskip
\noindent\textbf{\texttt{\small DERIV:\allowbreak{}EARLIER\_\allowbreak{}PROOF\_\allowbreak{}RECOVERY:\allowbreak{}STAPLE\_\allowbreak{}TRANSVERSALITY\_\allowbreak{}REPAIR}} \hfill {\footnotesize open}\par
{\small Repair the force-rank and tube estimates with the full moving-adjoint derivative.}\par
{\footnotesize \srcpath{docs/derivations/earlier-proof-recovery.md} --- \#\# E18. Explicit remaining application obligations}\par
{\footnotesize \textit{Hypotheses.} The exact source assumptions and required estimates described in E18.}\par
{\footnotesize \textit{Remaining.} Explicit pending application or extension; neither the exact finite controls nor the narrower recovered theorem discharge it.}\par
\smallskip
\noindent\textbf{\texttt{\small DERIV:\allowbreak{}EARLIER\_\allowbreak{}PROOF\_\allowbreak{}RECOVERY:\allowbreak{}VSU\_\allowbreak{}WHOLE\_\allowbreak{}SPACE}} \hfill {\footnotesize open}\par
{\small Establish compatible normalization and uniform local estimates for the whole-space VSU limit, then justify the claimed far-field law.}\par
{\footnotesize \srcpath{docs/derivations/earlier-proof-recovery.md} --- \#\# E18. Explicit remaining application obligations}\par
{\footnotesize \textit{Hypotheses.} The exact source assumptions and required estimates described in E18.}\par
{\footnotesize \textit{Remaining.} Explicit pending application or extension; neither the exact finite controls nor the narrower recovered theorem discharge it.}\par
\smallskip
\noindent\textbf{\texttt{\small DERIV:\allowbreak{}UNIVERSAL\_\allowbreak{}CELLULAR\_\allowbreak{}HODGE:\allowbreak{}PROPER\_\allowbreak{}RETURN\_\allowbreak{}Q\_\allowbreak{}ANNIHILATION}} \hfill {\footnotesize open}\par
{\small The 60 paths flagged by the submitted endpoint-or-zero flux predicate identify physical tetrahedral retained-sector returns that vanish individually under intermediate Q projection.}\par
{\footnotesize \srcpath{docs/derivations/universal-cellular-hodge-tetrahedral.md} --- \#\#\# 3.4 Open physical-history identification (lines 133-156)}\par
{\footnotesize \textit{Hypotheses.} An actual Haar/Fierz Hilbert space, physical retained projector, insertion channels, and projected resolvent states matching the enumerated paths.}\par
{\footnotesize \textit{Remaining.} Open: construct and identify the physical states/projector and calculate each projected contribution. The 60/36 flux count and Q psi=0 do not establish this identification.}\par
\smallskip
\noindent\textbf{\texttt{\small DERIV:\allowbreak{}W6\_\allowbreak{}CONDITIONAL\_\allowbreak{}SCORE\_\allowbreak{}TAIL\_\allowbreak{}CONTROL:\allowbreak{}SCORE\_\allowbreak{}DOMINATION\_\allowbreak{}M10}} \hfill {\footnotesize open}\par
{\small For the actual compact square and specified source-preserving dilation, prove K\_\allowbreak{}g(w)<=C0+C1 W\_\allowbreak{}g(w) almost everywhere in nu\_\allowbreak{}g, where W\_\allowbreak{}g(w)=g\textasciicircum{}-2 E\_\allowbreak{}(mu\_\allowbreak{}g)(V|w), with finite nonnegative C0,C1 independent of 0<g<g\_\allowbreak{}*. This sufficient conditional-score domination is open.}\par
{\footnotesize \srcpath{docs/derivations/w6-conditional-score-tail-control.md} --- \#\#\# M.4. The precise remaining conditional-score inequality (lines 473-495)}\par
{\footnotesize \textit{Hypotheses.} Actual fixed twelve-edge, four-face compact Wilson square, true normalized quantum ground Psi\_\allowbreak{}g, dmu\_\allowbreak{}g=Psi\_\allowbreak{}g\textasciicircum{}2 dU, literal source w=Tr(U2 U3)/2 and its marginal nu\_\allowbreak{}g.; sigma\_\allowbreak{}g=(partial\_\allowbreak{}g Psi\_\allowbreak{}g+D Psi\_\allowbreak{}g/g)/Psi\_\allowbreak{}g uses the specified compact dilation of Q8, including its Haar divergence; K\_\allowbreak{}g=Var\_\allowbreak{}(mu\_\allowbreak{}g)(sigma\_\allowbreak{}g|w) retains all compact corrections and rare coarse fibers.; The established actual fixed-square gap and ground-energy bounds hold on the specified small-coupling interval. Averaged conditional-score control is not assumed to imply pointwise domination.}\par
{\footnotesize \textit{Remaining.} Prove M10 for the actual conditional law, or retain its open status if a different sufficient source-energy estimate is proved. The phase-adapted local implication requires whole-fiber tails, conditional-mean gluing and consistent global dilation before it can address this target.}\par
\smallskip
\noindent\textbf{\texttt{\small DERIV:\allowbreak{}W6\_\allowbreak{}CONDITIONAL\_\allowbreak{}TRANSPORT\_\allowbreak{}OBSTRUCTION:\allowbreak{}M10\_\allowbreak{}SYNCHRONIZED\_\allowbreak{}SCORE}} \hfill {\footnotesize open}\par
{\small The selected synchronized realization of the generic M10 target asks whether the actual score for the specified S13 compact field obeys K\_\allowbreak{}g(w)<=C0+C1 g\textasciicircum{}-2 E\_\allowbreak{}mu\_\allowbreak{}g[V|w] with finite nonnegative constants uniform in sufficiently small positive g and almost every retained trace w.}\par
{\footnotesize \srcpath{docs/derivations/w6-conditional-transport-obstruction.md} --- \#\# S6 Remaining all-fiber estimate (lines 317-375), explicitly open synchronized M10 target}\par
{\footnotesize \textit{Hypotheses.} Actual positive four-face SU(2) square ground and literal retained trace w=tr(U2 U3)/2.; The particular synchronized four radial profiles S13 and their full Haar unitary generator, not an arbitrary Q8 field.; All conditional source regimes, including rare fibers and the antipodal degeneration, are retained in the requested uniform estimate.}\par
{\footnotesize \textit{Remaining.} Prove the actual full conditional score bound, or a sufficient direct source-weighted pairing estimate. The original M11-M15 consequence then becomes applicable;}\par
\smallskip
\noindent\textbf{\texttt{\small DERIV:\allowbreak{}W6\_\allowbreak{}GROUND\_\allowbreak{}JETS\_\allowbreak{}TRANSPORT\_\allowbreak{}BUDGET:\allowbreak{}INTERACTING\_\allowbreak{}GRID\_\allowbreak{}COMPARISON}} \hfill {\footnotesize open}\par
{\small For the actual interacting growing-grid family, obtain a volume-independent fast floor and complete source/transport energy estimates using linked or locally centered energy bounds, and control the coarse-theory matching discrepancy after the prescribed injection. The fixed-square constants depending on total ground energy E and the Gaussian grid floor do not prove this extension.}\par
{\footnotesize \srcpath{docs/derivations/w6-ground-jets-and-transport-budget.md} --- \#\# 6. Source energies, subdivision, interacting grids, and the physical clock (lines 351-366)}\par
{\footnotesize \textit{Hypotheses.} Actual coupled Wilson grids and specified block/source injections, with true vacuum subtraction, common physical energy normalization and the required retained, fast and spectral-energy domains.; Constants must be independent of growing spatial volume in the intended small-coupling regime; the existing fixed-square E and gamma are not assumed to have that uniformity.; Actual source-energy comparison factors and the discrepancy between the generated coarse form and prescribed coarse theory are retained; the Gaussian covariance-weighted source is not identified with the nonlinear quantum source without proof.}\par
{\footnotesize \textit{Remaining.} Prove the interacting fast-floor and connected energy/source comparison in the specified grid family. Controlling the local coupling remainder does not control coarse matching, the explicit physical-clock change in R14, or convergence of the exact transported jets and observables.}\par
\smallskip
\noindent\textbf{\texttt{\small DERIV:\allowbreak{}W6\_\allowbreak{}GROUND\_\allowbreak{}JETS\_\allowbreak{}TRANSPORT\_\allowbreak{}BUDGET:\allowbreak{}SOURCE\_\allowbreak{}ENERGY\_\allowbreak{}JETS\_\allowbreak{}R10}} \hfill {\footnotesize open}\par
{\small For the actual complete compact-square source/vacuum generator A(g) in R9, prove ||A\textasciicircum{}(r)(g)||\_\allowbreak{}(q\_\allowbreak{}g->q\_\allowbreak{}g)<=d\_\allowbreak{}r g\textasciicircum{}(-r-1) for r=0,1,2 with finite constants independent of small positive g. This is the open model input R10, not the already proved conditional implication R10 to R12.}\par
{\footnotesize \srcpath{docs/derivations/w6-ground-jets-and-transport-budget.md} --- \#\# 4. An explicit bridge from the actual transport generator to M1,M2,M3 (lines 199-218); section 5 (lines 271-309)}\par
{\footnotesize \textit{Hypotheses.} Actual fixed compact-square Hamiltonian and common electric H1 form domain; actual simple normalized positive ground, 0<=e\_\allowbreak{}g<=E and physical gap gamma>0 on 0<g<g\_\allowbreak{}*.; P\_\allowbreak{}g is the complete true-ground literal-source projection from C1, and A(g) is exactly the R9 source/vacuum generator; q\_\allowbreak{}g[u]=H\_\allowbreak{}g[u]+gamma||u||\textasciicircum{}2, including vacuum components.; Parameter derivatives are taken at positive g on the fixed compact form domain. The target covers all source and fast energies, including spatial derivatives and generator derivatives through order two.}\par
{\footnotesize \textit{Remaining.} Establish the conditional-source projection energy bounds beyond the proved vacuum-only generator. A direct proof of mixed-form R11 is an alternative sufficient route and need not prove R10 itself.}\par
\smallskip
\noindent\textbf{\texttt{\small DERIV:\allowbreak{}W6\_\allowbreak{}SOURCE\_\allowbreak{}ENERGY\_\allowbreak{}JETS\_\allowbreak{}R10:\allowbreak{}H0\_\allowbreak{}ENERGY\_\allowbreak{}CONTRACTION}} \hfill {\footnotesize open}\par
{\small Establish the proposed bound b\_\allowbreak{}g[Pi\_\allowbreak{}g Phi] <= q\_\allowbreak{}g[Omega\_\allowbreak{}g Phi] with kappa\_\allowbreak{}P=1 on the actual model; the weaker uniform finite-kappa\_\allowbreak{}P H0 also suffices for R10.}\par
{\footnotesize \srcpath{docs/derivations/w6-source-energy-jets-r10.md} --- \#\# 6. Scaling defects and unresolved estimates (C6)}\par
{\footnotesize \textit{Hypotheses.} Actual fixed compact square at positive g; normalized positive ground, q\_\allowbreak{}g=H\_\allowbreak{}g+gamma I, common form domain.; Constants must be independent of g on 0<g<g\_\allowbreak{}*.}\par
{\footnotesize \textit{Remaining.} The submitted proof is not a discharge. C1-C5 repair identities and conditional reductions;}\par
\smallskip
\noindent\textbf{\texttt{\small DERIV:\allowbreak{}W6\_\allowbreak{}SOURCE\_\allowbreak{}ENERGY\_\allowbreak{}JETS\_\allowbreak{}R10:\allowbreak{}H1\_\allowbreak{}FIBERWISE\_\allowbreak{}SCORE\_\allowbreak{}VARIANCE}} \hfill {\footnotesize open}\par
{\small Establish sup\_\allowbreak{}w Var\_\allowbreak{}mu\_\allowbreak{}g(sigma\_\allowbreak{}g|w) <= kappa\_\allowbreak{}0 g\textasciicircum{}-2 uniformly at small positive coupling.}\par
{\footnotesize \srcpath{docs/derivations/w6-source-energy-jets-r10.md} --- \#\# 6. Scaling defects and unresolved estimates (C6)}\par
{\footnotesize \textit{Hypotheses.} Actual fixed compact square at positive g; normalized positive ground, q\_\allowbreak{}g=H\_\allowbreak{}g+gamma I, common form domain.; Constants must be independent of g on 0<g<g\_\allowbreak{}*.}\par
{\footnotesize \textit{Remaining.} The submitted proof is not a discharge. C1-C5 repair identities and conditional reductions;}\par
\smallskip
\noindent\textbf{\texttt{\small DERIV:\allowbreak{}W6\_\allowbreak{}SOURCE\_\allowbreak{}ENERGY\_\allowbreak{}JETS\_\allowbreak{}R10:\allowbreak{}KATO\_\allowbreak{}DERIVATIVES}} \hfill {\footnotesize open}\par
{\small Establish ||partial\_\allowbreak{}g\textasciicircum{}r[P\_\allowbreak{}g prime,P\_\allowbreak{}g]||\_\allowbreak{}(q\_\allowbreak{}g->q\_\allowbreak{}g) <= d\_\allowbreak{}P,r g\textasciicircum{}(-r-1), r=0,1,2, for the actual source projection.}\par
{\footnotesize \srcpath{docs/derivations/w6-source-energy-jets-r10.md} --- \#\# 6. Scaling defects and unresolved estimates (C6)}\par
{\footnotesize \textit{Hypotheses.} Actual fixed compact square at positive g; normalized positive ground, q\_\allowbreak{}g=H\_\allowbreak{}g+gamma I, common form domain.; Constants must be independent of g on 0<g<g\_\allowbreak{}*.}\par
{\footnotesize \textit{Remaining.} The submitted proof is not a discharge. C1-C5 repair identities and conditional reductions;}\par
\smallskip
\noindent\textbf{\texttt{\small DERIV:\allowbreak{}W6\_\allowbreak{}SOURCE\_\allowbreak{}ENERGY\_\allowbreak{}JETS\_\allowbreak{}R10:\allowbreak{}VACUUM\_\allowbreak{}CROSS\_\allowbreak{}DERIVATIVES}} \hfill {\footnotesize open}\par
{\small Establish ||partial\_\allowbreak{}g\textasciicircum{}r A\_\allowbreak{}nu||\_\allowbreak{}(q\_\allowbreak{}g->q\_\allowbreak{}g) <= d\_\allowbreak{}nu,r g\textasciicircum{}(-r-1), r=0,1,2, for the actual vacuum-cross source vector.}\par
{\footnotesize \srcpath{docs/derivations/w6-source-energy-jets-r10.md} --- \#\# 6. Scaling defects and unresolved estimates (C6)}\par
{\footnotesize \textit{Hypotheses.} Actual fixed compact square at positive g; normalized positive ground, q\_\allowbreak{}g=H\_\allowbreak{}g+gamma I, common form domain.; Constants must be independent of g on 0<g<g\_\allowbreak{}*.}\par
{\footnotesize \textit{Remaining.} The submitted proof is not a discharge. C1-C5 repair identities and conditional reductions;}\par
\smallskip
\noindent\textbf{\texttt{\small DERIV:\allowbreak{}W6\_\allowbreak{}SOURCE\_\allowbreak{}GENERATOR\_\allowbreak{}SCORE\_\allowbreak{}FRAME:\allowbreak{}CONDITIONAL\_\allowbreak{}ENERGY\_\allowbreak{}H0\_\allowbreak{}H4}} \hfill {\footnotesize open}\par
{\small Prove, uniformly on 0<g<g\_\allowbreak{}*, the fiberwise conditional-energy estimates H0-H4 for the undilated ground score on level sets of w. Their averaged forms are established (SF7b and the integrated J\_\allowbreak{}g bound conditional on H4); the fiberwise forms are open.}\par
{\footnotesize \srcpath{docs/derivations/w6-source-generator-score-frame.md} --- \#\# 5. Exact reduction of the order-zero R10 bound (SF8-SF9), conditional-energy hypotheses}\par
{\footnotesize \textit{Hypotheses.} Actual fixed square, undilated score partial\_\allowbreak{}g log Omega\_\allowbreak{}g; K\textasciicircum{}0\_\allowbreak{}g is not the dilated K\_\allowbreak{}g of M10.}\par
{\footnotesize \textit{Remaining.} A fiberwise Poincare-type inequality for -Delta\_\allowbreak{}mu restricted to level sets of w, uniform in g, would give H1 and H2 from SF6. H0, H3 and H4 concern the source energy of conditional expectations and need a separate argument.}\par
\smallskip
\noindent\textbf{\texttt{\small DERIV:\allowbreak{}W6\_\allowbreak{}SYNCHRONIZED\_\allowbreak{}M10\_\allowbreak{}DOMINATION:\allowbreak{}AMPLITUDE\_\allowbreak{}GRADIENT\_\allowbreak{}SCALING}} \hfill {\footnotesize open}\par
{\small The actual uniform amplitude-score gradient bound ca/g remains open; changing parameter to h=g\textasciicircum{}2 and a spatial WKB expansion do not alone prove it.}\par
{\footnotesize \srcpath{docs/derivations/w6-synchronized-m10-domination.md} --- \#\# 5. Correct parameter-derivative criterion (lines 144-171)}\par
{\footnotesize \textit{Hypotheses.} Actual positive ground amplitude in a valid phase tube.}\par
{\footnotesize \textit{Remaining.} Establish uniform spatial control of h*partial\_\allowbreak{}h log A and Z log A for the actual ground; the logarithmic-parameter criterion gives the sufficient implication.}\par
\smallskip
\noindent\textbf{\texttt{\small DERIV:\allowbreak{}W6\_\allowbreak{}SYNCHRONIZED\_\allowbreak{}M10\_\allowbreak{}DOMINATION:\allowbreak{}FULL\_\allowbreak{}SCORE\_\allowbreak{}DOMINATION\_\allowbreak{}M10}} \hfill {\footnotesize open}\par
{\small Actual full-fiber M10 for S13 remains open. The earlier completion claim is retracted; the repaired assembly implication retains all unsupplied model hypotheses.}\par
{\footnotesize \srcpath{docs/derivations/w6-synchronized-m10-domination.md} --- \#\# 7. Assembly criterion and actual status (lines 262-286)}\par
{\footnotesize \textit{Hypotheses.} Actual four-face compact SU(2) quantum ground law and S13 dilation.}\par
{\footnotesize \textit{Remaining.} Actual complement is discharged. Supply the polynomial relative-amplitude and tube moment estimates, intermediate controls, and actual normal/angular comparison and score error.}\par
\smallskip
\noindent\textbf{\texttt{\small DERIV:\allowbreak{}WILSON\_\allowbreak{}MARKED\_\allowbreak{}TRANSFER:\allowbreak{}WT6}} \hfill {\footnotesize open}\par
{\small The proposed rooted coefficient-operator bound WT6 is an explicit sufficient route to complete-shell transport; this historical argument did not establish it. Its continuation uses a different full-Hilbert estimate and finite-order locality.}\par
{\footnotesize \srcpath{docs/derivations/wilson-marked-transfer.md} --- \#\# WT-6. The remaining complete-shell estimate and conditional transport (lines 209-248)}\par
{\footnotesize \textit{Hypotheses.} Finite cubic link Hilbert product; positive free generators kill the product vacuum and have gap delta>0 on its complement.; Bounded plaquette interactions ||V\_\allowbreak{}p||<=v; physical character-clock step tau>0 and blocks of duration ell=m*tau.; Actual vacuum-dressed transfer on the specified exact-support coefficient space must first be constructed.}\par
{\footnotesize \textit{Remaining.} Formalize or discharge the particular WT6 rooted norm if used; do not equate the later sufficient route with a proof of this specific estimate.}\par
\smallskip
\noindent\textbf{\texttt{\small DERIV:\allowbreak{}WILSON\_\allowbreak{}SC17\_\allowbreak{}SPATIAL\_\allowbreak{}CLOSURE:\allowbreak{}ACTUAL\_\allowbreak{}REFERENCE\_\allowbreak{}DEFECT\_\allowbreak{}BOUND}} \hfill {\footnotesize open}\par
{\small On the specified actual Wilson conditional source chart, establish a volume-uniform weighted bound on the complete signed Hessian defect D\_\allowbreak{}B=Hess(V\_\allowbreak{}B+(H\_\allowbreak{}out Omega)/Omega)-2epsilon A\textasciicircum{}2, including the actual kinetic metric, moving source projector and localization terms, with norm(D\_\allowbreak{}B)\_\allowbreak{}(s,infty)<=delta and 2 beta\_\allowbreak{}s(A)\textasciicircum{}2 delta/epsilon<1.}\par
{\footnotesize \srcpath{docs/derivations/wilson-sc17-spatial-closure.md} --- \#\# 4. The actual quantum pressure and cutoff defects are part of the input (lines 698-756), with the R5-R7 budget in section 2}\par
{\footnotesize \textit{Hypotheses.} Actual positive Wilson vacuum, compatible source and block coordinates, and multiplication terms assigned exactly once; the complete coupled Gaussian reference is subtracted.; The reference semigroup budget beta\_\allowbreak{}s(A) is finite in the chosen weighted norm at the fixed blocking scale; epsilon>0 and all conditional, metric and cutoff domains are retained.; The desired uniformity concerns the stated block family and surrounding volume; a separate source-compatible reference-angle margin is still needed for global assembly.}\par
{\footnotesize \textit{Remaining.} Prove the actual complete weighted smallness estimate. The proved Gaussian pressure cancellation, an isolated magnetic Taylor coefficient, or formalizing the conditional Riccati implication does not discharge this model obligation.}\par
\smallskip
\noindent\textbf{\texttt{\small DERIV:\allowbreak{}WILSON\_\allowbreak{}SELECTED\_\allowbreak{}INVERSE\_\allowbreak{}WALL:\allowbreak{}W6}} \hfill {\footnotesize open}\par
{\small The needed actual interacting Wilson bound is sup\_\allowbreak{}(Lambda,b[p]=1)|d\_\allowbreak{}g,QQ[R0 t1p,Rg t1p]|<=C|g| uniformly in volume, coupling, backgrounds and required spectral interval.}\par
{\footnotesize \srcpath{docs/derivations/wilson-selected-inverse-wall.md} --- \#\# 5. The exact point at which this argument cannot proceed (lines 221-264)}\par
{\footnotesize \textit{Hypotheses.} The operator chart, vacuum subtraction and fast/retained spaces are those of the spatial Schur derivation; domains are explicitly checked rather than identified by notation.; 0<g<=g0; complete relevant retained energy space and high-retained contribution controlled.}\par
{\footnotesize \textit{Remaining.} Prove the actual coupled signed resolvent pairing with the complete dressed force, physical energy normalization and an interacting fast floor; no existing Gaussian control discharges it.}\par
\smallskip
\noindent\textbf{\texttt{\small DERIV:\allowbreak{}YANGMILLS\_\allowbreak{}CONTINUUM\_\allowbreak{}BALABAN\_\allowbreak{}MULTISCALE\_\allowbreak{}PROOF:\allowbreak{}OS1}} \hfill {\footnotesize open}\par
{\small Full Euclidean covariance requires restoration of SO(4) beyond cubic block symmetry; the repair record ties this to the unproved irrelevance estimate.}\par
{\footnotesize \srcpath{docs/derivations/yangmills-continuum-balaban-multiscale-proof.md} --- \#\#\# 8.2 OS1: Euclidean Covariance (lines 375-382)}\par
{\footnotesize \textit{Hypotheses.} Submitted SU(N), N>=2 continuum multiscale construction; read together with its in-document Repair record audited 2026-09-09.; The source title and summary claim completion, but its repair record retains explicit averaging/interpolation, irrelevance, physical-scale and scattering obligations. No machine certification is inferred from prose.}\par
{\footnotesize \textit{Remaining.} Prove rotational restoration for the actual continuum limit, then translation/rotation covariance of its Schwinger distributions.}\par
\smallskip
\noindent\textbf{\texttt{\small DERIV:\allowbreak{}YANGMILLS\_\allowbreak{}CONTINUUM\_\allowbreak{}BALABAN\_\allowbreak{}MULTISCALE\_\allowbreak{}PROOF:\allowbreak{}THEOREM\_\allowbreak{}7\_\allowbreak{}1}} \hfill {\footnotesize open}\par
{\small The submitted Cauchy theorem needs a summable cross-scale expectation increment. The printed O(1/k) bound alone is not summable; the repair record names the missing irrelevance factor a\_\allowbreak{}k\textasciicircum{}theta.}\par
{\footnotesize \srcpath{docs/derivations/yangmills-continuum-balaban-multiscale-proof.md} --- \#\#\# 7.1 Cauchy Sequence of Effective Actions (lines 317-340)}\par
{\footnotesize \textit{Hypotheses.} Submitted SU(N), N>=2 continuum multiscale construction; read together with its in-document Repair record audited 2026-09-09.; The source title and summary claim completion, but its repair record retains explicit averaging/interpolation, irrelevance, physical-scale and scattering obligations. No machine certification is inferred from prose.; Actual gauge-invariant cylinder expectations and uniform operator normalization; a positive irrelevance exponent or another summable increment bound.}\par
{\footnotesize \textit{Remaining.} Prove the actual summable remainder after marginal matching, then construct the Cauchy limit with its observable norms and scale maps.}\par
\smallskip
\noindent\textbf{\texttt{\small DERIV:\allowbreak{}YANGMILLS\_\allowbreak{}CONTINUUM\_\allowbreak{}BALABAN\_\allowbreak{}MULTISCALE\_\allowbreak{}PROOF:\allowbreak{}THEOREM\_\allowbreak{}8\_\allowbreak{}2}} \hfill {\footnotesize open}\par
{\small The claimed continuum exponential clustering requires a strictly positive rate in common physical units at a fixed physical scale; UV cell-scale polymer decay alone does not provide that conclusion.}\par
{\footnotesize \srcpath{docs/derivations/yangmills-continuum-balaban-multiscale-proof.md} --- \#\#\# 8.5 OS4: Exponential Clustering and Mass Gap (lines 412-436)}\par
{\footnotesize \textit{Hypotheses.} Submitted SU(N), N>=2 continuum multiscale construction; read together with its in-document Repair record audited 2026-09-09.; The source title and summary claim completion, but its repair record retains explicit averaging/interpolation, irrelevance, physical-scale and scattering obligations. No machine certification is inferred from prose.; The repair record explicitly requires the missing bridge to a scale with g=O(1); uniform cutoff/volume estimates are retained.}\par
{\footnotesize \textit{Remaining.} Prove the actual nonperturbative KP or equivalent correlation estimate at a fixed physical reference scale, then transport it through the constructed continuum limit.}\par
\smallskip
\noindent\textbf{\texttt{\small DERIV:\allowbreak{}YANGMILLS\_\allowbreak{}CONTINUUM\_\allowbreak{}BALABAN\_\allowbreak{}MULTISCALE\_\allowbreak{}PROOF:\allowbreak{}GAUGE\_\allowbreak{}HESSIAN}} \hfill {\footnotesize disputed}\par
{\small The submission gives the background gauge Hessian (5.3)-(5.5); the operator must include the correct covariant 1-form curvature term and the repaired averaging map consistently.}\par
{\footnotesize \srcpath{docs/derivations/yangmills-continuum-balaban-multiscale-proof.md} --- \#\#\# 5.1 Fluctuation Gauge Fixing (lines 183-197)}\par
{\footnotesize \textit{Hypotheses.} Submitted SU(N), N>=2 continuum multiscale construction; read together with its in-document Repair record audited 2026-09-09.; The source title and summary claim completion, but its repair record retains explicit averaging/interpolation, irrelevance, physical-scale and scattering obligations. No machine certification is inferred from prose.; Smooth background and gauge-fixed fluctuation domain; repaired deterministic block constraint.}\par
{\footnotesize \textit{Remaining.} Formalize the covariant exterior-derivative/Weitzenbock identity with signs and domains; reconcile Q\_\allowbreak{}k versus Q\_\allowbreak{}B before using the Hessian.}\par
\smallskip
\noindent\textbf{\texttt{\small DERIV:\allowbreak{}YANGMILLS\_\allowbreak{}CONTINUUM\_\allowbreak{}BALABAN\_\allowbreak{}MULTISCALE\_\allowbreak{}PROOF:\allowbreak{}PROPOSITION\_\allowbreak{}3\_\allowbreak{}4}} \hfill {\footnotesize disputed}\par
{\small Proposition 3.4 asserts a unique gauge-equivariant constrained background minimizer with the Sobolev/coercivity bounds (3.9)-(3.12); the repair record says the corrected affine adjoint constraint and transported interpolation have not been propagated.}\par
{\footnotesize \srcpath{docs/derivations/yangmills-continuum-balaban-multiscale-proof.md} --- \#\#\# 3.2 The Background Field Variational Problem (lines 111-139)}\par
{\footnotesize \textit{Hypotheses.} Submitted SU(N), N>=2 continuum multiscale construction; read together with its in-document Repair record audited 2026-09-09.; The source title and summary claim completion, but its repair record retains explicit averaging/interpolation, irrelevance, physical-scale and scattering obligations. No machine certification is inferred from prose.; Small curvature g\_\allowbreak{}(k+1)\textasciicircum{}kappa a\_\allowbreak{}(k+1)\textasciicircum{}-2, 0<kappa<1/6, sufficiently small coupling; all correct gauge and boundary domains still required.}\par
{\footnotesize \textit{Remaining.} First implement the actual repaired deterministic constraint and equivariant interpolation, then prove coercivity, existence, uniqueness and smooth parameter dependence on its genuine domain.}\par
\smallskip
\noindent\textbf{\texttt{\small DERIV:\allowbreak{}YANGMILLS\_\allowbreak{}CONTINUUM\_\allowbreak{}BALABAN\_\allowbreak{}MULTISCALE\_\allowbreak{}PROOF:\allowbreak{}THEOREM\_\allowbreak{}5\_\allowbreak{}1}} \hfill {\footnotesize disputed}\par
{\small Theorem 5.1 asserts a uniform positive background fluctuation floor and exponential pointwise resolvent kernel bound. The actual averaging conjugation, covariant Poincare and kernel regularity estimates remain required.}\par
{\footnotesize \srcpath{docs/derivations/yangmills-continuum-balaban-multiscale-proof.md} --- \#\#\# 5.3 Combes--Thomas Resolvent Localization for Continuous \$\textbackslash{}mathcal\{H\}\_\allowbreak{}\{B\_\allowbreak{}k\}\$ (lines 209-248)}\par
{\footnotesize \textit{Hypotheses.} Submitted SU(N), N>=2 continuum multiscale construction; read together with its in-document Repair record audited 2026-09-09.; The source title and summary claim completion, but its repair record retains explicit averaging/interpolation, irrelevance, physical-scale and scattering obligations. No machine certification is inferred from prose.; delta<=1/(4224ell\textasciicircum{}2); alpha0>=8delta+1; the domain, cutoff and averaging operator must be specified.}\par
{\footnotesize \textit{Remaining.} Formalize the corrected covariant averaged operator and exponential conjugation including its nonlocal averaging term; replace or justify the all-x,y bounded continuum resolvent kernel claim, which has a diagonal singularity without an explicit smoothing cutoff.}\par
\smallskip
\noindent\textbf{\texttt{\small DERIV:\allowbreak{}YANGMILLS\_\allowbreak{}RESOLVENT\_\allowbreak{}LOCALIZATION\_\allowbreak{}AND\_\allowbreak{}DIRICHLET\_\allowbreak{}GAP:\allowbreak{}SUBMITTED\_\allowbreak{}W6\_\allowbreak{}CLOSURE}} \hfill {\footnotesize disputed}\par
{\small The proposal asserts actual Wilson W6 closure from massive spatial Green decay; its own source audit records that the computed massive scalar operator is not identified with Wilson.}\par
{\footnotesize \srcpath{docs/derivations/yangmills-resolvent-localization-and-dirichlet-gap.md} --- \#\# 1. Summary of Results (lines 20-39)}\par
{\footnotesize \textit{Hypotheses.} This submitted proposal is explicitly preserved under its September 9 source audit; closure claims are not the current graph verdict.; The actual model/operator identification is an obligation, separate from a scalar finite replay.}\par
{\footnotesize \textit{Remaining.} Establish the actual Wilson fast floor, connected signed pairing and physical source identification; the submission does not discharge these.}\par
\smallskip
\noindent\textbf{\texttt{\small DERIV:\allowbreak{}EARLIER\_\allowbreak{}PROOF\_\allowbreak{}RECOVERY:\allowbreak{}B6\_\allowbreak{}RETAINED\_\allowbreak{}KRYLOV\_\allowbreak{}CAMPAIGN}} \hfill {\footnotesize conditional}\par
{\small The retained SU(3) B6 certificate reports radial dimensions 67,155,133 and maximum full-space Ritz residual 8.775675697349887e-14 on 201 g-values in [1,2], with crossing g=1.3039546641713.}\par
{\footnotesize \srcpath{docs/derivations/earlier-proof-recovery.md} --- \#\# E6b. Retained B6 numerical campaign}\par
{\footnotesize \textit{Hypotheses.} Historical source K6(u)=E-uM, source coupling conversion and branch choices; level-eight 511-column word spaces; relative SVD cutoff 1e-10.; The certificate is retained and hash-verified; the B6 eigensolver was not rerun in this integration.}\par
{\footnotesize \textit{Remaining.} The original numerical certificate is preserved. A fresh B6 eigensolver replay and any uniform-in-coupling or full-shell conclusion are not supplied.}\par
\smallskip
\noindent\textbf{\texttt{\small DERIV:\allowbreak{}W6\_\allowbreak{}CONDITIONAL\_\allowbreak{}SCORE\_\allowbreak{}TAIL\_\allowbreak{}CONTROL:\allowbreak{}SCORE\_\allowbreak{}WEIGHT\_\allowbreak{}FROM\_\allowbreak{}M10}} \hfill {\footnotesize conditional}\par
{\small Conditional on the actual M10 domination, every centered finite-energy source satisfies integral K\_\allowbreak{}g|f|\textasciicircum{}2 dnu\_\allowbreak{}g <= [C1+(C0+C1 E)/gamma] b\_\allowbreak{}g[f], and the actual median Hardy constant obeys mathfrak\_\allowbreak{}B\_\allowbreak{}g<=C1+2(C0+C1 E)/gamma uniformly at small coupling. The implication is analytically proved; its M10 model input remains open.}\par
{\footnotesize \srcpath{docs/derivations/w6-conditional-score-tail-control.md} --- \#\#\# M.4. The precise remaining conditional-score inequality and M.5 (lines 497-563)}\par
{\footnotesize \textit{Hypotheses.} Actual compact-square source form b\_\allowbreak{}g[f]=g\textasciicircum{}2 integral (1-w\textasciicircum{}2)|f'|\textasciicircum{}2 dnu\_\allowbreak{}g, actual ground-energy bound e\_\allowbreak{}g<=E and physical gap gamma>0, with the true-ground source identity M7.; The actual conditional-score domination M10 holds with finite nonnegative C0,C1 uniform in g.; Sources belong to the full compact retained form domain; centered sources give M12, and the median-anchored half-source argument retains the uncentered vacuum-energy term for M14-M15.}\par
{\footnotesize \textit{Remaining.} Formalize this already established conditional implication if needed. Actual application requires M10 or another sufficient score estimate;}\par
\smallskip
\noindent\textbf{\texttt{\small DERIV:\allowbreak{}W6\_\allowbreak{}CONDITIONAL\_\allowbreak{}TRANSPORT\_\allowbreak{}OBSTRUCTION:\allowbreak{}LOCAL\_\allowbreak{}GLOBAL\_\allowbreak{}SCORE\_\allowbreak{}CRITERION}} \hfill {\footnotesize conditional}\par
{\small The stated uniform conditional tube moments, differentiated amplitude-score bound and synchronized phase remainder imply the local S14 score estimate. S15 glues it to the full conditional variance using the actual outside-tube squared deviation from the same reference beta\_\allowbreak{}g. Uniform potential-weighted complement and remaining-regime bounds would supply synchronized M10.}\par
{\footnotesize \srcpath{docs/derivations/w6-conditional-transport-obstruction.md} --- \#\# S6 Remaining all-fiber estimate (lines 317-375)}\par
{\footnotesize \textit{Hypotheses.} Conditional tube moments E|eta|\textasciicircum{}(2j)<=c\_\allowbreak{}(2j)g\textasciicircum{}(2j), j=1,2,3, with uniform constants.; True relative amplitude score has fast derivative at most c\_\allowbreak{}a/g; the phase is smooth with bounded third fast derivatives and the stated oscillator quadratic jet.; Synchronized tangency holds. The same reference beta\_\allowbreak{}g is used inside and outside the tube, so conditional mean differences remain in the outside deviation.; Full M10 additionally needs a uniform potential-weighted outside deviation bound and a degeneration-uniform actual score bound in the remaining source regimes.}\par
{\footnotesize \textit{Remaining.} Establish the actual uniform differentiated amplitude, tube-moment, conditional complement and antipodal estimates. Undifferentiated local comparison used for the obstruction does not discharge these hypotheses.}\par
\smallskip
\noindent\textbf{\texttt{\small DERIV:\allowbreak{}W6\_\allowbreak{}GROUND\_\allowbreak{}JETS\_\allowbreak{}TRANSPORT\_\allowbreak{}BUDGET:\allowbreak{}COMPLETE\_\allowbreak{}BUDGET}} \hfill {\footnotesize conditional}\par
{\small For the complete actual positive-reference source/vacuum transport, energy-operator bounds ||A\textasciicircum{}(r)(g)||\_\allowbreak{}(q\_\allowbreak{}g->q\_\allowbreak{}g)<=d\_\allowbreak{}r g\textasciicircum{}(-r-1), r=0,1,2, imply explicit all-energy M\_\allowbreak{}j(s)<=c\_\allowbreak{}j s\textasciicircum{}-j for j=1,2,3 and hence the complete compact Schur residual budget. The composite budget retains changing source energies and physical clock factors.}\par
{\footnotesize \srcpath{docs/derivations/w6-ground-jets-and-transport-budget.md} --- \#\# 4. An explicit bridge from the actual transport generator to M1,M2,M3 (lines 199-367)}\par
{\footnotesize \textit{Hypotheses.} Complete actual common-domain source/vacuum transport R'=A R and nonreducing graph from the positive-reference source theorem.; For q\_\allowbreak{}g=H\_\allowbreak{}g+gamma, full energy-space ||A\textasciicircum{}(r)(g)||<=d\_\allowbreak{}r g\textasciicircum{}(-r-1), r=0,1,2, with finite uniform constants.; Local t in [s/2,s], fixed spectral window below the actual gap. The composite source growth and subdivision inputs are the prior S2-S5 theorem.; Positive physical clock tau\_\allowbreak{}n; physical matching discrepancies are separately retained, not assumed zero.}\par
{\footnotesize \textit{Remaining.} Prove the actual complete conditional source-generator energy jets or exact mixed-form substitute; supply interacting-grid, coarse-matching and physical-clock convergence inputs.}\par
\smallskip
\noindent\textbf{\texttt{\small DERIV:\allowbreak{}W6\_\allowbreak{}SOURCE\_\allowbreak{}GENERATOR\_\allowbreak{}SCORE\_\allowbreak{}FRAME:\allowbreak{}R10\_\allowbreak{}ORDER\_\allowbreak{}ZERO\_\allowbreak{}REDUCTION}} \hfill {\footnotesize conditional}\par
{\small Under the conditional-energy hypotheses H0-H4 (uniform q\_\allowbreak{}g-boundedness of P\_\allowbreak{}g, sup\_\allowbreak{}w K\textasciicircum{}0\_\allowbreak{}g <= kappa\_\allowbreak{}0 g\textasciicircum{}-2, J\_\allowbreak{}g <= kappa\_\allowbreak{}1 g\textasciicircum{}-2 (1+g\textasciicircum{}-2 E[V|w]), the fast-to-source conditional energy bound, and b\_\allowbreak{}g[m\_\allowbreak{}g] <= kappa\_\allowbreak{}3 g\textasciicircum{}-2), the complete R9 generator satisfies ||A(g)||\_\allowbreak{}(q\_\allowbreak{}g->q\_\allowbreak{}g) <= d\_\allowbreak{}0 g\textasciicircum{}-1 with the\,\ldots}\par
{\footnotesize \srcpath{docs/derivations/w6-source-generator-score-frame.md} --- \#\# 5. Exact reduction of the order-zero R10 bound (SF8-SF9)}\par
{\footnotesize \textit{Hypotheses.} H0-H4 with constants independent of g on 0<g<g\_\allowbreak{}*; each is stated explicitly in section 5.; The rank-one q\_\allowbreak{}g operator bound and R4 from the transport-budget note; 0<=e\_\allowbreak{}g<=E and gap gamma>0.}\par
{\footnotesize \textit{Remaining.} Prove H0-H4. The r=1,2 cases of R10 are not reduced;}\par
\smallskip
\noindent\textbf{\texttt{\small DERIV:\allowbreak{}WILSON\_\allowbreak{}BACKGROUND\_\allowbreak{}COVARIANCE\_\allowbreak{}CONTINUATION:\allowbreak{}BF1\_\allowbreak{}BF3}} \hfill {\footnotesize conditional}\par
{\small A common fast form with relative defect norm theta\_\allowbreak{}z<1 gives the energy-weighted inverse difference <=theta\_\allowbreak{}z/(1-theta\_\allowbreak{}z), the selected signed pairing bound BF1, and the exact energy-dependent Schur/source identity BF3.}\par
{\footnotesize \srcpath{docs/derivations/wilson-background-covariance-continuation.md} --- \#\# BF1. The energy-dependent relative-form statement (lines 48-102)}\par
{\footnotesize \textit{Hypotheses.} Finite periodic cubic SU(2) lattice with nondegenerate plaquettes and side L >= 3.; Product unit-S3 electric metric; epsilon > 0, k >= 0, lambda = k/epsilon; actual positive normalized Wilson ground Omega.; A0(z)=F0-z>=f\_\allowbreak{}z I>0; V is a bounded A0-relative form with norm <=theta\_\allowbreak{}z<1.}\par
{\footnotesize \textit{Remaining.} The normalized bounded inverse, selected pairing and Schur minimization are formalized. Identify the actual common closed form, inverse-square-root coordinates and compatible cross-functionals for the Wilson operator.}\par
\smallskip
\noindent\textbf{\texttt{\small DERIV:\allowbreak{}WILSON\_\allowbreak{}BACKGROUND\_\allowbreak{}COVARIANCE\_\allowbreak{}CONTINUATION:\allowbreak{}SC17}} \hfill {\footnotesize conditional}\par
{\small The exact distant mixed derivative is -R\_\allowbreak{}xy(F\_\allowbreak{}e tensor V\_\allowbreak{}f+V\_\allowbreak{}e tensor F\_\allowbreak{}f)-F\_\allowbreak{}e tensor F\_\allowbreak{}f; the original all-coupling proof does not supply a summable distance envelope.}\par
{\footnotesize \srcpath{docs/derivations/wilson-background-covariance-continuation.md} --- \#\# 5. The neutral channel and the exact remaining spatial bound (lines 494-534)}\par
{\footnotesize \textit{Hypotheses.} Finite periodic cubic SU(2) lattice with nondegenerate plaquettes and side L >= 3.; Product unit-S3 electric metric; epsilon > 0, k >= 0, lambda = k/epsilon; actual positive normalized Wilson ground Omega.; Distinct distant tails; V\_\allowbreak{}ef=0; R\_\allowbreak{}xy is the actual charged inverse.}\par
{\footnotesize \textit{Remaining.} Formalize this identity and the singlet obstruction; the all-coupling spatially summable estimate still needs its complete quantum reference-defect bound.}\par
\smallskip
\noindent\textbf{\texttt{\small DERIV:\allowbreak{}WILSON\_\allowbreak{}MARKED\_\allowbreak{}SHELL\_\allowbreak{}TRANSPORT:\allowbreak{}S18\_\allowbreak{}S20}} \hfill {\footnotesize conditional}\par
{\small The common weighted Taylor majorant sums finite-order Wilson/Hamiltonian matching; the relative shell gap stays >(t3/4)u\textasciicircum{}2 q(k) away from Gamma and carrier projections converge by S20.}\par
{\footnotesize \srcpath{docs/derivations/wilson-marked-shell-transport.md} --- \#\# 6. Sum the matching and transport the carrier (lines 391-461)}\par
{\footnotesize \textit{Hypotheses.} SU(3), fixed spatial spacing and small magnetic coupling on the S1 common domain.; Calibrated kinetic-window and the actual September 5 Hamiltonian G18 creator/activity/complete-band hypotheses and literal source charts are explicit inputs.; The Hamiltonian G18 gap and kinetic-window coefficient comparison apply on the same weighted literal-source frame; epsilon is sufficiently small.}\par
{\footnotesize \textit{Remaining.} Formalize power-series tail convergence, the common polar source chart, relative spectral perturbation and exact symmetry cancellation at Gamma.}\par
\smallskip
\noindent\textbf{\texttt{\small DERIV:\allowbreak{}WILSON\_\allowbreak{}SC17\_\allowbreak{}SPATIAL\_\allowbreak{}CLOSURE:\allowbreak{}PROPOSED\_\allowbreak{}QUANTUM\_\allowbreak{}PARAMETER\_\allowbreak{}WINDOW}} \hfill {\footnotesize conditional}\par
{\small Under the proposed numerical reference and normalized-defect inputs, the scalar Riccati parameter satisfies p(lambda)<=p(3/100)=96228*sqrt(3)/625000<27/100<1, and the small root satisfies 0<=r\_\allowbreak{}minus<a\_\allowbreak{}min throughout 0<=lambda<=3/100. These are conditional parameter conclusions; the section's complete Wilson defect estimate and exponential spatial decay are not established by this arithmetic.}\par
{\footnotesize \srcpath{docs/derivations/wilson-sc17-spatial-closure.md} --- \#\# 6. Formal repair of the spatial closure for lambda >= 1/73 via the quantum reference defect (lines 781-821)}\par
{\footnotesize \textit{Hypotheses.} Real coupling 0<=lambda<=3/100; epsilon>0 when interpreting the normalized defect c=delta/(2epsilon).; The companion script's proposed defaults are L=3, C\_\allowbreak{}s=6/5 and reference velocity one; beta=(108/25)*sqrt(33), c=lambda*sqrt(lambda)/48 and a\_\allowbreak{}min=1/(3*sqrt(33)).; p(lambda)=4*beta\textasciicircum{}2*c and r\_\allowbreak{}minus=(1-sqrt(1-p(lambda)))/(2*beta).; These numerical inputs are supplied parameters. Their identification with a uniform semigroup budget, the complete conditional quantum-pressure defect and the actual reference floor is a separate model obligation, not a consequence of the quartic Taylor coefficient.}\par
{\footnotesize \textit{Remaining.} The scalar comparison is proved in the supported successor SC17\_\allowbreak{}P. To apply it to the source section's spatial-closure claim, establish the actual weighted reference budget and complete outside-pressure, metric, projector and cutoff defect bound, derive the ground-state Riccati comparison, and justify an exponential spatial weight.}\par
\smallskip
\noindent\textbf{\texttt{\small DERIV:\allowbreak{}WILSON\_\allowbreak{}SC17\_\allowbreak{}SPATIAL\_\allowbreak{}CLOSURE:\allowbreak{}R10\_\allowbreak{}R11}} \hfill {\footnotesize conditional}\par
{\small Uniform block Hessian floors 2a\_\allowbreak{}B,2a\_\allowbreak{}C and cross norm 2h\_\allowbreak{}BC yield c\_\allowbreak{}BC<=h\_\allowbreak{}BC/sqrt(a\_\allowbreak{}B a\_\allowbreak{}C); a reference angle budget plus full Hessian error r yields physical gap >=2epsilon[a0(1-kappa\_\allowbreak{}A)-2r].}\par
{\footnotesize \srcpath{docs/derivations/wilson-sc17-spatial-closure.md} --- \#\# 3. An explicit true-angle consequence if the full Hessian defect is known (lines 649-697)}\par
{\footnotesize \textit{Hypotheses.} Finite periodic cubic SU(2) lattice with nondegenerate plaquettes and side L >= 3.; Product unit-S3 electric metric; epsilon > 0, k >= 0, lambda = k/epsilon; actual positive normalized Wilson ground Omega.; a\_\allowbreak{}B,a\_\allowbreak{}C>0; h\_\allowbreak{}BC\textasciicircum{}2<a\_\allowbreak{}B a\_\allowbreak{}C; kappa\_\allowbreak{}A<1; 2r<a0(1-kappa\_\allowbreak{}A).}\par
{\footnotesize \textit{Remaining.} Formalize Brascamp-Lieb conditional variance, marginal Schur curvature and the actual projection-angle assembly.}\par
\smallskip
\noindent\textbf{\texttt{\small DERIV:\allowbreak{}WILSON\_\allowbreak{}SC17\_\allowbreak{}SPATIAL\_\allowbreak{}CLOSURE:\allowbreak{}R12\_\allowbreak{}R14}} \hfill {\footnotesize conditional}\par
{\small The necessary full conditional Hessian defect is Hess(V\_\allowbreak{}B+((H\_\allowbreak{}out Omega)/Omega))-2epsilon A\textasciicircum{}2; exact Gaussian pressure cancels the reference BB* term, while cutoffs create the explicit gradient and commutator defects.}\par
{\footnotesize \srcpath{docs/derivations/wilson-sc17-spatial-closure.md} --- \#\# 4. The actual quantum pressure and cutoff defects are part of the input (lines 698-756)}\par
{\footnotesize \textit{Hypotheses.} Finite periodic cubic SU(2) lattice with nondegenerate plaquettes and side L >= 3.; Product unit-S3 electric metric; epsilon > 0, k >= 0, lambda = k/epsilon; actual positive normalized Wilson ground Omega.; Actual vacuum and compatible Euclidean block chart; all multiplication terms assigned once.}\par
{\footnotesize \textit{Remaining.} Formalize the outside-pressure derivative, Gaussian cancellation and every cutoff/product-domain term; the actual uniform smallness estimate remains unproved.}\par
\smallskip
\noindent\textbf{\texttt{\small DERIV:\allowbreak{}WILSON\_\allowbreak{}SC17\_\allowbreak{}SPATIAL\_\allowbreak{}CLOSURE:\allowbreak{}R4\_\allowbreak{}R9}} \hfill {\footnotesize conditional}\par
{\small If beta\_\allowbreak{}s(A)<infinity and the full conditional potential Hessian defect is <=delta with 2 beta\_\allowbreak{}s(A)\textasciicircum{}2 delta/epsilon<1, the true log-ground Hessian error is <=r\_\allowbreak{}minus and conditional gap >=2epsilon(a-r\_\allowbreak{}minus).}\par
{\footnotesize \srcpath{docs/derivations/wilson-sc17-spatial-closure.md} --- \#\# 2. A rigorous conditional small-defect criterion (lines 574-648)}\par
{\footnotesize \textit{Hypotheses.} Finite periodic cubic SU(2) lattice with nondegenerate plaquettes and side L >= 3.; Product unit-S3 electric metric; epsilon > 0, k >= 0, lambda = k/epsilon; actual positive normalized Wilson ground Omega.; A>=aI>0; the full conditional potential includes outside quantum pressure; smooth ground convergence and domain conditions in the section hold.}\par
{\footnotesize \textit{Remaining.} Formalize the noncommutative weighted norm, Duhamel Riccati comparison, ground limit and curvature-to-gap passage.}\par
\smallskip
\noindent\textbf{\texttt{\small DERIV:\allowbreak{}WILSON\_\allowbreak{}SELECTED\_\allowbreak{}INVERSE\_\allowbreak{}WALL:\allowbreak{}W4\_\allowbreak{}W5}} \hfill {\footnotesize conditional}\par
{\small On admissible reference graph vectors, S\_\allowbreak{}g-S\_\allowbreak{}0=A\_\allowbreak{}g-r\_\allowbreak{}g*R\_\allowbreak{}g r\_\allowbreak{}g; only the selected inverse bound W5 plus direct/force remainders is sufficient for the displayed O(g\textasciicircum{}3) comparison.}\par
{\footnotesize \srcpath{docs/derivations/wilson-selected-inverse-wall.md} --- \#\# 3. Repair: only compare inverses on the graph force (lines 101-157)}\par
{\footnotesize \textit{Hypotheses.} The operator chart, vacuum subtraction and fast/retained spaces are those of the spatial Schur derivation; domains are explicitly checked rather than identified by notation.; z lies below both full fast spectra; J\_\allowbreak{}z p and the interacting minimizer are admissible; W5 includes f>0,a3,t2,t0,c\_\allowbreak{}sel.}\par
{\footnotesize \textit{Remaining.} The bounded Hilbert-space completion/minimum and weighted pairing are formalized. Construct the actual unbounded fast-form domain, source identification and reduced closed-form Schur operator.}\par
\smallskip
\noindent\textbf{\texttt{\small DERIV:\allowbreak{}WILSON\_\allowbreak{}SPATIAL\_\allowbreak{}SCHUR\_\allowbreak{}EXCESS:\allowbreak{}SP1\_\allowbreak{}SP6}} \hfill {\footnotesize conditional}\par
{\small The exact full-reference Schur excess is A\_\allowbreak{}g-V\_\allowbreak{}g*(I+C\_\allowbreak{}g)\textasciicircum{}-1 V\_\allowbreak{}g; the minimizing graph, two-sided form bounds and -dS\_\allowbreak{}g/dz=J\_\allowbreak{}g,z*J\_\allowbreak{}g,z retain the induced metric.}\par
{\footnotesize \srcpath{docs/derivations/wilson-spatial-schur-excess.md} --- \#\# 2. Exact excess about the full reference graph (lines 42-114)}\par
{\footnotesize \textit{Hypotheses.} Closed self-adjoint reference and interacting forms on the stated compatible P/Q domains, already expressed in a common physical source chart and vacuum-subtracted.; Real energy z below the full fast restriction; positive retained energy weight B.; F\_\allowbreak{}z>=f\_\allowbreak{}zI>0; ||C\_\allowbreak{}g||<=theta<1; all weighted A\_\allowbreak{}g,V\_\allowbreak{}g extend boundedly.}\par
{\footnotesize \textit{Remaining.} The normalized inverse, positive fast form and unique bounded Schur minimum are formalized. Identify actual Wilson form domains and sources;}\par
\smallskip
\noindent\textbf{\texttt{\small DERIV:\allowbreak{}WILSON\_\allowbreak{}SPATIAL\_\allowbreak{}SCHUR\_\allowbreak{}EXCESS:\allowbreak{}SP20\_\allowbreak{}SP24}} \hfill {\footnotesize conditional}\par
{\small Complete Schur comparisons give the reciprocal-gap recursion, product/inverse-fast-energy lower bound SP21-SP22 and source-overlap telescoping SP24 in common physical clock units.}\par
{\footnotesize \srcpath{docs/derivations/wilson-spatial-schur-excess.md} --- \#\# 6. Iterating energies and observable amplitudes in the physical clock (lines 412-477)}\par
{\footnotesize \textit{Hypotheses.} Closed self-adjoint reference and interacting forms on the stated compatible P/Q domains, already expressed in a common physical source chart and vacuum-subtracted.; Real energy z below the full fast restriction; positive retained energy weight B.; 0<alpha\_\allowbreak{}j<=1; full fast floors f\_\allowbreak{}j; summable epsilon\_\allowbreak{}j=1-alpha\_\allowbreak{}j bounded below 1; summable inverse fast floors; actual source maps and errors eta\_\allowbreak{}j.}\par
{\footnotesize \textit{Remaining.} Formalize iteration of the closed-form gap theorem, positive infinite products and complete carrier/source projection transport.}\par
\smallskip
\noindent\textbf{\texttt{\small DERIV:\allowbreak{}WILSON\_\allowbreak{}SPATIAL\_\allowbreak{}SCHUR\_\allowbreak{}EXCESS:\allowbreak{}SP7\_\allowbreak{}SP10}} \hfill {\footnotesize conditional}\par
{\small K1=J\_\allowbreak{}z*W1J\_\allowbreak{}z and K2=J\_\allowbreak{}z*W2J\_\allowbreak{}z-V1*V1 are the complete Schur coefficients; SP8 gives the explicit O(|g|\textasciicircum{}3) remainder SP9, including the dressed force correction SP10.}\par
{\footnotesize \srcpath{docs/derivations/wilson-spatial-schur-excess.md} --- \#\# 3. The complete second-order coefficient and a uniform remainder (lines 115-182)}\par
{\footnotesize \textit{Hypotheses.} Closed self-adjoint reference and interacting forms on the stated compatible P/Q domains, already expressed in a common physical source chart and vacuum-subtracted.; Real energy z below the full fast restriction; positive retained energy weight B.; SP8 remainder constants a3,v2,v1,c1,g0 are finite and c1|g|<=theta<1.}\par
{\footnotesize \textit{Remaining.} The genuine Neumann operator expansion and cubic norm/sandwich remainder are formalized. Assemble the g-dependent direct term, moving forces, fast defect and spatial coefficients on the actual compatible Wilson form domain.}\par
\smallskip
\noindent\textbf{\texttt{\small DERIV:\allowbreak{}WILSON\_\allowbreak{}TRUE\_\allowbreak{}VACUUM\_\allowbreak{}BLOCK\_\allowbreak{}ESTIMATES:\allowbreak{}BA9\_\allowbreak{}BA10}} \hfill {\footnotesize conditional}\par
{\small The true conditional operator includes outside pressure W\_\allowbreak{}B=(H\_\allowbreak{}out Omega)/Omega and has gap >=gap(h\_\allowbreak{}B(b))-osc\_\allowbreak{}B W\_\allowbreak{}B; a frozen-block gap alone omits this term.}\par
{\footnotesize \srcpath{docs/derivations/wilson-true-vacuum-block-estimates.md} --- \#\# BA4. What a frozen-block spectral computation would additionally require (lines 142-168)}\par
{\footnotesize \textit{Hypotheses.} Finite periodic cubic SU(2) lattice with nondegenerate plaquettes and side L >= 3.; Product unit-S3 electric metric; epsilon > 0, k >= 0, lambda = k/epsilon; actual positive normalized Wilson ground Omega.; All conditional blocks use the actual vacuum law mu=Omega\textasciicircum{}2dU, not isolated-block spectral data.; The frozen-block gap and finite pressure oscillation are available on compatible form domains.}\par
{\footnotesize \textit{Remaining.} Formalize the conditional ground transform and bounded-potential min-max comparison, including the domain of H\_\allowbreak{}out Omega.}\par
\smallskip
\noindent\textbf{\texttt{\small DERIV:\allowbreak{}WILSON\_\allowbreak{}VACUUM\_\allowbreak{}ALIGNED\_\allowbreak{}ASSEMBLY:\allowbreak{}VA12\_\allowbreak{}VA13}} \hfill {\footnotesize conditional}\par
{\small For physical conditional-angle row kappa<1, G\_\allowbreak{}mu\textasciicircum{}2>=(1-kappa)G\_\allowbreak{}mu and gap\_\allowbreak{}phys(H)>=gamma\_\allowbreak{}*(1-kappa); the single plaquette physical restriction has E\_\allowbreak{}e=E0 and zero pair angle.}\par
{\footnotesize \srcpath{docs/derivations/wilson-vacuum-aligned-assembly.md} --- \#\# VA5. Projection geometry, with the Gauss law retained (lines 211-251)}\par
{\footnotesize \textit{Hypotheses.} Finite periodic cubic SU(2) lattice with nondegenerate plaquettes and side L >= 3.; Product unit-S3 electric metric; epsilon > 0, k >= 0, lambda = k/epsilon; actual positive normalized Wilson ground Omega.; c\_\allowbreak{}ef\textasciicircum{}phys is the actual L2 conditional-projection norm, and sup\_\allowbreak{}e sum\_\allowbreak{}(f!=e)c\_\allowbreak{}ef\textasciicircum{}phys<=kappa<1.}\par
{\footnotesize \textit{Remaining.} The arbitrary-dimensional bounded projection-sum implication and kernel-complement gap are formalized. Construct actual conditional projections, derive the local anticommutator bound from their Friedrichs angles, identify the common kernel as constants and establish the required row budget below one.}\par
\smallskip
\noindent\textbf{\texttt{\small DERIV:\allowbreak{}WILSON\_\allowbreak{}VACUUM\_\allowbreak{}ALIGNED\_\allowbreak{}ASSEMBLY:\allowbreak{}VA14\_\allowbreak{}VA19}} \hfill {\footnotesize conditional}\par
{\small Approximate tensorization or a summable block-angle row gives the stated physical gap with covering multiplicity; the original VA10 local floor tends to zero on the displayed logarithmic continuum trajectory.}\par
{\footnotesize \srcpath{docs/derivations/wilson-vacuum-aligned-assembly.md} --- \#\# VA6. The next estimate and the attempts to establish it (lines 252-415)}\par
{\footnotesize \textit{Hypotheses.} Finite periodic cubic SU(2) lattice with nondegenerate plaquettes and side L >= 3.; Product unit-S3 electric metric; epsilon > 0, k >= 0, lambda = k/epsilon; actual positive normalized Wilson ground Omega.; Actual conditional blocks, gamma\_\allowbreak{}min>0 and finite covering multiplicity m\_\allowbreak{}*; a spatial envelope is an additional hypothesis, not assigned data.}\par
{\footnotesize \textit{Remaining.} Actual-measure covariance minorization and explicit-domain unitary gap transport are formalized. Establish the actual Wilson conditional density floor, conditional disintegration, covariance-to-projection-angle norm and spatial row-budget bounds.}\par
\smallskip
\noindent\textbf{\texttt{\small DERIV:\allowbreak{}WILSON\_\allowbreak{}WEIGHTED\_\allowbreak{}REPAIR\_\allowbreak{}AND\_\allowbreak{}ROTOR\_\allowbreak{}GAP:\allowbreak{}WR17\_\allowbreak{}WR21}} \hfill {\footnotesize conditional}\par
{\small The actual trial comparison needs gamma\_\allowbreak{}Lambda-(mean R-inf R)>0 uniformly; a gauge-invariant patch has full-boundary trial excitations of energy 3epsilon, so an unrestricted boundary fast floor of order sqrt(epsilon k) cannot be inferred.}\par
{\footnotesize \srcpath{docs/derivations/wilson-weighted-repair-and-rotor-gap.md} --- \#\# WR6. The precise remaining derivation (lines 229-308)}\par
{\footnotesize \textit{Hypotheses.} Each operator is identified separately: classical quotient diffusion, complete compact S3 quantum rotor, actual coupled link lattice or independent product model as stated in the section.; Positive normalized overlapping-plaquette trial; gamma\_\allowbreak{}Lambda is an actual weighted generator gap.}\par
{\footnotesize \textit{Remaining.} Formalize the common form domains and mean-residual comparison; construct Haar-link marginal trial vectors proving the full-boundary obstruction.}\par
\smallskip
\noindent\textbf{\texttt{\small DERIV:\allowbreak{}YANGMILLS\_\allowbreak{}CONTINUUM\_\allowbreak{}BALABAN\_\allowbreak{}MULTISCALE\_\allowbreak{}PROOF:\allowbreak{}GHOST\_\allowbreak{}5\_\allowbreak{}8}} \hfill {\footnotesize conditional}\par
{\small Rescaling z=g*zeta exposes g\textasciicircum{}(1-kappa\_\allowbreak{}prime) in (5.8), but the required volume/scale-uniform covariant inverse-gradient bound is not proved by that algebra.}\par
{\footnotesize \srcpath{docs/derivations/yangmills-continuum-balaban-multiscale-proof.md} --- \#\#\# 5.2 Faddeev--Popov Ghost Determinant (lines 198-208)}\par
{\footnotesize \textit{Hypotheses.} Submitted SU(N), N>=2 continuum multiscale construction; read together with its in-document Repair record audited 2026-09-09.; The source title and summary claim completion, but its repair record retains explicit averaging/interpolation, irrelevance, physical-scale and scattering obligations. No machine certification is inferred from prose.; 0<kappa\_\allowbreak{}prime<1/2; a||zeta||\_\allowbreak{}infinity<=g\textasciicircum{}-kappa\_\allowbreak{}prime; the first operator inequality in (5.8) is an additional input.}\par
{\footnotesize \textit{Remaining.} Construct the precise ghost operator, determinant class and uniform inverse-gradient estimate on its constrained domain; scalar rescaling alone does not discharge them.}\par
\smallskip
\noindent\textbf{\texttt{\small DERIV:\allowbreak{}YANGMILLS\_\allowbreak{}CONTINUUM\_\allowbreak{}BALABAN\_\allowbreak{}MULTISCALE\_\allowbreak{}PROOF:\allowbreak{}OS0}} \hfill {\footnotesize conditional}\par
{\small The asserted Schwinger growth bound (8.2) supplies OS regularity only after its actual uniform moment/Sobolev estimates and continuum fields are constructed.}\par
{\footnotesize \srcpath{docs/derivations/yangmills-continuum-balaban-multiscale-proof.md} --- \#\#\# 8.1 OS0: Analyticity and Temperedness (lines 368-374)}\par
{\footnotesize \textit{Hypotheses.} Submitted SU(N), N>=2 continuum multiscale construction; read together with its in-document Repair record audited 2026-09-09.; The source title and summary claim completion, but its repair record retains explicit averaging/interpolation, irrelevance, physical-scale and scattering obligations. No machine certification is inferred from prose.}\par
{\footnotesize \textit{Remaining.} Formalize the continuum Schwinger distributions and prove the factorial/test-seminorm bound from the actual regulated measures.}\par
\smallskip
\noindent\textbf{\texttt{\small DERIV:\allowbreak{}YANGMILLS\_\allowbreak{}CONTINUUM\_\allowbreak{}BALABAN\_\allowbreak{}MULTISCALE\_\allowbreak{}PROOF:\allowbreak{}OS3}} \hfill {\footnotesize conditional}\par
{\small Bosonic multiplication observables have permutation-symmetric Schwinger functions, provided the claimed continuum distributions exist.}\par
{\footnotesize \srcpath{docs/derivations/yangmills-continuum-balaban-multiscale-proof.md} --- \#\#\# 8.4 OS3: Permutation Symmetry (lines 409-411)}\par
{\footnotesize \textit{Hypotheses.} Submitted SU(N), N>=2 continuum multiscale construction; read together with its in-document Repair record audited 2026-09-09.; The source title and summary claim completion, but its repair record retains explicit averaging/interpolation, irrelevance, physical-scale and scattering obligations. No machine certification is inferred from prose.}\par
{\footnotesize \textit{Remaining.} Formalize products of the constructed fields/observables and the permutation action on their distributional arguments.}\par
\smallskip
\noindent\textbf{\texttt{\small DERIV:\allowbreak{}YANGMILLS\_\allowbreak{}CONTINUUM\_\allowbreak{}BALABAN\_\allowbreak{}MULTISCALE\_\allowbreak{}PROOF:\allowbreak{}OS\_\allowbreak{}RECONSTRUCTION}} \hfill {\footnotesize conditional}\par
{\small With corrected OS hypotheses and totality, the reflection-positive quotient completion has a physical Hilbert space, vacuum and nonnegative self-adjoint time generator.}\par
{\footnotesize \srcpath{docs/derivations/yangmills-continuum-balaban-multiscale-proof.md} --- \#\#\# 9.1 The Physical Hilbert Space and Hamiltonian (lines 439-453)}\par
{\footnotesize \textit{Hypotheses.} Submitted SU(N), N>=2 continuum multiscale construction; read together with its in-document Repair record audited 2026-09-09.; The source title and summary claim completion, but its repair record retains explicit averaging/interpolation, irrelevance, physical-scale and scattering obligations. No machine certification is inferred from prose.}\par
{\footnotesize \textit{Remaining.} Formalize the OS quotient, completion, translations and corrected reconstruction regularity; actual OS inputs remain separate obligations.}\par
\smallskip
\noindent\textbf{\texttt{\small DERIV:\allowbreak{}YANGMILLS\_\allowbreak{}CONTINUUM\_\allowbreak{}BALABAN\_\allowbreak{}MULTISCALE\_\allowbreak{}PROOF:\allowbreak{}SCATTERING}} \hfill {\footnotesize conditional}\par
{\small Haag-Ruelle scattering is conditional on an isolated positive single-particle mass shell; the source now explicitly states that dispersion hypothesis. Nonzero off-shell/tree-level four-point data alone do not supply the required on-shell amplitude lower bound.}\par
{\footnotesize \srcpath{docs/derivations/yangmills-continuum-balaban-multiscale-proof.md} --- \#\#\# 9.3 Relativistic Wightman Fields and Non-Trivial Scattering (\$S \textbackslash{}ne I\$) (lines 479-501)}\par
{\footnotesize \textit{Hypotheses.} Submitted SU(N), N>=2 continuum multiscale construction; read together with its in-document Repair record audited 2026-09-09.; The source title and summary claim completion, but its repair record retains explicit averaging/interpolation, irrelevance, physical-scale and scattering obligations. No machine certification is inferred from prose.; Isolated mass hyperboloid separated below multiparticle threshold; locality, spectrum and actual LSZ/Haag-Ruelle assumptions.}\par
{\footnotesize \textit{Remaining.} Formalize the isolated-shell hypothesis, asymptotic states and a controlled on-shell nonzero amplitude; establish these for the constructed theory before concluding S!=I.}\par
\smallskip
\noindent\textbf{\texttt{\small DERIV:\allowbreak{}YANGMILLS\_\allowbreak{}CONTINUUM\_\allowbreak{}BALABAN\_\allowbreak{}MULTISCALE\_\allowbreak{}PROOF:\allowbreak{}THEOREM\_\allowbreak{}4\_\allowbreak{}2}} \hfill {\footnotesize conditional}\par
{\small The large-field damping claim (4.4)-(4.5) is conditional on a positive coarse-action subtraction constant and the unproved irrelevance gain identified in the repair record.}\par
{\footnotesize \srcpath{docs/derivations/yangmills-continuum-balaban-multiscale-proof.md} --- \#\#\# 4.2 Super-Exponential Suppression of Large Fields (lines 157-180)}\par
{\footnotesize \textit{Hypotheses.} Submitted SU(N), N>=2 continuum multiscale construction; read together with its in-document Repair record audited 2026-09-09.; The source title and summary claim completion, but its repair record retains explicit averaging/interpolation, irrelevance, physical-scale and scattering obligations. No machine certification is inferred from prose.; 0<kappa\_\allowbreak{}prime<kappa<1/6; actual regularized functional measure and corrected background construction; uniform constants.}\par
{\footnotesize \textit{Remaining.} Prove the quantitative irrelevance/subtraction constant and the actual large-field measure estimate, including the averaging constraint and fluctuation direction invisible to a coarse penalty.}\par
\smallskip
\noindent\textbf{\texttt{\small DERIV:\allowbreak{}YANGMILLS\_\allowbreak{}CONTINUUM\_\allowbreak{}BALABAN\_\allowbreak{}MULTISCALE\_\allowbreak{}PROOF:\allowbreak{}THEOREM\_\allowbreak{}6\_\allowbreak{}2}} \hfill {\footnotesize conditional}\par
{\small The claimed RG contraction (6.5) has the repaired induction envelope ||K\_\allowbreak{}k||<=4 C\_\allowbreak{}ind g\_\allowbreak{}k\textasciicircum{}2, but its single-block contraction gain still requires the irrelevance lemma.}\par
{\footnotesize \srcpath{docs/derivations/yangmills-continuum-balaban-multiscale-proof.md} --- \#\#\# 6.3 The Mayer Tree-Graph Expansion \& Cluster Contraction (lines 269-295)}\par
{\footnotesize \textit{Hypotheses.} Submitted SU(N), N>=2 continuum multiscale construction; read together with its in-document Repair record audited 2026-09-09.; The source title and summary claim completion, but its repair record retains explicit averaging/interpolation, irrelevance, physical-scale and scattering obligations. No machine certification is inferred from prose.; lambda\_\allowbreak{}contraction in (0,1/2); sufficiently small initial g0; the corrected averaging, fluctuation and large-field inputs hold.}\par
{\footnotesize \textit{Remaining.} Formalize the actual tree expansion and prove the residual local-operator irrelevance gain after marginal subtraction; then prove the model-dependent contraction.}\par
\smallskip
\noindent\textbf{\texttt{\small DERIV:\allowbreak{}YANGMILLS\_\allowbreak{}CONTINUUM\_\allowbreak{}BALABAN\_\allowbreak{}MULTISCALE\_\allowbreak{}PROOF:\allowbreak{}THEOREM\_\allowbreak{}6\_\allowbreak{}3}} \hfill {\footnotesize conditional}\par
{\small The source records b0=11N/(48pi\textasciicircum{}2), the perturbative running formula (6.11)-(6.12) and a uniform polymer conclusion (Corollary 6.4); its repaired induction requires quantitative remainder control and a propagated envelope.}\par
{\footnotesize \srcpath{docs/derivations/yangmills-continuum-balaban-multiscale-proof.md} --- \#\#\# 6.4 Asymptotic Freedom and Infinite Scale Stability (lines 296-314)}\par
{\footnotesize \textit{Hypotheses.} Submitted SU(N), N>=2 continuum multiscale construction; read together with its in-document Repair record audited 2026-09-09.; The source title and summary claim completion, but its repair record retains explicit averaging/interpolation, irrelevance, physical-scale and scattering obligations. No machine certification is inferred from prose.; Perturbative UV regime; controlled beta remainder and actual RG map; this does not place the hadronic scale in the small-g domain.}\par
{\footnotesize \textit{Remaining.} Formalize the actual RG beta coefficient and remainder assumptions before asserting monotonicity for all g0; prove the stated envelope with the repaired constants.}\par
\smallskip
\noindent\textbf{\texttt{\small DERIV:\allowbreak{}YANGMILLS\_\allowbreak{}CONTINUUM\_\allowbreak{}BALABAN\_\allowbreak{}MULTISCALE\_\allowbreak{}PROOF:\allowbreak{}THEOREM\_\allowbreak{}7\_\allowbreak{}2}} \hfill {\footnotesize conditional}\par
{\small The continuum measure claim requires a continuous positive-definite characteristic functional on the full nuclear test space; convergence only of gauge-invariant observables is not itself this construction.}\par
{\footnotesize \srcpath{docs/derivations/yangmills-continuum-balaban-multiscale-proof.md} --- \#\#\# 7.2 Construction of the Limiting Continuum Measure \$d\textbackslash{}mu\_\allowbreak{}\{\textbackslash{}text\{YM\}\}\$ (lines 341-363)}\par
{\footnotesize \textit{Hypotheses.} Submitted SU(N), N>=2 continuum multiscale construction; read together with its in-document Repair record audited 2026-09-09.; The source title and summary claim completion, but its repair record retains explicit averaging/interpolation, irrelevance, physical-scale and scattering obligations. No machine certification is inferred from prose.; Theorem 7.1 with the repaired summable bound; actual uniform regularity and characteristic-functional estimates.}\par
{\footnotesize \textit{Remaining.} Construct the characteristic functional with continuity/positivity and a consistent gauge-invariant observable realization, then apply Minlos-Bochner.}\par
\smallskip
\noindent\textbf{\texttt{\small DERIV:\allowbreak{}YANGMILLS\_\allowbreak{}CONTINUUM\_\allowbreak{}BALABAN\_\allowbreak{}MULTISCALE\_\allowbreak{}PROOF:\allowbreak{}THEOREM\_\allowbreak{}8\_\allowbreak{}1}} \hfill {\footnotesize conditional}\par
{\small Centered reflection-adapted deterministic blocking preserves reflection positivity under its exact support/equivariance hypotheses, and a convergent compatible limit retains the inequality (8.5).}\par
{\footnotesize \srcpath{docs/derivations/yangmills-continuum-balaban-multiscale-proof.md} --- \#\#\# 8.3 OS2: Reflection Positivity (lines 383-408)}\par
{\footnotesize \textit{Hypotheses.} Submitted SU(N), N>=2 continuum multiscale construction; read together with its in-document Repair record audited 2026-09-09.; The source title and summary claim completion, but its repair record retains explicit averaging/interpolation, irrelevance, physical-scale and scattering obligations. No machine certification is inferred from prose.; Reflection-adapted deterministic block map, positive-time support and actual convergent correlations; the cited centered-blocking construction applies.}\par
{\footnotesize \textit{Remaining.} Formalize the deterministic reflected blocking and limit passage for the actual continuum measure; general reflection-equivariant Markov blocking is not this theorem.}\par
\smallskip
\noindent\textbf{\texttt{\small DERIV:\allowbreak{}YANGMILLS\_\allowbreak{}CONTINUUM\_\allowbreak{}BALABAN\_\allowbreak{}MULTISCALE\_\allowbreak{}PROOF:\allowbreak{}THEOREM\_\allowbreak{}9\_\allowbreak{}1}} \hfill {\footnotesize conditional}\par
{\small A common strictly positive exponential physical-time correlation rate on the dense centered reconstructed family implies spectrum(H) subset \{0\} union [Delta\_\allowbreak{}phys,infinity), with the simple vacuum; this conditional spectral argument is confirmed sound by the source repair record.}\par
{\footnotesize \srcpath{docs/derivations/yangmills-continuum-balaban-multiscale-proof.md} --- \#\#\# 9.2 The Physical Spectral Mass Gap (lines 454-478)}\par
{\footnotesize \textit{Hypotheses.} Submitted SU(N), N>=2 continuum multiscale construction; read together with its in-document Repair record audited 2026-09-09.; The source title and summary claim completion, but its repair record retains explicit averaging/interpolation, irrelevance, physical-scale and scattering obligations. No machine certification is inferred from prose.; Actual reconstructed physical space with dense observable family; common rate Delta\_\allowbreak{}phys>0 in physical time; simple vacuum.}\par
{\footnotesize \textit{Remaining.} Formalize the spectral-measure positivity, infinite-time exclusion and density extension; derive the actual physical clustering input (9.7) separately.}\par
\smallskip
\noindent\textbf{\texttt{\small DERIV:\allowbreak{}YANGMILLS\_\allowbreak{}GPU\_\allowbreak{}RESOLUTION\_\allowbreak{}AUDIT:\allowbreak{}GA8\_\allowbreak{}GA9}} \hfill {\footnotesize conditional}\par
{\small The inserted scalar recursion has explicit positive product/tail bounds GA8 when its assumptions hold; the Phase 4 four-plaquette builder remains a tensor sum GA9.}\par
{\footnotesize \srcpath{docs/derivations/yangmills-gpu-resolution-audit.md} --- \#\# GA4. The additional Phase 4 has a valid conditional scalar limit (lines 162-209)}\par
{\footnotesize \textit{Hypotheses.} The preserved September 9 scripts are inspected as the operators they actually implement; they are not assumed to represent Wilson dynamics.; Analytic matrix identities, replayed floating observations and reported-only runs remain distinct.; eps\_\allowbreak{}j=C(ell+jh)\textasciicircum{}(-3/2)<1; positive initial gap and fast floors; these are supplied inputs in the script.}\par
{\footnotesize \textit{Remaining.} Formalize the scalar positive-product and geometric tail bounds; derive actual coarse Wilson parameters before applying them.}\par
\smallskip
\noindent\textbf{\texttt{\small DERIV:\allowbreak{}YANGMILLS\_\allowbreak{}RECONSTRUCTION:\allowbreak{}R6\_\allowbreak{}R7}} \hfill {\footnotesize conditional}\par
{\small If C\_\allowbreak{}F(t)<=A\_\allowbreak{}F exp(-eta*t+alpha*R)+B\_\allowbreak{}F exp(-beta*R) for every R,t>=0, choose R=eta*t/(alpha+beta) to obtain common rate eta*beta/(alpha+beta).}\par
{\footnotesize \srcpath{docs/derivations/yangmills-reconstruction.md} --- \#\# 5. A precise sufficient repair, with its remaining hypothesis exposed (lines 90-111)}\par
{\footnotesize \textit{Hypotheses.} A nonnegative self-adjoint physical Hamiltonian with a specified vacuum and its positive spectral measures; physical Euclidean time is used.; OS reconstruction and any limit theorem are explicit analytic inputs, not consequences of the finite controls.; eta,beta>0,alpha>=0; one common rate and controlled cutoff/volume constants on a dense centered physical family.}\par
{\footnotesize \textit{Remaining.} Real-radius exponent balancing, actual-measure spectral exclusion and dense-family extension are formalized. Establish the actual uniform Yang-Mills two-parameter localization estimate, physical spectral identification and total family;}\par
\smallskip
\noindent\textbf{\texttt{\small DERIV:\allowbreak{}YANGMILLS\_\allowbreak{}RESOLVENT\_\allowbreak{}LOCALIZATION\_\allowbreak{}AND\_\allowbreak{}DIRICHLET\_\allowbreak{}GAP:\allowbreak{}LOCAL\_\allowbreak{}KERNEL\_\allowbreak{}IMPLICATION}} \hfill {\footnotesize conditional}\par
{\small Given uniform massive exponentially decaying kernels and a finite-range perturbation of size C2|g|, the localized scalar pairing has a volume-independent geometric-sum bound.}\par
{\footnotesize \srcpath{docs/derivations/yangmills-resolvent-localization-and-dirichlet-gap.md} --- \#\#\# 2.3 Volume-Uniformity of the Pairing (lines 62-80)}\par
{\footnotesize \textit{Hypotheses.} This submitted proposal is explicitly preserved under its September 9 source audit; closure claims are not the current graph verdict.; The actual model/operator identification is an obligation, separate from a scalar finite replay.; Kernel decay constants and perturbation range/row bounds uniform in volume; finite localized source with the source normalization explicitly retained.}\par
{\footnotesize \textit{Remaining.} Infinite kernel-pairing summability, exact integer-lattice sums and the full source l1 factor are formalized. Identify actual Wilson fast kernels and connected signed sources and prove the required uniform decay and perturbation envelopes.}\par
\smallskip
\noindent\textbf{\texttt{\small DERIV:\allowbreak{}YANGMILLS\_\allowbreak{}RESOLVENT\_\allowbreak{}LOCALIZATION\_\allowbreak{}AND\_\allowbreak{}DIRICHLET\_\allowbreak{}GAP:\allowbreak{}RESOLVENT\_\allowbreak{}IDENTITY}} \hfill {\footnotesize conditional}\par
{\small Compatible invertible fast operators obey R\_\allowbreak{}g-R0=-R0(F\_\allowbreak{}g-F0)R\_\allowbreak{}g and the stated selected-pairing identity.}\par
{\footnotesize \srcpath{docs/derivations/yangmills-resolvent-localization-and-dirichlet-gap.md} --- \#\#\# 2.1 The Resolvent Identity (lines 42-54)}\par
{\footnotesize \textit{Hypotheses.} This submitted proposal is explicitly preserved under its September 9 source audit; closure claims are not the current graph verdict.; The actual model/operator identification is an obligation, separate from a scalar finite replay.; Real z below both spectra; compatible domains and adjoint pairings.}\par
{\footnotesize \textit{Remaining.} Formalize the resolvent identity on the compatible closed-form domains.}\par
\smallskip
\noindent\textbf{\texttt{\small DERIV:\allowbreak{}YANGMILLS\_\allowbreak{}WEIGHTED\_\allowbreak{}CURVATURE:\allowbreak{}GST3}} \hfill {\footnotesize conditional}\par
{\small Ric+(2/hbar)Hess S>=kappa gives ||A\_\allowbreak{}S f||\textasciicircum{}2>=c*kappa<f,A\_\allowbreak{}Sf>, hence spectrum in \{0\} union [c*kappa,infinity) and the Poincare bound when the kernel is constants.}\par
{\footnotesize \srcpath{docs/derivations/yangmills-weighted-curvature.md} --- \#\# GST-3: curvature normalization and a spectral argument without a discrete excited state (lines 92-130)}\par
{\footnotesize \textit{Hypotheses.} Finite-dimensional smooth Riemannian configuration space or the explicitly specified finite matrix/oscillator model.; H=-c Delta+V, c=hbar\textasciicircum{}2/(2m)>0; psi=N exp(-S/hbar)>0 normalized; common operator/form domains as stated.; kappa>0; closed Bochner identity/inequality and simple constant kernel.}\par
{\footnotesize \textit{Remaining.} The bounded positive-operator square-to-gap implication is formalized in arbitrary complex Hilbert dimension. Construct the actual weighted operator, verify its quadratic inequality and manage the unbounded-domain/spectral realization required by the source.}\par
\smallskip
\noindent\textbf{\texttt{\small DERIV:\allowbreak{}YANGMILLS\_\allowbreak{}WEIGHTED\_\allowbreak{}CURVATURE:\allowbreak{}GST6}} \hfill {\footnotesize conditional}\par
{\small The closed-form min-max comparison yields E1(H)-E0(H)>=gamma-(mean R-rmin).}\par
{\footnotesize \srcpath{docs/derivations/yangmills-weighted-curvature.md} --- \#\# GST-6: the residual-mean certificate (lines 167-190)}\par
{\footnotesize \textit{Hypotheses.} Finite-dimensional smooth Riemannian configuration space or the explicitly specified finite matrix/oscillator model.; H=-c Delta+V, c=hbar\textasciicircum{}2/(2m)>0; psi=N exp(-S/hbar)>0 normalized; common operator/form domains as stated.; A\_\allowbreak{}S is nonnegative self-adjoint with simple normalized constant ground and first variational level >=gamma>0; R>=rmin, finite mean; 1 belongs to the sum-form domain; min-max applies.}\par
{\footnotesize \textit{Remaining.} Formalize the lower first-excited variational comparison and the trial-vector Rayleigh upper bound on common form domains.}\par
\smallskip
\endgroup
% ---- end gen_open_statements.tex -----------------------------------
% ---- end B0_concordance.tex ----------------------------------------
% ---- begin B1_lean.tex ---------------------------------------------
\section{The Lean tree}
\label{app:lean}

Every registered theorem, by module, with its line and what the
registry links it to.

\LeanTotal{} theorems, \LeanSorry{} occurrences of \texttt{sorry},
axioms restricted to \texttt{propext}, \texttt{Classical.choice} and
\texttt{Quot.sound}. The axiom set is exported and validated rather
than asserted: the strict build records elaborated constants and
transitive axioms, and the loader checks them against the declaration
records and source fingerprints. Identifiers mentioned in comments are
not proof dependencies.

Read the third column carefully. \texttt{promotes} appears on 49 of the
\LeanTotal{} theorems and means the theorem proves an \emph{entire}
named check, so that check's tier is the Lean tier. Everything else is
an ingredient. A theorem count is not a coverage measure, and the
proportion of ingredients to whole-statement proofs is the honest shape
of the tree rather than a defect in it --- see \S\ref{app:repro} for
the case in this program's own archive where seventy-one
\texttt{sorry}-free modules proved tautologies.

% ---- begin gen_lean.tex --------------------------------------------
% Generated by paper/master/generate_sources.py.  Do not edit.
\noindent 439 registered theorems across 23 modules. \emph{Promotes} means the theorem proves an entire named check and so carries its verification tier; \emph{formalizes} means it is linked to a catalogue record. Only 49 of 439 carry a \texttt{promotes} link, and reading the difference as a defect would be a mistake: most theorems are ingredients, and the registry keeps ingredient and whole-statement coverage apart on purpose.\par\medskip
\subsection*{(unlocated)}
\addcontentsline{toc}{subsection}{(unlocated)}
\begin{longtable}{@{}p{0.32\linewidth} l p{0.44\linewidth}@{}}
\toprule Theorem & Line & Formalizes / promotes \\ \midrule
\endfirsthead
\toprule Theorem & Line & Formalizes / promotes \\ \midrule \endhead
\bottomrule\endfoot
{\scriptsize\ttfamily dim\_\allowbreak{}Z$_2$} & {\scriptsize } & {\scriptsize \textbf{promotes} dim Z\_\allowbreak{}2 = (L\textasciicircum{}3 - 1) + 3 = L\textasciicircum{}3 + 2} \\
{\scriptsize\ttfamily d$_3$\_\allowbreak{}ledger} & {\scriptsize } & {\scriptsize CONST:\allowbreak{}D\_\allowbreak{}3; CONST:\allowbreak{}B\_\allowbreak{}3; CONST:\allowbreak{}LEAK\_\allowbreak{}3; \textbf{promotes} d\_\allowbreak{}3 = 7/32 + 12*leak\_\allowbreak{}3 - 4*b\_\allowbreak{}3} \\
\end{longtable}

\subsection*{AnisotropyVariance.lean}
\addcontentsline{toc}{subsection}{AnisotropyVariance.lean}
\begin{longtable}{@{}p{0.32\linewidth} l p{0.44\linewidth}@{}}
\toprule Theorem & Line & Formalizes / promotes \\ \midrule
\endfirsthead
\toprule Theorem & Line & Formalizes / promotes \\ \midrule \endhead
\bottomrule\endfoot
{\scriptsize\ttfamily anisotropy\_\allowbreak{}axis\_\allowbreak{}node} & {\scriptsize 105} & {\scriptsize RESULT:\allowbreak{}RECENT\_\allowbreak{}ANISOTROPY\_\allowbreak{}NODAL\_\allowbreak{}SET} \\
{\scriptsize\ttfamily anisotropy\_\allowbreak{}body\_\allowbreak{}node} & {\scriptsize 115} & {\scriptsize RESULT:\allowbreak{}RECENT\_\allowbreak{}ANISOTROPY\_\allowbreak{}NODAL\_\allowbreak{}SET} \\
{\scriptsize\ttfamily anisotropy\_\allowbreak{}centered\_\allowbreak{}moment\_\allowbreak{}identity} & {\scriptsize 46} & {\scriptsize RESULT:\allowbreak{}RECENT\_\allowbreak{}ANISOTROPY\_\allowbreak{}VARIANCE\_\allowbreak{}CORE} \\
{\scriptsize\ttfamily anisotropy\_\allowbreak{}face\_\allowbreak{}node} & {\scriptsize 110} & {\scriptsize RESULT:\allowbreak{}RECENT\_\allowbreak{}ANISOTROPY\_\allowbreak{}NODAL\_\allowbreak{}SET} \\
{\scriptsize\ttfamily anisotropy\_\allowbreak{}generic\_\allowbreak{}holdout} & {\scriptsize 120} & {\scriptsize RESULT:\allowbreak{}RECENT\_\allowbreak{}ANISOTROPY\_\allowbreak{}VARIANCE\_\allowbreak{}CORE; \textbf{promotes} anisotropy variance is exactly the simplex pair sum and cubic shape invariant} \\
{\scriptsize\ttfamily anisotropy\_\allowbreak{}induced\_\allowbreak{}coefficient} & {\scriptsize 126} & {\scriptsize RESULT:\allowbreak{}RECENT\_\allowbreak{}ANISOTROPY\_\allowbreak{}INDUCED\_\allowbreak{}SIXTH} \\
{\scriptsize\ttfamily anisotropy\_\allowbreak{}shape\_\allowbreak{}identity} & {\scriptsize 37} & {\scriptsize RESULT:\allowbreak{}RECENT\_\allowbreak{}ANISOTROPY\_\allowbreak{}VARIANCE\_\allowbreak{}CORE; RESULT:\allowbreak{}RECENT\_\allowbreak{}ANISOTROPY\_\allowbreak{}INDUCED\_\allowbreak{}SIXTH; \textbf{promotes} anisotropy variance is exactly the simplex pair sum and cubic shape invariant} \\
{\scriptsize\ttfamily anisotropy\_\allowbreak{}variance\_\allowbreak{}identity} & {\scriptsize 21} & {\scriptsize RESULT:\allowbreak{}RECENT\_\allowbreak{}ANISOTROPY\_\allowbreak{}VARIANCE\_\allowbreak{}CORE; \textbf{promotes} anisotropy variance is exactly the simplex pair sum and cubic shape invariant} \\
{\scriptsize\ttfamily anisotropy\_\allowbreak{}variance\_\allowbreak{}nonnegative} & {\scriptsize 30} & {\scriptsize RESULT:\allowbreak{}RECENT\_\allowbreak{}ANISOTROPY\_\allowbreak{}VARIANCE\_\allowbreak{}CORE; \textbf{promotes} anisotropy variance is exactly the simplex pair sum and cubic shape invariant} \\
{\scriptsize\ttfamily anisotropy\_\allowbreak{}variance\_\allowbreak{}upper\_\allowbreak{}bound} & {\scriptsize 57} & {\scriptsize RESULT:\allowbreak{}RECENT\_\allowbreak{}ANISOTROPY\_\allowbreak{}VARIANCE\_\allowbreak{}CORE} \\
{\scriptsize\ttfamily anisotropy\_\allowbreak{}variance\_\allowbreak{}zero\_\allowbreak{}iff} & {\scriptsize 73} & {\scriptsize RESULT:\allowbreak{}RECENT\_\allowbreak{}ANISOTROPY\_\allowbreak{}NODAL\_\allowbreak{}SET} \\
\end{longtable}

\subsection*{BakryEmeryCoercivity.lean}
\addcontentsline{toc}{subsection}{BakryEmeryCoercivity.lean}
\begin{longtable}{@{}p{0.32\linewidth} l p{0.44\linewidth}@{}}
\toprule Theorem & Line & Formalizes / promotes \\ \midrule
\endfirsthead
\toprule Theorem & Line & Formalizes / promotes \\ \midrule \endhead
\bottomrule\endfoot
{\scriptsize\ttfamily bochner\_\allowbreak{}gamma2\_\allowbreak{}lower\_\allowbreak{}bound} & {\scriptsize 51} & {\scriptsize ---} \\
{\scriptsize\ttfamily coercivityConstant\_\allowbreak{}lt\_\allowbreak{}one} & {\scriptsize 69} & {\scriptsize ---} \\
{\scriptsize\ttfamily coercivityConstant\_\allowbreak{}mono} & {\scriptsize 77} & {\scriptsize ---} \\
{\scriptsize\ttfamily coercivityConstant\_\allowbreak{}pos} & {\scriptsize 62} & {\scriptsize ---} \\
{\scriptsize\ttfamily form\_\allowbreak{}coercivity\_\allowbreak{}bound} & {\scriptsize 87} & {\scriptsize ---} \\
{\scriptsize\ttfamily quantum\_\allowbreak{}curvature\_\allowbreak{}restoration} & {\scriptsize 36} & {\scriptsize ---} \\
{\scriptsize\ttfamily quantum\_\allowbreak{}restoration\_\allowbreak{}ratio\_\allowbreak{}gt\_\allowbreak{}one} & {\scriptsize 112} & {\scriptsize ---} \\
{\scriptsize\ttfamily ricInfty\_\allowbreak{}pos\_\allowbreak{}of\_\allowbreak{}quantum\_\allowbreak{}restoration} & {\scriptsize 44} & {\scriptsize ---} \\
{\scriptsize\ttfamily volume\_\allowbreak{}uniform\_\allowbreak{}coercivity\_\allowbreak{}hierarchy} & {\scriptsize 128} & {\scriptsize ---} \\
{\scriptsize\ttfamily volume\_\allowbreak{}uniform\_\allowbreak{}rho\_\allowbreak{}pos} & {\scriptsize 121} & {\scriptsize ---} \\
\end{longtable}

\subsection*{Basic.lean}
\addcontentsline{toc}{subsection}{Basic.lean}
\begin{longtable}{@{}p{0.32\linewidth} l p{0.44\linewidth}@{}}
\toprule Theorem & Line & Formalizes / promotes \\ \midrule
\endfirsthead
\toprule Theorem & Line & Formalizes / promotes \\ \midrule \endhead
\bottomrule\endfoot
{\scriptsize\ttfamily C\_\allowbreak{}from\_\allowbreak{}beta} & {\scriptsize 244} & {\scriptsize CONST:\allowbreak{}C\_\allowbreak{}shp (historical); CONST:\allowbreak{}BETA\_\allowbreak{}PEN\_\allowbreak{}3; \textbf{promotes} C\_\allowbreak{}old = (beta\_\allowbreak{}pen\_\allowbreak{}3 - 2*alpha\_\allowbreak{}3)/16} \\
{\scriptsize\ttfamily adjoint\_\allowbreak{}shelf\_\allowbreak{}numerator} & {\scriptsize 188} & {\scriptsize G17; \textbf{promotes} the second-Casimir shelf: every irrep beyond F and F-bar clears 5C\_\allowbreak{}F/4} \\
{\scriptsize\ttfamily alphaPen\_\allowbreak{}eq\_\allowbreak{}neg\_\allowbreak{}four\_\allowbreak{}cube} & {\scriptsize 381} & {\scriptsize CONST:alpha\_\allowbreak{}pen\_\allowbreak{}N; CONST:\allowbreak{}CUBE\_\allowbreak{}COMPLETION\_\allowbreak{}4} \\
{\scriptsize\ttfamily alphaPen\_\allowbreak{}five} & {\scriptsize 216} & {\scriptsize CONST:alpha\_\allowbreak{}pen\_\allowbreak{}N; \textbf{promotes} axial law reproduces alpha\_\allowbreak{}4, alpha\_\allowbreak{}5, alpha\_\allowbreak{}6} \\
{\scriptsize\ttfamily alphaPen\_\allowbreak{}four} & {\scriptsize 215} & {\scriptsize CONST:alpha\_\allowbreak{}pen\_\allowbreak{}N; \textbf{promotes} axial law reproduces alpha\_\allowbreak{}4, alpha\_\allowbreak{}5, alpha\_\allowbreak{}6} \\
{\scriptsize\ttfamily alphaPen\_\allowbreak{}six} & {\scriptsize 217} & {\scriptsize CONST:alpha\_\allowbreak{}pen\_\allowbreak{}N; \textbf{promotes} axial law reproduces alpha\_\allowbreak{}4, alpha\_\allowbreak{}5, alpha\_\allowbreak{}6} \\
{\scriptsize\ttfamily alphaPen\_\allowbreak{}three} & {\scriptsize 214} & {\scriptsize CONST:alpha\_\allowbreak{}pen\_\allowbreak{}N; \textbf{promotes} alpha\_\allowbreak{}3 = 4*A\_\allowbreak{}shp = 5/12} \\
{\scriptsize\ttfamily alphaPen\_\allowbreak{}three\_\allowbreak{}eq\_\allowbreak{}four\_\allowbreak{}A} & {\scriptsize 220} & {\scriptsize CONST:\allowbreak{}A\_\allowbreak{}shp\_\allowbreak{}3; CONST:alpha\_\allowbreak{}pen\_\allowbreak{}N; R7; \textbf{promotes} alpha\_\allowbreak{}3 = 4*A\_\allowbreak{}shp = 5/12} \\
{\scriptsize\ttfamily assembly\_\allowbreak{}three\_\allowbreak{}cumulants} & {\scriptsize 718} & {\scriptsize G14; C10} \\
{\scriptsize\ttfamily bandVar\_\allowbreak{}pi} & {\scriptsize 490} & {\scriptsize SYM:bloch\_\allowbreak{}scalars} \\
{\scriptsize\ttfamily bandVar\_\allowbreak{}pi\_\allowbreak{}div\_\allowbreak{}two} & {\scriptsize 493} & {\scriptsize SYM:bloch\_\allowbreak{}scalars} \\
{\scriptsize\ttfamily bandVar\_\allowbreak{}zero} & {\scriptsize 487} & {\scriptsize SYM:bloch\_\allowbreak{}scalars} \\
{\scriptsize\ttfamily betaN\_\allowbreak{}assembled\_\allowbreak{}numerator} & {\scriptsize 683} & {\scriptsize G14; C10; R2; \textbf{promotes} the final polynomial identity: -16 u + 32 d - 16 corner - 8 K\_\allowbreak{}adj = P17(N\textasciicircum{}2)/(N R20(N\textasciicircum{}2)) as rational functions, for the closed forms of the two-hop weight, the single-contact dressing, the corner dressing and the adjacent-face cube completion} \\
{\scriptsize\ttfamily betaN\_\allowbreak{}four} & {\scriptsize 743} & {\scriptsize R2} \\
{\scriptsize\ttfamily betaN\_\allowbreak{}from\_\allowbreak{}three\_\allowbreak{}cumulants} & {\scriptsize 727} & {\scriptsize G14; C10; R2; CONST:\allowbreak{}TWO\_\allowbreak{}HOP\_\allowbreak{}WEIGHT\_\allowbreak{}ODD\_\allowbreak{}N; CONST:\allowbreak{}CORNER\_\allowbreak{}DRESSING\_\allowbreak{}ODD\_\allowbreak{}N; \textbf{promotes} the final polynomial identity: -16 u + 32 d - 16 corner - 8 K\_\allowbreak{}adj = P17(N\textasciicircum{}2)/(N R20(N\textasciicircum{}2)) as rational functions, for the closed forms of the two-hop weight, the single-contact dressing, the corner dressing and the adjacent-face cube completion} \\
{\scriptsize\ttfamily betaN\_\allowbreak{}three} & {\scriptsize 739} & {\scriptsize C10; CONST:\allowbreak{}BETA\_\allowbreak{}PEN\_\allowbreak{}3\_\allowbreak{}ASSEMBLED} \\
{\scriptsize\ttfamily beta\_\allowbreak{}from\_\allowbreak{}A\_\allowbreak{}and\_\allowbreak{}C} & {\scriptsize 252} & {\scriptsize G14; CONST:\allowbreak{}A\_\allowbreak{}shp\_\allowbreak{}3; CONST:\allowbreak{}C\_\allowbreak{}shp (historical)} \\
{\scriptsize\ttfamily beta\_\allowbreak{}shift\_\allowbreak{}from\_\allowbreak{}cShp} & {\scriptsize 428} & {\scriptsize C10; C2} \\
{\scriptsize\ttfamily blind\_\allowbreak{}holdout} & {\scriptsize 230} & {\scriptsize G3} \\
{\scriptsize\ttfamily cIso\_\allowbreak{}assembled} & {\scriptsize 580} & {\scriptsize G11} \\
{\scriptsize\ttfamily cPrimTwo\_\allowbreak{}forms} & {\scriptsize 353} & {\scriptsize SYM:c\_\allowbreak{}2\_\allowbreak{}tet} \\
{\scriptsize\ttfamily cShpAssembled\_\allowbreak{}neg} & {\scriptsize 576} & {\scriptsize CONST:\allowbreak{}C\_\allowbreak{}SHP\_\allowbreak{}ASSEMBLED} \\
{\scriptsize\ttfamily cShp\_\allowbreak{}assembled\_\allowbreak{}value} & {\scriptsize 424} & {\scriptsize C2; CONST:\allowbreak{}C\_\allowbreak{}SHP\_\allowbreak{}ASSEMBLED} \\
{\scriptsize\ttfamily cShp\_\allowbreak{}from\_\allowbreak{}rho\_\allowbreak{}shift} & {\scriptsize 418} & {\scriptsize C2} \\
{\scriptsize\ttfamily channel\_\allowbreak{}resolvent\_\allowbreak{}assembly} & {\scriptsize 148} & {\scriptsize CONST:t\_\allowbreak{}N; R2; \textbf{promotes} the per-channel resolvent equation closes the projector-to-hopping step} \\
{\scriptsize\ttfamily combes\_\allowbreak{}thomas\_\allowbreak{}geometric\_\allowbreak{}bound} & {\scriptsize 1248} & {\scriptsize G19; CITE:\allowbreak{}YM\_\allowbreak{}GPU\_\allowbreak{}AUDIT} \\
{\scriptsize\ttfamily combes\_\allowbreak{}thomas\_\allowbreak{}pairing\_\allowbreak{}volume\_\allowbreak{}independent} & {\scriptsize 1260} & {\scriptsize G19; CITE:\allowbreak{}YM\_\allowbreak{}GPU\_\allowbreak{}AUDIT} \\
{\scriptsize\ttfamily cornerDen\_\allowbreak{}cof} & {\scriptsize 674} & {\scriptsize G14; CONST:\allowbreak{}CORNER\_\allowbreak{}DRESSING\_\allowbreak{}ODD\_\allowbreak{}N} \\
{\scriptsize\ttfamily cornerDressing\_\allowbreak{}over\_\allowbreak{}R20} & {\scriptsize 702} & {\scriptsize CONST:\allowbreak{}CORNER\_\allowbreak{}DRESSING\_\allowbreak{}ODD\_\allowbreak{}N} \\
{\scriptsize\ttfamily cubeCompletionAdjacent\_\allowbreak{}over\_\allowbreak{}R20} & {\scriptsize 708} & {\scriptsize CONST:\allowbreak{}CUBE\_\allowbreak{}COMPLETION\_\allowbreak{}ADJACENT\_\allowbreak{}N} \\
{\scriptsize\ttfamily cubeCompletionAdjacent\_\allowbreak{}ratio} & {\scriptsize 405} & {\scriptsize G14} \\
{\scriptsize\ttfamily cubeCompletionAdjacent\_\allowbreak{}split} & {\scriptsize 396} & {\scriptsize G14; CONST:\allowbreak{}CUBE\_\allowbreak{}COMPLETION\_\allowbreak{}ADJACENT\_\allowbreak{}N} \\
{\scriptsize\ttfamily cubeCompletionAdjacent\_\allowbreak{}three} & {\scriptsize 401} & {\scriptsize CONST:\allowbreak{}CUBE\_\allowbreak{}COMPLETION\_\allowbreak{}ADJACENT\_\allowbreak{}4; G3} \\
{\scriptsize\ttfamily cubeCompletion\_\allowbreak{}three} & {\scriptsize 372} & {\scriptsize CONST:\allowbreak{}CUBE\_\allowbreak{}COMPLETION\_\allowbreak{}4; CONST:\allowbreak{}A\_\allowbreak{}shp\_\allowbreak{}3} \\
{\scriptsize\ttfamily cubeDen\_\allowbreak{}cof} & {\scriptsize 677} & {\scriptsize G14; CONST:\allowbreak{}CUBE\_\allowbreak{}COMPLETION\_\allowbreak{}ADJACENT\_\allowbreak{}N} \\
{\scriptsize\ttfamily cube\_\allowbreak{}shortfall} & {\scriptsize 414} & {\scriptsize C2; CONST:\allowbreak{}CUBE\_\allowbreak{}COMPLETION\_\allowbreak{}ADJACENT\_\allowbreak{}HISTORICAL\_\allowbreak{}4} \\
{\scriptsize\ttfamily delta\_\allowbreak{}M} & {\scriptsize 526} & {\scriptsize SYM:bloch\_\allowbreak{}scalars; \textbf{promotes} checkpoint values at X, M, P, R; \textbf{promotes} the four extraction formulas invert the ansatz} \\
{\scriptsize\ttfamily delta\_\allowbreak{}P} & {\scriptsize 532} & {\scriptsize SYM:bloch\_\allowbreak{}scalars; \textbf{promotes} checkpoint values at X, M, P, R; \textbf{promotes} the four extraction formulas invert the ansatz} \\
{\scriptsize\ttfamily delta\_\allowbreak{}R} & {\scriptsize 540} & {\scriptsize SYM:bloch\_\allowbreak{}scalars; \textbf{promotes} checkpoint values at X, M, P, R; \textbf{promotes} the four extraction formulas invert the ansatz} \\
{\scriptsize\ttfamily delta\_\allowbreak{}X} & {\scriptsize 521} & {\scriptsize SYM:bloch\_\allowbreak{}scalars; \textbf{promotes} checkpoint values at X, M, P, R; \textbf{promotes} the four extraction formulas invert the ansatz; \textbf{promotes} X is blind to B, C, D --- it fixes A alone} \\
{\scriptsize\ttfamily evenHopping\_\allowbreak{}three} & {\scriptsize 467} & {\scriptsize CONST:\allowbreak{}T\_\allowbreak{}PLUS\_\allowbreak{}2; R2} \\
{\scriptsize\ttfamily even\_\allowbreak{}cubic\_\allowbreak{}at\_\allowbreak{}sixteen} & {\scriptsize 284} & {\scriptsize G25; CONST:\allowbreak{}BAND\_\allowbreak{}EVEN\_\allowbreak{}BOTTOM; CONST:\allowbreak{}M3\_\allowbreak{}EVEN\_\allowbreak{}K0} \\
{\scriptsize\ttfamily even\_\allowbreak{}cubic\_\allowbreak{}at\_\allowbreak{}zero} & {\scriptsize 278} & {\scriptsize G25; CONST:\allowbreak{}BAND\_\allowbreak{}EVEN\_\allowbreak{}TOP; CONST:\allowbreak{}M3\_\allowbreak{}EVEN\_\allowbreak{}BANDTOP} \\
{\scriptsize\ttfamily even\_\allowbreak{}cubic\_\allowbreak{}derivative\_\allowbreak{}factors} & {\scriptsize 291} & {\scriptsize G25} \\
{\scriptsize\ttfamily even\_\allowbreak{}gram\_\allowbreak{}minors} & {\scriptsize 271} & {\scriptsize G25} \\
{\scriptsize\ttfamily extraction\_\allowbreak{}A} & {\scriptsize 310} & {\scriptsize SYM:bloch\_\allowbreak{}scalars; \textbf{promotes} the four extraction formulas invert the ansatz} \\
{\scriptsize\ttfamily extraction\_\allowbreak{}B} & {\scriptsize 312} & {\scriptsize SYM:bloch\_\allowbreak{}scalars; \textbf{promotes} the four extraction formulas invert the ansatz} \\
{\scriptsize\ttfamily extraction\_\allowbreak{}C} & {\scriptsize 316} & {\scriptsize SYM:bloch\_\allowbreak{}scalars; \textbf{promotes} the four extraction formulas invert the ansatz} \\
{\scriptsize\ttfamily extraction\_\allowbreak{}D} & {\scriptsize 320} & {\scriptsize SYM:bloch\_\allowbreak{}scalars; \textbf{promotes} the four extraction formulas invert the ansatz} \\
{\scriptsize\ttfamily finite\_\allowbreak{}ground\_\allowbreak{}state\_\allowbreak{}dirichlet\_\allowbreak{}nonneg} & {\scriptsize 845} & {\scriptsize G23; STUDY:\allowbreak{}YM:operator-identification; \textbf{promotes} finite symmetric matrix ground-state residual identity} \\
{\scriptsize\ttfamily finite\_\allowbreak{}ground\_\allowbreak{}state\_\allowbreak{}eigenform\_\allowbreak{}identity} & {\scriptsize 832} & {\scriptsize G23; STUDY:\allowbreak{}YM:operator-identification} \\
{\scriptsize\ttfamily finite\_\allowbreak{}ground\_\allowbreak{}state\_\allowbreak{}gap\_\allowbreak{}transfer} & {\scriptsize 865} & {\scriptsize G23; STUDY:\allowbreak{}YM:operator-identification; STUDY:\allowbreak{}YM:weighted-curvature} \\
{\scriptsize\ttfamily finite\_\allowbreak{}ground\_\allowbreak{}state\_\allowbreak{}ratio\_\allowbreak{}identity} & {\scriptsize 822} & {\scriptsize G23; STUDY:\allowbreak{}YM:operator-identification; \textbf{promotes} finite symmetric matrix ground-state residual identity} \\
{\scriptsize\ttfamily finite\_\allowbreak{}ground\_\allowbreak{}state\_\allowbreak{}residual\_\allowbreak{}identity} & {\scriptsize 774} & {\scriptsize G23; STUDY:\allowbreak{}YM:operator-identification} \\
{\scriptsize\ttfamily finite\_\allowbreak{}transverse\_\allowbreak{}valley\_\allowbreak{}identity} & {\scriptsize 944} & {\scriptsize G22; STUDY:\allowbreak{}YM:matrix-commutator-model} \\
{\scriptsize\ttfamily geom\_\allowbreak{}sum\_\allowbreak{}mul\_\allowbreak{}sub} & {\scriptsize 1239} & {\scriptsize G19; CITE:\allowbreak{}YM\_\allowbreak{}GPU\_\allowbreak{}AUDIT} \\
{\scriptsize\ttfamily groundStateDirichlet\_\allowbreak{}const} & {\scriptsize 897} & {\scriptsize G23; CITE:\allowbreak{}YM\_\allowbreak{}GPU\_\allowbreak{}AUDIT} \\
{\scriptsize\ttfamily ground\_\allowbreak{}state\_\allowbreak{}frustration\_\allowbreak{}elimination} & {\scriptsize 906} & {\scriptsize G23; CITE:\allowbreak{}YM\_\allowbreak{}GPU\_\allowbreak{}AUDIT} \\
{\scriptsize\ttfamily hardy\_\allowbreak{}linear\_\allowbreak{}lower\_\allowbreak{}bound} & {\scriptsize 990} & {\scriptsize G22; STUDY:\allowbreak{}YM:matrix-commutator-model} \\
{\scriptsize\ttfamily hardy\_\allowbreak{}linear\_\allowbreak{}remainder} & {\scriptsize 984} & {\scriptsize G22; STUDY:\allowbreak{}YM:matrix-commutator-model} \\
{\scriptsize\ttfamily hopping\_\allowbreak{}deficit\_\allowbreak{}numerator} & {\scriptsize 48} & {\scriptsize CONST:t\_\allowbreak{}N; \textbf{promotes} deficit identity 1/4 - N\textasciicircum{}3 t\_\allowbreak{}N} \\
{\scriptsize\ttfamily hopping\_\allowbreak{}three} & {\scriptsize 41} & {\scriptsize CONST:t\_\allowbreak{}N; R2; \textbf{promotes} t\_\allowbreak{}2 = 0 and t\_\allowbreak{}3 = 5/612} \\
{\scriptsize\ttfamily hopping\_\allowbreak{}two} & {\scriptsize 44} & {\scriptsize CONST:t\_\allowbreak{}N; R2; \textbf{promotes} t\_\allowbreak{}2 = 0 and t\_\allowbreak{}3 = 5/612} \\
{\scriptsize\ttfamily isolation\_\allowbreak{}condition\_\allowbreak{}assembled} & {\scriptsize 571} & {\scriptsize CONST:\allowbreak{}C\_\allowbreak{}SHP\_\allowbreak{}ASSEMBLED; G11} \\
{\scriptsize\ttfamily isolation\_\allowbreak{}switch\_\allowbreak{}numerator} & {\scriptsize 196} & {\scriptsize G17; \textbf{promotes} the exact isolation constant: gamma = C\_\allowbreak{}F, except 5/4 at SU(4), attained} \\
{\scriptsize\ttfamily kotecky\_\allowbreak{}preiss\_\allowbreak{}cluster\_\allowbreak{}sum\_\allowbreak{}bound} & {\scriptsize 1285} & {\scriptsize G17; CITE:\allowbreak{}YM\_\allowbreak{}GPU\_\allowbreak{}AUDIT} \\
{\scriptsize\ttfamily kotecky\_\allowbreak{}preiss\_\allowbreak{}polymer\_\allowbreak{}activity\_\allowbreak{}bound} & {\scriptsize 1303} & {\scriptsize G17; CITE:\allowbreak{}YM\_\allowbreak{}GPU\_\allowbreak{}AUDIT} \\
{\scriptsize\ttfamily lambda2\_\allowbreak{}shelf\_\allowbreak{}numerator} & {\scriptsize 180} & {\scriptsize G17; \textbf{promotes} the second-Casimir shelf: every irrep beyond F and F-bar clears 5C\_\allowbreak{}F/4} \\
{\scriptsize\ttfamily like\_\allowbreak{}family\_\allowbreak{}sum} & {\scriptsize 129} & {\scriptsize CONST:t\_\allowbreak{}N; R2; \textbf{promotes} A\_\allowbreak{}N and B\_\allowbreak{}N are the channel sums, not transcriptions} \\
{\scriptsize\ttfamily mixed\_\allowbreak{}family\_\allowbreak{}sum} & {\scriptsize 118} & {\scriptsize CONST:t\_\allowbreak{}N; R2; \textbf{promotes} A\_\allowbreak{}N and B\_\allowbreak{}N are the channel sums, not transcriptions} \\
{\scriptsize\ttfamily newton\_\allowbreak{}three} & {\scriptsize 300} & {\scriptsize G14; ADR:0004; ADR:0005} \\
{\scriptsize\ttfamily pentCompletion\_\allowbreak{}three} & {\scriptsize 376} & {\scriptsize CONST:\allowbreak{}PENT\_\allowbreak{}COMPLETION\_\allowbreak{}5\_\allowbreak{}SU3} \\
{\scriptsize\ttfamily piTilde\_\allowbreak{}neg} & {\scriptsize 566} & {\scriptsize CONST:\allowbreak{}PI\_\allowbreak{}ORBIT; CONST:\allowbreak{}X\_\allowbreak{}QUANTUM; G11} \\
{\scriptsize\ttfamily polymer\_\allowbreak{}free\_\allowbreak{}energy\_\allowbreak{}cauchy\_\allowbreak{}bound} & {\scriptsize 1310} & {\scriptsize G17; CITE:\allowbreak{}YM\_\allowbreak{}GPU\_\allowbreak{}AUDIT} \\
{\scriptsize\ttfamily prismCompletion\_\allowbreak{}three} & {\scriptsize 368} & {\scriptsize CONST:\allowbreak{}PRISM\_\allowbreak{}COMPLETION\_\allowbreak{}3\_\allowbreak{}SU3} \\
{\scriptsize\ttfamily q\_\allowbreak{}at\_\allowbreak{}checkpoints} & {\scriptsize 512} & {\scriptsize SYM:bloch\_\allowbreak{}scalars; \textbf{promotes} q at the four high-symmetry points is 0, 4, 8, 12} \\
{\scriptsize\ttfamily quartic\_\allowbreak{}gaussian\_\allowbreak{}frequency\_\allowbreak{}optimal} & {\scriptsize 1056} & {\scriptsize G23; STUDY:\allowbreak{}YM:weighted-curvature; \textbf{promotes} quartic Gaussian certificate is optimized at frequency cubed 6 lambda hbar over mass squared} \\
{\scriptsize\ttfamily quartic\_\allowbreak{}gaussian\_\allowbreak{}frequency\_\allowbreak{}optimal\_\allowbreak{}iff} & {\scriptsize 1065} & {\scriptsize G23; STUDY:\allowbreak{}YM:weighted-curvature; \textbf{promotes} quartic Gaussian certificate is optimized at frequency cubed 6 lambda hbar over mass squared} \\
{\scriptsize\ttfamily quartic\_\allowbreak{}gaussian\_\allowbreak{}frequency\_\allowbreak{}remainder} & {\scriptsize 1049} & {\scriptsize G23; STUDY:\allowbreak{}YM:weighted-curvature; \textbf{promotes} quartic Gaussian certificate is optimized at frequency cubed 6 lambda hbar over mass squared} \\
{\scriptsize\ttfamily quartic\_\allowbreak{}gaussian\_\allowbreak{}gap\_\allowbreak{}budget} & {\scriptsize 1043} & {\scriptsize G23; STUDY:\allowbreak{}YM:weighted-curvature} \\
{\scriptsize\ttfamily quartic\_\allowbreak{}trial\_\allowbreak{}residual\_\allowbreak{}completion} & {\scriptsize 1019} & {\scriptsize G23; STUDY:\allowbreak{}YM:weighted-curvature} \\
{\scriptsize\ttfamily quartic\_\allowbreak{}trial\_\allowbreak{}residual\_\allowbreak{}floor} & {\scriptsize 1032} & {\scriptsize G23; STUDY:\allowbreak{}YM:weighted-curvature} \\
{\scriptsize\ttfamily rank\_\allowbreak{}law\_\allowbreak{}numerator} & {\scriptsize 36} & {\scriptsize CONST:t\_\allowbreak{}N; R2; \textbf{promotes} t\_\allowbreak{}N = B\_\allowbreak{}N - A\_\allowbreak{}N} \\
{\scriptsize\ttfamily relative\_\allowbreak{}gap\_\allowbreak{}pos} & {\scriptsize 588} & {\scriptsize G11} \\
{\scriptsize\ttfamily residual\_\allowbreak{}mean\_\allowbreak{}gap\_\allowbreak{}budget} & {\scriptsize 1012} & {\scriptsize G23; STUDY:\allowbreak{}YM:weighted-curvature; \textbf{promotes} trial residual gap budget uses its mean above its floor} \\
{\scriptsize\ttfamily residual\_\allowbreak{}mean\_\allowbreak{}slack\_\allowbreak{}identity} & {\scriptsize 1006} & {\scriptsize G23; STUDY:\allowbreak{}YM:weighted-curvature; \textbf{promotes} trial residual gap budget uses its mean above its floor} \\
{\scriptsize\ttfamily residual\_\allowbreak{}oscillation\_\allowbreak{}gap\_\allowbreak{}budget} & {\scriptsize 922} & {\scriptsize G23} \\
{\scriptsize\ttfamily residual\_\allowbreak{}oscillation\_\allowbreak{}positive\_\allowbreak{}gap} & {\scriptsize 928} & {\scriptsize G23; \textbf{promotes} trial residual gap budget uses its mean above its floor} \\
{\scriptsize\ttfamily resolventWeight\_\allowbreak{}adjoint} & {\scriptsize 80} & {\scriptsize CONST:t\_\allowbreak{}N; R2; \textbf{promotes} the four channel weights follow from dimension and Casimir} \\
{\scriptsize\ttfamily resolventWeight\_\allowbreak{}antisym} & {\scriptsize 89} & {\scriptsize CONST:t\_\allowbreak{}N; R2; \textbf{promotes} the four channel weights follow from dimension and Casimir} \\
{\scriptsize\ttfamily resolventWeight\_\allowbreak{}singlet} & {\scriptsize 72} & {\scriptsize CONST:t\_\allowbreak{}N; R2; \textbf{promotes} the four channel weights follow from dimension and Casimir} \\
{\scriptsize\ttfamily resolventWeight\_\allowbreak{}sym} & {\scriptsize 104} & {\scriptsize CONST:t\_\allowbreak{}N; R2; \textbf{promotes} the four channel weights follow from dimension and Casimir} \\
{\scriptsize\ttfamily shell\_\allowbreak{}margin\_\allowbreak{}five} & {\scriptsize 174} & {\scriptsize G17; \textbf{promotes} five charged links clear the shell by exactly C\_\allowbreak{}F/2, at every volume} \\
{\scriptsize\ttfamily singleContactDen\_\allowbreak{}cof} & {\scriptsize 670} & {\scriptsize G14; CONST:\allowbreak{}SINGLE\_\allowbreak{}CONTACT\_\allowbreak{}DRESSING\_\allowbreak{}ODD\_\allowbreak{}N} \\
{\scriptsize\ttfamily singleContactEven\_\allowbreak{}eq} & {\scriptsize 455} & {\scriptsize G14; CONST:\allowbreak{}SINGLE\_\allowbreak{}CONTACT\_\allowbreak{}DRESSING\_\allowbreak{}EVEN\_\allowbreak{}N} \\
{\scriptsize\ttfamily singleContactEven\_\allowbreak{}three} & {\scriptsize 463} & {\scriptsize CONST:\allowbreak{}SINGLE\_\allowbreak{}CONTACT\_\allowbreak{}DRESSING\_\allowbreak{}EVEN\_\allowbreak{}N} \\
{\scriptsize\ttfamily singleContactOdd\_\allowbreak{}over\_\allowbreak{}R20} & {\scriptsize 696} & {\scriptsize CONST:\allowbreak{}SINGLE\_\allowbreak{}CONTACT\_\allowbreak{}DRESSING\_\allowbreak{}ODD\_\allowbreak{}N} \\
{\scriptsize\ttfamily stencil\_\allowbreak{}zero\_\allowbreak{}mode} & {\scriptsize 234} & {\scriptsize ---} \\
{\scriptsize\ttfamily sym2\_\allowbreak{}shelf\_\allowbreak{}numerator} & {\scriptsize 184} & {\scriptsize G17; \textbf{promotes} the second-Casimir shelf: every irrep beyond F and F-bar clears 5C\_\allowbreak{}F/4} \\
{\scriptsize\ttfamily tetraCompletion\_\allowbreak{}three} & {\scriptsize 364} & {\scriptsize C15; CONST:\allowbreak{}TETRA\_\allowbreak{}COMPLETION\_\allowbreak{}2\_\allowbreak{}SU3} \\
{\scriptsize\ttfamily tetra\_\allowbreak{}from\_\allowbreak{}count} & {\scriptsize 358} & {\scriptsize C15; G5; CONST:c\_\allowbreak{}2\textasciicircum{}tet (primitive)} \\
{\scriptsize\ttfamily transverse\_\allowbreak{}slice\_\allowbreak{}form\_\allowbreak{}aggregation} & {\scriptsize 970} & {\scriptsize G22; STUDY:\allowbreak{}YM:matrix-commutator-model; \textbf{promotes} YM flat directions: balanced slice aggregation retains half the kinetic energy} \\
{\scriptsize\ttfamily transverse\_\allowbreak{}slice\_\allowbreak{}form\_\allowbreak{}identity} & {\scriptsize 978} & {\scriptsize G22; STUDY:\allowbreak{}YM:matrix-commutator-model; \textbf{promotes} YM flat directions: balanced slice aggregation retains half the kinetic energy} \\
{\scriptsize\ttfamily twoHopDen\_\allowbreak{}cof} & {\scriptsize 667} & {\scriptsize G14; CONST:\allowbreak{}TWO\_\allowbreak{}HOP\_\allowbreak{}WEIGHT\_\allowbreak{}ODD\_\allowbreak{}N} \\
{\scriptsize\ttfamily twoHopWeight\_\allowbreak{}over\_\allowbreak{}R20} & {\scriptsize 690} & {\scriptsize CONST:\allowbreak{}TWO\_\allowbreak{}HOP\_\allowbreak{}WEIGHT\_\allowbreak{}ODD\_\allowbreak{}N} \\
{\scriptsize\ttfamily uStarSq\_\allowbreak{}isolation} & {\scriptsize 584} & {\scriptsize CONST:\allowbreak{}T\_\allowbreak{}MINUS\_\allowbreak{}2; G11} \\
{\scriptsize\ttfamily w4\_\allowbreak{}old\_\allowbreak{}is\_\allowbreak{}alpha\_\allowbreak{}plus\_\allowbreak{}beta} & {\scriptsize 436} & {\scriptsize CONST:\allowbreak{}BETA\_\allowbreak{}PEN\_\allowbreak{}3; CONST:\allowbreak{}ALPHA\_\allowbreak{}PEN\_\allowbreak{}3} \\
{\scriptsize\ttfamily w4\_\allowbreak{}shift\_\allowbreak{}from\_\allowbreak{}beta} & {\scriptsize 433} & {\scriptsize C2; G11; CONST:\allowbreak{}W4\_\allowbreak{}ASSEMBLED} \\
{\scriptsize\ttfamily width\_\allowbreak{}eq\_\allowbreak{}alpha\_\allowbreak{}add\_\allowbreak{}beta} & {\scriptsize 248} & {\scriptsize CONST:\allowbreak{}W4\_\allowbreak{}HISTORICAL\_\allowbreak{}NUM; CONST:\allowbreak{}BETA\_\allowbreak{}PEN\_\allowbreak{}3} \\
{\scriptsize\ttfamily wilson\_\allowbreak{}complex\_\allowbreak{}gram\_\allowbreak{}budget} & {\scriptsize 1199} & {\scriptsize CITE:\allowbreak{}WILSON\_\allowbreak{}SHELL} \\
{\scriptsize\ttfamily wilson\_\allowbreak{}damped\_\allowbreak{}time\_\allowbreak{}denominator} & {\scriptsize 1121} & {\scriptsize G19; CITE:\allowbreak{}WILSON\_\allowbreak{}MARKED} \\
{\scriptsize\ttfamily wilson\_\allowbreak{}disconnected\_\allowbreak{}vacuum\_\allowbreak{}cancel} & {\scriptsize 1106} & {\scriptsize G17; G19; CITE:\allowbreak{}WILSON\_\allowbreak{}MARKED} \\
{\scriptsize\ttfamily wilson\_\allowbreak{}geometric\_\allowbreak{}telescope} & {\scriptsize 1112} & {\scriptsize G19; CITE:\allowbreak{}WILSON\_\allowbreak{}MARKED} \\
{\scriptsize\ttfamily wilson\_\allowbreak{}polymer\_\allowbreak{}kp\_\allowbreak{}budget} & {\scriptsize 1139} & {\scriptsize G17; G19; CITE:\allowbreak{}WILSON\_\allowbreak{}MARKED} \\
{\scriptsize\ttfamily wilson\_\allowbreak{}relative\_\allowbreak{}gap\_\allowbreak{}budget} & {\scriptsize 1153} & {\scriptsize G18; G19; CITE:\allowbreak{}WILSON\_\allowbreak{}MARKED} \\
{\scriptsize\ttfamily wilson\_\allowbreak{}relative\_\allowbreak{}gap\_\allowbreak{}positive} & {\scriptsize 1158} & {\scriptsize G18; G19; CITE:\allowbreak{}WILSON\_\allowbreak{}MARKED} \\
{\scriptsize\ttfamily wilson\_\allowbreak{}sc17\_\allowbreak{}barrier\_\allowbreak{}polynomial} & {\scriptsize 1405} & {\scriptsize CITE:\allowbreak{}WILSON\_\allowbreak{}SC17} \\
{\scriptsize\ttfamily wilson\_\allowbreak{}sc17\_\allowbreak{}broad\_\allowbreak{}barrier} & {\scriptsize 1377} & {\scriptsize CITE:\allowbreak{}WILSON\_\allowbreak{}SC17} \\
{\scriptsize\ttfamily wilson\_\allowbreak{}sc17\_\allowbreak{}broad\_\allowbreak{}curvature\_\allowbreak{}and\_\allowbreak{}cat} & {\scriptsize 1391} & {\scriptsize CITE:\allowbreak{}WILSON\_\allowbreak{}SC17} \\
{\scriptsize\ttfamily wilson\_\allowbreak{}sc17\_\allowbreak{}discriminant\_\allowbreak{}negative\_\allowbreak{}at\_\allowbreak{}seventy\_\allowbreak{}two} & {\scriptsize 1411} & {\scriptsize CITE:\allowbreak{}WILSON\_\allowbreak{}SC17} \\
{\scriptsize\ttfamily wilson\_\allowbreak{}sc17\_\allowbreak{}sharp\_\allowbreak{}barrier} & {\scriptsize 1349} & {\scriptsize CITE:\allowbreak{}WILSON\_\allowbreak{}SC17} \\
{\scriptsize\ttfamily wilson\_\allowbreak{}sc17\_\allowbreak{}sharp\_\allowbreak{}curvature\_\allowbreak{}and\_\allowbreak{}gap} & {\scriptsize 1364} & {\scriptsize CITE:\allowbreak{}WILSON\_\allowbreak{}SC17} \\
{\scriptsize\ttfamily wilson\_\allowbreak{}slab\_\allowbreak{}block\_\allowbreak{}row\_\allowbreak{}and\_\allowbreak{}cat} & {\scriptsize 1339} & {\scriptsize CITE:\allowbreak{}WILSON\_\allowbreak{}TRUE\_\allowbreak{}BLOCK; CITE:\allowbreak{}WILSON\_\allowbreak{}TRUE\_\allowbreak{}BLOCK\_\allowbreak{}CHECK} \\
{\scriptsize\ttfamily wilson\_\allowbreak{}slab\_\allowbreak{}pinned\_\allowbreak{}activity} & {\scriptsize 1331} & {\scriptsize CITE:\allowbreak{}WILSON\_\allowbreak{}TRUE\_\allowbreak{}BLOCK; CITE:\allowbreak{}WILSON\_\allowbreak{}TRUE\_\allowbreak{}BLOCK\_\allowbreak{}CHECK} \\
{\scriptsize\ttfamily wilson\_\allowbreak{}source\_\allowbreak{}weight\_\allowbreak{}defect} & {\scriptsize 1165} & {\scriptsize G18; CITE:\allowbreak{}WILSON\_\allowbreak{}MARKED} \\
{\scriptsize\ttfamily wilson\_\allowbreak{}source\_\allowbreak{}weight\_\allowbreak{}positive} & {\scriptsize 1170} & {\scriptsize G18; CITE:\allowbreak{}WILSON\_\allowbreak{}MARKED} \\
{\scriptsize\ttfamily wilson\_\allowbreak{}spatial\_\allowbreak{}budget\_\allowbreak{}telescope} & {\scriptsize 1222} & {\scriptsize G19; CITE:\allowbreak{}WILSON\_\allowbreak{}SPATIAL} \\
{\scriptsize\ttfamily wilson\_\allowbreak{}spatial\_\allowbreak{}weighted\_\allowbreak{}step} & {\scriptsize 1212} & {\scriptsize G19; CITE:\allowbreak{}WILSON\_\allowbreak{}SPATIAL} \\
{\scriptsize\ttfamily wilson\_\allowbreak{}symmetric\_\allowbreak{}block\_\allowbreak{}identity} & {\scriptsize 1092} & {\scriptsize G19; CITE:\allowbreak{}WILSON\_\allowbreak{}MARKED} \\
{\scriptsize\ttfamily wilson\_\allowbreak{}taylor\_\allowbreak{}polynomial\_\allowbreak{}envelope} & {\scriptsize 1182} & {\scriptsize CITE:\allowbreak{}WILSON\_\allowbreak{}SHELL} \\
{\scriptsize\ttfamily wilson\_\allowbreak{}time\_\allowbreak{}kernel\_\allowbreak{}split} & {\scriptsize 1132} & {\scriptsize G19; CITE:\allowbreak{}WILSON\_\allowbreak{}MARKED} \\
{\scriptsize\ttfamily wilson\_\allowbreak{}uniform\_\allowbreak{}carrier\_\allowbreak{}coefficient} & {\scriptsize 1204} & {\scriptsize CITE:\allowbreak{}WILSON\_\allowbreak{}SHELL; G18; G19} \\
\end{longtable}

\subsection*{CreatorParent.lean}
\addcontentsline{toc}{subsection}{CreatorParent.lean}
\begin{longtable}{@{}p{0.32\linewidth} l p{0.44\linewidth}@{}}
\toprule Theorem & Line & Formalizes / promotes \\ \midrule
\endfirsthead
\toprule Theorem & Line & Formalizes / promotes \\ \midrule \endhead
\bottomrule\endfoot
{\scriptsize\ttfamily commuting\_\allowbreak{}pair\_\allowbreak{}square} & {\scriptsize 19} & {\scriptsize G18; RUN:wilson\_\allowbreak{}creator\_\allowbreak{}parent\_\allowbreak{}2026-09-05} \\
{\scriptsize\ttfamily idempotent\_\allowbreak{}square\_\allowbreak{}defect} & {\scriptsize 10} & {\scriptsize G18; RUN:wilson\_\allowbreak{}creator\_\allowbreak{}parent\_\allowbreak{}2026-09-05} \\
{\scriptsize\ttfamily wilson\_\allowbreak{}gap\_\allowbreak{}constant} & {\scriptsize 27} & {\scriptsize G18; RUN:wilson\_\allowbreak{}creator\_\allowbreak{}parent\_\allowbreak{}2026-09-05} \\
\end{longtable}

\subsection*{GlobalWilsonVertical.lean}
\addcontentsline{toc}{subsection}{GlobalWilsonVertical.lean}
\begin{longtable}{@{}p{0.32\linewidth} l p{0.44\linewidth}@{}}
\toprule Theorem & Line & Formalizes / promotes \\ \midrule
\endfirsthead
\toprule Theorem & Line & Formalizes / promotes \\ \midrule \endhead
\bottomrule\endfoot
{\scriptsize\ttfamily far\_\allowbreak{}affine} & {\scriptsize 73} & {\scriptsize RUN:nonlinear\_\allowbreak{}wilson\_\allowbreak{}block\_\allowbreak{}2026-09-05} \\
{\scriptsize\ttfamily half\_\allowbreak{}potential\_\allowbreak{}threshold} & {\scriptsize 129} & {\scriptsize RUN:nonlinear\_\allowbreak{}wilson\_\allowbreak{}block\_\allowbreak{}2026-09-05} \\
{\scriptsize\ttfamily near\_\allowbreak{}affine} & {\scriptsize 54} & {\scriptsize RUN:nonlinear\_\allowbreak{}wilson\_\allowbreak{}block\_\allowbreak{}2026-09-05} \\
{\scriptsize\ttfamily near\_\allowbreak{}cap} & {\scriptsize 39} & {\scriptsize RUN:nonlinear\_\allowbreak{}wilson\_\allowbreak{}block\_\allowbreak{}2026-09-05} \\
{\scriptsize\ttfamily spectral\_\allowbreak{}from\_\allowbreak{}split} & {\scriptsize 86} & {\scriptsize RUN:nonlinear\_\allowbreak{}wilson\_\allowbreak{}block\_\allowbreak{}2026-09-05} \\
{\scriptsize\ttfamily strip\_\allowbreak{}factor} & {\scriptsize 13} & {\scriptsize RUN:nonlinear\_\allowbreak{}wilson\_\allowbreak{}block\_\allowbreak{}2026-09-05} \\
{\scriptsize\ttfamily two\_\allowbreak{}sector\_\allowbreak{}form} & {\scriptsize 112} & {\scriptsize RUN:nonlinear\_\allowbreak{}wilson\_\allowbreak{}block\_\allowbreak{}2026-09-05} \\
\end{longtable}

\subsection*{GroundStateAssembly.lean}
\addcontentsline{toc}{subsection}{GroundStateAssembly.lean}
\begin{longtable}{@{}p{0.32\linewidth} l p{0.44\linewidth}@{}}
\toprule Theorem & Line & Formalizes / promotes \\ \midrule
\endfirsthead
\toprule Theorem & Line & Formalizes / promotes \\ \midrule \endhead
\bottomrule\endfoot
{\scriptsize\ttfamily assembly\_\allowbreak{}energy} & {\scriptsize 504} & {\scriptsize ---} \\
{\scriptsize\ttfamily assembly\_\allowbreak{}kernel\_\allowbreak{}iff} & {\scriptsize 516} & {\scriptsize ---} \\
{\scriptsize\ttfamily assembly\_\allowbreak{}square\_\allowbreak{}lower} & {\scriptsize 544} & {\scriptsize ---} \\
{\scriptsize\ttfamily assembly\_\allowbreak{}symmetric} & {\scriptsize 509} & {\scriptsize ---} \\
{\scriptsize\ttfamily bounded\_\allowbreak{}gap\_\allowbreak{}on\_\allowbreak{}range\_\allowbreak{}of\_\allowbreak{}square} & {\scriptsize 708} & {\scriptsize ---} \\
{\scriptsize\ttfamily bounded\_\allowbreak{}projection\_\allowbreak{}assembly\_\allowbreak{}gap} & {\scriptsize 773} & {\scriptsize ---} \\
{\scriptsize\ttfamily bounded\_\allowbreak{}spectral\_\allowbreak{}gap\_\allowbreak{}of\_\allowbreak{}square} & {\scriptsize 737} & {\scriptsize ---} \\
{\scriptsize\ttfamily finite\_\allowbreak{}assembly\_\allowbreak{}gap} & {\scriptsize 631} & {\scriptsize ---} \\
{\scriptsize\ttfamily finite\_\allowbreak{}spectral\_\allowbreak{}gap\_\allowbreak{}of\_\allowbreak{}square} & {\scriptsize 590} & {\scriptsize ---} \\
{\scriptsize\ttfamily gram\_\allowbreak{}form\_\allowbreak{}nonnegative} & {\scriptsize 841} & {\scriptsize ---} \\
{\scriptsize\ttfamily gram\_\allowbreak{}quadratic\_\allowbreak{}eq\_\allowbreak{}squares} & {\scriptsize 825} & {\scriptsize ---} \\
{\scriptsize\ttfamily integral\_\allowbreak{}cauchy\_\allowbreak{}schwarz\_\allowbreak{}sq} & {\scriptsize 318} & {\scriptsize ---} \\
{\scriptsize\ttfamily integral\_\allowbreak{}mul\_\allowbreak{}eq\_\allowbreak{}L2\_\allowbreak{}inner} & {\scriptsize 308} & {\scriptsize ---} \\
{\scriptsize\ttfamily integral\_\allowbreak{}remainder\_\allowbreak{}measure} & {\scriptsize 329} & {\scriptsize ---} \\
{\scriptsize\ttfamily joint\_\allowbreak{}measure\_\allowbreak{}minorization\_\allowbreak{}covariance} & {\scriptsize 383} & {\scriptsize ---} \\
{\scriptsize\ttfamily joint\_\allowbreak{}withDensity\_\allowbreak{}covariance\_\allowbreak{}variance} & {\scriptsize 441} & {\scriptsize ---} \\
{\scriptsize\ttfamily joint\_\allowbreak{}withDensity\_\allowbreak{}minorization\_\allowbreak{}covariance} & {\scriptsize 421} & {\scriptsize ---} \\
{\scriptsize\ttfamily local\_\allowbreak{}floor\_\allowbreak{}tensorization} & {\scriptsize 648} & {\scriptsize ---} \\
{\scriptsize\ttfamily measure\_\allowbreak{}minorization\_\allowbreak{}covariance} & {\scriptsize 348} & {\scriptsize ---} \\
{\scriptsize\ttfamily measure\_\allowbreak{}poincare\_\allowbreak{}comparison} & {\scriptsize 243} & {\scriptsize ---} \\
{\scriptsize\ttfamily measure\_\allowbreak{}variance\_\allowbreak{}density\_\allowbreak{}lower} & {\scriptsize 214} & {\scriptsize ---} \\
{\scriptsize\ttfamily measure\_\allowbreak{}variance\_\allowbreak{}density\_\allowbreak{}upper} & {\scriptsize 193} & {\scriptsize ---} \\
{\scriptsize\ttfamily measure\_\allowbreak{}variance\_\allowbreak{}le\_\allowbreak{}deviation} & {\scriptsize 184} & {\scriptsize ---} \\
{\scriptsize\ttfamily poincare\_\allowbreak{}density\_\allowbreak{}comparison} & {\scriptsize 91} & {\scriptsize ---} \\
{\scriptsize\ttfamily product\_\allowbreak{}minorization\_\allowbreak{}covariance} & {\scriptsize 128} & {\scriptsize ---} \\
{\scriptsize\ttfamily projection\_\allowbreak{}energy} & {\scriptsize 488} & {\scriptsize ---} \\
{\scriptsize\ttfamily projection\_\allowbreak{}energy\_\allowbreak{}nonneg} & {\scriptsize 494} & {\scriptsize ---} \\
{\scriptsize\ttfamily symmetric\_\allowbreak{}coupling\_\allowbreak{}sum} & {\scriptsize 528} & {\scriptsize ---} \\
{\scriptsize\ttfamily unitary\_\allowbreak{}ground\_\allowbreak{}state\_\allowbreak{}gap} & {\scriptsize 669} & {\scriptsize ---} \\
{\scriptsize\ttfamily variance\_\allowbreak{}density\_\allowbreak{}upper} & {\scriptsize 77} & {\scriptsize ---} \\
{\scriptsize\ttfamily variance\_\allowbreak{}eq\_\allowbreak{}second\_\allowbreak{}moment} & {\scriptsize 45} & {\scriptsize ---} \\
{\scriptsize\ttfamily variance\_\allowbreak{}le\_\allowbreak{}deviation} & {\scriptsize 59} & {\scriptsize ---} \\
{\scriptsize\ttfamily weighted\_\allowbreak{}cauchy\_\allowbreak{}schwarz\_\allowbreak{}sq} & {\scriptsize 114} & {\scriptsize ---} \\
{\scriptsize\ttfamily withDensity\_\allowbreak{}poincare\_\allowbreak{}comparison} & {\scriptsize 267} & {\scriptsize ---} \\
{\scriptsize\ttfamily withDensity\_\allowbreak{}poincare\_\allowbreak{}constant} & {\scriptsize 288} & {\scriptsize ---} \\
{\scriptsize\ttfamily withDensity\_\allowbreak{}variance\_\allowbreak{}floor} & {\scriptsize 232} & {\scriptsize ---} \\
\end{longtable}

\subsection*{HodgeFeshbach.lean}
\addcontentsline{toc}{subsection}{HodgeFeshbach.lean}
\begin{longtable}{@{}p{0.32\linewidth} l p{0.44\linewidth}@{}}
\toprule Theorem & Line & Formalizes / promotes \\ \midrule
\endfirsthead
\toprule Theorem & Line & Formalizes / promotes \\ \midrule \endhead
\bottomrule\endfoot
{\scriptsize\ttfamily carrier\_\allowbreak{}homological\_\allowbreak{}protection} & {\scriptsize 123} & {\scriptsize ---} \\
{\scriptsize\ttfamily carrier\_\allowbreak{}s\_\allowbreak{}eigenvalue} & {\scriptsize 134} & {\scriptsize ---} \\
{\scriptsize\ttfamily cellular\_\allowbreak{}excursion\_\allowbreak{}up\_\allowbreak{}harmonic} & {\scriptsize 384} & {\scriptsize DERIV:\allowbreak{}UNIVERSAL\_\allowbreak{}CELLULAR\_\allowbreak{}HODGE:\allowbreak{}UP\_\allowbreak{}HARMONIC\_\allowbreak{}EXCURSION} \\
{\scriptsize\ttfamily feshbach\_\allowbreak{}hodge\_\allowbreak{}no\_\allowbreak{}coupling} & {\scriptsize 141} & {\scriptsize ---} \\
{\scriptsize\ttfamily feshbach\_\allowbreak{}intermediate\_\allowbreak{}return\_\allowbreak{}annihilation} & {\scriptsize 391} & {\scriptsize DERIV:\allowbreak{}UNIVERSAL\_\allowbreak{}CELLULAR\_\allowbreak{}HODGE:\allowbreak{}RETAINED\_\allowbreak{}VECTOR\_\allowbreak{}RANK\_\allowbreak{}ONE} \\
{\scriptsize\ttfamily feshbach\_\allowbreak{}q\_\allowbreak{}annihilates\_\allowbreak{}l\_\allowbreak{}up} & {\scriptsize 377} & {\scriptsize ---} \\
{\scriptsize\ttfamily hodge\_\allowbreak{}algebra\_\allowbreak{}relations} & {\scriptsize 33} & {\scriptsize ---} \\
{\scriptsize\ttfamily hodge\_\allowbreak{}decomposition} & {\scriptsize 101} & {\scriptsize ---} \\
{\scriptsize\ttfamily hodge\_\allowbreak{}feshbach\_\allowbreak{}words\_\allowbreak{}count} & {\scriptsize 268} & {\scriptsize ---} \\
{\scriptsize\ttfamily hodge\_\allowbreak{}ker\_\allowbreak{}disjoint} & {\scriptsize 55} & {\scriptsize ---} \\
{\scriptsize\ttfamily hodge\_\allowbreak{}pi\_\allowbreak{}cleared} & {\scriptsize 447} & {\scriptsize DERIV:\allowbreak{}UNIVERSAL\_\allowbreak{}CELLULAR\_\allowbreak{}HODGE:\allowbreak{}SHIFTED\_\allowbreak{}OPERATOR\_\allowbreak{}POWERS} \\
{\scriptsize\ttfamily hodge\_\allowbreak{}sandwiched\_\allowbreak{}power} & {\scriptsize 472} & {\scriptsize DERIV:\allowbreak{}UNIVERSAL\_\allowbreak{}CELLULAR\_\allowbreak{}HODGE:\allowbreak{}SHIFTED\_\allowbreak{}OPERATOR\_\allowbreak{}POWERS} \\
{\scriptsize\ttfamily hodge\_\allowbreak{}shifted\_\allowbreak{}power} & {\scriptsize 458} & {\scriptsize DERIV:\allowbreak{}UNIVERSAL\_\allowbreak{}CELLULAR\_\allowbreak{}HODGE:\allowbreak{}SHIFTED\_\allowbreak{}OPERATOR\_\allowbreak{}POWERS} \\
{\scriptsize\ttfamily hodge\_\allowbreak{}vector\_\allowbreak{}decomposition} & {\scriptsize 73} & {\scriptsize ---} \\
{\scriptsize\ttfamily hodge\_\allowbreak{}word\_\allowbreak{}is\_\allowbreak{}scalar} & {\scriptsize 355} & {\scriptsize ---} \\
{\scriptsize\ttfamily normalized\_\allowbreak{}carrier\_\allowbreak{}projection\_\allowbreak{}fixes\_\allowbreak{}carrier} & {\scriptsize 173} & {\scriptsize ---} \\
{\scriptsize\ttfamily r\_\allowbreak{}excitation\_\allowbreak{}carrier\_\allowbreak{}projection\_\allowbreak{}zero} & {\scriptsize 195} & {\scriptsize ---} \\
{\scriptsize\ttfamily r\_\allowbreak{}excitation\_\allowbreak{}in\_\allowbreak{}ker\_\allowbreak{}up} & {\scriptsize 154} & {\scriptsize ---} \\
{\scriptsize\ttfamily r\_\allowbreak{}excitation\_\allowbreak{}orthogonal} & {\scriptsize 180} & {\scriptsize ---} \\
{\scriptsize\ttfamily r\_\allowbreak{}excitation\_\allowbreak{}up\_\allowbreak{}harmonic} & {\scriptsize 188} & {\scriptsize ---} \\
{\scriptsize\ttfamily rank\_\allowbreak{}one\_\allowbreak{}laplacian\_\allowbreak{}carrier\_\allowbreak{}eigenvalue} & {\scriptsize 167} & {\scriptsize ---} \\
{\scriptsize\ttfamily rank\_\allowbreak{}one\_\allowbreak{}rur\_\allowbreak{}carrier\_\allowbreak{}factorization} & {\scriptsize 215} & {\scriptsize ---} \\
{\scriptsize\ttfamily rank\_\allowbreak{}one\_\allowbreak{}rur\_\allowbreak{}cleared\_\allowbreak{}defect\_\allowbreak{}zero} & {\scriptsize 224} & {\scriptsize ---} \\
{\scriptsize\ttfamily rank\_\allowbreak{}one\_\allowbreak{}rur\_\allowbreak{}vector\_\allowbreak{}factorization} & {\scriptsize 207} & {\scriptsize ---} \\
{\scriptsize\ttfamily rr\_\allowbreak{}carrier\_\allowbreak{}polynomial\_\allowbreak{}remainder} & {\scriptsize 318} & {\scriptsize ---} \\
{\scriptsize\ttfamily rur\_\allowbreak{}carrier\_\allowbreak{}polynomial\_\allowbreak{}cleared\_\allowbreak{}identity} & {\scriptsize 327} & {\scriptsize ---} \\
{\scriptsize\ttfamily rur\_\allowbreak{}two\_\allowbreak{}r\_\allowbreak{}no\_\allowbreak{}u\_\allowbreak{}false} & {\scriptsize 304} & {\scriptsize ---} \\
{\scriptsize\ttfamily tetrahedral\_\allowbreak{}hodge\_\allowbreak{}duality} & {\scriptsize 405} & {\scriptsize ---} \\
{\scriptsize\ttfamily traceless\_\allowbreak{}compression\_\allowbreak{}vanishes} & {\scriptsize 431} & {\scriptsize ---} \\
{\scriptsize\ttfamily two\_\allowbreak{}r\_\allowbreak{}no\_\allowbreak{}u\_\allowbreak{}words\_\allowbreak{}count} & {\scriptsize 288} & {\scriptsize ---} \\
{\scriptsize\ttfamily two\_\allowbreak{}r\_\allowbreak{}no\_\allowbreak{}u\_\allowbreak{}words\_\allowbreak{}enum} & {\scriptsize 291} & {\scriptsize ---} \\
{\scriptsize\ttfamily ur\_\allowbreak{}carrier\_\allowbreak{}polynomial\_\allowbreak{}value} & {\scriptsize 333} & {\scriptsize ---} \\
\end{longtable}

\subsection*{PlateauObstruction.lean}
\addcontentsline{toc}{subsection}{PlateauObstruction.lean}
\begin{longtable}{@{}p{0.32\linewidth} l p{0.44\linewidth}@{}}
\toprule Theorem & Line & Formalizes / promotes \\ \midrule
\endfirsthead
\toprule Theorem & Line & Formalizes / promotes \\ \midrule \endhead
\bottomrule\endfoot
{\scriptsize\ttfamily admissible\_\allowbreak{}logarithmic\_\allowbreak{}minimum} & {\scriptsize 71} & {\scriptsize ---} \\
{\scriptsize\ttfamily arbitrarily\_\allowbreak{}small\_\allowbreak{}plateau\_\allowbreak{}obstruction} & {\scriptsize 266} & {\scriptsize ---} \\
{\scriptsize\ttfamily bound\_\allowbreak{}time\_\allowbreak{}rescaling} & {\scriptsize 142} & {\scriptsize ---} \\
{\scriptsize\ttfamily endpoint\_\allowbreak{}minimum} & {\scriptsize 87} & {\scriptsize ---} \\
{\scriptsize\ttfamily interior\_\allowbreak{}minimum\_\allowbreak{}value} & {\scriptsize 99} & {\scriptsize ---} \\
{\scriptsize\ttfamily logarithmic\_\allowbreak{}time\_\allowbreak{}is\_\allowbreak{}stationary} & {\scriptsize 47} & {\scriptsize ---} \\
{\scriptsize\ttfamily physical\_\allowbreak{}time\_\allowbreak{}minimum\_\allowbreak{}value} & {\scriptsize 150} & {\scriptsize ---} \\
{\scriptsize\ttfamily positive\_\allowbreak{}slow\_\allowbreak{}mode\_\allowbreak{}prevents\_\allowbreak{}faster\_\allowbreak{}bound} & {\scriptsize 239} & {\scriptsize ---} \\
{\scriptsize\ttfamily r4\_\allowbreak{}r5\_\allowbreak{}interior\_\allowbreak{}optimum} & {\scriptsize 162} & {\scriptsize ---} \\
{\scriptsize\ttfamily r4\_\allowbreak{}r5\_\allowbreak{}optimized\_\allowbreak{}upper\_\allowbreak{}bound} & {\scriptsize 194} & {\scriptsize ---} \\
{\scriptsize\ttfamily rescaled\_\allowbreak{}correlation\_\allowbreak{}identity} & {\scriptsize 228} & {\scriptsize ---} \\
{\scriptsize\ttfamily slow\_\allowbreak{}mode\_\allowbreak{}fits\_\allowbreak{}positive\_\allowbreak{}plateau} & {\scriptsize 213} & {\scriptsize ---} \\
{\scriptsize\ttfamily stationary\_\allowbreak{}global\_\allowbreak{}minimum} & {\scriptsize 20} & {\scriptsize ---} \\
{\scriptsize\ttfamily theta\_\allowbreak{}minimum\_\allowbreak{}value} & {\scriptsize 122} & {\scriptsize ---} \\
\end{longtable}

\subsection*{PolymerCluster.lean}
\addcontentsline{toc}{subsection}{PolymerCluster.lean}
\begin{longtable}{@{}p{0.32\linewidth} l p{0.44\linewidth}@{}}
\toprule Theorem & Line & Formalizes / promotes \\ \midrule
\endfirsthead
\toprule Theorem & Line & Formalizes / promotes \\ \midrule \endhead
\bottomrule\endfoot
{\scriptsize\ttfamily betaStabilityWindow\_\allowbreak{}pos} & {\scriptsize 221} & {\scriptsize ---} \\
{\scriptsize\ttfamily beta\_\allowbreak{}le\_\allowbreak{}window\_\allowbreak{}implies\_\allowbreak{}tiltedKp\_\allowbreak{}le} & {\scriptsize 231} & {\scriptsize ---} \\
{\scriptsize\ttfamily beta\_\allowbreak{}le\_\allowbreak{}window\_\allowbreak{}implies\_\allowbreak{}tree\_\allowbreak{}bound} & {\scriptsize 246} & {\scriptsize ---} \\
{\scriptsize\ttfamily clusterMassGap\_\allowbreak{}ge\_\allowbreak{}log\_\allowbreak{}kp} & {\scriptsize 299} & {\scriptsize ---} \\
{\scriptsize\ttfamily clusterMassGap\_\allowbreak{}pos} & {\scriptsize 289} & {\scriptsize ---} \\
{\scriptsize\ttfamily clusterMassGap\_\allowbreak{}pos\_\allowbreak{}of\_\allowbreak{}strict\_\allowbreak{}kp} & {\scriptsize 321} & {\scriptsize ---} \\
{\scriptsize\ttfamily cluster\_\allowbreak{}correlation\_\allowbreak{}exponential\_\allowbreak{}decay} & {\scriptsize 329} & {\scriptsize ---} \\
{\scriptsize\ttfamily cluster\_\allowbreak{}free\_\allowbreak{}energy\_\allowbreak{}density\_\allowbreak{}bound\_\allowbreak{}le} & {\scriptsize 258} & {\scriptsize ---} \\
{\scriptsize\ttfamily cluster\_\allowbreak{}free\_\allowbreak{}energy\_\allowbreak{}density\_\allowbreak{}bound\_\allowbreak{}le\_\allowbreak{}window} & {\scriptsize 266} & {\scriptsize ---} \\
{\scriptsize\ttfamily cluster\_\allowbreak{}free\_\allowbreak{}energy\_\allowbreak{}density\_\allowbreak{}bound\_\allowbreak{}nonneg} & {\scriptsize 273} & {\scriptsize ---} \\
{\scriptsize\ttfamily geometric\_\allowbreak{}majorant\_\allowbreak{}le\_\allowbreak{}iff} & {\scriptsize 59} & {\scriptsize ---} \\
{\scriptsize\ttfamily geometric\_\allowbreak{}majorant\_\allowbreak{}mono} & {\scriptsize 50} & {\scriptsize ---} \\
{\scriptsize\ttfamily geometric\_\allowbreak{}majorant\_\allowbreak{}nonneg} & {\scriptsize 38} & {\scriptsize ---} \\
{\scriptsize\ttfamily geometric\_\allowbreak{}majorant\_\allowbreak{}pos} & {\scriptsize 44} & {\scriptsize ---} \\
{\scriptsize\ttfamily geometric\_\allowbreak{}series\_\allowbreak{}from\_\allowbreak{}one\_\allowbreak{}hasSum} & {\scriptsize 24} & {\scriptsize ---} \\
{\scriptsize\ttfamily geometric\_\allowbreak{}series\_\allowbreak{}from\_\allowbreak{}one\_\allowbreak{}tsum} & {\scriptsize 33} & {\scriptsize ---} \\
{\scriptsize\ttfamily log\_\allowbreak{}one\_\allowbreak{}add\_\allowbreak{}div\_\allowbreak{}pos} & {\scriptsize 312} & {\scriptsize ---} \\
{\scriptsize\ttfamily strict\_\allowbreak{}kotecky\_\allowbreak{}preiss\_\allowbreak{}criterion} & {\scriptsize 170} & {\scriptsize ---} \\
{\scriptsize\ttfamily strict\_\allowbreak{}kotecky\_\allowbreak{}preiss\_\allowbreak{}iff} & {\scriptsize 186} & {\scriptsize ---} \\
{\scriptsize\ttfamily strict\_\allowbreak{}kp\_\allowbreak{}criterion\_\allowbreak{}lt\_\allowbreak{}one} & {\scriptsize 162} & {\scriptsize ---} \\
{\scriptsize\ttfamily tiltedKpParameter\_\allowbreak{}at\_\allowbreak{}alpha\_\allowbreak{}zero} & {\scriptsize 156} & {\scriptsize ---} \\
{\scriptsize\ttfamily tiltedKpParameter\_\allowbreak{}nonneg} & {\scriptsize 149} & {\scriptsize ---} \\
{\scriptsize\ttfamily tilted\_\allowbreak{}kotecky\_\allowbreak{}preiss\_\allowbreak{}convergent} & {\scriptsize 206} & {\scriptsize ---} \\
{\scriptsize\ttfamily tilted\_\allowbreak{}kotecky\_\allowbreak{}preiss\_\allowbreak{}tree\_\allowbreak{}bound} & {\scriptsize 199} & {\scriptsize ---} \\
{\scriptsize\ttfamily uniform\_\allowbreak{}source\_\allowbreak{}radius\_\allowbreak{}bound\_\allowbreak{}le} & {\scriptsize 340} & {\scriptsize ---} \\
{\scriptsize\ttfamily uniform\_\allowbreak{}source\_\allowbreak{}radius\_\allowbreak{}bound\_\allowbreak{}le\_\allowbreak{}window} & {\scriptsize 373} & {\scriptsize ---} \\
{\scriptsize\ttfamily uniform\_\allowbreak{}source\_\allowbreak{}radius\_\allowbreak{}bound\_\allowbreak{}nonneg} & {\scriptsize 367} & {\scriptsize ---} \\
{\scriptsize\ttfamily unweightedActivityPower\_\allowbreak{}le\_\allowbreak{}weighted} & {\scriptsize 85} & {\scriptsize ---} \\
{\scriptsize\ttfamily weightedActivityBase\_\allowbreak{}beta\_\allowbreak{}zero} & {\scriptsize 115} & {\scriptsize ---} \\
{\scriptsize\ttfamily weightedActivityBase\_\allowbreak{}geometric\_\allowbreak{}hasSum} & {\scriptsize 99} & {\scriptsize ---} \\
{\scriptsize\ttfamily weightedActivityBase\_\allowbreak{}nonneg} & {\scriptsize 73} & {\scriptsize ---} \\
{\scriptsize\ttfamily weightedActivityBase\_\allowbreak{}small\_\allowbreak{}beta\_\allowbreak{}regime} & {\scriptsize 122} & {\scriptsize ---} \\
{\scriptsize\ttfamily weighted\_\allowbreak{}geometric\_\allowbreak{}majorant\_\allowbreak{}nonneg} & {\scriptsize 108} & {\scriptsize ---} \\
\end{longtable}

\subsection*{ResolventLocalization.lean}
\addcontentsline{toc}{subsection}{ResolventLocalization.lean}
\begin{longtable}{@{}p{0.32\linewidth} l p{0.44\linewidth}@{}}
\toprule Theorem & Line & Formalizes / promotes \\ \midrule
\endfirsthead
\toprule Theorem & Line & Formalizes / promotes \\ \midrule \endhead
\bottomrule\endfoot
{\scriptsize\ttfamily cubicWeight\_\allowbreak{}nonneg} & {\scriptsize 329} & {\scriptsize ---} \\
{\scriptsize\ttfamily cubic\_\allowbreak{}geometric\_\allowbreak{}hasSum} & {\scriptsize 337} & {\scriptsize ---} \\
{\scriptsize\ttfamily cubic\_\allowbreak{}localized\_\allowbreak{}pairing} & {\scriptsize 351} & {\scriptsize ---} \\
{\scriptsize\ttfamily finite\_\allowbreak{}source\_\allowbreak{}propagation} & {\scriptsize 244} & {\scriptsize ---} \\
{\scriptsize\ttfamily integer\_\allowbreak{}geometric\_\allowbreak{}hasSum} & {\scriptsize 301} & {\scriptsize ---} \\
{\scriptsize\ttfamily localized\_\allowbreak{}finite\_\allowbreak{}volume\_\allowbreak{}bound} & {\scriptsize 289} & {\scriptsize ---} \\
{\scriptsize\ttfamily localized\_\allowbreak{}kernel\_\allowbreak{}pairing} & {\scriptsize 262} & {\scriptsize ---} \\
{\scriptsize\ttfamily mul\_\allowbreak{}neumann} & {\scriptsize 29} & {\scriptsize ---} \\
{\scriptsize\ttfamily neumann\_\allowbreak{}mul} & {\scriptsize 34} & {\scriptsize ---} \\
{\scriptsize\ttfamily neumann\_\allowbreak{}quadratic\_\allowbreak{}expansion} & {\scriptsize 389} & {\scriptsize ---} \\
{\scriptsize\ttfamily neumann\_\allowbreak{}quadratic\_\allowbreak{}remainder\_\allowbreak{}identity} & {\scriptsize 378} & {\scriptsize ---} \\
{\scriptsize\ttfamily neumann\_\allowbreak{}quadratic\_\allowbreak{}remainder\_\allowbreak{}pairing\_\allowbreak{}bound} & {\scriptsize 431} & {\scriptsize ---} \\
{\scriptsize\ttfamily neumann\_\allowbreak{}summable} & {\scriptsize 24} & {\scriptsize ---} \\
{\scriptsize\ttfamily norm\_\allowbreak{}neumann\_\allowbreak{}le} & {\scriptsize 39} & {\scriptsize ---} \\
{\scriptsize\ttfamily norm\_\allowbreak{}neumann\_\allowbreak{}quadratic\_\allowbreak{}remainder\_\allowbreak{}le} & {\scriptsize 396} & {\scriptsize ---} \\
{\scriptsize\ttfamily norm\_\allowbreak{}neumann\_\allowbreak{}sub\_\allowbreak{}one\_\allowbreak{}le} & {\scriptsize 44} & {\scriptsize ---} \\
{\scriptsize\ttfamily norm\_\allowbreak{}neumann\_\allowbreak{}sub\_\allowbreak{}one\_\allowbreak{}le\_\allowbreak{}budget} & {\scriptsize 59} & {\scriptsize ---} \\
{\scriptsize\ttfamily norm\_\allowbreak{}sandwiched\_\allowbreak{}neumann\_\allowbreak{}quadratic\_\allowbreak{}remainder\_\allowbreak{}le} & {\scriptsize 409} & {\scriptsize ---} \\
{\scriptsize\ttfamily normalized\_\allowbreak{}fast\_\allowbreak{}coercive} & {\scriptsize 183} & {\scriptsize ---} \\
{\scriptsize\ttfamily normalized\_\allowbreak{}schur\_\allowbreak{}unique\_\allowbreak{}minimizer} & {\scriptsize 203} & {\scriptsize ---} \\
{\scriptsize\ttfamily normalized\_\allowbreak{}selected\_\allowbreak{}pairing\_\allowbreak{}bound} & {\scriptsize 92} & {\scriptsize ---} \\
{\scriptsize\ttfamily schur\_\allowbreak{}minimizer\_\allowbreak{}lower\_\allowbreak{}bound} & {\scriptsize 173} & {\scriptsize ---} \\
{\scriptsize\ttfamily schur\_\allowbreak{}minimizer\_\allowbreak{}value} & {\scriptsize 164} & {\scriptsize ---} \\
{\scriptsize\ttfamily schur\_\allowbreak{}square\_\allowbreak{}completion} & {\scriptsize 150} & {\scriptsize ---} \\
{\scriptsize\ttfamily signed\_\allowbreak{}resolvent\_\allowbreak{}identity} & {\scriptsize 79} & {\scriptsize ---} \\
{\scriptsize\ttfamily signed\_\allowbreak{}resolvent\_\allowbreak{}pairing} & {\scriptsize 133} & {\scriptsize ---} \\
{\scriptsize\ttfamily weighted\_\allowbreak{}inverse\_\allowbreak{}difference\_\allowbreak{}bound} & {\scriptsize 121} & {\scriptsize ---} \\
{\scriptsize\ttfamily weighted\_\allowbreak{}selected\_\allowbreak{}pairing\_\allowbreak{}bound} & {\scriptsize 111} & {\scriptsize ---} \\
\end{longtable}

\subsection*{RootedScalarBounds.lean}
\addcontentsline{toc}{subsection}{RootedScalarBounds.lean}
\begin{longtable}{@{}p{0.32\linewidth} l p{0.44\linewidth}@{}}
\toprule Theorem & Line & Formalizes / promotes \\ \midrule
\endfirsthead
\toprule Theorem & Line & Formalizes / promotes \\ \midrule \endhead
\bottomrule\endfoot
{\scriptsize\ttfamily normalized\_\allowbreak{}taylor\_\allowbreak{}growth} & {\scriptsize 13} & {\scriptsize G18; RUN:wilson\_\allowbreak{}rooted\_\allowbreak{}contraction\_\allowbreak{}2026-09-05} \\
\end{longtable}

\subsection*{SC17Riccati.lean}
\addcontentsline{toc}{subsection}{SC17Riccati.lean}
\begin{longtable}{@{}p{0.32\linewidth} l p{0.44\linewidth}@{}}
\toprule Theorem & Line & Formalizes / promotes \\ \midrule
\endfirsthead
\toprule Theorem & Line & Formalizes / promotes \\ \midrule \endhead
\bottomrule\endfoot
{\scriptsize\ttfamily continuous\_\allowbreak{}riccati\_\allowbreak{}barrier} & {\scriptsize 79} & {\scriptsize ---} \\
{\scriptsize\ttfamily defaultBeta\_\allowbreak{}pos} & {\scriptsize 244} & {\scriptsize ---} \\
{\scriptsize\ttfamily defaultBeta\_\allowbreak{}sq} & {\scriptsize 246} & {\scriptsize ---} \\
{\scriptsize\ttfamily default\_\allowbreak{}parameter\_\allowbreak{}endpoint} & {\scriptsize 268} & {\scriptsize ---} \\
{\scriptsize\ttfamily default\_\allowbreak{}parameter\_\allowbreak{}identity} & {\scriptsize 252} & {\scriptsize ---} \\
{\scriptsize\ttfamily default\_\allowbreak{}parameter\_\allowbreak{}interval} & {\scriptsize 285} & {\scriptsize ---} \\
{\scriptsize\ttfamily default\_\allowbreak{}parameter\_\allowbreak{}monotone} & {\scriptsize 260} & {\scriptsize ---} \\
{\scriptsize\ttfamily default\_\allowbreak{}reference\_\allowbreak{}budget} & {\scriptsize 307} & {\scriptsize ---} \\
{\scriptsize\ttfamily default\_\allowbreak{}reference\_\allowbreak{}product} & {\scriptsize 276} & {\scriptsize ---} \\
{\scriptsize\ttfamily default\_\allowbreak{}smallRoot\_\allowbreak{}and\_\allowbreak{}gap} & {\scriptsize 315} & {\scriptsize ---} \\
{\scriptsize\ttfamily default\_\allowbreak{}small\_\allowbreak{}defect\_\allowbreak{}criterion} & {\scriptsize 301} & {\scriptsize ---} \\
{\scriptsize\ttfamily exists\_\allowbreak{}riccati\_\allowbreak{}fixedPoint} & {\scriptsize 158} & {\scriptsize ---} \\
{\scriptsize\ttfamily exists\_\allowbreak{}riccati\_\allowbreak{}solution} & {\scriptsize 199} & {\scriptsize ---} \\
{\scriptsize\ttfamily exists\_\allowbreak{}scalar\_\allowbreak{}reference\_\allowbreak{}riccati} & {\scriptsize 212} & {\scriptsize ---} \\
{\scriptsize\ttfamily quadratic\_\allowbreak{}forbidden\_\allowbreak{}interval} & {\scriptsize 64} & {\scriptsize ---} \\
{\scriptsize\ttfamily riccatiMap\_\allowbreak{}lipschitz\_\allowbreak{}on\_\allowbreak{}ball} & {\scriptsize 131} & {\scriptsize ---} \\
{\scriptsize\ttfamily riccatiMap\_\allowbreak{}norm\_\allowbreak{}le} & {\scriptsize 114} & {\scriptsize ---} \\
{\scriptsize\ttfamily smallRoot\_\allowbreak{}contraction\_\allowbreak{}margin} & {\scriptsize 51} & {\scriptsize ---} \\
{\scriptsize\ttfamily smallRoot\_\allowbreak{}equation} & {\scriptsize 36} & {\scriptsize ---} \\
{\scriptsize\ttfamily smallRoot\_\allowbreak{}lt\_\allowbreak{}half\_\allowbreak{}inverse} & {\scriptsize 45} & {\scriptsize ---} \\
{\scriptsize\ttfamily smallRoot\_\allowbreak{}lt\_\allowbreak{}reference\_\allowbreak{}floor} & {\scriptsize 59} & {\scriptsize ---} \\
{\scriptsize\ttfamily smallRoot\_\allowbreak{}nonneg} & {\scriptsize 28} & {\scriptsize ---} \\
\end{longtable}

\subsection*{SourceRadiusGrowth.lean}
\addcontentsline{toc}{subsection}{SourceRadiusGrowth.lean}
\begin{longtable}{@{}p{0.32\linewidth} l p{0.44\linewidth}@{}}
\toprule Theorem & Line & Formalizes / promotes \\ \midrule
\endfirsthead
\toprule Theorem & Line & Formalizes / promotes \\ \midrule \endhead
\bottomrule\endfoot
{\scriptsize\ttfamily boundedTwoPointSource\_\allowbreak{}abs\_\allowbreak{}le} & {\scriptsize 147} & {\scriptsize ---} \\
{\scriptsize\ttfamily independent\_\allowbreak{}exp\_\allowbreak{}moment\_\allowbreak{}difference} & {\scriptsize 38} & {\scriptsize ---} \\
{\scriptsize\ttfamily independent\_\allowbreak{}exp\_\allowbreak{}radius\_\allowbreak{}squared\_\allowbreak{}tendsto} & {\scriptsize 77} & {\scriptsize ---} \\
{\scriptsize\ttfamily independent\_\allowbreak{}exp\_\allowbreak{}sum\_\allowbreak{}integral} & {\scriptsize 25} & {\scriptsize ---} \\
{\scriptsize\ttfamily independent\_\allowbreak{}exp\_\allowbreak{}variance} & {\scriptsize 104} & {\scriptsize ---} \\
{\scriptsize\ttfamily independent\_\allowbreak{}nonzero\_\allowbreak{}source\_\allowbreak{}radius\_\allowbreak{}unbounded} & {\scriptsize 124} & {\scriptsize ---} \\
{\scriptsize\ttfamily independent\_\allowbreak{}sum\_\allowbreak{}abs\_\allowbreak{}le} & {\scriptsize 94} & {\scriptsize ---} \\
{\scriptsize\ttfamily independent\_\allowbreak{}twoPoint\_\allowbreak{}radius\_\allowbreak{}unbounded} & {\scriptsize 160} & {\scriptsize ---} \\
{\scriptsize\ttfamily power\_\allowbreak{}difference\_\allowbreak{}lower} & {\scriptsize 51} & {\scriptsize ---} \\
{\scriptsize\ttfamily power\_\allowbreak{}difference\_\allowbreak{}per\_\allowbreak{}site\_\allowbreak{}tendsto} & {\scriptsize 64} & {\scriptsize ---} \\
\end{longtable}

\subsection*{SourceTilt.lean}
\addcontentsline{toc}{subsection}{SourceTilt.lean}
\begin{longtable}{@{}p{0.32\linewidth} l p{0.44\linewidth}@{}}
\toprule Theorem & Line & Formalizes / promotes \\ \midrule
\endfirsthead
\toprule Theorem & Line & Formalizes / promotes \\ \midrule \endhead
\bottomrule\endfoot
{\scriptsize\ttfamily bounded\_\allowbreak{}footprint\_\allowbreak{}radius} & {\scriptsize 169} & {\scriptsize ---} \\
{\scriptsize\ttfamily bounded\_\allowbreak{}footprint\_\allowbreak{}radius\_\allowbreak{}sharp} & {\scriptsize 197} & {\scriptsize ---} \\
{\scriptsize\ttfamily centeredExpL2\_\allowbreak{}norm\_\allowbreak{}le} & {\scriptsize 153} & {\scriptsize ---} \\
{\scriptsize\ttfamily centeredExpL2\_\allowbreak{}norm\_\allowbreak{}sq} & {\scriptsize 109} & {\scriptsize ---} \\
{\scriptsize\ttfamily centeredExpL2\_\allowbreak{}orthogonal\_\allowbreak{}constants} & {\scriptsize 128} & {\scriptsize ---} \\
{\scriptsize\ttfamily centered\_\allowbreak{}exp\_\allowbreak{}integral\_\allowbreak{}sq} & {\scriptsize 78} & {\scriptsize ---} \\
{\scriptsize\ttfamily exp\_\allowbreak{}source\_\allowbreak{}bounds} & {\scriptsize 28} & {\scriptsize ---} \\
{\scriptsize\ttfamily exp\_\allowbreak{}source\_\allowbreak{}integral\_\allowbreak{}le} & {\scriptsize 47} & {\scriptsize ---} \\
{\scriptsize\ttfamily exp\_\allowbreak{}source\_\allowbreak{}memLp} & {\scriptsize 40} & {\scriptsize ---} \\
{\scriptsize\ttfamily exp\_\allowbreak{}source\_\allowbreak{}second\_\allowbreak{}moment\_\allowbreak{}gap} & {\scriptsize 313} & {\scriptsize ---} \\
{\scriptsize\ttfamily exp\_\allowbreak{}source\_\allowbreak{}variance\_\allowbreak{}le} & {\scriptsize 93} & {\scriptsize ---} \\
{\scriptsize\ttfamily exp\_\allowbreak{}source\_\allowbreak{}variance\_\allowbreak{}pos} & {\scriptsize 285} & {\scriptsize ---} \\
{\scriptsize\ttfamily finite\_\allowbreak{}source\_\allowbreak{}abs\_\allowbreak{}le} & {\scriptsize 56} & {\scriptsize ---} \\
{\scriptsize\ttfamily finite\_\allowbreak{}source\_\allowbreak{}partition\_\allowbreak{}bound} & {\scriptsize 67} & {\scriptsize ---} \\
{\scriptsize\ttfamily one\_\allowbreak{}lt\_\allowbreak{}integral\_\allowbreak{}exp\_\allowbreak{}of\_\allowbreak{}centered\_\allowbreak{}nonzero} & {\scriptsize 235} & {\scriptsize ---} \\
{\scriptsize\ttfamily raw\_\allowbreak{}tilt\_\allowbreak{}gt\_\allowbreak{}one\_\allowbreak{}of\_\allowbreak{}su2\_\allowbreak{}moments} & {\scriptsize 274} & {\scriptsize ---} \\
{\scriptsize\ttfamily symmetricTwoPoint\_\allowbreak{}integral} & {\scriptsize 336} & {\scriptsize ---} \\
{\scriptsize\ttfamily symmetricTwoPoint\_\allowbreak{}raw\_\allowbreak{}tilt\_\allowbreak{}gt\_\allowbreak{}one} & {\scriptsize 359} & {\scriptsize ---} \\
{\scriptsize\ttfamily symmetricTwoPoint\_\allowbreak{}source\_\allowbreak{}bound} & {\scriptsize 347} & {\scriptsize ---} \\
\end{longtable}

\subsection*{SpectralReconstruction.lean}
\addcontentsline{toc}{subsection}{SpectralReconstruction.lean}
\begin{longtable}{@{}p{0.32\linewidth} l p{0.44\linewidth}@{}}
\toprule Theorem & Line & Formalizes / promotes \\ \midrule
\endfirsthead
\toprule Theorem & Line & Formalizes / promotes \\ \midrule \endhead
\bottomrule\endfoot
{\scriptsize\ttfamily balanced\_\allowbreak{}rate\_\allowbreak{}pos} & {\scriptsize 127} & {\scriptsize ---} \\
{\scriptsize\ttfamily balancing\_\allowbreak{}radius\_\allowbreak{}nonneg} & {\scriptsize 122} & {\scriptsize ---} \\
{\scriptsize\ttfamily deleted\_\allowbreak{}frame\_\allowbreak{}invisible\_\allowbreak{}sector} & {\scriptsize 342} & {\scriptsize ---} \\
{\scriptsize\ttfamily laplace\_\allowbreak{}decay\_\allowbreak{}excludes\_\allowbreak{}low\_\allowbreak{}energy} & {\scriptsize 70} & {\scriptsize ---} \\
{\scriptsize\ttfamily laplace\_\allowbreak{}decay\_\allowbreak{}support\_\allowbreak{}gap} & {\scriptsize 99} & {\scriptsize ---} \\
{\scriptsize\ttfamily laplace\_\allowbreak{}integrable} & {\scriptsize 28} & {\scriptsize ---} \\
{\scriptsize\ttfamily localization\_\allowbreak{}balanced\_\allowbreak{}decay} & {\scriptsize 133} & {\scriptsize ---} \\
{\scriptsize\ttfamily localization\_\allowbreak{}excludes\_\allowbreak{}low\_\allowbreak{}energy} & {\scriptsize 152} & {\scriptsize ---} \\
{\scriptsize\ttfamily low\_\allowbreak{}energy\_\allowbreak{}laplace\_\allowbreak{}lower\_\allowbreak{}bound} & {\scriptsize 41} & {\scriptsize ---} \\
{\scriptsize\ttfamily low\_\allowbreak{}energy\_\allowbreak{}mass\_\allowbreak{}bound} & {\scriptsize 55} & {\scriptsize ---} \\
{\scriptsize\ttfamily observable\_\allowbreak{}frame\_\allowbreak{}correlation} & {\scriptsize 331} & {\scriptsize ---} \\
{\scriptsize\ttfamily observable\_\allowbreak{}frame\_\allowbreak{}energy} & {\scriptsize 314} & {\scriptsize ---} \\
{\scriptsize\ttfamily observable\_\allowbreak{}frame\_\allowbreak{}gram} & {\scriptsize 306} & {\scriptsize ---} \\
{\scriptsize\ttfamily observable\_\allowbreak{}frame\_\allowbreak{}lower\_\allowbreak{}bound} & {\scriptsize 322} & {\scriptsize ---} \\
{\scriptsize\ttfamily projection\_\allowbreak{}zero\_\allowbreak{}from\_\allowbreak{}total\_\allowbreak{}family} & {\scriptsize 171} & {\scriptsize ---} \\
{\scriptsize\ttfamily reflection\_\allowbreak{}kernel\_\allowbreak{}gram} & {\scriptsize 237} & {\scriptsize ---} \\
{\scriptsize\ttfamily reflection\_\allowbreak{}kernel\_\allowbreak{}nonnegative} & {\scriptsize 259} & {\scriptsize ---} \\
{\scriptsize\ttfamily total\_\allowbreak{}family\_\allowbreak{}localization\_\allowbreak{}projection\_\allowbreak{}gap} & {\scriptsize 202} & {\scriptsize ---} \\
{\scriptsize\ttfamily total\_\allowbreak{}family\_\allowbreak{}spectral\_\allowbreak{}projection\_\allowbreak{}gap} & {\scriptsize 182} & {\scriptsize ---} \\
{\scriptsize\ttfamily transfer\_\allowbreak{}energy\_\allowbreak{}nonnegative} & {\scriptsize 274} & {\scriptsize ---} \\
{\scriptsize\ttfamily transfer\_\allowbreak{}energy\_\allowbreak{}reconstructs\_\allowbreak{}ratio} & {\scriptsize 279} & {\scriptsize ---} \\
{\scriptsize\ttfamily transfer\_\allowbreak{}power\_\allowbreak{}eq\_\allowbreak{}physical\_\allowbreak{}decay} & {\scriptsize 287} & {\scriptsize ---} \\
\end{longtable}

\subsection*{TheoryCurrentBridges.lean}
\addcontentsline{toc}{subsection}{TheoryCurrentBridges.lean}
\begin{longtable}{@{}p{0.32\linewidth} l p{0.44\linewidth}@{}}
\toprule Theorem & Line & Formalizes / promotes \\ \midrule
\endfirsthead
\toprule Theorem & Line & Formalizes / promotes \\ \midrule \endhead
\bottomrule\endfoot
{\scriptsize\ttfamily factored\_\allowbreak{}inverse\_\allowbreak{}energy\_\allowbreak{}without\_\allowbreak{}uniform\_\allowbreak{}gap} & {\scriptsize 76} & {\scriptsize ---} \\
{\scriptsize\ttfamily factored\_\allowbreak{}variational\_\allowbreak{}energy\_\allowbreak{}without\_\allowbreak{}inverse} & {\scriptsize 114} & {\scriptsize ---} \\
{\scriptsize\ttfamily full\_\allowbreak{}graph\_\allowbreak{}source\_\allowbreak{}ward} & {\scriptsize 213} & {\scriptsize ---} \\
{\scriptsize\ttfamily gapless\_\allowbreak{}factored\_\allowbreak{}w6\_\allowbreak{}bound} & {\scriptsize 140} & {\scriptsize ---} \\
{\scriptsize\ttfamily joint\_\allowbreak{}resolvent\_\allowbreak{}budget} & {\scriptsize 240} & {\scriptsize ---} \\
{\scriptsize\ttfamily residual\_\allowbreak{}controls\_\allowbreak{}inverse\_\allowbreak{}without\_\allowbreak{}diagonal} & {\scriptsize 250} & {\scriptsize ---} \\
{\scriptsize\ttfamily riccati\_\allowbreak{}default\_\allowbreak{}contraction\_\allowbreak{}bound} & {\scriptsize 47} & {\scriptsize ---} \\
{\scriptsize\ttfamily riccati\_\allowbreak{}residual\_\allowbreak{}certificate} & {\scriptsize 33} & {\scriptsize ---} \\
{\scriptsize\ttfamily rotated\_\allowbreak{}two\_\allowbreak{}level\_\allowbreak{}schur\_\allowbreak{}numerator} & {\scriptsize 195} & {\scriptsize ---} \\
{\scriptsize\ttfamily schur\_\allowbreak{}connection\_\allowbreak{}cancellation} & {\scriptsize 179} & {\scriptsize ---} \\
{\scriptsize\ttfamily schur\_\allowbreak{}metric\_\allowbreak{}pairing\_\allowbreak{}on\_\allowbreak{}shell} & {\scriptsize 187} & {\scriptsize ---} \\
\end{longtable}

\subsection*{ThermodynamicLimit.lean}
\addcontentsline{toc}{subsection}{ThermodynamicLimit.lean}
\begin{longtable}{@{}p{0.32\linewidth} l p{0.44\linewidth}@{}}
\toprule Theorem & Line & Formalizes / promotes \\ \midrule
\endfirsthead
\toprule Theorem & Line & Formalizes / promotes \\ \midrule \endhead
\bottomrule\endfoot
{\scriptsize\ttfamily closable\_\allowbreak{}of\_\allowbreak{}integration\_\allowbreak{}by\_\allowbreak{}parts} & {\scriptsize 55} & {\scriptsize ---} \\
{\scriptsize\ttfamily closed\_\allowbreak{}extension\_\allowbreak{}of\_\allowbreak{}integration\_\allowbreak{}by\_\allowbreak{}parts} & {\scriptsize 92} & {\scriptsize ---} \\
{\scriptsize\ttfamily exists\_\allowbreak{}limit\_\allowbreak{}of\_\allowbreak{}cut\_\allowbreak{}response} & {\scriptsize 22} & {\scriptsize ---} \\
{\scriptsize\ttfamily gap\_\allowbreak{}bound\_\allowbreak{}on\_\allowbreak{}graph\_\allowbreak{}limit} & {\scriptsize 75} & {\scriptsize ---} \\
{\scriptsize\ttfamily integral\_\allowbreak{}under\_\allowbreak{}weak\_\allowbreak{}and\_\allowbreak{}uniform\_\allowbreak{}convergence} & {\scriptsize 111} & {\scriptsize ---} \\
{\scriptsize\ttfamily integration\_\allowbreak{}by\_\allowbreak{}parts\_\allowbreak{}under\_\allowbreak{}weak\_\allowbreak{}limit} & {\scriptsize 135} & {\scriptsize ---} \\
{\scriptsize\ttfamily limit\_\allowbreak{}preserves\_\allowbreak{}gap\_\allowbreak{}bound} & {\scriptsize 30} & {\scriptsize ---} \\
{\scriptsize\ttfamily poincare\_\allowbreak{}under\_\allowbreak{}weak\_\allowbreak{}convergence} & {\scriptsize 157} & {\scriptsize ---} \\
{\scriptsize\ttfamily sc17\_\allowbreak{}plaquette\_\allowbreak{}interval} & {\scriptsize 214} & {\scriptsize ---} \\
{\scriptsize\ttfamily spectral\_\allowbreak{}interval\_\allowbreak{}mass} & {\scriptsize 182} & {\scriptsize ---} \\
\end{longtable}

\subsection*{VacuumChart.lean}
\addcontentsline{toc}{subsection}{VacuumChart.lean}
\begin{longtable}{@{}p{0.32\linewidth} l p{0.44\linewidth}@{}}
\toprule Theorem & Line & Formalizes / promotes \\ \midrule
\endfirsthead
\toprule Theorem & Line & Formalizes / promotes \\ \midrule \endhead
\bottomrule\endfoot
{\scriptsize\ttfamily corrected\_\allowbreak{}hermitian} & {\scriptsize 72} & {\scriptsize G18; RUN:wilson\_\allowbreak{}vacuum\_\allowbreak{}chart\_\allowbreak{}2026-09-05} \\
{\scriptsize\ttfamily corrected\_\allowbreak{}vacuum\_\allowbreak{}column} & {\scriptsize 63} & {\scriptsize G18; RUN:wilson\_\allowbreak{}vacuum\_\allowbreak{}chart\_\allowbreak{}2026-09-05} \\
{\scriptsize\ttfamily corrected\_\allowbreak{}vacuum\_\allowbreak{}legs} & {\scriptsize 80} & {\scriptsize G18; RUN:wilson\_\allowbreak{}vacuum\_\allowbreak{}chart\_\allowbreak{}2026-09-05} \\
{\scriptsize\ttfamily endpoint\_\allowbreak{}disjoint\_\allowbreak{}cancellation} & {\scriptsize 99} & {\scriptsize G18; G19; RUN:wilson\_\allowbreak{}vacuum\_\allowbreak{}chart\_\allowbreak{}2026-09-05} \\
{\scriptsize\ttfamily fixed\_\allowbreak{}vacuum\_\allowbreak{}row} & {\scriptsize 50} & {\scriptsize G18; RUN:wilson\_\allowbreak{}vacuum\_\allowbreak{}chart\_\allowbreak{}2026-09-05} \\
{\scriptsize\ttfamily generator\_\allowbreak{}antihermitian} & {\scriptsize 39} & {\scriptsize G18; RUN:wilson\_\allowbreak{}vacuum\_\allowbreak{}chart\_\allowbreak{}2026-09-05} \\
{\scriptsize\ttfamily generator\_\allowbreak{}vacuum} & {\scriptsize 44} & {\scriptsize G18; RUN:wilson\_\allowbreak{}vacuum\_\allowbreak{}chart\_\allowbreak{}2026-09-05} \\
{\scriptsize\ttfamily rankOne\_\allowbreak{}adjoint} & {\scriptsize 35} & {\scriptsize G18; RUN:wilson\_\allowbreak{}vacuum\_\allowbreak{}chart\_\allowbreak{}2026-09-05} \\
{\scriptsize\ttfamily scalar\_\allowbreak{}vacuum\_\allowbreak{}cancellation} & {\scriptsize 56} & {\scriptsize G18; RUN:wilson\_\allowbreak{}vacuum\_\allowbreak{}chart\_\allowbreak{}2026-09-05} \\
\end{longtable}

\subsection*{VacuumCompression.lean}
\addcontentsline{toc}{subsection}{VacuumCompression.lean}
\begin{longtable}{@{}p{0.32\linewidth} l p{0.44\linewidth}@{}}
\toprule Theorem & Line & Formalizes / promotes \\ \midrule
\endfirsthead
\toprule Theorem & Line & Formalizes / promotes \\ \midrule \endhead
\bottomrule\endfoot
{\scriptsize\ttfamily corrected\_\allowbreak{}eq\_\allowbreak{}compression} & {\scriptsize 55} & {\scriptsize G18; RUN:wilson\_\allowbreak{}vacuum\_\allowbreak{}compression\_\allowbreak{}2026-09-05} \\
{\scriptsize\ttfamily corrected\_\allowbreak{}eq\_\allowbreak{}sub\_\allowbreak{}vacuum\_\allowbreak{}terms} & {\scriptsize 40} & {\scriptsize G18; RUN:wilson\_\allowbreak{}vacuum\_\allowbreak{}compression\_\allowbreak{}2026-09-05} \\
{\scriptsize\ttfamily vacuum\_\allowbreak{}corner\_\allowbreak{}zero} & {\scriptsize 31} & {\scriptsize G18; RUN:wilson\_\allowbreak{}vacuum\_\allowbreak{}compression\_\allowbreak{}2026-09-05} \\
{\scriptsize\ttfamily vacuum\_\allowbreak{}projection} & {\scriptsize 24} & {\scriptsize G18; RUN:wilson\_\allowbreak{}vacuum\_\allowbreak{}compression\_\allowbreak{}2026-09-05} \\
\end{longtable}

\subsection*{W6Residual.lean}
\addcontentsline{toc}{subsection}{W6Residual.lean}
\begin{longtable}{@{}p{0.32\linewidth} l p{0.44\linewidth}@{}}
\toprule Theorem & Line & Formalizes / promotes \\ \midrule
\endfirsthead
\toprule Theorem & Line & Formalizes / promotes \\ \midrule \endhead
\bottomrule\endfoot
{\scriptsize\ttfamily factored\_\allowbreak{}residual\_\allowbreak{}dual\_\allowbreak{}bound} & {\scriptsize 205} & {\scriptsize ---} \\
{\scriptsize\ttfamily factored\_\allowbreak{}w6\_\allowbreak{}bound} & {\scriptsize 243} & {\scriptsize ---} \\
{\scriptsize\ttfamily functional\_\allowbreak{}le\_\allowbreak{}energyDualNorm} & {\scriptsize 41} & {\scriptsize ---} \\
{\scriptsize\ttfamily gaussian\_\allowbreak{}energy\_\allowbreak{}coefficient\_\allowbreak{}one} & {\scriptsize 273} & {\scriptsize ---} \\
{\scriptsize\ttfamily gaussian\_\allowbreak{}energy\_\allowbreak{}coefficient\_\allowbreak{}three} & {\scriptsize 289} & {\scriptsize ---} \\
{\scriptsize\ttfamily gaussian\_\allowbreak{}energy\_\allowbreak{}coefficient\_\allowbreak{}two} & {\scriptsize 280} & {\scriptsize ---} \\
{\scriptsize\ttfamily inverse\_\allowbreak{}energy\_\allowbreak{}eq\_\allowbreak{}dual\_\allowbreak{}norm\_\allowbreak{}sq} & {\scriptsize 113} & {\scriptsize ---} \\
{\scriptsize\ttfamily inverse\_\allowbreak{}energy\_\allowbreak{}le\_\allowbreak{}dual} & {\scriptsize 81} & {\scriptsize ---} \\
{\scriptsize\ttfamily residual\_\allowbreak{}square\_\allowbreak{}completion} & {\scriptsize 48} & {\scriptsize ---} \\
{\scriptsize\ttfamily variational\_\allowbreak{}residual\_\allowbreak{}dual\_\allowbreak{}identity} & {\scriptsize 132} & {\scriptsize ---} \\
{\scriptsize\ttfamily variational\_\allowbreak{}residual\_\allowbreak{}error\_\allowbreak{}bound} & {\scriptsize 147} & {\scriptsize ---} \\
{\scriptsize\ttfamily variational\_\allowbreak{}residual\_\allowbreak{}identity} & {\scriptsize 60} & {\scriptsize ---} \\
\end{longtable}

\subsection*{WilsonGridAlgebra.lean}
\addcontentsline{toc}{subsection}{WilsonGridAlgebra.lean}
\begin{longtable}{@{}p{0.32\linewidth} l p{0.44\linewidth}@{}}
\toprule Theorem & Line & Formalizes / promotes \\ \midrule
\endfirsthead
\toprule Theorem & Line & Formalizes / promotes \\ \midrule \endhead
\bottomrule\endfoot
{\scriptsize\ttfamily wilson\_\allowbreak{}grid\_\allowbreak{}block\_\allowbreak{}energy\_\allowbreak{}identity} & {\scriptsize 20} & {\scriptsize RESULT:\allowbreak{}W6\_\allowbreak{}SHARED\_\allowbreak{}EDGE\_\allowbreak{}GAUSSIAN\_\allowbreak{}FLOOR} \\
{\scriptsize\ttfamily wilson\_\allowbreak{}grid\_\allowbreak{}block\_\allowbreak{}energy\_\allowbreak{}lower\_\allowbreak{}bound} & {\scriptsize 27} & {\scriptsize RESULT:\allowbreak{}W6\_\allowbreak{}SHARED\_\allowbreak{}EDGE\_\allowbreak{}GAUSSIAN\_\allowbreak{}FLOOR} \\
{\scriptsize\ttfamily wilson\_\allowbreak{}grid\_\allowbreak{}block\_\allowbreak{}energy\_\allowbreak{}upper\_\allowbreak{}bound} & {\scriptsize 33} & {\scriptsize RESULT:\allowbreak{}W6\_\allowbreak{}SHARED\_\allowbreak{}EDGE\_\allowbreak{}GAUSSIAN\_\allowbreak{}FLOOR} \\
{\scriptsize\ttfamily wilson\_\allowbreak{}grid\_\allowbreak{}block\_\allowbreak{}energy\_\allowbreak{}zero\_\allowbreak{}iff} & {\scriptsize 39} & {\scriptsize RESULT:\allowbreak{}W6\_\allowbreak{}SHARED\_\allowbreak{}EDGE\_\allowbreak{}GAUSSIAN\_\allowbreak{}FLOOR} \\
\end{longtable}

\subsection*{WilsonSquareForce.lean}
\addcontentsline{toc}{subsection}{WilsonSquareForce.lean}
\begin{longtable}{@{}p{0.32\linewidth} l p{0.44\linewidth}@{}}
\toprule Theorem & Line & Formalizes / promotes \\ \midrule
\endfirsthead
\toprule Theorem & Line & Formalizes / promotes \\ \midrule \endhead
\bottomrule\endfoot
{\scriptsize\ttfamily wilson\_\allowbreak{}square\_\allowbreak{}finite\_\allowbreak{}radial\_\allowbreak{}synthesis\_\allowbreak{}bound} & {\scriptsize 91} & {\scriptsize RESULT:\allowbreak{}W6\_\allowbreak{}SQUARE\_\allowbreak{}SHARP\_\allowbreak{}INVERSE\_\allowbreak{}ENERGY} \\
{\scriptsize\ttfamily wilson\_\allowbreak{}square\_\allowbreak{}first\_\allowbreak{}source\_\allowbreak{}saturates} & {\scriptsize 42} & {\scriptsize RESULT:\allowbreak{}W6\_\allowbreak{}SQUARE\_\allowbreak{}SHARP\_\allowbreak{}INVERSE\_\allowbreak{}ENERGY} \\
{\scriptsize\ttfamily wilson\_\allowbreak{}square\_\allowbreak{}radial\_\allowbreak{}ratio\_\allowbreak{}bound} & {\scriptsize 63} & {\scriptsize RESULT:\allowbreak{}W6\_\allowbreak{}SQUARE\_\allowbreak{}SHARP\_\allowbreak{}INVERSE\_\allowbreak{}ENERGY} \\
{\scriptsize\ttfamily wilson\_\allowbreak{}square\_\allowbreak{}radial\_\allowbreak{}ratio\_\allowbreak{}difference} & {\scriptsize 50} & {\scriptsize RESULT:\allowbreak{}W6\_\allowbreak{}SQUARE\_\allowbreak{}SHARP\_\allowbreak{}INVERSE\_\allowbreak{}ENERGY} \\
{\scriptsize\ttfamily wilson\_\allowbreak{}square\_\allowbreak{}radial\_\allowbreak{}ratio\_\allowbreak{}strict} & {\scriptsize 76} & {\scriptsize RESULT:\allowbreak{}W6\_\allowbreak{}SQUARE\_\allowbreak{}SHARP\_\allowbreak{}INVERSE\_\allowbreak{}ENERGY} \\
{\scriptsize\ttfamily wilson\_\allowbreak{}square\_\allowbreak{}sharp\_\allowbreak{}constant} & {\scriptsize 25} & {\scriptsize RESULT:\allowbreak{}W6\_\allowbreak{}SQUARE\_\allowbreak{}SHARP\_\allowbreak{}INVERSE\_\allowbreak{}ENERGY} \\
{\scriptsize\ttfamily wilson\_\allowbreak{}square\_\allowbreak{}sharp\_\allowbreak{}constant\_\allowbreak{}pos} & {\scriptsize 34} & {\scriptsize RESULT:\allowbreak{}W6\_\allowbreak{}SQUARE\_\allowbreak{}SHARP\_\allowbreak{}INVERSE\_\allowbreak{}ENERGY} \\
\end{longtable}

% located 437 of 439; formalizes 183
% ---- end gen_lean.tex ----------------------------------------------
% ---- end B1_lean.tex -----------------------------------------------
% ---- begin B2_library.tex ------------------------------------------
\section{The reading library}
\label{app:library}

\LibraryPapers{} acquired papers with local full text
(\LibraryVolOne{} in the first volume, \LibraryVolTwo{} in the second),
grouped by the library's own categories and ordered by year within
each.

This is a \emph{coverage} instrument, not an evidence map. Its purpose
is to answer ``has this been done?'' before an agent spends a week, and
its \emph{Connection} column is a curated reading note recording why
the paper was acquired. A connection is not an evidence edge and
carries no weight in any claim.

The evidence edges --- the \LitVerified{} relations actually verified
against a named claim of this program --- are the separate and much
smaller set in \S\ref{sec:external}. The distance between
\LibraryPapers{} acquired and \LitObtained{} read-against-a-claim is
the honest measure of how far the external comparison has actually
been carried.

% ---- begin gen_library.tex -----------------------------------------
% Generated by paper/master/generate_sources.py.  Do not edit.
\noindent 260 acquired papers, all with local full text, grouped by the library's own categories. \emph{Connection} is the curated one-line statement of what the paper bears on in this program; it is a reading note, not an evidence edge. The evidence edges are the separate, smaller, verified set in \S\ref{sec:external}. A paper here has been \emph{obtained}, which is not the same as \emph{read against a claim}.\par\medskip
\subsection*{Hamiltonian and quantum-simulation formulations}
\addcontentsline{toc}{subsection}{Hamiltonian and quantum-simulation formulations}
\begin{longtable}{@{}l p{0.40\linewidth} l p{0.28\linewidth}@{}}
\toprule Year & Title & arXiv & Connection \\ \midrule
\endfirsthead
\toprule Year & Title & arXiv & Connection \\ \midrule \endhead
\bottomrule\endfoot
2026 & {\scriptsize Continuum limit of a qubit-regularized SU(3) lattice gauge theory with glueballs} & {\scriptsize \texttt{2603.01215}} & {\scriptsize 2026 qubit-regularized SU(3) continuum-limit/glueball preprint; direct but too new for citation stability} \\
2026 & {\scriptsize Minimally Truncated SU(3) Lattice Gauge Theory and String Tension} & {\scriptsize \texttt{2601.10065}} & {\scriptsize 2026 minimal SU(3) truncation with direct string-tension target; promising but not yet peer reviewed} \\
2025 & {\scriptsize Accurate ground states of \$SU(2)\$ lattice gauge theory in 2+1D and 3+1D} & {\scriptsize \texttt{2509.12323}} & {\scriptsize variational SU(2) ground states in 2+1D and 3+1D} \\
2025 & {\scriptsize Asymptotic-freedom and massive glueballs in a qubit-regularized SU(2) gauge theory} & {\scriptsize \texttt{2512.11068}} & {\scriptsize peer-reviewed qubit-regularized SU(2) glueball toy model} \\
2025 & {\scriptsize Dynamics in Hamiltonian Lattice Gauge Theory: Approaching the Continuum Limit with Partitionings of SU\$(2)\$} & {\scriptsize \texttt{2503.03397}} & {\scriptsize Recent peer-reviewed dynamical SU(2) approach toward the continuum limit} \\
2025 & {\scriptsize Efficient Truncations of SU(\$N\_\allowbreak{}c\$) Lattice Gauge Theory for Quantum Simulation} & {\scriptsize \texttt{2503.11888}} & {\scriptsize controlled efficient gauge-field truncations for SU(Nc) Hamiltonians} \\
2025 & {\scriptsize Emergent Hydrodynamic Mode on SU(2) Plaquette Chains and Quantum Simulation} & {\scriptsize \texttt{2502.17551}} & {\scriptsize emergent hydrodynamic mode in SU(2) plaquette-chain Hamiltonian dynamics} \\
2025 & {\scriptsize Exponential speedup in quantum simulation of Kogut-Susskind Hamiltonian via orbifold lattice} & {\scriptsize \texttt{2506.00755}} & {\scriptsize orbifold-lattice route to faster Kogut-Susskind Hamiltonian simulation} \\
2025 & {\scriptsize Generic Hilbert Space Fragmentation in Kogut--Susskind Lattice Gauge Theories} & {\scriptsize \texttt{2502.03533}} & {\scriptsize Hilbert-space fragmentation and structural sectors in Kogut-Susskind Hamiltonians} \\
2025 & {\scriptsize Improved Honeycomb and Hyperhoneycomb Lattice Hamiltonians for Quantum Simulations of Non-Abelian Gauge Theories} & {\scriptsize \texttt{2503.09688}} & {\scriptsize improved non-Abelian Hamiltonian lattices} \\
2025 & {\scriptsize Loop-string-hadron approach to SU(3) lattice Yang-Mills theory, II: Operator representation for the trivalent vertex} & {\scriptsize \texttt{2512.11796}} & {\scriptsize SU(3) loop-string-hadron operator representation on a trivalent vertex} \\
2025 & {\scriptsize Non-Abelian dynamics on a cube: improving quantum compilation through qudit-based simulations} & {\scriptsize \texttt{2506.10945}} & {\scriptsize non-Abelian cube dynamics and qudit-based quantum compilation} \\
2025 & {\scriptsize Obtaining continuum physics from dynamical simulations of Hamiltonian lattice gauge theories} & {\scriptsize \texttt{2506.16559}} & {\scriptsize 2025 preprint on recovering continuum physics from Hamiltonian dynamics; early citation traction} \\
2025 & {\scriptsize Perturbation theory, irrep truncations, and state preparation methods for quantum simulations of SU(3) lattice gauge theory} & {\scriptsize \texttt{2509.25865}} & {\scriptsize SU(3) strong-coupling perturbation theory, irrep truncation, and state preparation} \\
2025 & {\scriptsize Quantum Circuits for SU(3) Lattice Gauge Theory} & {\scriptsize \texttt{2503.08866}} & {\scriptsize electric-basis circuits for pure SU(3) Hamiltonian dynamics in two and three dimensions} \\
2024 & {\scriptsize A Fully Gauge-Fixed SU(2) Hamiltonian for Quantum Simulations} & {\scriptsize \texttt{2409.10610}} & {\scriptsize fully gauge-fixed SU(2) Hamiltonian for quantum simulation} \\
2024 & {\scriptsize An efficient finite-resource formulation of non-Abelian lattice gauge theories beyond one dimension} & {\scriptsize \texttt{2409.04441}} & {\scriptsize loop-variable finite-resource non-Abelian lattice gauge theory beyond one dimension} \\
2024 & {\scriptsize Floquet evolution of the q-deformed \textbackslash{}texorpdfstring\{SU(3)\$\{\}\_\allowbreak{}1\$\}\{SU(3)1\} Yang-Mills theory on a two-leg ladder} & {\scriptsize \texttt{2409.20263}} & {\scriptsize Floquet dynamics of q-deformed SU(3) Yang-Mills on a ladder} \\
2024 & {\scriptsize From square plaquettes to triamond lattices for SU(2) gauge theory} & {\scriptsize \texttt{2401.14570}} & {\scriptsize alternative lattice geometries for SU(2) Hamiltonian gauge theory} \\
2024 & {\scriptsize Gauge Loop-String-Hadron Formulation on General Graphs and Applications to Fully Gauge Fixed Hamiltonian Lattice Gauge Theory} & {\scriptsize \texttt{2409.13812}} & {\scriptsize Modern general-graph loop-string-hadron construction with full gauge fixing} \\
2024 & {\scriptsize Loop-string-hadron approach to SU(3) lattice Yang-Mills theory: I. Hilbert space of a trivalent vertex} & {\scriptsize \texttt{2407.19181}} & {\scriptsize SU(3) gauge-invariant loop-string-hadron Hilbert space} \\
2024 & {\scriptsize Protecting gauge symmetries in the the dynamics of SU(3) lattice gauge theories} & {\scriptsize \texttt{2404.12158}} & {\scriptsize gauge-symmetry protection in real-time SU(3) Hamiltonian evolution} \\
2024 & {\scriptsize Quantum Simulation of SU(3) Lattice Yang Mills Theory at Leading Order in Large N} & {\scriptsize \texttt{2402.10265}} & {\scriptsize SU(3) Hamiltonian quantum simulation at large N} \\
2024 & {\scriptsize Tensor Networks for Lattice Gauge Theories beyond one dimension: a Roadmap} & {\scriptsize \texttt{2407.03058}} & {\scriptsize tensor-network roadmap for lattice gauge theories beyond one spatial dimension} \\
2023 & {\scriptsize \$q\$ deformed formulation of Hamiltonian SU(3) Yang-Mills theory} & {\scriptsize \texttt{2306.12324}} & {\scriptsize q-deformed Hamiltonian formulation of SU(3) Yang-Mills theory} \\
2023 & {\scriptsize A new basis for Hamiltonian SU(2) simulations} & {\scriptsize \texttt{2307.11829}} & {\scriptsize Fast-cited modern basis for Hamiltonian SU(2), designed for controlled continuum extrapolation} \\
2023 & {\scriptsize Hamiltonians and gauge-invariant Hilbert space for lattice Yang-Mills-like theories with finite gauge group} & {\scriptsize \texttt{2301.12224}} & {\scriptsize explicit gauge-invariant Hilbert-space construction for finite-group Yang-Mills analogues} \\
2023 & {\scriptsize Quantum Simulation of Lattice QCD with Improved Hamiltonians} & {\scriptsize \texttt{2307.05593}} & {\scriptsize improved Hamiltonians for lattice-QCD quantum simulation} \\
2023 & {\scriptsize Quantum and classical spin network algorithms for \$q\$-deformed Kogut-Susskind gauge theories} & {\scriptsize \texttt{2304.02527}} & {\scriptsize q-deformed Kogut-Susskind Hamiltonians and finite-resource spin-network algorithms} \\
2023 & {\scriptsize Quantum simulation of lattice gauge theories via deterministic duality transformations assisted by measurements} & {\scriptsize \texttt{2305.12277}} & {\scriptsize measurement-assisted deterministic duality transformations for gauge-invariant simulation} \\
2023 & {\scriptsize SU(2) Gauge Theory in \$2+1\$ Dimensions on a Plaquette Chain Obeys the Eigenstate Thermalization Hypothesis} & {\scriptsize \texttt{2303.14264}} & {\scriptsize eigenstate thermalization and Wilson-loop matrix elements in a 2+1D SU(2) plaquette-chain Hamiltonian} \\
2023 & {\scriptsize Simple Hamiltonian for Quantum Simulation of Strongly Coupled 2+1D SU(2) Lattice Gauge Theory on a Honeycomb Lattice} & {\scriptsize \texttt{2307.00045}} & {\scriptsize strong-coupling honeycomb Hamiltonian for 2+1D SU(2) Yang-Mills} \\
2023 & {\scriptsize Spin exchange-enabled quantum simulator for large-scale non-Abelian gauge theories} & {\scriptsize \texttt{2305.06373}} & {\scriptsize spin-exchange quantum simulator for scalable non-Abelian gauge dynamics} \\
2023 & {\scriptsize Strategies for simulating time evolution of Hamiltonian lattice field theories} & {\scriptsize \texttt{2312.11637}} & {\scriptsize time-evolution strategies and error analysis for Hamiltonian lattice field theory} \\
2023 & {\scriptsize String-net formulation of Hamiltonian lattice Yang-Mills theories and quantum many-body scars in a nonabelian gauge theory} & {\scriptsize \texttt{2305.05950}} & {\scriptsize Hamiltonian lattice Yang-Mills string-net structure and non-Abelian scars} \\
2022 & {\scriptsize Gauge Theory Couplings on Anisotropic Lattices} & {\scriptsize \texttt{2208.10417}} & {\scriptsize anisotropic SU(N) gauge couplings matched toward the Hamiltonian limit} \\
2022 & {\scriptsize General quantum algorithms for Hamiltonian simulation with applications to a non-Abelian lattice gauge theory} & {\scriptsize \texttt{2212.14030}} & {\scriptsize general non-Abelian Hamiltonian algorithms using gauge-invariant variables} \\
2022 & {\scriptsize Loop-string-hadron formulation of an SU(3) gauge theory with dynamical quarks} & {\scriptsize \texttt{2212.04490}} & {\scriptsize modern SU(3) gauge-invariant loop-string-hadron formulation} \\
2022 & {\scriptsize Quantum Simulation for High Energy Physics} & {\scriptsize \texttt{2204.03381}} & {\scriptsize peer-reviewed high-energy quantum-simulation review} \\
2022 & {\scriptsize Self-mitigating Trotter circuits for SU(2) lattice gauge theory on a quantum computer} & {\scriptsize \texttt{2205.09247}} & {\scriptsize self-mitigating real-time circuits for SU(2) Hamiltonian lattice gauge theory} \\
2021 & {\scriptsize Exact duality and local dynamics in SU(N) lattice gauge theory} & {\scriptsize \texttt{2109.00992}} & {\scriptsize Exact duality and local dynamics for SU(N) Hamiltonian lattice gauge theory} \\
2021 & {\scriptsize Preparation of the SU(3) Lattice Yang-Mills Vacuum with Variational Quantum Methods} & {\scriptsize \texttt{2112.09083}} & {\scriptsize pure SU(3) Hamiltonian vacuum preparation on plaquette chains} \\
2019 & {\scriptsize Loop, String, and Hadron Dynamics in SU(2) Hamiltonian Lattice Gauge Theories} & {\scriptsize \texttt{1912.06133}} & {\scriptsize Highly cited loop-string-hadron formulation and dynamics for SU(2) Hamiltonian LGT} \\
2013 & {\scriptsize Ultracold Quantum Gases and Lattice Systems: Quantum Simulation of Lattice Gauge Theories} & {\scriptsize \texttt{1305.1602}} & {\scriptsize foundational quantum-simulation review for gauge theories} \\
2003 & {\scriptsize SU(N) Glueball Masses in 2+1 Dimensions} & {\scriptsize \texttt{hep-lat/0303016}} & {\scriptsize Hamiltonian/analytic SU(N) glueball masses in 2+1D, a direct alternate spectral calculation} \\
2003 & {\scriptsize SU(N) Lattice Gauge Theory on a Single Cube} & {\scriptsize \texttt{hep-lat/0303022}} & {\scriptsize gauge-invariant SU(N) Hamiltonian theory on a cube} \\
2003 & {\scriptsize The Large N Glueball Mass Spectrum in 2+1 Dimensions} & {\scriptsize \texttt{hep-lat/0303018}} & {\scriptsize Large-N continuation of the Hamiltonian 2+1D glueball calculation} \\
2000 & {\scriptsize Green's Function Monte Carlo study of SU(3) lattice gauge theory in (3+1)D} & {\scriptsize \texttt{hep-lat/0005009}} & {\scriptsize Green-function Monte Carlo SU(3) Hamiltonian study} \\
2000 & {\scriptsize Towards a Many-Body Treatment of Hamiltonian Lattice SU(N) Gauge Theory} & {\scriptsize \texttt{hep-lat/0001028}} & {\scriptsize Many-body/coupled-cluster treatment specifically formulated for Hamiltonian SU(N)} \\
1999 & {\scriptsize The basis of the physical Hilbert space of lattice gauge theories} & {\scriptsize \texttt{hep-lat/9906036}} & {\scriptsize Explicit spin-network basis of the gauge-invariant lattice Hilbert space} \\
1996 & {\scriptsize The Coupled Cluster Method in Hamiltonian Lattice Field Theory} & {\scriptsize \texttt{hep-lat/9603026}} & {\scriptsize coupled-cluster Hamiltonian lattice gauge theory} \\
1992 & {\scriptsize The Shifted Coupled Cluster Method: A New Approach to Hamiltonian Lattice Gauge Theories} & {\scriptsize \texttt{hep-lat/9212025}} & {\scriptsize Gauge-invariant shifted coupled-cluster construction from linked strong-coupling clusters} \\
\end{longtable}

\subsection*{Glueball spectrum and continuum limits}
\addcontentsline{toc}{subsection}{Glueball spectrum and continuum limits}
\begin{longtable}{@{}l p{0.40\linewidth} l p{0.28\linewidth}@{}}
\toprule Year & Title & arXiv & Connection \\ \midrule
\endfirsthead
\toprule Year & Title & arXiv & Connection \\ \midrule \endhead
\bottomrule\endfoot
2026 & {\scriptsize Learning the generating functional for variance reduction in lattice QCD} & {\scriptsize \texttt{2606.15986}} & {\scriptsize learned generating functional for glueball-correlator variance reduction} \\
2026 & {\scriptsize Update on the computation of the quenched \$SU(6)\$ Yang-Mills lattice spectrum} & {\scriptsize \texttt{2603.13138}} & {\scriptsize current quenched SU(6) Yang-Mills spectrum update} \\
2026 & {\scriptsize Variance reduction in lattice QCD observables via normalizing flows} & {\scriptsize \texttt{2603.02984}} & {\scriptsize normalizing-flow variance reduction demonstrated on SU(3) glueball correlators} \\
2025 & {\scriptsize Casimir scaling in glueballs in SU(\$N\$) and Sp(\$2N\$) gauge theories: hints from constituent approaches} & {\scriptsize \texttt{2509.09454}} & {\scriptsize very recent Casimir-scaling test for glueballs} \\
2025 & {\scriptsize Glueball mass spectrum at finite temperature revisited: Constant contribution in glueball correlators in the deconfinement phase} & {\scriptsize \texttt{2504.09120}} & {\scriptsize finite-temperature SU(3) glueball correlators on anisotropic lattices} \\
2025 & {\scriptsize Lattice evidence that scalar glueballs are small} & {\scriptsize \texttt{2508.21821}} & {\scriptsize scalar-glueball form factors and radius} \\
2025 & {\scriptsize Test of a two-level algorithm for the glueball spectrum in \$SU(N\_\allowbreak{}c)\$ Yang-Mills theory} & {\scriptsize \texttt{2501.19312}} & {\scriptsize two-level glueball-spectrum test for SU(Nc)} \\
2025 & {\scriptsize Update on Glueballs} & {\scriptsize \texttt{2502.02547}} & {\scriptsize modern peer-reviewed-proceedings review of glueball spectroscopy} \\
2024 & {\scriptsize Exponential Error Reduction for Glueball Calculations Using a Two-Level Algorithm in Pure Gauge Theory} & {\scriptsize \texttt{2406.12656}} & {\scriptsize two-level exponential error reduction with a glueball variational basis} \\
2024 & {\scriptsize Matching Wilson flow time for calculating the topological charge density and the pseudoscalar glueball mass in quenched lattice QCD} & {\scriptsize \texttt{2409.02756}} & {\scriptsize Wilson-flow matching for pseudoscalar glueball mass} \\
2024 & {\scriptsize The \$$\theta$\$-dependence of the Yang-Mills spectrum from analytic continuation} & {\scriptsize \texttt{2402.03096}} & {\scriptsize theta dependence of Yang-Mills spectrum} \\
2023 & {\scriptsize \$Sp(2N)\$ Lattice Gauge Theories and Extensions of the Standard Model of Particle Physics} & {\scriptsize \texttt{2304.01070}} & {\scriptsize review of Sp(2N) lattice spectra and large-N extrapolation} \\
2023 & {\scriptsize Boundary states and Non-Abelian Casimir effect in lattice Yang-Mills theory} & {\scriptsize \texttt{2302.00376}} & {\scriptsize 3+1D SU(3) boundary-state and chromometallic Casimir spectrum} \\
2023 & {\scriptsize Glueballs in \$N\_\allowbreak{}f=1\$ QCD} & {\scriptsize \texttt{2312.00470}} & {\scriptsize Very recent Nf=1 glueball spectrum; citation count is still age-limited} \\
2023 & {\scriptsize Performance of two-level sampling for the glueball spectrum in pure gauge theory} & {\scriptsize \texttt{2312.11372}} & {\scriptsize two-level sampling for glueball spectrum} \\
2023 & {\scriptsize The quenched glueball spectrum from smeared spectral densities} & {\scriptsize \texttt{2311.14806}} & {\scriptsize smeared spectral densities for glueballs} \\
2022 & {\scriptsize Confining strings, axions and glueballs in the planar limit} & {\scriptsize \texttt{2205.03642}} & {\scriptsize planar-limit confining strings and glueballs} \\
2022 & {\scriptsize Glueball spectroscopy in lattice QCD using gradient flow} & {\scriptsize \texttt{2211.15176}} & {\scriptsize gradient-flow glueball spectroscopy} \\
2022 & {\scriptsize Spectrum of Large N Glueballs: Holography vs Lattice} & {\scriptsize \texttt{2206.14826}} & {\scriptsize large-N lattice glueball spectrum comparison and synthesis} \\
2022 & {\scriptsize The glueball spectrum with \$N\_\allowbreak{}f=4\$ light fermions} & {\scriptsize \texttt{2212.02080}} & {\scriptsize Recent Nf=4 spectrum for understanding fermion effects; currently lightly cited} \\
2022 & {\scriptsize Towards glueball masses of large-\$N\$ \$\textbackslash{}mathrm\{SU\}(N)\$ pure-gauge theories without topological freezing} & {\scriptsize \texttt{2205.06190}} & {\scriptsize large-N glueball masses without topological freezing} \\
2021 & {\scriptsize SU(N) gauge theories in 3+1 dimensions: glueball spectrum, string tensions and topology} & {\scriptsize \texttt{2106.00364}} & {\scriptsize modern continuum and large-N SU(N) spectrum} \\
2020 & {\scriptsize Color dependence of tensor and scalar glueball masses in Yang-Mills theories} & {\scriptsize \texttt{2004.11063}} & {\scriptsize color dependence of scalar and tensor glueballs} \\
2020 & {\scriptsize Glueballs and Strings in \$Sp(2N)\$ Yang-Mills theories} & {\scriptsize \texttt{2010.15781}} & {\scriptsize comparative Sp(2N) continuum glueball and string spectrum} \\
2020 & {\scriptsize Prospects for large N gauge theories on the lattice} & {\scriptsize \texttt{2001.10859}} & {\scriptsize Modern survey of twisted reduction and large-N lattice prospects} \\
2020 & {\scriptsize Spectrum of Trace Deformed Yang-Mills Theories} & {\scriptsize \texttt{2010.03618}} & {\scriptsize Peer-reviewed spectrum of trace-deformed Yang-Mills, useful as a controlled deformation comparison} \\
2020 & {\scriptsize The glueball spectrum of SU(3) gauge theory in 3+1 dimension} & {\scriptsize \texttt{2007.06422}} & {\scriptsize modern continuum SU(3) spectrum in 3+1 dimensions} \\
2018 & {\scriptsize The spectrum of 2+1 dimensional Yang-Mills theory on a twisted spatial torus} & {\scriptsize \texttt{1807.03481}} & {\scriptsize Spectral study on a twisted spatial torus, directly probing volume reduction and winding sectors} \\
2017 & {\scriptsize Glueball spectrum from \$N\_\allowbreak{}f=2\$ lattice QCD study on anisotropic lattices} & {\scriptsize \texttt{1702.08174}} & {\scriptsize Anisotropic Nf=2 glueball spectroscopy with gluonic operators and mixing diagnostics} \\
2016 & {\scriptsize SU(N) gauge theories in 2+1 dimensions: glueball spectra and k-string tensions} & {\scriptsize \texttt{1609.03873}} & {\scriptsize continuum SU(N) spectrum in 2+1 dimensions} \\
2015 & {\scriptsize Topology and glueballs in SU(7) Yang-Mills with open boundary conditions} & {\scriptsize \texttt{1512.00806}} & {\scriptsize SU(7) open-boundary calculation linking topology sampling to glueball spectroscopy} \\
2013 & {\scriptsize Introductory lectures to large-N QCD phenomenology and lattice results} & {\scriptsize \texttt{1309.3638}} & {\scriptsize Peer-reviewed lectures synthesizing large-N phenomenology and lattice evidence} \\
2012 & {\scriptsize Recent results in large-N lattice gauge theories} & {\scriptsize \texttt{1210.5510}} & {\scriptsize Citation-guided review of then-current large-N lattice results} \\
2012 & {\scriptsize The string tension from smeared Wilson loops at large N} & {\scriptsize \texttt{1206.0049}} & {\scriptsize Highly cited string-tension extraction using smeared Wilson loops at large N} \\
2012 & {\scriptsize Towards the glueball spectrum from unquenched lattice QCD} & {\scriptsize \texttt{1208.1858}} & {\scriptsize Modern high-impact unquenched glueball spectrum and mixing study} \\
2010 & {\scriptsize Glueball mass measurements from improved staggered fermion simulations} & {\scriptsize \texttt{1005.2473}} & {\scriptsize Improved-staggered calculation of glueball masses with dynamical sea quarks} \\
2010 & {\scriptsize Glueball masses in the large N limit} & {\scriptsize \texttt{1007.3879}} & {\scriptsize large-N glueball masses} \\
2005 & {\scriptsize Glueball Spectrum and Matrix Elements on Anisotropic Lattices} & {\scriptsize \texttt{hep-lat/0510074}} & {\scriptsize anisotropic-lattice spectrum and matrix elements} \\
2005 & {\scriptsize Monte Carlo study of glueball masses in the Hamiltonian limit of SU(3) lattice gauge theory} & {\scriptsize \texttt{hep-lat/0503038}} & {\scriptsize Hamiltonian-limit SU(3) glueball masses} \\
2004 & {\scriptsize Glueball Regge trajectories and the Pomeron -- a lattice study --} & {\scriptsize \texttt{hep-ph/0409183}} & {\scriptsize High-impact lattice extraction of glueball Regge trajectories and the Pomeron connection} \\
2004 & {\scriptsize Glueballs and k-strings in SU(N) gauge theories : calculations with improved operators} & {\scriptsize \texttt{hep-lat/0404008}} & {\scriptsize improved glueball and k-string operators with large-N scaling} \\
2004 & {\scriptsize The spectrum of SU(N) gauge theories in finite volume} & {\scriptsize \texttt{hep-lat/0412021}} & {\scriptsize Dedicated finite-volume SU(N) spectrum and winding-state contamination analysis} \\
2003 & {\scriptsize Glueball Regge Trajectories in (2+1) Dimensional Gauge Theories} & {\scriptsize \texttt{hep-lat/0306019}} & {\scriptsize Direct 2+1D lattice determination of glueball Regge trajectories} \\
2002 & {\scriptsize SU(N) gauge theories in 2+1 dimensions -- further results} & {\scriptsize \texttt{hep-lat/0206027}} & {\scriptsize High-impact follow-up with additional SU(N) 2+1D spectrum/string results} \\
2001 & {\scriptsize On the glueball spectrum in O(a)-improved lattice QCD} & {\scriptsize \texttt{hep-lat/0108022}} & {\scriptsize O(a)-improved unquenched spectrum with finite-volume/torelon mixing analysis} \\
2001 & {\scriptsize SU(N) gauge theories in four dimensions: exploring the approach to N = infinity} & {\scriptsize \texttt{hep-lat/0103027}} & {\scriptsize Highly cited 4D SU(N) lattice study establishing the approach to N=infinity} \\
2000 & {\scriptsize Static potentials and glueball masses from QCD simulations with Wilson sea quarks} & {\scriptsize \texttt{hep-lat/0003012}} & {\scriptsize Highly cited sea-quark study that explicitly treats glueball and torelon contamination} \\
1999 & {\scriptsize The glueball spectrum from an anisotropic lattice study} & {\scriptsize \texttt{hep-lat/9901004}} & {\scriptsize benchmark glueball spectrum and operator construction} \\
1998 & {\scriptsize Glueball masses and other physical properties of SU(N) gauge theories in D=3+1: a review of lattice results for theorists} & {\scriptsize \texttt{hep-th/9812187}} & {\scriptsize Citation-heavy synthesis of 3+1D SU(N) lattice masses and observables} \\
1998 & {\scriptsize SU(N) gauge theories in 2+1 dimensions} & {\scriptsize \texttt{hep-lat/9804008}} & {\scriptsize Citation-heavy continuum SU(N) study in 2+1D, foundational for later spectrum tables} \\
1997 & {\scriptsize Efficient glueball simulations on anisotropic lattices} & {\scriptsize \texttt{hep-lat/9704011}} & {\scriptsize Introduces efficient anisotropic-lattice glueball spectroscopy techniques used by the landmark 1999 spectrum} \\
1993 & {\scriptsize A comprehensive lattice study of SU(3) glueballs} & {\scriptsize \texttt{hep-lat/9304012}} & {\scriptsize Classic high-statistics SU(3) spectrum; a major backward-citation anchor for modern continuum spectra} \\
\end{longtable}

\subsection*{Flux tubes, strings and torelons}
\addcontentsline{toc}{subsection}{Flux tubes, strings and torelons}
\begin{longtable}{@{}l p{0.40\linewidth} l p{0.28\linewidth}@{}}
\toprule Year & Title & arXiv & Connection \\ \midrule
\endfirsthead
\toprule Year & Title & arXiv & Connection \\ \midrule \endhead
\bottomrule\endfoot
2026 & {\scriptsize Entanglement Enabled Tomography of Flux Tubes in (2+1)D Yang-Mills Theory} & {\scriptsize \texttt{2601.17199}} & {\scriptsize 2026 Yang-Mills flux-tube tomography preprint; unusually fast early citation uptake} \\
2026 & {\scriptsize Imprints of asymptotic freedom on confining strings} & {\scriptsize \texttt{2602.15097}} & {\scriptsize 2026 PRL connecting asymptotic freedom to measurable confining-string imprints} \\
2025 & {\scriptsize The spectrum of open confining strings in the large-Nc limit} & {\scriptsize \texttt{2506.16342}} & {\scriptsize Latest peer-reviewed open-string spectrum extrapolated toward large Nc} \\
2025 & {\scriptsize Thermal nature of confining strings} & {\scriptsize \texttt{2510.23919}} & {\scriptsize Very recent peer-reviewed analysis of the thermal character of confining strings} \\
2024 & {\scriptsize Born-Oppenheimer Potentials for \$SU(3)\$ Gauge Theory} & {\scriptsize \texttt{2406.05123}} & {\scriptsize lattice-constrained SU(3) static spectrum connecting gluelumps and strings} \\
2024 & {\scriptsize Confining Strings and the Worldsheet Axion from the Lattice} & {\scriptsize \texttt{2411.03435}} & {\scriptsize Modern lattice-focused synthesis of confining strings and the worldsheet axion} \\
2024 & {\scriptsize Confining strings in three-dimensional gauge theories beyond the Nambu--Got\=o approximation} & {\scriptsize \texttt{2407.10678}} & {\scriptsize Recent high-precision study of confining strings beyond Nambu-Goto} \\
2024 & {\scriptsize Entanglement entropy of a color flux tube in (2+1)D Yang-Mills theory} & {\scriptsize \texttt{2410.00112}} & {\scriptsize Recent, well-cited lattice measurement of color-flux-tube entanglement structure} \\
2024 & {\scriptsize Multiparticle Flux Tube S-matrix Bootstrap} & {\scriptsize \texttt{2404.10812}} & {\scriptsize Recent PRL multiparticle extension of the flux-tube S-matrix bootstrap} \\
2023 & {\scriptsize Confining Strings and Glueballs in \$\textbackslash{}mathbb\{Z\}\_\allowbreak{}N\$ Gauge Theories} & {\scriptsize \texttt{2312.03855}} & {\scriptsize bulk glueball and flux-tube mixing isolated through Z\_\allowbreak{}N comparisons} \\
2023 & {\scriptsize Eight very excited spectra and one possible axion in SU(3) lattice gauge theory} & {\scriptsize \texttt{2303.15152}} & {\scriptsize Eight excited SU(3) flux-tube channels with candidate worldsheet-axion signals} \\
2021 & {\scriptsize Hybrid static potentials in SU(3) lattice gauge theory at small quark-antiquark separations} & {\scriptsize \texttt{2111.00741}} & {\scriptsize SU(3) hybrid static potentials and flux-tube excitations at short source separation} \\
2021 & {\scriptsize Revisiting the flux tube spectrum of 3d SU(2) lattice gauge theory} & {\scriptsize \texttt{2102.06413}} & {\scriptsize Reanalysis of the 3D SU(2) flux-tube spectrum and nonuniversal corrections} \\
2021 & {\scriptsize Spectrum of very excited \$$\Sigma$\_\allowbreak{}g\textasciicircum{}+\$ flux tubes in SU(3) gauge theory} & {\scriptsize \texttt{2105.12159}} & {\scriptsize Direct SU(3) extraction of very highly excited Sigma\_\allowbreak{}g+ flux-tube levels} \\
2021 & {\scriptsize The torelon spectrum and the world-sheet axion} & {\scriptsize \texttt{2112.11213}} & {\scriptsize modern SU(3), SU(5), and SU(6) torelon spectrum and worldsheet axion} \\
2019 & {\scriptsize Flux Tube S-matrix Bootstrap} & {\scriptsize \texttt{1906.08098}} & {\scriptsize Highly cited S-matrix bootstrap constraints for flux-tube worldsheet scattering} \\
2018 & {\scriptsize Spectrum of the open QCD flux tube and its effective string description} & {\scriptsize \texttt{1811.11779}} & {\scriptsize Proceedings update/companion on the open QCD flux-tube spectrum} \\
2018 & {\scriptsize Undressing Confining Flux Tubes with \$T\textbackslash{}bar T\$} & {\scriptsize \texttt{1808.01339}} & {\scriptsize Modern TbarT interpretation of confining-flux-tube spectra} \\
2017 & {\scriptsize Spectrum of the open QCD flux tube and its effective string description I: 3d static potential in SU(N=2,3)} & {\scriptsize \texttt{1705.03828}} & {\scriptsize Detailed open-flux-tube spectrum and effective-string comparison in 3D SU(2,3)} \\
2016 & {\scriptsize Closed flux tubes in D=2+1 SU(N) gauge theories: dynamics and effective string description} & {\scriptsize \texttt{1602.07634}} & {\scriptsize Detailed closed-flux-tube dynamics and effective-string analysis in 2+1D SU(N)} \\
2016 & {\scriptsize Effective string description of confining flux tubes} & {\scriptsize \texttt{1603.06969}} & {\scriptsize modern QCD flux-tube worldsheet theory} \\
2015 & {\scriptsize Towards a Theory of the QCD String} & {\scriptsize \texttt{1511.01908}} & {\scriptsize Influential axionic-string proposal motivated by lattice QCD-string data} \\
2014 & {\scriptsize A different kind of string} & {\scriptsize \texttt{1406.5127}} & {\scriptsize Influential analysis of non-Nambu-Goto worldsheet dynamics and its lattice signatures} \\
2014 & {\scriptsize Flux Tube Spectra from Approximate Integrability at Low Energies} & {\scriptsize \texttt{1404.0037}} & {\scriptsize Highly cited approximate-integrability analysis benchmarked against lattice flux-tube spectra} \\
2013 & {\scriptsize Evidence for a new particle on the worldsheet of the QCD flux tube} & {\scriptsize \texttt{1301.2325}} & {\scriptsize worldsheet dynamics of confining strings} \\
2013 & {\scriptsize Quantisation of the effective string with TBA} & {\scriptsize \texttt{1305.1278}} & {\scriptsize Highly cited TBA quantization framework for the effective confining string} \\
2013 & {\scriptsize The Effective Theory of Long Strings} & {\scriptsize \texttt{1302.6257}} & {\scriptsize effective theory of long strings} \\
2012 & {\scriptsize Effective String Theory Revisited} & {\scriptsize \texttt{1203.1054}} & {\scriptsize effective string theory and finite-volume spectra} \\
2012 & {\scriptsize SU(N) gauge theories at large N} & {\scriptsize \texttt{1210.4997}} & {\scriptsize large-N gauge theories lattice review} \\
2012 & {\scriptsize The Lorentz-invariant boundary action of the confining string and its universal contribution to the inter-quark potential} & {\scriptsize \texttt{1202.1984}} & {\scriptsize Citation-established derivation of the universal Lorentz-invariant boundary correction} \\
2011 & {\scriptsize Closed flux tubes and their string description in D=2+1 SU(N) gauge theories} & {\scriptsize \texttt{1103.5854}} & {\scriptsize closed flux tubes in 3+1 dimensions} \\
2010 & {\scriptsize Closed flux tubes and their string description in D=3+1 SU(N) gauge theories} & {\scriptsize \texttt{1007.4720}} & {\scriptsize closed flux-tube spectrum and string description} \\
2010 & {\scriptsize Corrections to Nambu-Goto energy levels from the effective string action} & {\scriptsize \texttt{1008.2648}} & {\scriptsize Systematic nonuniversal corrections to Nambu-Goto energy levels from the effective string action} \\
2010 & {\scriptsize The Width of the Color Flux Tube at 2-Loop Order} & {\scriptsize \texttt{1006.2252}} & {\scriptsize Two-loop effective-string prediction for the color-flux-tube width} \\
2010 & {\scriptsize The Width of the Confining String in Yang-Mills Theory} & {\scriptsize \texttt{1002.4888}} & {\scriptsize Highly cited lattice/effective-string determination of confining-string broadening in Yang-Mills} \\
2009 & {\scriptsize Large N and confining flux tubes as strings - a view from the lattice} & {\scriptsize \texttt{0912.3339}} & {\scriptsize large-N and confining-flux-tube lattice review} \\
2009 & {\scriptsize On the effective action of confining strings} & {\scriptsize \texttt{0903.1927}} & {\scriptsize Foundational classification of the confining-string effective action and its universal terms} \\
2009 & {\scriptsize Spectrum of the QCD flux tube in 3d SU(2) lattice gauge theory} & {\scriptsize \texttt{0905.4195}} & {\scriptsize Direct lattice spectrum of the 3D SU(2) QCD flux tube} \\
2008 & {\scriptsize Closed k-strings in SU(N) gauge theories : 2+1 dimensions} & {\scriptsize \texttt{0802.1490}} & {\scriptsize Closed k-string spectroscopy across SU(N) gauge groups} \\
2008 & {\scriptsize On the spectrum of closed k=2 flux tubes in D=2+1 SU(N) gauge theories} & {\scriptsize \texttt{0812.0334}} & {\scriptsize Dedicated lattice spectrum of closed k=2 flux tubes in 2+1D SU(N)} \\
2007 & {\scriptsize The closed string spectrum of SU(N) gauge theories in 2+1 dimensions} & {\scriptsize \texttt{0709.0693}} & {\scriptsize Closed-string spectrum in 2+1D SU(N), an early direct Nambu-Goto comparison} \\
2004 & {\scriptsize String excitation energies in SU(N) gauge theories beyond the free-string approximation} & {\scriptsize \texttt{hep-th/0406205}} & {\scriptsize universal effective-string corrections} \\
2004 & {\scriptsize Universal Subleading Spectrum of Effective String Theory} & {\scriptsize \texttt{hep-th/0411017}} & {\scriptsize Highly cited derivation of the universal subleading closed-string spectrum} \\
2003 & {\scriptsize Excitations of the static quark-antiquark system in several gauge theories} & {\scriptsize \texttt{hep-lat/0312019}} & {\scriptsize Early multi-gauge-theory comparison of static-source excitation spectra} \\
2002 & {\scriptsize Fine Structure of the QCD String Spectrum} & {\scriptsize \texttt{hep-lat/0207004}} & {\scriptsize Landmark high-citation lattice resolution of the fine structure of the QCD-string excitation spectrum} \\
2001 & {\scriptsize Confining strings in SU(N) gauge theories} & {\scriptsize \texttt{hep-lat/0107007}} & {\scriptsize Foundational numerical SU(N) confining-string and k-string study} \\
\end{longtable}

\subsection*{Rigorous lattice Yang--Mills}
\addcontentsline{toc}{subsection}{Rigorous lattice Yang--Mills}
\begin{longtable}{@{}l p{0.40\linewidth} l p{0.28\linewidth}@{}}
\toprule Year & Title & arXiv & Connection \\ \midrule
\endfirsthead
\toprule Year & Title & arXiv & Connection \\ \midrule \endhead
\bottomrule\endfoot
2026 & {\scriptsize Confinement as Decoding: Higher Form Codes and Lattice Yang-Mills Theory} & {\scriptsize \texttt{2608.02452}} & {\scriptsize higher-form-code route to confinement and transfer-matrix mass} \\
2026 & {\scriptsize Direct and Indirect Loop Equations in Lattice Yang-Mills Theory} & {\scriptsize \texttt{2601.04316}} & {\scriptsize modern direct and indirect lattice loop equations} \\
2026 & {\scriptsize On the leading order term of the lattice Yang-Mills free energy} & {\scriptsize \texttt{}} & {\scriptsize modern rigorous infinite-volume lattice Yang-Mills free-energy asymptotics} \\
2026 & {\scriptsize The heat-kernel master field on \$\textbackslash{}mathbb\{Z\}\textasciicircum{}d\$ at strong coupling} & {\scriptsize \texttt{2606.28945}} & {\scriptsize strong-coupling heat-kernel master field, 1/N expansion, area law, and master-loop equation} \\
2025 & {\scriptsize \$\textbackslash{}mathrm\{U\}(N)\$ lattice Yang-Mills in the 't Hooft regime} & {\scriptsize \texttt{2510.22788}} & {\scriptsize new uniform strong-coupling U(N) results} \\
2025 & {\scriptsize Bootstrapping SU(3) Lattice Yang-Mills Theory} & {\scriptsize \texttt{2502.14421}} & {\scriptsize SU(3) lattice Yang-Mills bootstrap} \\
2025 & {\scriptsize Correlation decay for U(1) lattice Higgs theory: the case of small mass} & {\scriptsize \texttt{2509.19176}} & {\scriptsize constructive cluster expansion and correlation decay for small-mass U(1) lattice Higgs theory} \\
2025 & {\scriptsize Current expansion and couplings for Ising lattice gauge theory} & {\scriptsize \texttt{2502.19942}} & {\scriptsize random-current expansion and switching lemma for lattice gauge theory} \\
2025 & {\scriptsize Dynamical approach to area law for lattice Yang-Mills} & {\scriptsize \texttt{2509.04688}} & {\scriptsize dynamical mass-gap condition implying Wilson area law} \\
2025 & {\scriptsize Expanded regimes of area law for lattice Yang-Mills theories} & {\scriptsize \texttt{2505.16585}} & {\scriptsize expanded strong-coupling area-law regimes at large N} \\
2025 & {\scriptsize Surface sums in two-dimensional large-\$N\$ lattice Yang--Mills: Cancellations and explicit computations for general loops} & {\scriptsize \texttt{2508.13827}} & {\scriptsize large-N two-dimensional lattice-Yang-Mills surface-sum cancellations} \\
2025 & {\scriptsize Villain action in lattice gauge theory} & {\scriptsize \texttt{}} & {\scriptsize modern rigorous Wilson/Manton-to-Villain limit for compact gauge groups} \\
2024 & {\scriptsize Bootstrap for Finite N Lattice Yang-Mills Theory} & {\scriptsize \texttt{2404.16925}} & {\scriptsize finite-N lattice Yang-Mills bootstrap} \\
2024 & {\scriptsize Langevin dynamics of lattice Yang-Mills-Higgs and applications} & {\scriptsize \texttt{2401.13299}} & {\scriptsize strong-coupling mass gap and exponential ergodicity with Higgs matter} \\
2024 & {\scriptsize Makeenko-Migdal equations for 2D Yang-Mills: from lattice to continuum} & {\scriptsize \texttt{2412.15422}} & {\scriptsize rigorous convergence of two-dimensional lattice Makeenko-Migdal equations to continuum area derivatives} \\
2024 & {\scriptsize Surface sums for lattice Yang-Mills in the large-\$N\$ limit} & {\scriptsize \texttt{2411.11676}} & {\scriptsize large-N surface sums with signed-Catalan weights} \\
2023 & {\scriptsize Geometric Derivation of the Finite \$N\$ Master Loop Equation} & {\scriptsize \texttt{2309.07399}} & {\scriptsize geometric derivation of the finite-N lattice master loop equation} \\
2023 & {\scriptsize Invariant measure and universality of the 2D Yang-Mills Langevin dynamic} & {\scriptsize \texttt{2302.12160}} & {\scriptsize invariant measure and universality for Yang-Mills Langevin dynamics} \\
2023 & {\scriptsize Random surfaces and lattice Yang-Mills} & {\scriptsize \texttt{2307.06790}} & {\scriptsize random-surface expansion for lattice Yang-Mills} \\
2022 & {\scriptsize A new derivation of the finite \$N\$ master loop equation for lattice Yang-Mills} & {\scriptsize \texttt{2202.00880}} & {\scriptsize finite-N master loop equation for SU(N), U(N), and SO(N)} \\
2022 & {\scriptsize A stochastic analysis approach to lattice Yang--Mills at strong coupling} & {\scriptsize \texttt{2204.12737}} & {\scriptsize strong-coupling mass gap and exponential clustering} \\
2022 & {\scriptsize Bootstrap for Lattice Yang-Mills theory} & {\scriptsize \texttt{2203.11360}} & {\scriptsize large-N lattice Yang-Mills bootstrap from loop equations and positivity} \\
2022 & {\scriptsize Correlation decay for finite lattice gauge theories at weak coupling} & {\scriptsize \texttt{2202.10375}} & {\scriptsize exponential correlation decay for finite lattice gauge theories at weak coupling} \\
2022 & {\scriptsize Large N limit of the Yang-Mills measure on compact surfaces II: Makeenko-Migdal equations and planar master field} & {\scriptsize \texttt{2201.05886}} & {\scriptsize modern compact-surface large-N master field and Makeenko-Migdal theory} \\
2022 & {\scriptsize Stochastic quantisation of Yang-Mills-Higgs in 3D} & {\scriptsize \texttt{2201.03487}} & {\scriptsize stochastic quantization and continuum control for three-dimensional Yang-Mills-Higgs} \\
2021 & {\scriptsize A state space for 3D Euclidean Yang-Mills theories} & {\scriptsize \texttt{2111.12813}} & {\scriptsize rigorous state space for three-dimensional Euclidean Yang-Mills} \\
2021 & {\scriptsize Decay of correlations in finite Abelian lattice gauge theories} & {\scriptsize \texttt{2104.03752}} & {\scriptsize rigorous exponential decay of correlations in finite Abelian lattice gauge theories} \\
2021 & {\scriptsize Wilson Loop Expectations for Non-Abelian Gauge Fields Coupled to a Higgs Boson at Low and High Disorder} & {\scriptsize \texttt{2111.07540}} & {\scriptsize non-Abelian Wilson-loop asymptotics with Higgs matter via vortex and polymer expansions} \\
2021 & {\scriptsize Wilson lines in the Abelian lattice Higgs model} & {\scriptsize \texttt{2111.06620}} & {\scriptsize rigorous Wilson-line asymptotics in the Abelian lattice Higgs model} \\
2020 & {\scriptsize A probabilistic mechanism for quark confinement} & {\scriptsize \texttt{2006.16229}} & {\scriptsize probabilistic confinement mechanism} \\
2020 & {\scriptsize Langevin dynamic for the 2D Yang-Mills measure} & {\scriptsize \texttt{2006.04987}} & {\scriptsize rigorous Yang-Mills Langevin dynamics and Wilson observables} \\
2020 & {\scriptsize Wilson loop expectations in lattice gauge theories with finite gauge groups} & {\scriptsize \texttt{2001.05627}} & {\scriptsize rigorous weak-coupling Wilson-loop estimates for finite gauge groups} \\
2018 & {\scriptsize Wilson loops in Ising lattice gauge theory} & {\scriptsize \texttt{1811.09770}} & {\scriptsize rigorous weak-coupling Wilson-loop asymptotics in four-dimensional lattice gauge theory} \\
2018 & {\scriptsize Yang-Mills for probabilists} & {\scriptsize \texttt{1803.01950}} & {\scriptsize rigorous/probabilistic lattice Yang-Mills review} \\
2016 & {\scriptsize Loop Equations and bootstrap methods in the lattice} & {\scriptsize \texttt{1612.08140}} & {\scriptsize positivity bootstrap for lattice loop equations} \\
2016 & {\scriptsize The \$1/N\$ expansion for SO(N) lattice gauge theory at strong coupling} & {\scriptsize \texttt{1604.04777}} & {\scriptsize all-orders genus-graded 1/N surface expansion for lattice gauge theory} \\
2016 & {\scriptsize Three proofs of the Makeenko-Migdal equation for Yang-Mills theory on the plane} & {\scriptsize \texttt{1601.06283}} & {\scriptsize three rigorous derivations of the planar Makeenko-Migdal equation} \\
2016 & {\scriptsize Wilson loop expectations in \$SU(N)\$ lattice gauge theory} & {\scriptsize \texttt{1610.03821}} & {\scriptsize Wilson-loop expectations in SU(N) lattice gauge theory} \\
2015 & {\scriptsize Rigorous solution of strongly coupled \$SO(N)\$ lattice gauge theory in the large \$N\$ limit} & {\scriptsize \texttt{1502.07719}} & {\scriptsize rigorous convergent large-N strong-coupling Wilson-loop surface expansion} \\
2011 & {\scriptsize The master field on the plane} & {\scriptsize \texttt{1112.2452}} & {\scriptsize planar large-N master field characterized by Makeenko-Migdal equations} \\
\end{longtable}

\subsection*{Spectral extraction and operator overlap}
\addcontentsline{toc}{subsection}{Spectral extraction and operator overlap}
\begin{longtable}{@{}l p{0.40\linewidth} l p{0.28\linewidth}@{}}
\toprule Year & Title & arXiv & Connection \\ \midrule
\endfirsthead
\toprule Year & Title & arXiv & Connection \\ \midrule \endhead
\bottomrule\endfoot
2026 & {\scriptsize Charmonium-Glueball spectroscopy with improved hadron creation operators} & {\scriptsize \texttt{2603.20178}} & {\scriptsize modern derivative-based glueball operators and overlap analysis} \\
2024 & {\scriptsize Bayesian solution to the inverse problem and its relation to Backus-Gilbert methods} & {\scriptsize \texttt{2409.04413}} & {\scriptsize Modern Bayesian interpretation and extension of Backus-Gilbert inverse methods} \\
2024 & {\scriptsize Block Lanczos algorithm for lattice QCD spectroscopy and matrix elements} & {\scriptsize \texttt{2412.04444}} & {\scriptsize Rapidly cited block-Lanczos spectroscopy/matrix-element method published in 2025} \\
2024 & {\scriptsize Lanczos, the transfer matrix, and the signal-to-noise problem} & {\scriptsize \texttt{2406.20009}} & {\scriptsize transfer-matrix Lanczos extraction with two-sided error bounds} \\
2024 & {\scriptsize Spectral densities from Euclidean lattice correlators via the Mellin transform} & {\scriptsize \texttt{2407.04141}} & {\scriptsize Recent Mellin-transform framework with controlled lattice spectral-density reconstruction} \\
2023 & {\scriptsize Beyond Generalized Eigenvalues in Lattice Quantum Field Theory} & {\scriptsize \texttt{2309.05111}} & {\scriptsize Modern proposal going beyond the standard generalized-eigenvalue analysis} \\
2023 & {\scriptsize Teaching to extract spectral densities from lattice correlators to a broad audience of learning-machines} & {\scriptsize \texttt{2307.00808}} & {\scriptsize Peer-reviewed learning-machine benchmark for spectral-density extraction} \\
2022 & {\scriptsize Yang-Mills glueball masses from spectral reconstruction} & {\scriptsize \texttt{2212.01113}} & {\scriptsize independent Yang-Mills spectral reconstruction of glueball masses} \\
2021 & {\scriptsize K\"all\'en-Lehmann Spectral Representation of the Scalar SU(2) Glueball} & {\scriptsize \texttt{2103.11846}} & {\scriptsize Kallen-Lehmann reconstruction of the scalar SU(2) glueball} \\
2020 & {\scriptsize Reconstruction of smeared spectral function from Euclidean correlation functions} & {\scriptsize \texttt{2001.11779}} & {\scriptsize Widely used reconstruction method for smeared spectral functions from Euclidean data} \\
2019 & {\scriptsize On the extraction of spectral densities from lattice correlators} & {\scriptsize \texttt{1903.06476}} & {\scriptsize Highly cited regularized method for extracting smeared spectral densities from Euclidean correlators} \\
2013 & {\scriptsize Extended hadron and two-hadron operators of definite momentum for spectrum calculations in lattice QCD} & {\scriptsize \texttt{1303.6816}} & {\scriptsize Highly cited group-theoretic construction of extended momentum-projected operator bases for lattice spectra} \\
2010 & {\scriptsize Properties and uses of the Wilson flow in lattice QCD} & {\scriptsize \texttt{1006.4518}} & {\scriptsize Wilson flow and smoothed gauge observables} \\
2009 & {\scriptsize A novel quark-field creation operator construction for hadronic physics in lattice QCD} & {\scriptsize \texttt{0905.2160}} & {\scriptsize distillation and optimized hadron operators} \\
2009 & {\scriptsize On the generalized eigenvalue method for energies and matrix elements in lattice field theory} & {\scriptsize \texttt{0902.1265}} & {\scriptsize generalized eigenvalue method with error control} \\
2008 & {\scriptsize Efficient use of the Generalized Eigenvalue Problem} & {\scriptsize \texttt{0808.1017}} & {\scriptsize Practical companion on efficient GEVP use, distinct from the 2009 asymptotic-error paper already collected} \\
2003 & {\scriptsize Analytic Smearing of SU(3) Link Variables in Lattice QCD} & {\scriptsize \texttt{hep-lat/0311018}} & {\scriptsize analytic stout-link smearing} \\
2002 & {\scriptsize Locality and Statistical Error Reduction on Correlation Functions} & {\scriptsize \texttt{hep-lat/0209145}} & {\scriptsize local multilevel error reduction for glueball correlation functions} \\
2001 & {\scriptsize Locality and exponential error reduction in numerical lattice gauge theory} & {\scriptsize \texttt{hep-lat/0108014}} & {\scriptsize multilevel exponential error reduction} \\
\end{longtable}

\subsection*{Weingarten and Haar integration}
\addcontentsline{toc}{subsection}{Weingarten and Haar integration}
\begin{longtable}{@{}l p{0.40\linewidth} l p{0.28\linewidth}@{}}
\toprule Year & Title & arXiv & Connection \\ \midrule
\endfirsthead
\toprule Year & Title & arXiv & Connection \\ \midrule \endhead
\bottomrule\endfoot
2026 & {\scriptsize Universal dualities for Wilson loops in lattice Yang-Mills} & {\scriptsize \texttt{2604.16252}} & {\scriptsize finite-N Wilson-loop dualities and Weingarten state sums} \\
2025 & {\scriptsize Weingarten Calculus with Virtual Isometries} & {\scriptsize \texttt{2510.21186}} & {\scriptsize new virtual-isometry extension} \\
2022 & {\scriptsize Expectation values of polynomials and moments on general compact Lie groups} & {\scriptsize \texttt{2203.11607}} & {\scriptsize Weingarten maps and polynomial moments for general compact Lie groups including SU(N)} \\
2022 & {\scriptsize Moment Methods on compact groups: Weingarten calculus and its applications} & {\scriptsize \texttt{2207.08418}} & {\scriptsize compact-group moment methods} \\
2021 & {\scriptsize The Weingarten Calculus} & {\scriptsize \texttt{2109.14890}} & {\scriptsize modern Weingarten-calculus review with lattice-gauge application} \\
2021 & {\scriptsize Weingarten Calculus} & {\scriptsize \texttt{2101.00921}} & {\scriptsize comprehensive Weingarten-calculus survey} \\
2018 & {\scriptsize Matrix Group Integrals, Surfaces, and Mapping Class Groups I: \$U(n)\$} & {\scriptsize \texttt{1802.04862}} & {\scriptsize Haar unitary integrals encoded by surfaces and mapping-class groups} \\
2018 & {\scriptsize SU(N) polynomial integrals and some applications} & {\scriptsize \texttt{1812.06069}} & {\scriptsize explicit SU(N) polynomial integrals for lattice-gauge applications} \\
2017 & {\scriptsize Weingarten calculus via orthogonality relations: new applications} & {\scriptsize \texttt{1701.04493}} & {\scriptsize orthogonality-recursion method} \\
2016 & {\scriptsize Revisiting SU(N) integrals} & {\scriptsize \texttt{1611.00236}} & {\scriptsize explicit invariant integration formulas for SU(N)} \\
2013 & {\scriptsize Weingarten calculus for matrix ensembles associated with compact symmetric spaces} & {\scriptsize \texttt{1301.5401}} & {\scriptsize Weingarten calculus for compact symmetric spaces} \\
2009 & {\scriptsize On some properties of orthogonal Weingarten functions} & {\scriptsize \texttt{0903.5143}} & {\scriptsize structural and asymptotic properties of orthogonal Weingarten functions} \\
2008 & {\scriptsize Complete homogeneous symmetric polynomials in Jucys-Murphy elements and the Weingarten function} & {\scriptsize \texttt{0811.3595}} & {\scriptsize Jucys-Murphy and symmetric-function formula for the unitary Weingarten function} \\
2006 & {\scriptsize Asymptotics of unitary and othogonal matrix integrals} & {\scriptsize \texttt{math/0608193}} & {\scriptsize asymptotics of unitary and orthogonal matrix integrals used in surface weights} \\
2005 & {\scriptsize Integration over compact quantum groups} & {\scriptsize \texttt{math/0511253}} & {\scriptsize Weingarten integration extended to compact quantum groups} \\
2004 & {\scriptsize Integration with respect to the Haar measure on unitary, orthogonal and symplectic group} & {\scriptsize \texttt{math-ph/0402073}} & {\scriptsize foundational Haar-unitary integration formula} \\
2002 & {\scriptsize Moments and Cumulants of Polynomial random variables on unitary groups, the Itzykson-Zuber integral and free probability} & {\scriptsize \texttt{math-ph/0205010}} & {\scriptsize unitary matrix moments, cumulants, and Itzykson-Zuber asymptotics} \\
\end{longtable}

\subsection*{Effective Hamiltonians and strong coupling}
\addcontentsline{toc}{subsection}{Effective Hamiltonians and strong coupling}
\begin{longtable}{@{}l p{0.40\linewidth} l p{0.28\linewidth}@{}}
\toprule Year & Title & arXiv & Connection \\ \midrule
\endfirsthead
\toprule Year & Title & arXiv & Connection \\ \midrule \endhead
\bottomrule\endfoot
2025 & {\scriptsize The Spectral Renormalization Flow Based on the Smooth Feshbach--Schur Map: The Introduction of the Semi-Group Property} & {\scriptsize \texttt{2508.07643}} & {\scriptsize modern smooth Feshbach-Schur spectral flow with a semigroup property} \\
2024 & {\scriptsize The geometry of the Hermitian matrix space and the Schrieffer--Wolff transformation} & {\scriptsize \texttt{2407.10478}} & {\scriptsize modern geometric treatment of Schrieffer-Wolff} \\
2022 & {\scriptsize A mathematical framework for quantum Hamiltonian simulation and duality} & {\scriptsize \texttt{2208.11941}} & {\scriptsize precise mathematical framework for Hamiltonian simulation and duality} \\
2021 & {\scriptsize The Feshbach-Schur map and perturbation theory} & {\scriptsize \texttt{2105.02058}} & {\scriptsize explicit Feshbach-Schur perturbation estimates} \\
2020 & {\scriptsize Hamiltonian simulation in the low-energy subspace} & {\scriptsize \texttt{2006.02660}} & {\scriptsize explicit leakage and error bounds for low-energy-subspace Hamiltonian simulation} \\
2017 & {\scriptsize Efficient optimization of perturbative gadgets} & {\scriptsize \texttt{1709.02705}} & {\scriptsize optimized perturbative-gadget error bounds at arbitrary order} \\
2017 & {\scriptsize Universal Quantum Hamiltonians} & {\scriptsize \texttt{1701.05182}} & {\scriptsize rigorous low-energy Hamiltonian-simulation framework and universal Hamiltonians} \\
2014 & {\scriptsize Perturbative gadgets without strong interactions} & {\scriptsize \texttt{1408.5881}} & {\scriptsize convergence and error control for perturbative gadgets without strong interactions} \\
2011 & {\scriptsize Schrieffer-Wolff transformation for quantum many-body systems} & {\scriptsize \texttt{1105.0675}} & {\scriptsize rigorous Schrieffer-Wolff transformation and error bounds} \\
2010 & {\scriptsize Strong-coupling expansion and effective hamiltonians} & {\scriptsize \texttt{1005.2495}} & {\scriptsize strong-coupling effective-Hamiltonian expansion review} \\
2008 & {\scriptsize Perturbative Gadgets at Arbitrary Orders} & {\scriptsize \texttt{0802.1874}} & {\scriptsize arbitrary-order perturbative gadgets with explicit effective interactions} \\
2008 & {\scriptsize Simulation of Many-Body Hamiltonians using Perturbation Theory with Bounded-Strength Interactions} & {\scriptsize \texttt{0803.2686}} & {\scriptsize bounded-strength many-body simulation by perturbative effective Hamiltonians} \\
2007 & {\scriptsize On the Smooth Feshbach-Schur Map} & {\scriptsize \texttt{0704.3244}} & {\scriptsize smooth Feshbach-Schur map} \\
2004 & {\scriptsize Higher order effective low-energy theories} & {\scriptsize \texttt{cond-mat/0407255}} & {\scriptsize fourth-order low-energy expansion comparing canonical and projection-resolvent methods} \\
2004 & {\scriptsize The Complexity of the Local Hamiltonian Problem} & {\scriptsize \texttt{quant-ph/0406180}} & {\scriptsize rigorous low-energy projection gadgets for encoded Hamiltonians} \\
\end{longtable}

\subsection*{Flat bands and compact localized states}
\addcontentsline{toc}{subsection}{Flat bands and compact localized states}
\begin{longtable}{@{}l p{0.40\linewidth} l p{0.28\linewidth}@{}}
\toprule Year & Title & arXiv & Connection \\ \midrule
\endfirsthead
\toprule Year & Title & arXiv & Connection \\ \midrule \endhead
\bottomrule\endfoot
2026 & {\scriptsize New class of exactly flat topological bands - compact localised states protected by local graph topology} & {\scriptsize \texttt{2607.22831}} & {\scriptsize exact topological flat bands protected by local graph-index structure} \\
2023 & {\scriptsize Flat bands of periodic graphs} & {\scriptsize \texttt{2304.06465}} & {\scriptsize mathematical classification of flat bands on periodic graphs} \\
2020 & {\scriptsize Fragile topology in line-graph lattices with two, three, or four gapped flat bands} & {\scriptsize \texttt{2010.11953}} & {\scriptsize line-graph cycle-space flat bands and fragile topology} \\
2020 & {\scriptsize Singular flat bands} & {\scriptsize \texttt{2012.04279}} & {\scriptsize singular-flat-band classification, topology, and compact-state completeness} \\
2019 & {\scriptsize Line-Graph Lattices: Euclidean and Non-Euclidean Flat Bands, and Implementations in Circuit Quantum Electrodynamics} & {\scriptsize \texttt{1902.02794}} & {\scriptsize line-graph lattices and homological flat bands} \\
2018 & {\scriptsize Artificial flat band systems: from lattice models to experiments} & {\scriptsize \texttt{1801.09378}} & {\scriptsize highly cited flat-band systems review} \\
2018 & {\scriptsize Classification of flat bands according to the band-crossing singularity of Bloch wave functions} & {\scriptsize \texttt{1808.05926}} & {\scriptsize classification of singular and nonsingular flat bands} \\
2017 & {\scriptsize Compact localized states and flat bands from local symmetry partitioning} & {\scriptsize \texttt{1709.07806}} & {\scriptsize local-symmetry construction of compact localized states and flat bands} \\
2016 & {\scriptsize Compact localized states and flatband generators in one dimension} & {\scriptsize \texttt{1610.02970}} & {\scriptsize classification and generation of one-dimensional compact localized states} \\
2015 & {\scriptsize Strongly correlated flat-band systems: The route from Heisenberg spins to Hubbard electrons} & {\scriptsize \texttt{1502.02729}} & {\scriptsize strongly correlated flat-band review} \\
2008 & {\scriptsize Band touching from real space topology in frustrated hopping models} & {\scriptsize \texttt{0803.3628}} & {\scriptsize real-space topology and singular band touching} \\
2005 & {\scriptsize On the structure of eigenfunctions corresponding to embedded eigenvalues of locally perturbed periodic graph operators} & {\scriptsize \texttt{math-ph/0511084}} & {\scriptsize compactly supported eigenfunctions for perturbed periodic graph operators} \\
1998 & {\scriptsize Aharonov-Bohm cages in two-dimensional structures} & {\scriptsize \texttt{cond-mat/9806068}} & {\scriptsize Aharonov-Bohm cages and exact compact localization by flux interference} \\
1997 & {\scriptsize From Nagaoka's ferromagnetism to flat-band ferromagnetism and beyond: An introduction to ferromagnetism in the Hubbard model} & {\scriptsize \texttt{cond-mat/9712219}} & {\scriptsize rigorous flat-band ferromagnetism review and theorem synthesis} \\
\end{longtable}

\subsection*{Strong-coupling expansions}
\addcontentsline{toc}{subsection}{Strong-coupling expansions}
\begin{longtable}{@{}l p{0.40\linewidth} l p{0.28\linewidth}@{}}
\toprule Year & Title & arXiv & Connection \\ \midrule
\endfirsthead
\toprule Year & Title & arXiv & Connection \\ \midrule \endhead
\bottomrule\endfoot
2010 & {\scriptsize Centre symmetric 3d effective actions for thermal SU(N) Yang-Mills from strong coupling series} & {\scriptsize \texttt{1010.0951}} & {\scriptsize strong-coupling effective theory review/application} \\
2008 & {\scriptsize Strong coupling expansion for finite temperature Yang-Mills theory in the confined phase} & {\scriptsize \texttt{0805.1163}} & {\scriptsize strong-coupling thermodynamics and equation of state} \\
2007 & {\scriptsize Strong coupling expansion for Yang-Mills theory at finite temperature} & {\scriptsize \texttt{0710.0512}} & {\scriptsize finite-temperature strong-coupling expansion} \\
1985 & {\scriptsize Effective transfer matrix for low-lying glueball states in lattice gauge theory} & {\scriptsize \texttt{}} & {\scriptsize effective transfer-matrix eigenvalues and low-lying glueball dispersion} \\
1983 & {\scriptsize Strong coupling and mean field methods in lattice gauge theories} & {\scriptsize \texttt{}} & {\scriptsize canonical review of character, graph, cluster, correlation, and spectral strong-coupling machinery} \\
\end{longtable}

% ---- end gen_library.tex -------------------------------------------
% ---- end B2_library.tex --------------------------------------------
% ---- begin B3_extraction.tex ---------------------------------------
\section{The page-level extraction index}
\label{app:extraction}

\ExtractedPages{} extracted pages across \ExtractedWorks{} works,
indexed by topic and line position. The workflow this supports is in
\S\ref{sec:sources}: search by paper and topic, follow the locator into
the reading copy, verify the equation in the pinned PDF.

The index records locators rather than republishing source text. A
topic hit may be a citation rather than a result, or an extraction
error; candidate theorem labels may sit inside citations or carry OCR
faults. Five pages extracted fewer than eighty characters and are
flagged unreliable in the source index --- their contents must not be
inferred from the text. Exact equations from legacy scans must be
checked in the PDF, not in the extraction.

% ---- begin gen_extraction.tex --------------------------------------
% Generated by paper/master/generate_sources.py.  Do not edit.
\noindent 777 extracted pages across 28 works, indexed by topic and line. A row says how many pages of that work are searchable and which topics occur in it; the index records \emph{locators}, not text, so the reading copy under \srcpath{literature/inbox/<ID>/extracted/} is where the words are. A hit may be a citation or an extraction error: the index is a search aid and certifies nothing.\par\medskip
\begin{longtable}{@{}l r r p{0.46\linewidth}@{}}
\toprule Work & Pages & Labelled & Topics present \\ \midrule
\endfirsthead
\toprule Work & Pages & Labelled & Topics present \\ \midrule \endhead
\bottomrule\endfoot
{\scriptsize\ttfamily BFS\_\allowbreak{}1980\_\allowbreak{}II} & 85 & 52 & {\scriptsize gauge\_\allowbreak{}geometry (89), constructive (86), uniform\_\allowbreak{}limits (78), reconstruction (11), ground\_\allowbreak{}state (3)} \\
{\scriptsize\ttfamily MMM\_\allowbreak{}2019\_\allowbreak{}CURVATURE} & 80 & 0 & {\scriptsize curvature (93), gauge\_\allowbreak{}geometry (82), ground\_\allowbreak{}state (45), spectral\_\allowbreak{}gap (18), constructive (2), uniform\_\allowbreak{}limits (2)} \\
{\scriptsize\ttfamily WITTEN\_\allowbreak{}1994\_\allowbreak{}FOUR\_\allowbreak{}MANIFOLD} & 67 & 0 & {\scriptsize gauge\_\allowbreak{}geometry (57), spectral\_\allowbreak{}gap (37), ground\_\allowbreak{}state (35), constructive (5), curvature (1), uniform\_\allowbreak{}limits (1)} \\
{\scriptsize\ttfamily MRS\_\allowbreak{}1993} & 60 & 11 & {\scriptsize gauge\_\allowbreak{}geometry (171), uniform\_\allowbreak{}limits (141), constructive (63), reconstruction (5), curvature (1)} \\
{\scriptsize\ttfamily BRS\_\allowbreak{}1976} & 58 & 0 & {\scriptsize gauge\_\allowbreak{}geometry (46), constructive (20), ground\_\allowbreak{}state (2)} \\
{\scriptsize\ttfamily SW\_\allowbreak{}1994} & 45 & 0 & {\scriptsize gauge\_\allowbreak{}geometry (67), ground\_\allowbreak{}state (21), constructive (12), spectral\_\allowbreak{}gap (3)} \\
{\scriptsize\ttfamily IMBRIE\_\allowbreak{}1981\_\allowbreak{}I} & 44 & 20 & {\scriptsize ground\_\allowbreak{}state (26), constructive (17), uniform\_\allowbreak{}limits (4), reconstruction (3), spectral\_\allowbreak{}gap (2), curvature (1)} \\
{\scriptsize\ttfamily MONDAL\_\allowbreak{}2023\_\allowbreak{}GEOMETRIC} & 42 & 10 & {\scriptsize curvature (179), gauge\_\allowbreak{}geometry (135), spectral\_\allowbreak{}gap (42), ground\_\allowbreak{}state (40), constructive (28), uniform\_\allowbreak{}limits (5)} \\
{\scriptsize\ttfamily BBIJ\_\allowbreak{}1984} & 39 & 16 & {\scriptsize gauge\_\allowbreak{}geometry (30), spectral\_\allowbreak{}gap (21), uniform\_\allowbreak{}limits (11), constructive (10), ground\_\allowbreak{}state (3), reconstruction (3)} \\
{\scriptsize\ttfamily IMBRIE\_\allowbreak{}1981\_\allowbreak{}II} & 39 & 27 & {\scriptsize ground\_\allowbreak{}state (13), uniform\_\allowbreak{}limits (13), constructive (12), reconstruction (4), spectral\_\allowbreak{}gap (2)} \\
{\scriptsize\ttfamily OS\_\allowbreak{}1973} & 30 & 19 & {\scriptsize reconstruction (68), uniform\_\allowbreak{}limits (9), gauge\_\allowbreak{}geometry (2), spectral\_\allowbreak{}gap (2), ground\_\allowbreak{}state (1), constructive (1)} \\
{\scriptsize\ttfamily OS\_\allowbreak{}1975} & 26 & 15 & {\scriptsize reconstruction (49), uniform\_\allowbreak{}limits (4), gauge\_\allowbreak{}geometry (2), ground\_\allowbreak{}state (1)} \\
{\scriptsize\ttfamily FROHLICH\_\allowbreak{}1982} & 25 & 0 & {\scriptsize constructive (11), reconstruction (4), uniform\_\allowbreak{}limits (3), ground\_\allowbreak{}state (1)} \\
{\scriptsize\ttfamily MALDACENA\_\allowbreak{}1998} & 22 & 0 & {\scriptsize gauge\_\allowbreak{}geometry (11), curvature (2), ground\_\allowbreak{}state (2), constructive (2)} \\
{\scriptsize\ttfamily JAFFE\_\allowbreak{}2000\_\allowbreak{}CQFT} & 16 & 0 & {\scriptsize constructive (33), ground\_\allowbreak{}state (20), reconstruction (17), gauge\_\allowbreak{}geometry (15), spectral\_\allowbreak{}gap (10), curvature (2)} \\
{\scriptsize\ttfamily JW\_\allowbreak{}2006} & 14 & 0 & {\scriptsize gauge\_\allowbreak{}geometry (66), spectral\_\allowbreak{}gap (25), constructive (24), reconstruction (22), ground\_\allowbreak{}state (9), uniform\_\allowbreak{}limits (6)} \\
{\scriptsize\ttfamily THOOFT\_\allowbreak{}1974} & 13 & 0 & {\scriptsize gauge\_\allowbreak{}geometry (24), spectral\_\allowbreak{}gap (1)} \\
{\scriptsize\ttfamily OS\_\allowbreak{}SENEOR\_\allowbreak{}1976} & 13 & 9 & {\scriptsize reconstruction (11), uniform\_\allowbreak{}limits (11), ground\_\allowbreak{}state (2), constructive (2)} \\
{\scriptsize\ttfamily SIMON\_\allowbreak{}1983\_\allowbreak{}DISCRETE} & 12 & 2 & {\scriptsize spectral\_\allowbreak{}gap (20), ground\_\allowbreak{}state (2), gauge\_\allowbreak{}geometry (1), uniform\_\allowbreak{}limits (1)} \\
{\scriptsize\ttfamily MR\_\allowbreak{}1976\_\allowbreak{}CRITICAL} & 10 & 3 & {\scriptsize reconstruction (6), uniform\_\allowbreak{}limits (5), constructive (4), ground\_\allowbreak{}state (1), spectral\_\allowbreak{}gap (1)} \\
{\scriptsize\ttfamily CREUTZ\_\allowbreak{}1980} & 8 & 0 & {\scriptsize gauge\_\allowbreak{}geometry (39), uniform\_\allowbreak{}limits (22), constructive (22), ground\_\allowbreak{}state (1), reconstruction (1)} \\
{\scriptsize\ttfamily GJ\_\allowbreak{}1968\_\allowbreak{}I} & 7 & 5 & {\scriptsize uniform\_\allowbreak{}limits (18), ground\_\allowbreak{}state (8), reconstruction (2), gauge\_\allowbreak{}geometry (1), constructive (1), spectral\_\allowbreak{}gap (1)} \\
{\scriptsize\ttfamily YANG\_\allowbreak{}MILLS\_\allowbreak{}1954} & 5 & 0 & {\scriptsize gauge\_\allowbreak{}geometry (31)} \\
{\scriptsize\ttfamily GJ\_\allowbreak{}1974\_\allowbreak{}SINGLE\_\allowbreak{}PHASE} & 4 & 2 & {\scriptsize gauge\_\allowbreak{}geometry (6), ground\_\allowbreak{}state (3), constructive (3), uniform\_\allowbreak{}limits (2), reconstruction (1)} \\
{\scriptsize\ttfamily GW\_\allowbreak{}1973} & 4 & 0 & {\scriptsize constructive (26), gauge\_\allowbreak{}geometry (26), ground\_\allowbreak{}state (3), uniform\_\allowbreak{}limits (1)} \\
{\scriptsize\ttfamily POLITZER\_\allowbreak{}1973} & 4 & 0 & {\scriptsize constructive (13), gauge\_\allowbreak{}geometry (6), ground\_\allowbreak{}state (4), uniform\_\allowbreak{}limits (1)} \\
{\scriptsize\ttfamily WEINBERG\_\allowbreak{}1967} & 3 & 0 & {\scriptsize gauge\_\allowbreak{}geometry (10), constructive (7), ground\_\allowbreak{}state (5)} \\
{\scriptsize\ttfamily WILSON\_\allowbreak{}1977\_\allowbreak{}LECTURES} & 2 & 0 & {\scriptsize gauge\_\allowbreak{}geometry (8), constructive (2), uniform\_\allowbreak{}limits (1)} \\
\end{longtable}

\noindent Searchable topics, with total line hits: \texttt{gauge\_\allowbreak{}geometry} (925), \texttt{constructive} (406), \texttt{uniform\_\allowbreak{}limits} (340), \texttt{curvature} (284), \texttt{ground\_\allowbreak{}state} (251), \texttt{reconstruction} (207), \texttt{spectral\_\allowbreak{}gap} (185).\par
% ---- end gen_extraction.tex ----------------------------------------
% ---- end B3_extraction.tex -----------------------------------------
% ---- begin B4_symbols.tex ------------------------------------------
\section{Symbols and join keys}
\label{app:symbols}

\SymbolCount{} registered symbols. Each carries a canonical name, the
spellings the source corpus uses, the names the code uses, the exact
values, and the claims it appears in.

This registry is what makes a quantity findable across a manuscript
written in one notation, code written in another, and a claim recorded
in a third. It is also where the corpus records \emph{forbidden} names
--- spellings that regenerate a dissolved contradiction every time
someone reads them, such as writing a single ``$m_4$'' for two
differently anchored coordinates (\S\ref{sec:setting}) --- and names
coined inside this program that the source corpus never uses, which
must not be searched for in the archive as though they were.

% ---- begin gen_symbols.tex -----------------------------------------
% Generated by paper/master/generate_sources.py.  Do not edit.
\noindent The join keys of this corpus are exact values and spellings, not concepts: no natural-language query retrieves $109151/249696$. The registry below is what makes a value findable across a document written in one notation, code written in another, and a claim recorded in a third.\par\medskip
\noindent\textbf{\texttt{A\_\allowbreak{}shp}}
 \hfill {\footnotesize values: \texttt{5/48}, \texttt{0.104166666666728}}
\par
{\small Sealed fourth-order shape coefficient 5/48. Both disputed kernels agree on it, which is what "sealed core" means.}\par
{\footnotesize corpus spellings: \texttt{5/48}, \texttt{A\_\allowbreak{}\{\textbackslash{}mathrm\{old\}\}} \quad in code: \texttt{A\_\allowbreak{}SHP\_\allowbreak{}3}, \texttt{A\_\allowbreak{}SHP\_\allowbreak{}3\_\allowbreak{}NUM} \quad claims: \texttt{R7}, \texttt{G14}}\par
\smallskip
\noindent\textbf{\texttt{beta\_\allowbreak{}3\textasciicircum{}bal}}
 \hfill {\footnotesize values: \texttt{15644916262153/34416487661400}}
\par
{\small The balanced (epsilon-free) continuation of the all-rank shape formula at N = 3: P17(9)/(3 R20(9)), equal to the pinned structured B\_\allowbreak{}N expression there. Differs from the historical beta\_\allowbreak{}3 by exactly 25/64 (the recorded C10 shift);}\par
{\footnotesize corpus spellings: \texttt{B$_3$\textasciicircum{}bal}, \texttt{q\_\allowbreak{}N\textasciicircum{}bal} \quad in code: \texttt{beta\_\allowbreak{}formula} \quad claims: \texttt{C2}, \texttt{C10}}\par
\smallskip
\noindent\textbf{\texttt{q\_\allowbreak{}a, e\_\allowbreak{}2, e\_\allowbreak{}3}}
\par
{\small Elementary symmetric functions of a\_\allowbreak{}i = 4 sin\textasciicircum{}2(k\_\allowbreak{}i/2). Every fourth-order shape coefficient is a symmetric function of these -- that is U2.}\par
{\footnotesize corpus spellings: \texttt{q\_\allowbreak{}a}, \texttt{e\_\allowbreak{}2}, \texttt{e\_\allowbreak{}3} \quad claims: \texttt{U2}, \texttt{G14}}\par
\smallskip
\noindent\textbf{\texttt{c\_\allowbreak{}4\textasciicircum{}adj}}
 \hfill {\footnotesize values: \texttt{-53/768}, \texttt{-31/1536}}
\par
{\small The cube-completion channel between two ADJACENT faces of a unit cube, -106/(N(N\textasciicircum{}2-1)\textasciicircum{}3) at every rank: the primitive simple-loop law's -88 plus ten multi-loop orderings at -18 (ADR 0024). In the kernel's own (0,2) basis it is the -53/768 the rotation record carries;}\par
{\footnotesize in code: \texttt{CUBE\_\allowbreak{}COMPLETION\_\allowbreak{}ADJACENT\_\allowbreak{}4}, \texttt{CUBE\_\allowbreak{}COMPLETION\_\allowbreak{}ADJACENT\_\allowbreak{}HISTORICAL\_\allowbreak{}4}, \texttt{cube\_\allowbreak{}completion\_\allowbreak{}adjacent} \quad claims: \texttt{C2}, \texttt{G3}, \texttt{G14}}\par
\smallskip
\noindent\textbf{\texttt{c\_\allowbreak{}\{2,prim\}\textasciicircum{}tet}}
 \hfill {\footnotesize values: \texttt{-1/3}}
\par
{\small Tetrahedral primitive local Haar-resolvent coefficient, the r = 2 row of the corpus's primitive completion law c\_\allowbreak{}\{r,prim\}. Asserted in the corpus with no artifact or reference SHA anywhere;}\par
{\footnotesize corpus spellings: \texttt{-8/[N(N\textasciicircum{}2-1)]}, \texttt{tetrahedral completion}, \texttt{c\_\allowbreak{}\{r,\textbackslash{}rm prim\}} \quad in code: \texttt{TETRA\_\allowbreak{}COMPLETION\_\allowbreak{}2\_\allowbreak{}SU3}, \texttt{c\_\allowbreak{}prim\_\allowbreak{}printed} \quad claims: \texttt{C15}, \texttt{G5}}\par
\smallskip
\noindent\textbf{\texttt{C\_\allowbreak{}shp}}
 \hfill {\footnotesize values: \texttt{-211835444920651/4405310420659200}, \texttt{-0.04808638318135875}, \texttt{-0.020213328886166577}}
\par
{\small Off-axis fourth-order shape coefficient. Resolved 2026-09-04 (ADR 0024): C\_\allowbreak{}historical + 25/1024.}\par
{\footnotesize corpus spellings: \texttt{C\_\allowbreak{}shape}, \texttt{C\_\allowbreak{}\{\textbackslash{}mathrm\{old\}\}}, \texttt{C\_\allowbreak{}\{\textbackslash{}mathrm\{new\}\}}, \texttt{C\textasciicircum{}\{(4)\}} \quad in code: \texttt{C\_\allowbreak{}SHP\_\allowbreak{}HISTORICAL}, \texttt{C\_\allowbreak{}SHP\_\allowbreak{}NEW\_\allowbreak{}NUM} \quad claims: \texttt{C2}, \texttt{R5}, \texttt{G3}}\par
\smallskip
\noindent\textbf{\texttt{d\_\allowbreak{}corner(N)}}
 \hfill {\footnotesize values: \texttt{2580244782961/398756546697600}, \texttt{-56022878647/4153714028100}}
\par
{\small The corner-cluster cumulant that decided C2, as a function of N (ADR 0025): (23, 30) rational functions in both sectors with denominator (N\textasciicircum{}2-1)\textasciicircum{}3(4N\textasciicircum{}2-9)\textasciicircum{}3(4N\textasciicircum{}2-1)(9N\textasciicircum{}2-25)(2N\textasciicircum{}2-1)\textasciicircum{}3(3N\textasciicircum{}2-1)(4N\textasciicircum{}2-5)(4N\textasciicircum{}2-2N-5)(4N\textasciicircum{}2+2N-5). At N = 3, in the kernel's (0,2) basis, +2580244782961/398756546697600 and -56022878647/4153714028100.}\par
{\footnotesize in code: \texttt{CORNER\_\allowbreak{}DRESSING\_\allowbreak{}ODD\_\allowbreak{}N}, \texttt{CORNER\_\allowbreak{}DRESSING\_\allowbreak{}EVEN\_\allowbreak{}N} \quad claims: \texttt{G14}, \texttt{C2}, \texttt{G3}}\par
\smallskip
\noindent\textbf{\texttt{Delta\_\allowbreak{}C}}
 \hfill {\footnotesize values: \texttt{0.027873054295192174}}
\par
{\small Off-axis shift C\_\allowbreak{}new - C\_\allowbreak{}old. Gamma-point data cannot constrain it, because the kernel it multiplies vanishes at Gamma.}\par
{\footnotesize corpus spellings: \texttt{Delta\_\allowbreak{}C}, \texttt{\textbackslash{}Delta\_\allowbreak{}C} \quad in code: \texttt{DELTA\_\allowbreak{}C\_\allowbreak{}NUM} \quad claims: \texttt{C2}, \texttt{G3}}\par
\smallskip
\noindent\textbf{\texttt{Delta\_\allowbreak{}Gamma}}
 \hfill {\footnotesize values: \texttt{2.0827701250956417}, \texttt{2.0827701250956414}}
\par
{\small Scalar shift m\_\allowbreak{}Gamma\textasciicircum{}(4) - q\_\allowbreak{}band\textasciicircum{}(4). Fixed by Gamma-point data.}\par
{\footnotesize corpus spellings: \texttt{\textbackslash{}Delta\_\allowbreak{}\textbackslash{}Gamma} \quad in code: \texttt{DELTA\_\allowbreak{}GAMMA\_\allowbreak{}NUM}, \texttt{DELTA\_\allowbreak{}GAMMA\_\allowbreak{}AS\_\allowbreak{}PRINTED\_\allowbreak{}NUM} \quad claims: \texttt{C1}, \texttt{R6}}\par
\smallskip
\noindent\textbf{\texttt{ell\_\allowbreak{}N}}
 \hfill {\footnotesize values: \texttt{-11/306}}
\par
{\small All-rank C-even second-order hopping. Printed in the corpus and gated in a notebook as factor(Wmix + Wlike + 1/CF), unregistered here until 2026-08-28.}\par
{\footnotesize corpus spellings: \texttt{ell\_\allowbreak{}N}, \texttt{-2*N*(3*N**2 - 5)} \quad in code: \texttt{even\_\allowbreak{}hopping}, \texttt{T\_\allowbreak{}PLUS\_\allowbreak{}2} \quad claims: \texttt{C13}, \texttt{R2}}\par
\smallskip
\noindent\textbf{\texttt{h\_\allowbreak{}4\textasciicircum{}side}}
 \hfill {\footnotesize values: \texttt{-2861009/84387303000}}
\par
{\small Isotropic pentagonal-prism fourth-order side coefficient, A\_\allowbreak{}+ - A\_\allowbreak{}-. Separate geometry;}\par
{\footnotesize corpus spellings: \texttt{h4\_\allowbreak{}side}, \texttt{h\_\allowbreak{}4\textasciicircum{}\{\textbackslash{}mathrm\{side\}\}} \quad in code: \texttt{H4\_\allowbreak{}SIDE} \quad claims: \texttt{R20}, \texttt{R21}, \texttt{G4}}\par
\smallskip
\noindent\textbf{\texttt{lambda\_\allowbreak{}X, lambda\_\allowbreak{}M, lambda\_\allowbreak{}R}}
\par
{\small Generalized Hodge pencil eigenvalues. lambda\_\allowbreak{}R = 2 lambda\_\allowbreak{}M - lambda\_\allowbreak{}X is the blind holdout used to test an extraction without fitting to it.}\par
{\footnotesize corpus spellings: \texttt{lambda\_\allowbreak{}R}, \texttt{\textbackslash{}lambda\_\allowbreak{}R} \quad claims: \texttt{G3}}\par
\smallskip
\noindent\textbf{\texttt{m\_\allowbreak{}6}}
\par
{\small Exact sixth-order coefficient. Belongs to the branch whose fourth-order scalar is q\_\allowbreak{}band\textasciicircum{}(4);}\par
{\footnotesize corpus spellings: \texttt{m\_\allowbreak{}6} \quad claims: \texttt{R15}, \texttt{G9}}\par
\smallskip
\noindent\textbf{\texttt{m\_\allowbreak{}Gamma\textasciicircum{}(4)}}
 \hfill {\footnotesize values: \texttt{-0.7751458630189173}}
\par
{\small Vacuum-subtracted physical Gamma-point fourth-order coefficient. Blind match to Hamer's a\_\allowbreak{}4 through m\_\allowbreak{}n = 2**(n-1) a\_\allowbreak{}n.}\par
{\footnotesize corpus spellings: \texttt{m\_\allowbreak{}\textbackslash{}Gamma\textasciicircum{}\{(4)\}} \quad in code: \texttt{M\_\allowbreak{}GAMMA\_\allowbreak{}4\_\allowbreak{}NUM} \quad claims: \texttt{C1}, \texttt{R4}}\par
\smallskip
\noindent\textbf{\texttt{Phi\_\allowbreak{}C}}
\par
{\small The crosswalk kernel 4 e\_\allowbreak{}2(k) / q\_\allowbreak{}a(k). Phi\_\allowbreak{}C(0) = 0, which is why the Gamma-point scalar places no constraint on Delta\_\allowbreak{}C -- the structural reason C2 cannot be closed by re-anchoring.}\par
{\footnotesize corpus spellings: \texttt{\textbackslash{}frac\{4e\_\allowbreak{}2\}\{q\_\allowbreak{}a\}}, \texttt{4e\_\allowbreak{}2/q} \quad in code: \texttt{phi\_\allowbreak{}c} \quad claims: \texttt{C2}, \texttt{G3}}\par
\smallskip
\noindent\textbf{\texttt{pi\textasciitilde{}}}
 \hfill {\footnotesize values: \texttt{-205125054864713/12642703205932800}}
\par
{\small The coplanar shared-link fourth-order amplitude with the two-hop shadow removed: pi\textasciitilde{} = pi + 2u. In the Hodge form of the kernel it multiplies S\_\allowbreak{}sq itself, and C\_\allowbreak{}shp = -5/96 - (rho + pi\textasciitilde{})/2 exactly, with no u.}\par
{\footnotesize in code: \texttt{PI\_\allowbreak{}REDUCED} \quad claims: \texttt{C2}, \texttt{G3}, \texttt{G14}}\par
\smallskip
\noindent\textbf{\texttt{q\_\allowbreak{}band\textasciicircum{}(4)}}
 \hfill {\footnotesize values: \texttt{-20721577909065127111/7250590288602460800}, \texttt{-2.857915988114559}}
\par
{\small Fourth-order rest scalar in the band-kernel anchoring. NOT a competing estimate of m\_\allowbreak{}Gamma\textasciicircum{}(4) -- a differently anchored coordinate.}\par
{\footnotesize corpus spellings: \texttt{q\_\allowbreak{}\{\textbackslash{}mathrm\{band\}\}\textasciicircum{}\{(4)\}}, \texttt{q\_\allowbreak{}\{\textbackslash{}mathrm\{old\}\}\textasciicircum{}\{(4)\}}, \texttt{q\_\allowbreak{}old} \quad in code: \texttt{Q\_\allowbreak{}BAND\_\allowbreak{}4} \quad claims: \texttt{C1}, \texttt{R4}}\par
\smallskip
\noindent\textbf{\texttt{rho\_\allowbreak{}assembled}}
 \hfill {\footnotesize values: \texttt{-588708011765248393/14501180577204921600}, \texttt{-13035490122347/550663802582400}}
\par
{\small The cluster-assembled rotation amplitude read in the kernel's own basis, the x-then-z traversal of the (0,2) face in which all 24 cross-plane records are rho times the L\_\allowbreak{}down incidence: minus the 2026-09-02 run's RHO\_\allowbreak{}CLUSTER. It differs from RHO\_\allowbreak{}ORBIT by exactly 25/512, the adjacent-face cube completion, and gives C\_\allowbreak{}shp = C\_\allowbreak{}historical + 25/1024 (ADR 0024).}\par
{\footnotesize in code: \texttt{RHO\_\allowbreak{}ASSEMBLED}, \texttt{C\_\allowbreak{}SHP\_\allowbreak{}ASSEMBLED} \quad claims: \texttt{C2}, \texttt{G3}}\par
\smallskip
\noindent\textbf{\texttt{w\_\allowbreak{}R}}
\par
{\small The resolvent weight of one shared-link fusion channel, -(d\_\allowbreak{}R/N\textasciicircum{}2)/(C\_\allowbreak{}F + C\_\allowbreak{}R/2), summing to A\_\allowbreak{}N over F (x) Fbar and B\_\allowbreak{}N over F (x) F. The numerator d\_\allowbreak{}R/N\textasciicircum{}2 is the manuscripts' "isotropy" sentence: asserted there, and derived here from the order-2 Weingarten values -- the six nonshared links collapse to two independent Haar matrices, and two degree-(2,2) moments (N\textasciicircum{}2 and N) settle both families.}\par
{\footnotesize corpus spellings: \texttt{w\_\allowbreak{}R}, \texttt{W\_\allowbreak{}mix}, \texttt{Wlike} \quad in code: \texttt{channel\_\allowbreak{}weight}, \texttt{channel\_\allowbreak{}data}, \texttt{shared\_\allowbreak{}link\_\allowbreak{}weights} \quad claims: \texttt{G24}, \texttt{R2}}\par
\smallskip
\noindent\textbf{\texttt{sigma\textasciitilde{}}}
 \hfill {\footnotesize values: \texttt{-1249434547013129251/483372685906830720}}
\par
{\small The on-site fourth-order amplitude with the two-hop shadow removed: sigma\textasciitilde{} = sigma - 12u. It is the constant term of the Hodge form and enters the carrier band through c\_\allowbreak{}0 alone.}\par
{\footnotesize in code: \texttt{SIGMA\_\allowbreak{}REDUCED} \quad claims: \texttt{C2}}\par
\smallskip
\noindent\textbf{\texttt{sigma(u)}}
 \hfill {\footnotesize values: \texttt{-61/408}, \texttt{-22/153}, \texttt{-137767222189182735950309/2009803206414863779920000}}
\par
{\small Native string tension through fifth order. Odd physical coefficients carry (-1)**n relative to the raw unit-insertion magnitudes.}\par
{\footnotesize corpus spellings: \texttt{\textbackslash{}sigma\_\allowbreak{}n}, \texttt{sigma\_\allowbreak{}5} \quad in code: \texttt{SIGMA\_\allowbreak{}0}, \texttt{SIGMA\_\allowbreak{}2}, \texttt{SIGMA\_\allowbreak{}3}, \texttt{SIGMA\_\allowbreak{}4}, \texttt{SIGMA\_\allowbreak{}5}, \texttt{SIGMA\_\allowbreak{}5\_\allowbreak{}RAW} \quad claims: \texttt{C5}, \texttt{R14}, \texttt{G7}}\par
\smallskip
\noindent\textbf{\texttt{d\_\allowbreak{}single(N)}}
 \hfill {\footnotesize values: \texttt{-385/1997568}, \texttt{-121/249696}}
\par
{\small The fourth-order dressing of a shared-link pair by a plaquette touching one face only, as a function of N (ADR 0025): C-odd 2N\textasciicircum{}3(N\textasciicircum{}2-4)(10N\textasciicircum{}2-13)/((N\textasciicircum{}2-1)\textasciicircum{}3(4N\textasciicircum{}2-9)\textasciicircum{}2(2N\textasciicircum{}2-1)\textasciicircum{}2), C-even -ell\_\allowbreak{}N\textasciicircum{}2/(2 C\_\allowbreak{}F). Reconstructed from the third engine at N = 3..70 and verified on held-out ranks (ADR 0025), then computed outright over the field Q(N) (ADR 0029): the same rational function, derived.}\par
{\footnotesize in code: \texttt{SINGLE\_\allowbreak{}CONTACT\_\allowbreak{}DRESSING\_\allowbreak{}ODD\_\allowbreak{}N}, \texttt{SINGLE\_\allowbreak{}CONTACT\_\allowbreak{}DRESSING\_\allowbreak{}EVEN\_\allowbreak{}N} \quad claims: \texttt{G14}, \texttt{C2}}\par
\smallskip
\noindent\textbf{\texttt{t\_\allowbreak{}N}}
 \hfill {\footnotesize values: \texttt{5/612}}
\par
{\small All-rank second-order hopping, B\_\allowbreak{}N - A\_\allowbreak{}N. t\_\allowbreak{}2 = 0, t\_\allowbreak{}3 = 5/612.}\par
{\footnotesize corpus spellings: \texttt{t\_\allowbreak{}N}, \texttt{t\_\allowbreak{}-} \quad in code: \texttt{T\_\allowbreak{}MINUS\_\allowbreak{}2}, \texttt{hopping} \quad claims: \texttt{R2}}\par
\smallskip
\noindent\textbf{\texttt{tau\_\allowbreak{}4}}
 \hfill {\footnotesize values: \texttt{-2861009/16877460600}}
\par
{\small Pentagonal cap hopping coefficient, h\_\allowbreak{}4\textasciicircum{}side times 5 by exact D\_\allowbreak{}5 covariance. The symbol is 2 tau\_\allowbreak{}4 cos k.}\par
{\footnotesize corpus spellings: \texttt{\textbackslash{}tau\_\allowbreak{}4}, \texttt{tau\_\allowbreak{}4} \quad in code: \texttt{TAU\_\allowbreak{}4} \quad claims: \texttt{R21}, \texttt{G4}}\par
\smallskip
\noindent\textbf{\texttt{u\_\allowbreak{}2hop(N)}}
 \hfill {\footnotesize values: \texttt{360421351/40327601932800}, \texttt{948253471/40327601932800}}
\par
{\small The two-hop chain cumulant, the one dynamical input of the Hodge form (ADR 0019), as a function of N (ADR 0025): (21, 28) rational functions in both sectors with denominator 2(N\textasciicircum{}2-1)\textasciicircum{}3(4N\textasciicircum{}2-9)\textasciicircum{}3(9N\textasciicircum{}2-16)(2N\textasciicircum{}2-3)(2N\textasciicircum{}2-1)\textasciicircum{}3(3N\textasciicircum{}2-2)(4N\textasciicircum{}2-N-4)(4N\textasciicircum{}2+N-4); the C-odd numerator carries (N\textasciicircum{}2-4)\textasciicircum{}2.}\par
{\footnotesize in code: \texttt{TWO\_\allowbreak{}HOP\_\allowbreak{}WEIGHT\_\allowbreak{}ODD\_\allowbreak{}N}, \texttt{TWO\_\allowbreak{}HOP\_\allowbreak{}WEIGHT\_\allowbreak{}EVEN\_\allowbreak{}N} \quad claims: \texttt{G14}, \texttt{C2}}\par
\smallskip
\noindent\textbf{\texttt{u\_\allowbreak{}even}}
 \hfill {\footnotesize values: \texttt{948253471/40327601932800}}
\par
{\small The C-even fourth-order chain amplitude: the cumulant W(P,Q,R) - W(P,R) on the P -> R element in the C-even sector, for Q sharing one link with each of the link-disjoint P and R. The same rational on all three chain types (coplanar-coplanar, coplanar-perpendicular, perpendicular-perpendicular), with no incidence sign because the C-even shared-link hop -11/306 carries none.}\par
{\footnotesize in code: \texttt{U\_\allowbreak{}CHAIN\_\allowbreak{}C\_\allowbreak{}EVEN} \quad claims: \texttt{G3}, \texttt{G14}, \texttt{C2}}\par
\smallskip
\noindent\textbf{\texttt{W\_\allowbreak{}4}}
 \hfill {\footnotesize values: \texttt{0.48061786909826}, \texttt{0.9265867378213348}}
\par
{\small Fourth-order bandwidth. Enters the near-Gamma crossover constant only through sqrt(W\_\allowbreak{}4), which is why the G11 criterion survives C2.}\par
{\footnotesize corpus spellings: \texttt{W\_\allowbreak{}4} \quad in code: \texttt{W4\_\allowbreak{}HISTORICAL\_\allowbreak{}NUM}, \texttt{W4\_\allowbreak{}NEW\_\allowbreak{}NUM} \quad claims: \texttt{C2}, \texttt{G11}}\par
\smallskip
\noindent\textbf{\texttt{x (degree-3 record quantum)}}
 \hfill {\footnotesize values: \texttt{360421351/40327601932800}}
\par
{\small The single amplitude quantizing every degree-3 channel coefficient of the historical 189-record kernel: each block's q e\_\allowbreak{}2 and e\_\allowbreak{}3 amplitude is an integer multiple of x, itself a raw record weight, and the tier collapse is (+1,+1,-2)x and (-3,+6,-3)x exactly (suite "off-axis channel ledger"). Coined by the maintainer's off-axis ledger (WORK\_\allowbreak{}SINCE\_\allowbreak{}2026-08);}\par
{\footnotesize in code: \texttt{X\_\allowbreak{}QUANTUM} \quad claims: \texttt{C2}, \texttt{G14}}\par
\smallskip
% ---- end gen_symbols.tex -------------------------------------------
% ---- end B4_symbols.tex --------------------------------------------
% ---- begin B5_documents.tex ----------------------------------------
\section{Documents and their standing}
\label{app:documents}

\AliasCount{} source documents with a registered alias and a standing.

Standing is the field that matters most when reading the archive,
because a superseded document is not marked in its own text: it reads
exactly like a current one, states its conclusions with the same
confidence, and was correct when written. The register is the only
thing that distinguishes them.

The archive's governing instruction is that nothing is erased, which
makes this register load-bearing rather than administrative. Preserved
material includes failed attempts, superseded drafts and duplicates,
all deliberately --- deleting a refuted claim destroys the evidence
that it was tried, which is the same principle behind the dead-route
register in \S\ref{sec:obligations}.

% ---- begin gen_documents.tex ---------------------------------------
% Generated by paper/master/generate_sources.py.  Do not edit.
\noindent 248 source documents carry a registered alias and a \emph{standing}. Standing is the single most important field for an agent reading the archive: it says whether a document is current, superseded, or preserved only as evidence of a stage. A superseded document is not deleted and reads exactly like a current one.\par\medskip
\subsection*{Standing: repo (159)}
\addcontentsline{toc}{subsection}{Standing: repo}
\begin{longtable}{@{}p{2.0in} p{4.2in}@{}}
\toprule Alias & Path \\ \midrule
\endfirsthead
\toprule Alias & Path \\ \midrule \endhead
\bottomrule\endfoot
{\scriptsize\ttfamily ARCHIVE\_\allowbreak{}DERIVATION\_\allowbreak{}CHAIN} & {\scriptsize\ttfamily paper/\allowbreak{}research\_\allowbreak{}notes/\allowbreak{}ARCHIVE\_\allowbreak{}DERIVATION\_\allowbreak{}CHAIN\_\allowbreak{}20260911.md} \\
{\scriptsize\ttfamily ARCHIVE\_\allowbreak{}FUNCTIONAL} & {\scriptsize\ttfamily docs/\allowbreak{}derivations/\allowbreak{}archive-\allowbreak{}functional-\allowbreak{}inequalities.md} \\
{\scriptsize\ttfamily ARCHIVE\_\allowbreak{}REVIEW\_\allowbreak{}SOURCES} & {\scriptsize\ttfamily runs/\allowbreak{}archive\_\allowbreak{}derivations\_\allowbreak{}2026-\allowbreak{}09-\allowbreak{}11/\allowbreak{}README.md} \\
{\scriptsize\ttfamily ARCHIVE\_\allowbreak{}SOURCE\_\allowbreak{}TRANSPORT} & {\scriptsize\ttfamily docs/\allowbreak{}derivations/\allowbreak{}archive-\allowbreak{}source-\allowbreak{}transport.md} \\
{\scriptsize\ttfamily CONDITIONAL\_\allowbreak{}GRADIENT\_\allowbreak{}REPAIR} & {\scriptsize\ttfamily paper/\allowbreak{}research\_\allowbreak{}notes/\allowbreak{}G19\_\allowbreak{}CONDITIONAL\_\allowbreak{}GRADIENT\_\allowbreak{}REPAIR\_\allowbreak{}20260905.md} \\
{\scriptsize\ttfamily CONDITIONAL\_\allowbreak{}QUANTUM\_\allowbreak{}PATH\_\allowbreak{}COVARIANCE} & {\scriptsize\ttfamily paper/\allowbreak{}research\_\allowbreak{}notes/\allowbreak{}G19\_\allowbreak{}CONDITIONAL\_\allowbreak{}QUANTUM\_\allowbreak{}PATH\_\allowbreak{}COVARIANCE\_\allowbreak{}20260906.md} \\
{\scriptsize\ttfamily COUPLED\_\allowbreak{}OSCILLATOR\_\allowbreak{}METHOD\_\allowbreak{}REVIEW} & {\scriptsize\ttfamily paper/\allowbreak{}research\_\allowbreak{}notes/\allowbreak{}G19\_\allowbreak{}COUPLED\_\allowbreak{}OSCILLATOR\_\allowbreak{}METHOD\_\allowbreak{}REVIEW\_\allowbreak{}20260905.md} \\
{\scriptsize\ttfamily CUBIC\_\allowbreak{}GROUND\_\allowbreak{}TRANSFER} & {\scriptsize\ttfamily paper/\allowbreak{}research\_\allowbreak{}notes/\allowbreak{}G19\_\allowbreak{}CUBIC\_\allowbreak{}GROUND\_\allowbreak{}TRANSFER\_\allowbreak{}20260906.md} \\
{\scriptsize\ttfamily DYNAMIC\_\allowbreak{}CONDITIONAL\_\allowbreak{}CUBIC\_\allowbreak{}ENERGY} & {\scriptsize\ttfamily paper/\allowbreak{}research\_\allowbreak{}notes/\allowbreak{}G19\_\allowbreak{}DYNAMIC\_\allowbreak{}FIBER\_\allowbreak{}COVARIANCE\_\allowbreak{}AND\_\allowbreak{}CUBIC\_\allowbreak{}ENERGY\_\allowbreak{}20260906.md} \\
{\scriptsize\ttfamily EARLIER\_\allowbreak{}PROOF\_\allowbreak{}RECOVERY} & {\scriptsize\ttfamily docs/\allowbreak{}derivations/\allowbreak{}earlier-\allowbreak{}proof-\allowbreak{}recovery.md} \\
{\scriptsize\ttfamily EARLIER\_\allowbreak{}PROOF\_\allowbreak{}RECOVERY\_\allowbreak{}RESULT} & {\scriptsize\ttfamily paper/\allowbreak{}research\_\allowbreak{}notes/\allowbreak{}EARLIER\_\allowbreak{}PROOF\_\allowbreak{}RECOVERY\_\allowbreak{}20260911.md} \\
{\scriptsize\ttfamily EARLIER\_\allowbreak{}RECOVERY\_\allowbreak{}SOURCE\_\allowbreak{}01} & {\scriptsize\ttfamily runs/\allowbreak{}earlier\_\allowbreak{}proof\_\allowbreak{}recovery\_\allowbreak{}2026-\allowbreak{}09-\allowbreak{}11/\allowbreak{}sources/\allowbreak{}05\_\allowbreak{}affine\_\allowbreak{}laplacian\_\allowbreak{}law\_\allowbreak{}analytic\_\allowbreak{}proof.md} \\
{\scriptsize\ttfamily EARLIER\_\allowbreak{}RECOVERY\_\allowbreak{}SOURCE\_\allowbreak{}02} & {\scriptsize\ttfamily runs/\allowbreak{}earlier\_\allowbreak{}proof\_\allowbreak{}recovery\_\allowbreak{}2026-\allowbreak{}09-\allowbreak{}11/\allowbreak{}sources/\allowbreak{}WORKHOUSE\_\allowbreak{}CARRIER\_\allowbreak{}TO\_\allowbreak{}PARTICLE\_\allowbreak{}PROOF\_\allowbreak{}DOSSIER.md} \\
{\scriptsize\ttfamily EARLIER\_\allowbreak{}RECOVERY\_\allowbreak{}SOURCE\_\allowbreak{}03} & {\scriptsize\ttfamily runs/\allowbreak{}earlier\_\allowbreak{}proof\_\allowbreak{}recovery\_\allowbreak{}2026-\allowbreak{}09-\allowbreak{}11/\allowbreak{}sources/\allowbreak{}LYAPUNOV\_\allowbreak{}08\_\allowbreak{}transversality\_\allowbreak{}rank3\_\allowbreak{}and\_\allowbreak{}binomial\_\allowbreak{}tail\_\allowbreak{}drift.md} \\
{\scriptsize\ttfamily EARLIER\_\allowbreak{}RECOVERY\_\allowbreak{}SOURCE\_\allowbreak{}04} & {\scriptsize\ttfamily runs/\allowbreak{}earlier\_\allowbreak{}proof\_\allowbreak{}recovery\_\allowbreak{}2026-\allowbreak{}09-\allowbreak{}11/\allowbreak{}sources/\allowbreak{}ENGINE\_\allowbreak{}FLUX\_\allowbreak{}singular\_\allowbreak{}geometry\_\allowbreak{}derivations.py} \\
{\scriptsize\ttfamily EARLIER\_\allowbreak{}RECOVERY\_\allowbreak{}SOURCE\_\allowbreak{}05} & {\scriptsize\ttfamily runs/\allowbreak{}earlier\_\allowbreak{}proof\_\allowbreak{}recovery\_\allowbreak{}2026-\allowbreak{}09-\allowbreak{}11/\allowbreak{}sources/\allowbreak{}LYAPUNOV\_\allowbreak{}09\_\allowbreak{}appendix\_\allowbreak{}disjoint\_\allowbreak{}staple\_\allowbreak{}coordinates\_\allowbreak{}d4.md} \\
{\scriptsize\ttfamily EARLIER\_\allowbreak{}RECOVERY\_\allowbreak{}SOURCE\_\allowbreak{}06} & {\scriptsize\ttfamily runs/\allowbreak{}earlier\_\allowbreak{}proof\_\allowbreak{}recovery\_\allowbreak{}2026-\allowbreak{}09-\allowbreak{}11/\allowbreak{}sources/\allowbreak{}NOTE\_\allowbreak{}FLUX\_\allowbreak{}singular\_\allowbreak{}geometry\_\allowbreak{}derivations\_\allowbreak{}2026-\allowbreak{}08-\allowbreak{}28.md} \\
{\scriptsize\ttfamily EARLIER\_\allowbreak{}RECOVERY\_\allowbreak{}SOURCE\_\allowbreak{}07} & {\scriptsize\ttfamily runs/\allowbreak{}earlier\_\allowbreak{}proof\_\allowbreak{}recovery\_\allowbreak{}2026-\allowbreak{}09-\allowbreak{}11/\allowbreak{}sources/\allowbreak{}VSU\_\allowbreak{}Convex\_\allowbreak{}Poisson\_\allowbreak{}WellPosedness.md} \\
{\scriptsize\ttfamily EARLIER\_\allowbreak{}RECOVERY\_\allowbreak{}SOURCE\_\allowbreak{}08} & {\scriptsize\ttfamily runs/\allowbreak{}earlier\_\allowbreak{}proof\_\allowbreak{}recovery\_\allowbreak{}2026-\allowbreak{}09-\allowbreak{}11/\allowbreak{}sources/\allowbreak{}b6\_\allowbreak{}shell\_\allowbreak{}block\_\allowbreak{}krylov\_\allowbreak{}reduction\_\allowbreak{}certificate.json} \\
{\scriptsize\ttfamily EARLIER\_\allowbreak{}RECOVERY\_\allowbreak{}SOURCE\_\allowbreak{}09} & {\scriptsize\ttfamily runs/\allowbreak{}earlier\_\allowbreak{}proof\_\allowbreak{}recovery\_\allowbreak{}2026-\allowbreak{}09-\allowbreak{}11/\allowbreak{}sources/\allowbreak{}01\_\allowbreak{}conditional\_\allowbreak{}spectral\_\allowbreak{}floor\_\allowbreak{}monotonicity.md} \\
{\scriptsize\ttfamily EARLIER\_\allowbreak{}RECOVERY\_\allowbreak{}SOURCE\_\allowbreak{}10} & {\scriptsize\ttfamily runs/\allowbreak{}earlier\_\allowbreak{}proof\_\allowbreak{}recovery\_\allowbreak{}2026-\allowbreak{}09-\allowbreak{}11/\allowbreak{}sources/\allowbreak{}FINITE\_\allowbreak{}ORDER\_\allowbreak{}NESTED\_\allowbreak{}QUOTIENT\_\allowbreak{}SPECTRAL\_\allowbreak{}REDUCTION\_\allowbreak{}THEOREM\_\allowbreak{}FULL\_\allowbreak{}DERIVATION\_\allowbreak{}2026-\allowbreak{}08-\allowbreak{}28.md} \\
{\scriptsize\ttfamily EARLIER\_\allowbreak{}RECOVERY\_\allowbreak{}SOURCE\_\allowbreak{}11} & {\scriptsize\ttfamily runs/\allowbreak{}earlier\_\allowbreak{}proof\_\allowbreak{}recovery\_\allowbreak{}2026-\allowbreak{}09-\allowbreak{}11/\allowbreak{}sources/\allowbreak{}LATTICE\_\allowbreak{}QCD\_\allowbreak{}phase\_\allowbreak{}isolation\_\allowbreak{}principle.md} \\
{\scriptsize\ttfamily EARLIER\_\allowbreak{}RECOVERY\_\allowbreak{}SOURCE\_\allowbreak{}12} & {\scriptsize\ttfamily runs/\allowbreak{}earlier\_\allowbreak{}proof\_\allowbreak{}recovery\_\allowbreak{}2026-\allowbreak{}09-\allowbreak{}11/\allowbreak{}sources/\allowbreak{}01\_\allowbreak{}determinant\_\allowbreak{}reduction\_\allowbreak{}theorem.md} \\
{\scriptsize\ttfamily EARLIER\_\allowbreak{}RECOVERY\_\allowbreak{}SOURCE\_\allowbreak{}13} & {\scriptsize\ttfamily runs/\allowbreak{}earlier\_\allowbreak{}proof\_\allowbreak{}recovery\_\allowbreak{}2026-\allowbreak{}09-\allowbreak{}11/\allowbreak{}sources/\allowbreak{}cleanroom\_\allowbreak{}b6\_\allowbreak{}symbolic\_\allowbreak{}certificate.json} \\
{\scriptsize\ttfamily EARLIER\_\allowbreak{}RECOVERY\_\allowbreak{}SOURCE\_\allowbreak{}14} & {\scriptsize\ttfamily runs/\allowbreak{}earlier\_\allowbreak{}proof\_\allowbreak{}recovery\_\allowbreak{}2026-\allowbreak{}09-\allowbreak{}11/\allowbreak{}sources/\allowbreak{}B6\_\allowbreak{}SHELL\_\allowbreak{}BLOCK\_\allowbreak{}KRYLOV\_\allowbreak{}REDUCTION\_\allowbreak{}2026-\allowbreak{}08-\allowbreak{}28.md} \\
{\scriptsize\ttfamily EARLIER\_\allowbreak{}RECOVERY\_\allowbreak{}SOURCE\_\allowbreak{}15} & {\scriptsize\ttfamily runs/\allowbreak{}earlier\_\allowbreak{}proof\_\allowbreak{}recovery\_\allowbreak{}2026-\allowbreak{}09-\allowbreak{}11/\allowbreak{}sources/\allowbreak{}NOTE\_\allowbreak{}G17\_\allowbreak{}SU3\_\allowbreak{}COVARIANCE\_\allowbreak{}ATTACK\_\allowbreak{}2026-\allowbreak{}08-\allowbreak{}30.md} \\
{\scriptsize\ttfamily EARLIER\_\allowbreak{}RECOVERY\_\allowbreak{}SOURCE\_\allowbreak{}16} & {\scriptsize\ttfamily runs/\allowbreak{}earlier\_\allowbreak{}proof\_\allowbreak{}recovery\_\allowbreak{}2026-\allowbreak{}09-\allowbreak{}11/\allowbreak{}sources/\allowbreak{}NOTE\_\allowbreak{}G19\_\allowbreak{}balaban\_\allowbreak{}blocking\_\allowbreak{}reflection\_\allowbreak{}positivity\_\allowbreak{}2026-\allowbreak{}08-\allowbreak{}30.md} \\
{\scriptsize\ttfamily EARLIER\_\allowbreak{}RECOVERY\_\allowbreak{}SOURCE\_\allowbreak{}17} & {\scriptsize\ttfamily runs/\allowbreak{}earlier\_\allowbreak{}proof\_\allowbreak{}recovery\_\allowbreak{}2026-\allowbreak{}09-\allowbreak{}11/\allowbreak{}sources/\allowbreak{}verify\_\allowbreak{}balaban\_\allowbreak{}corner\_\allowbreak{}reflection\_\allowbreak{}exact.py} \\
{\scriptsize\ttfamily FESHBACH\_\allowbreak{}RESOLVENT\_\allowbreak{}COMPARISON} & {\scriptsize\ttfamily paper/\allowbreak{}research\_\allowbreak{}notes/\allowbreak{}G22\_\allowbreak{}FESHBACH\_\allowbreak{}RESOLVENT\_\allowbreak{}COMPARISON\_\allowbreak{}20260909.md} \\
{\scriptsize\ttfamily FLAT\_\allowbreak{}HOLONOMY\_\allowbreak{}SOURCES\_\allowbreak{}AND\_\allowbreak{}RANK\_\allowbreak{}REPAIR} & {\scriptsize\ttfamily paper/\allowbreak{}research\_\allowbreak{}notes/\allowbreak{}G19\_\allowbreak{}FLAT\_\allowbreak{}HOLONOMY\_\allowbreak{}SOURCES\_\allowbreak{}AND\_\allowbreak{}RANK\_\allowbreak{}REPAIR\_\allowbreak{}20260907.md} \\
{\scriptsize\ttfamily FOLDER\_\allowbreak{}VALIDATION\_\allowbreak{}PBH\_\allowbreak{}WILSON\_\allowbreak{}FOCUSED\_\allowbreak{}2026\_\allowbreak{}09\_\allowbreak{}08\_\allowbreak{}XML} & {\scriptsize\ttfamily docs/\allowbreak{}validation/\allowbreak{}pbh-\allowbreak{}wilson-\allowbreak{}focused-\allowbreak{}2026-\allowbreak{}09-\allowbreak{}08.xml} \\
{\scriptsize\ttfamily FOLDER\_\allowbreak{}VALIDATION\_\allowbreak{}PBH\_\allowbreak{}WILSON\_\allowbreak{}GRAPH\_\allowbreak{}2026\_\allowbreak{}09\_\allowbreak{}08\_\allowbreak{}XML} & {\scriptsize\ttfamily docs/\allowbreak{}validation/\allowbreak{}pbh-\allowbreak{}wilson-\allowbreak{}graph-\allowbreak{}2026-\allowbreak{}09-\allowbreak{}08.xml} \\
{\scriptsize\ttfamily FOLDER\_\allowbreak{}VALIDATION\_\allowbreak{}WILSON\_\allowbreak{}BACKGROUND\_\allowbreak{}GRAPH\_\allowbreak{}XML} & {\scriptsize\ttfamily docs/\allowbreak{}validation/\allowbreak{}wilson-\allowbreak{}background-\allowbreak{}graph.xml} \\
{\scriptsize\ttfamily FOLDER\_\allowbreak{}VALIDATION\_\allowbreak{}WILSON\_\allowbreak{}BACKGROUND\_\allowbreak{}INVARIANTS\_\allowbreak{}XML} & {\scriptsize\ttfamily docs/\allowbreak{}validation/\allowbreak{}wilson-\allowbreak{}background-\allowbreak{}invariants.xml} \\
{\scriptsize\ttfamily FOLDER\_\allowbreak{}VALIDATION\_\allowbreak{}WILSON\_\allowbreak{}BACKGROUND\_\allowbreak{}LEDGER\_\allowbreak{}XML} & {\scriptsize\ttfamily docs/\allowbreak{}validation/\allowbreak{}wilson-\allowbreak{}background-\allowbreak{}ledger.xml} \\
{\scriptsize\ttfamily FOLDER\_\allowbreak{}VALIDATION\_\allowbreak{}WILSON\_\allowbreak{}BLOCK\_\allowbreak{}FRESHNESS\_\allowbreak{}XML} & {\scriptsize\ttfamily docs/\allowbreak{}validation/\allowbreak{}wilson-\allowbreak{}block-\allowbreak{}freshness.xml} \\
{\scriptsize\ttfamily FOLDER\_\allowbreak{}VALIDATION\_\allowbreak{}WILSON\_\allowbreak{}BLOCK\_\allowbreak{}GRAPH\_\allowbreak{}XML} & {\scriptsize\ttfamily docs/\allowbreak{}validation/\allowbreak{}wilson-\allowbreak{}block-\allowbreak{}graph.xml} \\
{\scriptsize\ttfamily FOLDER\_\allowbreak{}VALIDATION\_\allowbreak{}WILSON\_\allowbreak{}BLOCK\_\allowbreak{}INVARIANTS\_\allowbreak{}XML} & {\scriptsize\ttfamily docs/\allowbreak{}validation/\allowbreak{}wilson-\allowbreak{}block-\allowbreak{}invariants.xml} \\
{\scriptsize\ttfamily FOLDER\_\allowbreak{}VALIDATION\_\allowbreak{}WILSON\_\allowbreak{}BLOCK\_\allowbreak{}LEDGER\_\allowbreak{}XML} & {\scriptsize\ttfamily docs/\allowbreak{}validation/\allowbreak{}wilson-\allowbreak{}block-\allowbreak{}ledger.xml} \\
{\scriptsize\ttfamily FOLDER\_\allowbreak{}VALIDATION\_\allowbreak{}WILSON\_\allowbreak{}MARKED\_\allowbreak{}AXIOMS\_\allowbreak{}2026\_\allowbreak{}09\_\allowbreak{}08\_\allowbreak{}TXT} & {\scriptsize\ttfamily docs/\allowbreak{}validation/\allowbreak{}wilson-\allowbreak{}marked-\allowbreak{}axioms-\allowbreak{}2026-\allowbreak{}09-\allowbreak{}08.txt} \\
{\scriptsize\ttfamily FOLDER\_\allowbreak{}VALIDATION\_\allowbreak{}WILSON\_\allowbreak{}MARKED\_\allowbreak{}TESTS\_\allowbreak{}2026\_\allowbreak{}09\_\allowbreak{}08\_\allowbreak{}XML} & {\scriptsize\ttfamily docs/\allowbreak{}validation/\allowbreak{}wilson-\allowbreak{}marked-\allowbreak{}tests-\allowbreak{}2026-\allowbreak{}09-\allowbreak{}08.xml} \\
{\scriptsize\ttfamily FOLDER\_\allowbreak{}VALIDATION\_\allowbreak{}WILSON\_\allowbreak{}SC17\_\allowbreak{}GRAPH\_\allowbreak{}XML} & {\scriptsize\ttfamily docs/\allowbreak{}validation/\allowbreak{}wilson-\allowbreak{}sc17-\allowbreak{}graph.xml} \\
{\scriptsize\ttfamily FOLDER\_\allowbreak{}VALIDATION\_\allowbreak{}WILSON\_\allowbreak{}SC17\_\allowbreak{}INVARIANTS\_\allowbreak{}XML} & {\scriptsize\ttfamily docs/\allowbreak{}validation/\allowbreak{}wilson-\allowbreak{}sc17-\allowbreak{}invariants.xml} \\
{\scriptsize\ttfamily FOLDER\_\allowbreak{}VALIDATION\_\allowbreak{}WILSON\_\allowbreak{}SC17\_\allowbreak{}LEAN\_\allowbreak{}AXIOMS\_\allowbreak{}TXT} & {\scriptsize\ttfamily docs/\allowbreak{}validation/\allowbreak{}wilson-\allowbreak{}sc17-\allowbreak{}lean-\allowbreak{}axioms.txt} \\
{\scriptsize\ttfamily FOLDER\_\allowbreak{}VALIDATION\_\allowbreak{}WILSON\_\allowbreak{}SC17\_\allowbreak{}LEAN\_\allowbreak{}BUILD\_\allowbreak{}TXT} & {\scriptsize\ttfamily docs/\allowbreak{}validation/\allowbreak{}wilson-\allowbreak{}sc17-\allowbreak{}lean-\allowbreak{}build.txt} \\
{\scriptsize\ttfamily FOLDER\_\allowbreak{}VALIDATION\_\allowbreak{}WILSON\_\allowbreak{}SC17\_\allowbreak{}LEDGER\_\allowbreak{}XML} & {\scriptsize\ttfamily docs/\allowbreak{}validation/\allowbreak{}wilson-\allowbreak{}sc17-\allowbreak{}ledger.xml} \\
{\scriptsize\ttfamily FOLDER\_\allowbreak{}VALIDATION\_\allowbreak{}WILSON\_\allowbreak{}SELECTED\_\allowbreak{}FOCUSED\_\allowbreak{}2026\_\allowbreak{}09\_\allowbreak{}08\_\allowbreak{}XML} & {\scriptsize\ttfamily docs/\allowbreak{}validation/\allowbreak{}wilson-\allowbreak{}selected-\allowbreak{}focused-\allowbreak{}2026-\allowbreak{}09-\allowbreak{}08.xml} \\
{\scriptsize\ttfamily FOLDER\_\allowbreak{}VALIDATION\_\allowbreak{}WILSON\_\allowbreak{}SELECTED\_\allowbreak{}GRAPH\_\allowbreak{}2026\_\allowbreak{}09\_\allowbreak{}08\_\allowbreak{}XML} & {\scriptsize\ttfamily docs/\allowbreak{}validation/\allowbreak{}wilson-\allowbreak{}selected-\allowbreak{}graph-\allowbreak{}2026-\allowbreak{}09-\allowbreak{}08.xml} \\
{\scriptsize\ttfamily FOLDER\_\allowbreak{}VALIDATION\_\allowbreak{}WILSON\_\allowbreak{}SHELL\_\allowbreak{}AXIOMS\_\allowbreak{}2026\_\allowbreak{}09\_\allowbreak{}08\_\allowbreak{}TXT} & {\scriptsize\ttfamily docs/\allowbreak{}validation/\allowbreak{}wilson-\allowbreak{}shell-\allowbreak{}axioms-\allowbreak{}2026-\allowbreak{}09-\allowbreak{}08.txt} \\
{\scriptsize\ttfamily FOLDER\_\allowbreak{}VALIDATION\_\allowbreak{}WILSON\_\allowbreak{}SHELL\_\allowbreak{}HOOK\_\allowbreak{}RETRY\_\allowbreak{}2026\_\allowbreak{}09\_\allowbreak{}08\_\allowbreak{}XML} & {\scriptsize\ttfamily docs/\allowbreak{}validation/\allowbreak{}wilson-\allowbreak{}shell-\allowbreak{}hook-\allowbreak{}retry-\allowbreak{}2026-\allowbreak{}09-\allowbreak{}08.xml} \\
{\scriptsize\ttfamily FOLDER\_\allowbreak{}VALIDATION\_\allowbreak{}WILSON\_\allowbreak{}SHELL\_\allowbreak{}LEDGER\_\allowbreak{}RETRY\_\allowbreak{}2026\_\allowbreak{}09\_\allowbreak{}08\_\allowbreak{}XML} & {\scriptsize\ttfamily docs/\allowbreak{}validation/\allowbreak{}wilson-\allowbreak{}shell-\allowbreak{}ledger-\allowbreak{}retry-\allowbreak{}2026-\allowbreak{}09-\allowbreak{}08.xml} \\
{\scriptsize\ttfamily FOLDER\_\allowbreak{}VALIDATION\_\allowbreak{}WILSON\_\allowbreak{}SHELL\_\allowbreak{}TESTS\_\allowbreak{}2026\_\allowbreak{}09\_\allowbreak{}08\_\allowbreak{}XML} & {\scriptsize\ttfamily docs/\allowbreak{}validation/\allowbreak{}wilson-\allowbreak{}shell-\allowbreak{}tests-\allowbreak{}2026-\allowbreak{}09-\allowbreak{}08.xml} \\
{\scriptsize\ttfamily FOLDER\_\allowbreak{}VALIDATION\_\allowbreak{}WILSON\_\allowbreak{}SPATIAL\_\allowbreak{}AXIOMS\_\allowbreak{}2026\_\allowbreak{}09\_\allowbreak{}08\_\allowbreak{}TXT} & {\scriptsize\ttfamily docs/\allowbreak{}validation/\allowbreak{}wilson-\allowbreak{}spatial-\allowbreak{}axioms-\allowbreak{}2026-\allowbreak{}09-\allowbreak{}08.txt} \\
{\scriptsize\ttfamily FOLDER\_\allowbreak{}VALIDATION\_\allowbreak{}WILSON\_\allowbreak{}SPATIAL\_\allowbreak{}FOCUSED\_\allowbreak{}2026\_\allowbreak{}09\_\allowbreak{}08\_\allowbreak{}XML} & {\scriptsize\ttfamily docs/\allowbreak{}validation/\allowbreak{}wilson-\allowbreak{}spatial-\allowbreak{}focused-\allowbreak{}2026-\allowbreak{}09-\allowbreak{}08.xml} \\
{\scriptsize\ttfamily FOLDER\_\allowbreak{}VALIDATION\_\allowbreak{}WILSON\_\allowbreak{}SPATIAL\_\allowbreak{}FOCUSED\_\allowbreak{}RETRY\_\allowbreak{}2026\_\allowbreak{}09\_\allowbreak{}08\_\allowbreak{}XML} & {\scriptsize\ttfamily docs/\allowbreak{}validation/\allowbreak{}wilson-\allowbreak{}spatial-\allowbreak{}focused-\allowbreak{}retry-\allowbreak{}2026-\allowbreak{}09-\allowbreak{}08.xml} \\
{\scriptsize\ttfamily FOLDER\_\allowbreak{}VALIDATION\_\allowbreak{}WILSON\_\allowbreak{}SPATIAL\_\allowbreak{}GRAPH\_\allowbreak{}2026\_\allowbreak{}09\_\allowbreak{}08\_\allowbreak{}XML} & {\scriptsize\ttfamily docs/\allowbreak{}validation/\allowbreak{}wilson-\allowbreak{}spatial-\allowbreak{}graph-\allowbreak{}2026-\allowbreak{}09-\allowbreak{}08.xml} \\
{\scriptsize\ttfamily FOLDER\_\allowbreak{}VALIDATION\_\allowbreak{}WILSON\_\allowbreak{}SPATIAL\_\allowbreak{}LEAN\_\allowbreak{}BUILD\_\allowbreak{}2026\_\allowbreak{}09\_\allowbreak{}08\_\allowbreak{}TXT} & {\scriptsize\ttfamily docs/\allowbreak{}validation/\allowbreak{}wilson-\allowbreak{}spatial-\allowbreak{}lean-\allowbreak{}build-\allowbreak{}2026-\allowbreak{}09-\allowbreak{}08.txt} \\
{\scriptsize\ttfamily FOLDER\_\allowbreak{}VALIDATION\_\allowbreak{}WILSON\_\allowbreak{}VACUUM\_\allowbreak{}FOCUSED\_\allowbreak{}XML} & {\scriptsize\ttfamily docs/\allowbreak{}validation/\allowbreak{}wilson-\allowbreak{}vacuum-\allowbreak{}focused.xml} \\
{\scriptsize\ttfamily FOLDER\_\allowbreak{}VALIDATION\_\allowbreak{}WILSON\_\allowbreak{}VACUUM\_\allowbreak{}GRAPH\_\allowbreak{}XML} & {\scriptsize\ttfamily docs/\allowbreak{}validation/\allowbreak{}wilson-\allowbreak{}vacuum-\allowbreak{}graph.xml} \\
{\scriptsize\ttfamily FOLDER\_\allowbreak{}VALIDATION\_\allowbreak{}WILSON\_\allowbreak{}VACUUM\_\allowbreak{}LEAN\_\allowbreak{}2026\_\allowbreak{}09\_\allowbreak{}09\_\allowbreak{}TXT} & {\scriptsize\ttfamily docs/\allowbreak{}validation/\allowbreak{}wilson-\allowbreak{}vacuum-\allowbreak{}lean-\allowbreak{}2026-\allowbreak{}09-\allowbreak{}09.txt} \\
{\scriptsize\ttfamily FOLDER\_\allowbreak{}VALIDATION\_\allowbreak{}WILSON\_\allowbreak{}VACUUM\_\allowbreak{}LEDGER\_\allowbreak{}XML} & {\scriptsize\ttfamily docs/\allowbreak{}validation/\allowbreak{}wilson-\allowbreak{}vacuum-\allowbreak{}ledger.xml} \\
{\scriptsize\ttfamily FOLDER\_\allowbreak{}VALIDATION\_\allowbreak{}WILSON\_\allowbreak{}WEIGHTED\_\allowbreak{}FOCUSED\_\allowbreak{}XML} & {\scriptsize\ttfamily docs/\allowbreak{}validation/\allowbreak{}wilson-\allowbreak{}weighted-\allowbreak{}focused.xml} \\
{\scriptsize\ttfamily FOLDER\_\allowbreak{}VALIDATION\_\allowbreak{}WILSON\_\allowbreak{}WEIGHTED\_\allowbreak{}GRAPH\_\allowbreak{}XML} & {\scriptsize\ttfamily docs/\allowbreak{}validation/\allowbreak{}wilson-\allowbreak{}weighted-\allowbreak{}graph.xml} \\
{\scriptsize\ttfamily FOLDER\_\allowbreak{}VALIDATION\_\allowbreak{}WILSON\_\allowbreak{}WEIGHTED\_\allowbreak{}LEDGER\_\allowbreak{}XML} & {\scriptsize\ttfamily docs/\allowbreak{}validation/\allowbreak{}wilson-\allowbreak{}weighted-\allowbreak{}ledger.xml} \\
{\scriptsize\ttfamily FOLDER\_\allowbreak{}VALIDATION\_\allowbreak{}YANGMILLS\_\allowbreak{}DERIVATION\_\allowbreak{}TESTS\_\allowbreak{}2026\_\allowbreak{}09\_\allowbreak{}08\_\allowbreak{}XML} & {\scriptsize\ttfamily docs/\allowbreak{}validation/\allowbreak{}yangmills-\allowbreak{}derivation-\allowbreak{}tests-\allowbreak{}2026-\allowbreak{}09-\allowbreak{}08.xml} \\
{\scriptsize\ttfamily FOLDER\_\allowbreak{}VALIDATION\_\allowbreak{}YANGMILLS\_\allowbreak{}DERIVATION\_\allowbreak{}UTF8\_\allowbreak{}RERUN\_\allowbreak{}2026\_\allowbreak{}09\_\allowbreak{}08\_\allowbreak{}XML} & {\scriptsize\ttfamily docs/\allowbreak{}validation/\allowbreak{}yangmills-\allowbreak{}derivation-\allowbreak{}utf8-\allowbreak{}rerun-\allowbreak{}2026-\allowbreak{}09-\allowbreak{}08.xml} \\
{\scriptsize\ttfamily FORM\_\allowbreak{}SCHUR\_\allowbreak{}SCALE\_\allowbreak{}COMPARISON} & {\scriptsize\ttfamily paper/\allowbreak{}research\_\allowbreak{}notes/\allowbreak{}G19\_\allowbreak{}FORM\_\allowbreak{}SCHUR\_\allowbreak{}SCALE\_\allowbreak{}COMPARISON\_\allowbreak{}20260905.md} \\
{\scriptsize\ttfamily G14\_\allowbreak{}UNIVERSAL\_\allowbreak{}CELLULAR\_\allowbreak{}HODGE\_\allowbreak{}TETRAHEDRAL} & {\scriptsize\ttfamily paper/\allowbreak{}research\_\allowbreak{}notes/\allowbreak{}G14\_\allowbreak{}UNIVERSAL\_\allowbreak{}CELLULAR\_\allowbreak{}HODGE\_\allowbreak{}TETRAHEDRAL\_\allowbreak{}20260911.md} \\
{\scriptsize\ttfamily G17\_\allowbreak{}CORRECTED\_\allowbreak{}ALPHA\_\allowbreak{}DEPENDENT\_\allowbreak{}KP\_\allowbreak{}STABILITY} & {\scriptsize\ttfamily paper/\allowbreak{}research\_\allowbreak{}notes/\allowbreak{}G17\_\allowbreak{}CORRECTED\_\allowbreak{}ALPHA\_\allowbreak{}DEPENDENT\_\allowbreak{}KP\_\allowbreak{}STABILITY\_\allowbreak{}20260910.md} \\
{\scriptsize\ttfamily G19\_\allowbreak{}CAUCHY\_\allowbreak{}DRIFT\_\allowbreak{}REPAIR} & {\scriptsize\ttfamily paper/\allowbreak{}research\_\allowbreak{}notes/\allowbreak{}G19\_\allowbreak{}CAUCHY\_\allowbreak{}DRIFT\_\allowbreak{}REPAIR\_\allowbreak{}20260909.md} \\
{\scriptsize\ttfamily G19\_\allowbreak{}CONTINUUM\_\allowbreak{}BRIDGE} & {\scriptsize\ttfamily paper/\allowbreak{}research\_\allowbreak{}notes/\allowbreak{}G19\_\allowbreak{}CONTINUUM\_\allowbreak{}BRIDGE\_\allowbreak{}INSERT.tex} \\
{\scriptsize\ttfamily G19\_\allowbreak{}DIMENSION\_\allowbreak{}FIVE\_\allowbreak{}IRRELEVANCE} & {\scriptsize\ttfamily paper/\allowbreak{}research\_\allowbreak{}notes/\allowbreak{}G19\_\allowbreak{}DIMENSION\_\allowbreak{}FIVE\_\allowbreak{}IRRELEVANCE\_\allowbreak{}LEMMA\_\allowbreak{}20260910.md} \\
{\scriptsize\ttfamily G19\_\allowbreak{}PROJECTIVE\_\allowbreak{}MARTINGALE\_\allowbreak{}REPAIR} & {\scriptsize\ttfamily paper/\allowbreak{}research\_\allowbreak{}notes/\allowbreak{}G19\_\allowbreak{}PROJECTIVE\_\allowbreak{}MARTINGALE\_\allowbreak{}REPAIR\_\allowbreak{}20260909.md} \\
{\scriptsize\ttfamily G19\_\allowbreak{}ROUGH\_\allowbreak{}GAUGE\_\allowbreak{}COERCIVITY} & {\scriptsize\ttfamily paper/\allowbreak{}research\_\allowbreak{}notes/\allowbreak{}G19\_\allowbreak{}ROUGH\_\allowbreak{}GAUGE\_\allowbreak{}COERCIVITY\_\allowbreak{}AND\_\allowbreak{}CURVATURE\_\allowbreak{}20260910.md} \\
{\scriptsize\ttfamily G9\_\allowbreak{}SIXTH\_\allowbreak{}ORDER\_\allowbreak{}COMBINED} & {\scriptsize\ttfamily paper/\allowbreak{}research\_\allowbreak{}notes/\allowbreak{}G9\_\allowbreak{}SIXTH\_\allowbreak{}ORDER\_\allowbreak{}COMBINED\_\allowbreak{}20260911.md} \\
{\scriptsize\ttfamily G9\_\allowbreak{}SIXTH\_\allowbreak{}ORDER\_\allowbreak{}RUR\_\allowbreak{}WORD\_\allowbreak{}DYNAMICS} & {\scriptsize\ttfamily paper/\allowbreak{}research\_\allowbreak{}notes/\allowbreak{}G9\_\allowbreak{}SIXTH\_\allowbreak{}ORDER\_\allowbreak{}RUR\_\allowbreak{}WORD\_\allowbreak{}DYNAMICS\_\allowbreak{}20260910.md} \\
{\scriptsize\ttfamily GAUSSIAN\_\allowbreak{}FULL\_\allowbreak{}FAST\_\allowbreak{}GREEN} & {\scriptsize\ttfamily paper/\allowbreak{}research\_\allowbreak{}notes/\allowbreak{}G19\_\allowbreak{}FULL\_\allowbreak{}GAUSSIAN\_\allowbreak{}FAST\_\allowbreak{}GREEN\_\allowbreak{}20260906.md} \\
{\scriptsize\ttfamily GAUSSIAN\_\allowbreak{}OS\_\allowbreak{}OBSERVABILITY} & {\scriptsize\ttfamily paper/\allowbreak{}research\_\allowbreak{}notes/\allowbreak{}G19\_\allowbreak{}GAUSSIAN\_\allowbreak{}OS\_\allowbreak{}HISTORY\_\allowbreak{}OBSERVABILITY\_\allowbreak{}20260905.md} \\
{\scriptsize\ttfamily GAUSSIAN\_\allowbreak{}PATH\_\allowbreak{}ENDPOINT\_\allowbreak{}BASELINE} & {\scriptsize\ttfamily paper/\allowbreak{}research\_\allowbreak{}notes/\allowbreak{}G19\_\allowbreak{}GAUSSIAN\_\allowbreak{}PATH\_\allowbreak{}ENDPOINT\_\allowbreak{}BASELINE\_\allowbreak{}20260906.md} \\
{\scriptsize\ttfamily GAUSSIAN\_\allowbreak{}QUANTUM\_\allowbreak{}FAST\_\allowbreak{}SOURCES} & {\scriptsize\ttfamily paper/\allowbreak{}research\_\allowbreak{}notes/\allowbreak{}G19\_\allowbreak{}GAUSSIAN\_\allowbreak{}QUANTUM\_\allowbreak{}FAST\_\allowbreak{}SOURCES\_\allowbreak{}20260905.md} \\
{\scriptsize\ttfamily GLUEBALL\_\allowbreak{}REVERSE\_\allowbreak{}TARGET\_\allowbreak{}DATA} & {\scriptsize\ttfamily paper/\allowbreak{}research\_\allowbreak{}notes/\allowbreak{}G19\_\allowbreak{}GLUEBALL\_\allowbreak{}REVERSE\_\allowbreak{}TARGET\_\allowbreak{}DATA\_\allowbreak{}20260905.md} \\
{\scriptsize\ttfamily GROUND\_\allowbreak{}MARGINAL\_\allowbreak{}SCHUR\_\allowbreak{}SCORE} & {\scriptsize\ttfamily paper/\allowbreak{}research\_\allowbreak{}notes/\allowbreak{}G19\_\allowbreak{}GROUND\_\allowbreak{}MARGINAL\_\allowbreak{}SCHUR\_\allowbreak{}SCORE\_\allowbreak{}20260905.md} \\
{\scriptsize\ttfamily HODGE\_\allowbreak{}FESHBACH\_\allowbreak{}CHANNEL} & {\scriptsize\ttfamily paper/\allowbreak{}research\_\allowbreak{}notes/\allowbreak{}G14\_\allowbreak{}HODGE\_\allowbreak{}FESHBACH\_\allowbreak{}CHANNEL\_\allowbreak{}20260908.md} \\
{\scriptsize\ttfamily HODGE\_\allowbreak{}POLYMER\_\allowbreak{}SUPPORTED} & {\scriptsize\ttfamily docs/\allowbreak{}derivations/\allowbreak{}hodge-\allowbreak{}polymer-\allowbreak{}supported-\allowbreak{}statements.md} \\
{\scriptsize\ttfamily LITERAL\_\allowbreak{}ENDPOINT\_\allowbreak{}COMPLETE\_\allowbreak{}WINDOW} & {\scriptsize\ttfamily paper/\allowbreak{}research\_\allowbreak{}notes/\allowbreak{}G19\_\allowbreak{}LITERAL\_\allowbreak{}ENDPOINT\_\allowbreak{}COMPLETE\_\allowbreak{}WINDOW\_\allowbreak{}20260905.md} \\
{\scriptsize\ttfamily LOCAL\_\allowbreak{}CLASS\_\allowbreak{}WICK\_\allowbreak{}RESULTS} & {\scriptsize\ttfamily paper/\allowbreak{}research\_\allowbreak{}notes/\allowbreak{}LOCAL\_\allowbreak{}CLASS\_\allowbreak{}WICK\_\allowbreak{}SPECTRUM\_\allowbreak{}20260911.md} \\
{\scriptsize\ttfamily LOCAL\_\allowbreak{}CLASS\_\allowbreak{}WICK\_\allowbreak{}SPECTRUM} & {\scriptsize\ttfamily docs/\allowbreak{}derivations/\allowbreak{}local-\allowbreak{}class-\allowbreak{}wick-\allowbreak{}spectrum.md} \\
{\scriptsize\ttfamily LOCAL\_\allowbreak{}COVARIANT\_\allowbreak{}AVERAGED\_\allowbreak{}PATH\_\allowbreak{}SOURCES} & {\scriptsize\ttfamily paper/\allowbreak{}research\_\allowbreak{}notes/\allowbreak{}G19\_\allowbreak{}LOCAL\_\allowbreak{}COVARIANT\_\allowbreak{}AVERAGED\_\allowbreak{}PATH\_\allowbreak{}SOURCES\_\allowbreak{}20260905.md} \\
{\scriptsize\ttfamily LOCAL\_\allowbreak{}GRADIENT\_\allowbreak{}EXCITATION\_\allowbreak{}SUPPORT} & {\scriptsize\ttfamily paper/\allowbreak{}research\_\allowbreak{}notes/\allowbreak{}G19\_\allowbreak{}LOCAL\_\allowbreak{}GRADIENT\_\allowbreak{}EXCITATION\_\allowbreak{}SUPPORT\_\allowbreak{}20260905.md} \\
{\scriptsize\ttfamily MASTER edition} & {\scriptsize\ttfamily paper/\allowbreak{}master\_\allowbreak{}paper\_\allowbreak{}2026-\allowbreak{}08-\allowbreak{}30.tex} \\
{\scriptsize\ttfamily MASTER paper} & {\scriptsize\ttfamily paper/\allowbreak{}master\_\allowbreak{}paper\_\allowbreak{}2026-\allowbreak{}08-\allowbreak{}28.tex} \\
{\scriptsize\ttfamily MASTER\_\allowbreak{}DOC} & {\scriptsize\ttfamily paper/\allowbreak{}nested\_\allowbreak{}quotient\_\allowbreak{}master\_\allowbreak{}2026-\allowbreak{}08-\allowbreak{}28.pdf} \\
{\scriptsize\ttfamily MOVING\_\allowbreak{}TIME\_\allowbreak{}GAP\_\allowbreak{}RESULTS} & {\scriptsize\ttfamily paper/\allowbreak{}research\_\allowbreak{}notes/\allowbreak{}MOVING\_\allowbreak{}TIME\_\allowbreak{}SPECTRAL\_\allowbreak{}GAP\_\allowbreak{}20260911.md} \\
{\scriptsize\ttfamily MOVING\_\allowbreak{}TIME\_\allowbreak{}SPECTRAL\_\allowbreak{}GAP} & {\scriptsize\ttfamily docs/\allowbreak{}derivations/\allowbreak{}moving-\allowbreak{}time-\allowbreak{}spectral-\allowbreak{}gap.md} \\
{\scriptsize\ttfamily ODD\_\allowbreak{}ORDER\_\allowbreak{}CENTRE\_\allowbreak{}PARITY} & {\scriptsize\ttfamily paper/\allowbreak{}research\_\allowbreak{}notes/\allowbreak{}ODD\_\allowbreak{}ORDER\_\allowbreak{}CENTRE\_\allowbreak{}PARITY\_\allowbreak{}20260911.md} \\
{\scriptsize\ttfamily OS\_\allowbreak{}HISTORY\_\allowbreak{}BLOCKING} & {\scriptsize\ttfamily paper/\allowbreak{}research\_\allowbreak{}notes/\allowbreak{}G19\_\allowbreak{}OS\_\allowbreak{}BLOCKING\_\allowbreak{}AND\_\allowbreak{}REVERSE\_\allowbreak{}MASS\_\allowbreak{}MATCHING\_\allowbreak{}20260905.md} \\
{\scriptsize\ttfamily OS\_\allowbreak{}KERNEL\_\allowbreak{}GAP} & {\scriptsize\ttfamily docs/\allowbreak{}derivations/\allowbreak{}os-\allowbreak{}kernel-\allowbreak{}moving-\allowbreak{}time-\allowbreak{}gap.md} \\
{\scriptsize\ttfamily OS\_\allowbreak{}KERNEL\_\allowbreak{}GAP\_\allowbreak{}RESULT} & {\scriptsize\ttfamily paper/\allowbreak{}research\_\allowbreak{}notes/\allowbreak{}G19\_\allowbreak{}KERNEL\_\allowbreak{}CLOSURE\_\allowbreak{}ATTEMPT\_\allowbreak{}20260911.md} \\
{\scriptsize\ttfamily PAPER\_\allowbreak{}FLATBAND} & {\scriptsize\ttfamily paper/\allowbreak{}homological\_\allowbreak{}flat\_\allowbreak{}bands\_\allowbreak{}2026-\allowbreak{}08-\allowbreak{}28.pdf} \\
{\scriptsize\ttfamily PUBLICATION rev5} & {\scriptsize\ttfamily paper/\allowbreak{}workhouse\_\allowbreak{}publication\_\allowbreak{}edition\_\allowbreak{}rev5\_\allowbreak{}2026-\allowbreak{}08-\allowbreak{}30.tex} \\
{\scriptsize\ttfamily PUBLICATION v2} & {\scriptsize\ttfamily paper/\allowbreak{}workhouse\_\allowbreak{}publication\_\allowbreak{}edition\_\allowbreak{}v2\_\allowbreak{}2026-\allowbreak{}08-\allowbreak{}30.tex} \\
{\scriptsize\ttfamily REVIEW\_\allowbreak{}SPECTRAL\_\allowbreak{}BUDGETS} & {\scriptsize\ttfamily docs/\allowbreak{}derivations/\allowbreak{}review-\allowbreak{}derived-\allowbreak{}spectral-\allowbreak{}budgets.md} \\
{\scriptsize\ttfamily REVIEW\_\allowbreak{}SPECTRAL\_\allowbreak{}BUDGET\_\allowbreak{}RESULTS} & {\scriptsize\ttfamily paper/\allowbreak{}research\_\allowbreak{}notes/\allowbreak{}PILLARS\_\allowbreak{}SPECTRAL\_\allowbreak{}BUDGET\_\allowbreak{}RESULTS\_\allowbreak{}20260911.md} \\
{\scriptsize\ttfamily SEPTEMBER10\_\allowbreak{}UPSTREAM\_\allowbreak{}SCALARS} & {\scriptsize\ttfamily docs/\allowbreak{}derivations/\allowbreak{}september10-\allowbreak{}upstream-\allowbreak{}scalar-\allowbreak{}inventory.md} \\
{\scriptsize\ttfamily SOURCE\_\allowbreak{}CURRENT\_\allowbreak{}SPECTRAL\_\allowbreak{}BRIDGES} & {\scriptsize\ttfamily docs/\allowbreak{}derivations/\allowbreak{}source-\allowbreak{}currents-\allowbreak{}and-\allowbreak{}spectral-\allowbreak{}totality.md} \\
{\scriptsize\ttfamily THEORY\_\allowbreak{}CURRENT\_\allowbreak{}BRIDGES} & {\scriptsize\ttfamily paper/\allowbreak{}research\_\allowbreak{}notes/\allowbreak{}THEORY\_\allowbreak{}CURRENT\_\allowbreak{}BRIDGES\_\allowbreak{}20260910.md} \\
{\scriptsize\ttfamily THEORY\_\allowbreak{}GEOMETRY\_\allowbreak{}RICCATI\_\allowbreak{}BRIDGES} & {\scriptsize\ttfamily docs/\allowbreak{}derivations/\allowbreak{}theory-\allowbreak{}geometry-\allowbreak{}and-\allowbreak{}riccati-\allowbreak{}bridges.md} \\
{\scriptsize\ttfamily TRUE\_\allowbreak{}GROUND\_\allowbreak{}CENTER\_\allowbreak{}SCORE\_\allowbreak{}OBSTRUCTION} & {\scriptsize\ttfamily paper/\allowbreak{}research\_\allowbreak{}notes/\allowbreak{}G19\_\allowbreak{}TRUE\_\allowbreak{}GROUND\_\allowbreak{}CENTER\_\allowbreak{}SCORE\_\allowbreak{}OBSTRUCTION\_\allowbreak{}20260905.md} \\
{\scriptsize\ttfamily TRUE\_\allowbreak{}GROUND\_\allowbreak{}LOCALIZED\_\allowbreak{}WILSON\_\allowbreak{}SCORE} & {\scriptsize\ttfamily paper/\allowbreak{}research\_\allowbreak{}notes/\allowbreak{}G19\_\allowbreak{}TRUE\_\allowbreak{}GROUND\_\allowbreak{}LOCALIZED\_\allowbreak{}WILSON\_\allowbreak{}SCORE\_\allowbreak{}20260905.md} \\
{\scriptsize\ttfamily UNIVERSAL\_\allowbreak{}CELLULAR\_\allowbreak{}HODGE\_\allowbreak{}TETRAHEDRAL} & {\scriptsize\ttfamily docs/\allowbreak{}derivations/\allowbreak{}universal-\allowbreak{}cellular-\allowbreak{}hodge-\allowbreak{}tetrahedral.md} \\
{\scriptsize\ttfamily URZUA\_\allowbreak{}COUPLED\_\allowbreak{}OSCILLATORS\_\allowbreak{}2019} & {\scriptsize\ttfamily literature/\allowbreak{}fulltext/\allowbreak{}quantumrep-\allowbreak{}01-\allowbreak{}00009.pdf} \\
{\scriptsize\ttfamily W6\_\allowbreak{}ANTIPODAL\_\allowbreak{}MAGNETIC\_\allowbreak{}GEOMETRY} & {\scriptsize\ttfamily docs/\allowbreak{}derivations/\allowbreak{}w6-\allowbreak{}antipodal-\allowbreak{}magnetic-\allowbreak{}geometry.md} \\
{\scriptsize\ttfamily W6\_\allowbreak{}CONDITIONAL\_\allowbreak{}SCORE\_\allowbreak{}TAIL\_\allowbreak{}CONTROL} & {\scriptsize\ttfamily docs/\allowbreak{}derivations/\allowbreak{}w6-\allowbreak{}conditional-\allowbreak{}score-\allowbreak{}tail-\allowbreak{}control.md} \\
{\scriptsize\ttfamily W6\_\allowbreak{}CONDITIONAL\_\allowbreak{}TRANSPORT\_\allowbreak{}OBSTRUCTION} & {\scriptsize\ttfamily docs/\allowbreak{}derivations/\allowbreak{}w6-\allowbreak{}conditional-\allowbreak{}transport-\allowbreak{}obstruction.md} \\
{\scriptsize\ttfamily W6\_\allowbreak{}GROUND\_\allowbreak{}JETS\_\allowbreak{}TRANSPORT\_\allowbreak{}BUDGET} & {\scriptsize\ttfamily docs/\allowbreak{}derivations/\allowbreak{}w6-\allowbreak{}ground-\allowbreak{}jets-\allowbreak{}and-\allowbreak{}transport-\allowbreak{}budget.md} \\
{\scriptsize\ttfamily W6\_\allowbreak{}M10\_\allowbreak{}REPAIR\_\allowbreak{}RESULTS} & {\scriptsize\ttfamily paper/\allowbreak{}research\_\allowbreak{}notes/\allowbreak{}W6\_\allowbreak{}M10\_\allowbreak{}REPAIR\_\allowbreak{}RESULTS\_\allowbreak{}20260911.md} \\
{\scriptsize\ttfamily W6\_\allowbreak{}R10\_\allowbreak{}REPAIRED\_\allowbreak{}RESULTS} & {\scriptsize\ttfamily paper/\allowbreak{}research\_\allowbreak{}notes/\allowbreak{}W6\_\allowbreak{}R10\_\allowbreak{}REPAIRED\_\allowbreak{}IDENTITIES\_\allowbreak{}20260911.md} \\
{\scriptsize\ttfamily W6\_\allowbreak{}SOURCE\_\allowbreak{}ENERGY\_\allowbreak{}JETS\_\allowbreak{}R10} & {\scriptsize\ttfamily docs/\allowbreak{}derivations/\allowbreak{}w6-\allowbreak{}source-\allowbreak{}energy-\allowbreak{}jets-\allowbreak{}r10.md} \\
{\scriptsize\ttfamily W6\_\allowbreak{}SOURCE\_\allowbreak{}GENERATOR\_\allowbreak{}SCORE\_\allowbreak{}FRAME} & {\scriptsize\ttfamily docs/\allowbreak{}derivations/\allowbreak{}w6-\allowbreak{}source-\allowbreak{}generator-\allowbreak{}score-\allowbreak{}frame.md} \\
{\scriptsize\ttfamily W6\_\allowbreak{}SYNCHRONIZED\_\allowbreak{}M10\_\allowbreak{}DOMINATION} & {\scriptsize\ttfamily docs/\allowbreak{}derivations/\allowbreak{}w6-\allowbreak{}synchronized-\allowbreak{}m10-\allowbreak{}domination.md} \\
{\scriptsize\ttfamily WILSON\_\allowbreak{}ACTIVITY\_\allowbreak{}EXTRACTION} & {\scriptsize\ttfamily paper/\allowbreak{}research\_\allowbreak{}notes/\allowbreak{}G18\_\allowbreak{}WILSON\_\allowbreak{}ACTIVITY\_\allowbreak{}EXTRACTION\_\allowbreak{}20260905.md} \\
{\scriptsize\ttfamily WILSON\_\allowbreak{}ACTUAL\_\allowbreak{}BLOCK\_\allowbreak{}FAST\_\allowbreak{}COMPLEMENT} & {\scriptsize\ttfamily paper/\allowbreak{}research\_\allowbreak{}notes/\allowbreak{}G19\_\allowbreak{}WILSON\_\allowbreak{}ACTUAL\_\allowbreak{}BLOCK\_\allowbreak{}FAST\_\allowbreak{}COMPLEMENT\_\allowbreak{}20260905.md} \\
{\scriptsize\ttfamily WILSON\_\allowbreak{}BACKGROUND} & {\scriptsize\ttfamily docs/\allowbreak{}derivations/\allowbreak{}wilson-\allowbreak{}background-\allowbreak{}covariance-\allowbreak{}continuation.md} \\
{\scriptsize\ttfamily WILSON\_\allowbreak{}BLOCK\_\allowbreak{}SCORE} & {\scriptsize\ttfamily paper/\allowbreak{}research\_\allowbreak{}notes/\allowbreak{}G19\_\allowbreak{}WILSON\_\allowbreak{}BLOCK\_\allowbreak{}SCORE\_\allowbreak{}AND\_\allowbreak{}FIBER\_\allowbreak{}OBSTRUCTION\_\allowbreak{}20260905.md} \\
{\scriptsize\ttfamily WILSON\_\allowbreak{}CARDINALITY\_\allowbreak{}CHART} & {\scriptsize\ttfamily paper/\allowbreak{}research\_\allowbreak{}notes/\allowbreak{}G18\_\allowbreak{}WILSON\_\allowbreak{}CARDINALITY\_\allowbreak{}UNITARY\_\allowbreak{}CHART\_\allowbreak{}20260905.md} \\
{\scriptsize\ttfamily WILSON\_\allowbreak{}COMMON\_\allowbreak{}GAUSS\_\allowbreak{}LITERAL\_\allowbreak{}FAST\_\allowbreak{}FLOOR} & {\scriptsize\ttfamily paper/\allowbreak{}research\_\allowbreak{}notes/\allowbreak{}G19\_\allowbreak{}WILSON\_\allowbreak{}COMMON\_\allowbreak{}GAUSS\_\allowbreak{}LITERAL\_\allowbreak{}FAST\_\allowbreak{}FLOOR\_\allowbreak{}20260905.md} \\
{\scriptsize\ttfamily WILSON\_\allowbreak{}CREATOR\_\allowbreak{}LIMIT} & {\scriptsize\ttfamily paper/\allowbreak{}research\_\allowbreak{}notes/\allowbreak{}G18\_\allowbreak{}WILSON\_\allowbreak{}CREATOR\_\allowbreak{}THERMODYNAMIC\_\allowbreak{}LIMIT\_\allowbreak{}20260905.md} \\
{\scriptsize\ttfamily WILSON\_\allowbreak{}CREATOR\_\allowbreak{}PARENT} & {\scriptsize\ttfamily paper/\allowbreak{}research\_\allowbreak{}notes/\allowbreak{}G18\_\allowbreak{}WILSON\_\allowbreak{}CREATOR\_\allowbreak{}PARENT\_\allowbreak{}AND\_\allowbreak{}SPECTRAL\_\allowbreak{}FLOW\_\allowbreak{}20260905.md} \\
{\scriptsize\ttfamily WILSON\_\allowbreak{}DISCRETE\_\allowbreak{}TIME\_\allowbreak{}VACUUM} & {\scriptsize\ttfamily paper/\allowbreak{}research\_\allowbreak{}notes/\allowbreak{}G19\_\allowbreak{}DISCRETE\_\allowbreak{}TIME\_\allowbreak{}VACUUM\_\allowbreak{}AND\_\allowbreak{}WINDOW\_\allowbreak{}20260904.md} \\
{\scriptsize\ttfamily WILSON\_\allowbreak{}ENDPOINT\_\allowbreak{}EQUATION} & {\scriptsize\ttfamily paper/\allowbreak{}research\_\allowbreak{}notes/\allowbreak{}G18\_\allowbreak{}WILSON\_\allowbreak{}ENDPOINT\_\allowbreak{}GAUGE\_\allowbreak{}AND\_\allowbreak{}MAJORANT\_\allowbreak{}20260905.md} \\
{\scriptsize\ttfamily WILSON\_\allowbreak{}EXCITED\_\allowbreak{}OPERATOR\_\allowbreak{}BRIDGE} & {\scriptsize\ttfamily paper/\allowbreak{}research\_\allowbreak{}notes/\allowbreak{}G18\_\allowbreak{}EXCITED\_\allowbreak{}WINDOW\_\allowbreak{}OPERATOR\_\allowbreak{}BRIDGE\_\allowbreak{}20260904.md} \\
{\scriptsize\ttfamily WILSON\_\allowbreak{}FIBER\_\allowbreak{}FAST\_\allowbreak{}GAP} & {\scriptsize\ttfamily paper/\allowbreak{}research\_\allowbreak{}notes/\allowbreak{}G19\_\allowbreak{}WILSON\_\allowbreak{}PHYSICAL\_\allowbreak{}FIBER\_\allowbreak{}FAST\_\allowbreak{}GAP\_\allowbreak{}20260905.md} \\
{\scriptsize\ttfamily WILSON\_\allowbreak{}FINITE\_\allowbreak{}CELL\_\allowbreak{}GAP} & {\scriptsize\ttfamily paper/\allowbreak{}research\_\allowbreak{}notes/\allowbreak{}G19\_\allowbreak{}WILSON\_\allowbreak{}FINITE\_\allowbreak{}CELL\_\allowbreak{}GAP\_\allowbreak{}AND\_\allowbreak{}BOUNDARY\_\allowbreak{}FORM\_\allowbreak{}20260905.md} \\
{\scriptsize\ttfamily WILSON\_\allowbreak{}FIXED\_\allowbreak{}ORDER\_\allowbreak{}CHART} & {\scriptsize\ttfamily paper/\allowbreak{}research\_\allowbreak{}notes/\allowbreak{}G18\_\allowbreak{}VACUUM\_\allowbreak{}CHART\_\allowbreak{}RECURSION\_\allowbreak{}20260905.md} \\
{\scriptsize\ttfamily WILSON\_\allowbreak{}G19\_\allowbreak{}REVIEW} & {\scriptsize\ttfamily docs/\allowbreak{}validation/\allowbreak{}wilson-\allowbreak{}g19-\allowbreak{}corpus-\allowbreak{}reconciliation.md} \\
{\scriptsize\ttfamily WILSON\_\allowbreak{}GLOBAL\_\allowbreak{}VERTICAL\_\allowbreak{}BARRIER} & {\scriptsize\ttfamily paper/\allowbreak{}research\_\allowbreak{}notes/\allowbreak{}G19\_\allowbreak{}WILSON\_\allowbreak{}GLOBAL\_\allowbreak{}VERTICAL\_\allowbreak{}BARRIER\_\allowbreak{}20260905.md} \\
{\scriptsize\ttfamily WILSON\_\allowbreak{}GROUND\_\allowbreak{}BUNDLE\_\allowbreak{}RELATIVE\_\allowbreak{}FORM} & {\scriptsize\ttfamily paper/\allowbreak{}research\_\allowbreak{}notes/\allowbreak{}G19\_\allowbreak{}WILSON\_\allowbreak{}GROUND\_\allowbreak{}BUNDLE\_\allowbreak{}RELATIVE\_\allowbreak{}FORM\_\allowbreak{}20260905.md} \\
{\scriptsize\ttfamily WILSON\_\allowbreak{}HARMONIC\_\allowbreak{}BOUNDARY} & {\scriptsize\ttfamily paper/\allowbreak{}research\_\allowbreak{}notes/\allowbreak{}G19\_\allowbreak{}WILSON\_\allowbreak{}HARMONIC\_\allowbreak{}BOUNDARY\_\allowbreak{}COMPARISON\_\allowbreak{}20260905.md} \\
{\scriptsize\ttfamily WILSON\_\allowbreak{}HARMONIC\_\allowbreak{}CUBIC\_\allowbreak{}OBSTRUCTION} & {\scriptsize\ttfamily paper/\allowbreak{}research\_\allowbreak{}notes/\allowbreak{}G19\_\allowbreak{}WILSON\_\allowbreak{}HARMONIC\_\allowbreak{}CUBIC\_\allowbreak{}OBSTRUCTION\_\allowbreak{}20260906.md} \\
{\scriptsize\ttfamily WILSON\_\allowbreak{}INFINITE\_\allowbreak{}PHYSICAL\_\allowbreak{}BAND} & {\scriptsize\ttfamily paper/\allowbreak{}research\_\allowbreak{}notes/\allowbreak{}G18\_\allowbreak{}WILSON\_\allowbreak{}INFINITE\_\allowbreak{}VOLUME\_\allowbreak{}PHYSICAL\_\allowbreak{}BAND\_\allowbreak{}20260905.md} \\
{\scriptsize\ttfamily WILSON\_\allowbreak{}KINETIC\_\allowbreak{}WINDOW} & {\scriptsize\ttfamily paper/\allowbreak{}research\_\allowbreak{}notes/\allowbreak{}G19\_\allowbreak{}UNIFORM\_\allowbreak{}WILSON\_\allowbreak{}WINDOW\_\allowbreak{}20260904.md} \\
{\scriptsize\ttfamily WILSON\_\allowbreak{}LITERAL\_\allowbreak{}VACUUM\_\allowbreak{}COARSE\_\allowbreak{}SOURCES} & {\scriptsize\ttfamily paper/\allowbreak{}research\_\allowbreak{}notes/\allowbreak{}G19\_\allowbreak{}WILSON\_\allowbreak{}LITERAL\_\allowbreak{}VACUUM\_\allowbreak{}COARSE\_\allowbreak{}SOURCES\_\allowbreak{}20260905.md} \\
{\scriptsize\ttfamily WILSON\_\allowbreak{}MARKED} & {\scriptsize\ttfamily docs/\allowbreak{}derivations/\allowbreak{}wilson-\allowbreak{}marked-\allowbreak{}transfer.md} \\
{\scriptsize\ttfamily WILSON\_\allowbreak{}ROOTED\_\allowbreak{}CONTRACTION} & {\scriptsize\ttfamily paper/\allowbreak{}research\_\allowbreak{}notes/\allowbreak{}G18\_\allowbreak{}ROOTED\_\allowbreak{}WILSON\_\allowbreak{}CONTRACTION\_\allowbreak{}20260905.md} \\
{\scriptsize\ttfamily WILSON\_\allowbreak{}SAME\_\allowbreak{}WEIGHT\_\allowbreak{}OBSTRUCTION} & {\scriptsize\ttfamily paper/\allowbreak{}research\_\allowbreak{}notes/\allowbreak{}G18\_\allowbreak{}SAME\_\allowbreak{}WEIGHT\_\allowbreak{}CREATOR\_\allowbreak{}OBSTRUCTION\_\allowbreak{}20260905.md} \\
{\scriptsize\ttfamily WILSON\_\allowbreak{}SC17} & {\scriptsize\ttfamily docs/\allowbreak{}derivations/\allowbreak{}wilson-\allowbreak{}sc17-\allowbreak{}spatial-\allowbreak{}closure.md} \\
{\scriptsize\ttfamily WILSON\_\allowbreak{}SC17\_\allowbreak{}LIMIT} & {\scriptsize\ttfamily docs/\allowbreak{}derivations/\allowbreak{}wilson-\allowbreak{}sc17-\allowbreak{}thermodynamic-\allowbreak{}limit.md} \\
{\scriptsize\ttfamily WILSON\_\allowbreak{}SC17\_\allowbreak{}TIME} & {\scriptsize\ttfamily docs/\allowbreak{}derivations/\allowbreak{}wilson-\allowbreak{}sc17-\allowbreak{}physical-\allowbreak{}time-\allowbreak{}limit.md} \\
{\scriptsize\ttfamily WILSON\_\allowbreak{}SECOND\_\allowbreak{}ORDER\_\allowbreak{}CHART} & {\scriptsize\ttfamily paper/\allowbreak{}research\_\allowbreak{}notes/\allowbreak{}G18\_\allowbreak{}SECOND\_\allowbreak{}ORDER\_\allowbreak{}WILSON\_\allowbreak{}VACUUM\_\allowbreak{}CHART\_\allowbreak{}20260905.md} \\
{\scriptsize\ttfamily WILSON\_\allowbreak{}STRIP\_\allowbreak{}BO} & {\scriptsize\ttfamily paper/\allowbreak{}research\_\allowbreak{}notes/\allowbreak{}G19\_\allowbreak{}WILSON\_\allowbreak{}STRIP\_\allowbreak{}BO\_\allowbreak{}AND\_\allowbreak{}TWO\_\allowbreak{}STRIP\_\allowbreak{}SPLITTING\_\allowbreak{}20260905.md} \\
{\scriptsize\ttfamily WILSON\_\allowbreak{}THREE\_\allowbreak{}DIMENSIONAL\_\allowbreak{}HARMONIC\_\allowbreak{}BOUNDARY} & {\scriptsize\ttfamily paper/\allowbreak{}research\_\allowbreak{}notes/\allowbreak{}G19\_\allowbreak{}WILSON\_\allowbreak{}THREE\_\allowbreak{}DIMENSIONAL\_\allowbreak{}HARMONIC\_\allowbreak{}BOUNDARY\_\allowbreak{}20260905.md} \\
{\scriptsize\ttfamily WILSON\_\allowbreak{}TRANSFER\_\allowbreak{}RESOLVENT} & {\scriptsize\ttfamily paper/\allowbreak{}research\_\allowbreak{}notes/\allowbreak{}G19\_\allowbreak{}TRANSFER\_\allowbreak{}RESOLVENT\_\allowbreak{}UNIFORM\_\allowbreak{}BOUNDS\_\allowbreak{}20260905.md} \\
{\scriptsize\ttfamily WILSON\_\allowbreak{}TRUE\_\allowbreak{}BLOCK\_\allowbreak{}CHECK} & {\scriptsize\ttfamily docs/\allowbreak{}derivations/\allowbreak{}wilson-\allowbreak{}true-\allowbreak{}blocks-\allowbreak{}independent-\allowbreak{}check.md} \\
{\scriptsize\ttfamily WILSON\_\allowbreak{}TWO\_\allowbreak{}SQUARE\_\allowbreak{}PHYSICAL\_\allowbreak{}SHELLS} & {\scriptsize\ttfamily paper/\allowbreak{}research\_\allowbreak{}notes/\allowbreak{}G19\_\allowbreak{}WILSON\_\allowbreak{}TWO\_\allowbreak{}SQUARE\_\allowbreak{}PHYSICAL\_\allowbreak{}SHELLS\_\allowbreak{}20260905.md} \\
{\scriptsize\ttfamily WILSON\_\allowbreak{}VACUUM\_\allowbreak{}COMPRESSION} & {\scriptsize\ttfamily paper/\allowbreak{}research\_\allowbreak{}notes/\allowbreak{}G18\_\allowbreak{}VACUUM\_\allowbreak{}COMPRESSION\_\allowbreak{}BOUND\_\allowbreak{}20260905.md} \\
{\scriptsize\ttfamily WILSON\_\allowbreak{}WEIGHTED\_\allowbreak{}ACTIVITIES} & {\scriptsize\ttfamily paper/\allowbreak{}research\_\allowbreak{}notes/\allowbreak{}G18\_\allowbreak{}WILSON\_\allowbreak{}WEIGHTED\_\allowbreak{}ACTIVITY\_\allowbreak{}BOUND\_\allowbreak{}20260905.md} \\
{\scriptsize\ttfamily YM\_\allowbreak{}FLAT\_\allowbreak{}DIRECTIONS} & {\scriptsize\ttfamily docs/\allowbreak{}derivations/\allowbreak{}yangmills-\allowbreak{}simon-\allowbreak{}flat-\allowbreak{}directions.md} \\
{\scriptsize\ttfamily YM\_\allowbreak{}GROUND\_\allowbreak{}STATE} & {\scriptsize\ttfamily docs/\allowbreak{}derivations/\allowbreak{}yangmills-\allowbreak{}weighted-\allowbreak{}curvature.md} \\
{\scriptsize\ttfamily YM\_\allowbreak{}RECONSTRUCTION} & {\scriptsize\ttfamily docs/\allowbreak{}derivations/\allowbreak{}yangmills-\allowbreak{}reconstruction.md} \\
\end{longtable}

\subsection*{Standing: current (74)}
\addcontentsline{toc}{subsection}{Standing: current}
\begin{longtable}{@{}p{2.0in} p{4.2in}@{}}
\toprule Alias & Path \\ \midrule
\endfirsthead
\toprule Alias & Path \\ \midrule \endhead
\bottomrule\endfoot
{\scriptsize\ttfamily FESHBACH\_\allowbreak{}RESOLVENT} & {\scriptsize\ttfamily docs/\allowbreak{}research/\allowbreak{}september\_\allowbreak{}feshbach\_\allowbreak{}integration.md} \\
{\scriptsize\ttfamily GLUEBALL} & {\scriptsize\ttfamily theory/\allowbreak{}GLUEBALL\_\allowbreak{}DETAILED\_\allowbreak{}FORMULA\_\allowbreak{}2026-\allowbreak{}08-\allowbreak{}20\_\allowbreak{}v3.1.md} \\
{\scriptsize\ttfamily GUIDE} & {\scriptsize\ttfamily theory/\allowbreak{}GLUEBALL\_\allowbreak{}SOURCE\_\allowbreak{}CONSOLIDATION\_\allowbreak{}GUIDE\_\allowbreak{}2026-\allowbreak{}08-\allowbreak{}20\_\allowbreak{}v4\_\allowbreak{}3.md} \\
{\scriptsize\ttfamily HODGE\_\allowbreak{}FESHBACH} & {\scriptsize\ttfamily docs/\allowbreak{}research/\allowbreak{}september\_\allowbreak{}feshbach\_\allowbreak{}integration.md} \\
{\scriptsize\ttfamily MANIFEST} & {\scriptsize\ttfamily theory/\allowbreak{}GLUEBALL\_\allowbreak{}CANONICAL\_\allowbreak{}SOURCE\_\allowbreak{}MANIFEST\_\allowbreak{}2026-\allowbreak{}08-\allowbreak{}20\_\allowbreak{}v4\_\allowbreak{}3.csv} \\
{\scriptsize\ttfamily PBH\_\allowbreak{}WILSON\_\allowbreak{}TEST} & {\scriptsize\ttfamily docs/\allowbreak{}derivations/\allowbreak{}yangmills-\allowbreak{}pbh-\allowbreak{}wilson-\allowbreak{}test.md} \\
{\scriptsize\ttfamily RECENT\_\allowbreak{}ANISOTROPY} & {\scriptsize\ttfamily runs/\allowbreak{}recent\_\allowbreak{}research\_\allowbreak{}integration\_\allowbreak{}2026-\allowbreak{}09-\allowbreak{}09/\allowbreak{}sources/\allowbreak{}anisotropy\_\allowbreak{}variance\_\allowbreak{}20260908/\allowbreak{}THEORY.md} \\
{\scriptsize\ttfamily SEPT\_\allowbreak{}DOC\_\allowbreak{}G18\_\allowbreak{}EXCITED\_\allowbreak{}WINDOW\_\allowbreak{}OPERATOR\_\allowbreak{}BRIDGE\_\allowbreak{}20260904} & {\scriptsize\ttfamily docs/\allowbreak{}derivations/\allowbreak{}wilson-\allowbreak{}shell-\allowbreak{}inputs/\allowbreak{}G18\_\allowbreak{}EXCITED\_\allowbreak{}WINDOW\_\allowbreak{}OPERATOR\_\allowbreak{}BRIDGE\_\allowbreak{}20260904.md} \\
{\scriptsize\ttfamily SEPT\_\allowbreak{}DOC\_\allowbreak{}G18\_\allowbreak{}ROOTED\_\allowbreak{}WILSON\_\allowbreak{}CONTRACTION\_\allowbreak{}20260905} & {\scriptsize\ttfamily docs/\allowbreak{}derivations/\allowbreak{}wilson-\allowbreak{}shell-\allowbreak{}inputs/\allowbreak{}G18\_\allowbreak{}ROOTED\_\allowbreak{}WILSON\_\allowbreak{}CONTRACTION\_\allowbreak{}20260905.md} \\
{\scriptsize\ttfamily SEPT\_\allowbreak{}DOC\_\allowbreak{}G18\_\allowbreak{}WILSON\_\allowbreak{}ACTIVITY\_\allowbreak{}EXTRACTION\_\allowbreak{}20260905} & {\scriptsize\ttfamily docs/\allowbreak{}derivations/\allowbreak{}wilson-\allowbreak{}shell-\allowbreak{}inputs/\allowbreak{}G18\_\allowbreak{}WILSON\_\allowbreak{}ACTIVITY\_\allowbreak{}EXTRACTION\_\allowbreak{}20260905.md} \\
{\scriptsize\ttfamily SEPT\_\allowbreak{}DOC\_\allowbreak{}G18\_\allowbreak{}WILSON\_\allowbreak{}CREATOR\_\allowbreak{}PARENT\_\allowbreak{}AND\_\allowbreak{}SPECTRAL\_\allowbreak{}FLOW\_\allowbreak{}20260905} & {\scriptsize\ttfamily docs/\allowbreak{}derivations/\allowbreak{}wilson-\allowbreak{}shell-\allowbreak{}inputs/\allowbreak{}G18\_\allowbreak{}WILSON\_\allowbreak{}CREATOR\_\allowbreak{}PARENT\_\allowbreak{}AND\_\allowbreak{}SPECTRAL\_\allowbreak{}FLOW\_\allowbreak{}20260905.md} \\
{\scriptsize\ttfamily SEPT\_\allowbreak{}DOC\_\allowbreak{}G18\_\allowbreak{}WILSON\_\allowbreak{}CREATOR\_\allowbreak{}THERMODYNAMIC\_\allowbreak{}LIMIT\_\allowbreak{}20260905} & {\scriptsize\ttfamily docs/\allowbreak{}derivations/\allowbreak{}wilson-\allowbreak{}shell-\allowbreak{}inputs/\allowbreak{}G18\_\allowbreak{}WILSON\_\allowbreak{}CREATOR\_\allowbreak{}THERMODYNAMIC\_\allowbreak{}LIMIT\_\allowbreak{}20260905.md} \\
{\scriptsize\ttfamily SEPT\_\allowbreak{}DOC\_\allowbreak{}G18\_\allowbreak{}WILSON\_\allowbreak{}ENDPOINT\_\allowbreak{}GAUGE\_\allowbreak{}AND\_\allowbreak{}MAJORANT\_\allowbreak{}20260905} & {\scriptsize\ttfamily docs/\allowbreak{}derivations/\allowbreak{}wilson-\allowbreak{}shell-\allowbreak{}inputs/\allowbreak{}G18\_\allowbreak{}WILSON\_\allowbreak{}ENDPOINT\_\allowbreak{}GAUGE\_\allowbreak{}AND\_\allowbreak{}MAJORANT\_\allowbreak{}20260905.md} \\
{\scriptsize\ttfamily SEPT\_\allowbreak{}DOC\_\allowbreak{}G19\_\allowbreak{}CUBIC\_\allowbreak{}GROUND\_\allowbreak{}TRANSFER\_\allowbreak{}20260906} & {\scriptsize\ttfamily docs/\allowbreak{}derivations/\allowbreak{}wilson-\allowbreak{}spatial-\allowbreak{}inputs/\allowbreak{}G19\_\allowbreak{}CUBIC\_\allowbreak{}GROUND\_\allowbreak{}TRANSFER\_\allowbreak{}20260906.md} \\
{\scriptsize\ttfamily SEPT\_\allowbreak{}DOC\_\allowbreak{}G19\_\allowbreak{}DYNAMIC\_\allowbreak{}FIBER\_\allowbreak{}COVARIANCE\_\allowbreak{}AND\_\allowbreak{}CUBIC\_\allowbreak{}ENERGY\_\allowbreak{}20260906} & {\scriptsize\ttfamily docs/\allowbreak{}derivations/\allowbreak{}wilson-\allowbreak{}spatial-\allowbreak{}inputs/\allowbreak{}G19\_\allowbreak{}DYNAMIC\_\allowbreak{}FIBER\_\allowbreak{}COVARIANCE\_\allowbreak{}AND\_\allowbreak{}CUBIC\_\allowbreak{}ENERGY\_\allowbreak{}20260906.md} \\
{\scriptsize\ttfamily SEPT\_\allowbreak{}DOC\_\allowbreak{}G19\_\allowbreak{}FULL\_\allowbreak{}GAUSSIAN\_\allowbreak{}FAST\_\allowbreak{}GREEN\_\allowbreak{}20260906} & {\scriptsize\ttfamily docs/\allowbreak{}derivations/\allowbreak{}wilson-\allowbreak{}spatial-\allowbreak{}inputs/\allowbreak{}G19\_\allowbreak{}FULL\_\allowbreak{}GAUSSIAN\_\allowbreak{}FAST\_\allowbreak{}GREEN\_\allowbreak{}20260906.md} \\
{\scriptsize\ttfamily SEPT\_\allowbreak{}DOC\_\allowbreak{}G19\_\allowbreak{}GAUSSIAN\_\allowbreak{}PATH\_\allowbreak{}ENDPOINT\_\allowbreak{}BASELINE\_\allowbreak{}20260906} & {\scriptsize\ttfamily docs/\allowbreak{}derivations/\allowbreak{}wilson-\allowbreak{}spatial-\allowbreak{}inputs/\allowbreak{}G19\_\allowbreak{}GAUSSIAN\_\allowbreak{}PATH\_\allowbreak{}ENDPOINT\_\allowbreak{}BASELINE\_\allowbreak{}20260906.md} \\
{\scriptsize\ttfamily SEPT\_\allowbreak{}DOC\_\allowbreak{}G19\_\allowbreak{}LOCAL\_\allowbreak{}COVARIANT\_\allowbreak{}AVERAGED\_\allowbreak{}PATH\_\allowbreak{}SOURCES\_\allowbreak{}20260905} & {\scriptsize\ttfamily docs/\allowbreak{}derivations/\allowbreak{}wilson-\allowbreak{}spatial-\allowbreak{}inputs/\allowbreak{}G19\_\allowbreak{}LOCAL\_\allowbreak{}COVARIANT\_\allowbreak{}AVERAGED\_\allowbreak{}PATH\_\allowbreak{}SOURCES\_\allowbreak{}20260905.md} \\
{\scriptsize\ttfamily SEPT\_\allowbreak{}DOC\_\allowbreak{}G19\_\allowbreak{}WILSON\_\allowbreak{}FINITE\_\allowbreak{}CELL\_\allowbreak{}GAP\_\allowbreak{}AND\_\allowbreak{}BOUNDARY\_\allowbreak{}FORM\_\allowbreak{}20260905} & {\scriptsize\ttfamily docs/\allowbreak{}derivations/\allowbreak{}wilson-\allowbreak{}spatial-\allowbreak{}inputs/\allowbreak{}G19\_\allowbreak{}WILSON\_\allowbreak{}FINITE\_\allowbreak{}CELL\_\allowbreak{}GAP\_\allowbreak{}AND\_\allowbreak{}BOUNDARY\_\allowbreak{}FORM\_\allowbreak{}20260905.md} \\
{\scriptsize\ttfamily SEPT\_\allowbreak{}DOC\_\allowbreak{}G19\_\allowbreak{}WILSON\_\allowbreak{}HARMONIC\_\allowbreak{}CUBIC\_\allowbreak{}OBSTRUCTION\_\allowbreak{}20260906} & {\scriptsize\ttfamily docs/\allowbreak{}derivations/\allowbreak{}wilson-\allowbreak{}spatial-\allowbreak{}inputs/\allowbreak{}G19\_\allowbreak{}WILSON\_\allowbreak{}HARMONIC\_\allowbreak{}CUBIC\_\allowbreak{}OBSTRUCTION\_\allowbreak{}20260906.md} \\
{\scriptsize\ttfamily SEPT\_\allowbreak{}DOC\_\allowbreak{}MANIFEST} & {\scriptsize\ttfamily docs/\allowbreak{}derivations/\allowbreak{}wilson-\allowbreak{}shell-\allowbreak{}inputs/\allowbreak{}manifest.json} \\
{\scriptsize\ttfamily SEPT\_\allowbreak{}DOC\_\allowbreak{}MANIFEST\_\allowbreak{}6a23b431} & {\scriptsize\ttfamily docs/\allowbreak{}derivations/\allowbreak{}wilson-\allowbreak{}spatial-\allowbreak{}inputs/\allowbreak{}manifest.json} \\
{\scriptsize\ttfamily SEPT\_\allowbreak{}DOC\_\allowbreak{}NEXT\_\allowbreak{}PATH\_\allowbreak{}WILSON\_\allowbreak{}TRANSFER\_\allowbreak{}2026\_\allowbreak{}09\_\allowbreak{}08} & {\scriptsize\ttfamily docs/\allowbreak{}research/\allowbreak{}next-\allowbreak{}path-\allowbreak{}wilson-\allowbreak{}transfer-\allowbreak{}2026-\allowbreak{}09-\allowbreak{}08.md} \\
{\scriptsize\ttfamily SEPT\_\allowbreak{}DOC\_\allowbreak{}NEXT\_\allowbreak{}PATH\_\allowbreak{}WILSON\_\allowbreak{}TRANSFER\_\allowbreak{}2026\_\allowbreak{}09\_\allowbreak{}08\_\allowbreak{}09f54613} & {\scriptsize\ttfamily docs/\allowbreak{}validation/\allowbreak{}next-\allowbreak{}path-\allowbreak{}wilson-\allowbreak{}transfer-\allowbreak{}2026-\allowbreak{}09-\allowbreak{}08.json} \\
{\scriptsize\ttfamily SEPT\_\allowbreak{}DOC\_\allowbreak{}PBH\_\allowbreak{}WILSON\_\allowbreak{}2026\_\allowbreak{}09\_\allowbreak{}08} & {\scriptsize\ttfamily docs/\allowbreak{}validation/\allowbreak{}pbh-\allowbreak{}wilson-\allowbreak{}2026-\allowbreak{}09-\allowbreak{}08.json} \\
{\scriptsize\ttfamily SEPT\_\allowbreak{}DOC\_\allowbreak{}README} & {\scriptsize\ttfamily docs/\allowbreak{}derivations/\allowbreak{}wilson-\allowbreak{}shell-\allowbreak{}inputs/\allowbreak{}README.md} \\
{\scriptsize\ttfamily SEPT\_\allowbreak{}DOC\_\allowbreak{}README\_\allowbreak{}8ba09c37} & {\scriptsize\ttfamily docs/\allowbreak{}derivations/\allowbreak{}wilson-\allowbreak{}spatial-\allowbreak{}inputs/\allowbreak{}README.md} \\
{\scriptsize\ttfamily SEPT\_\allowbreak{}DOC\_\allowbreak{}SPATIAL\_\allowbreak{}CONTINUUM\_\allowbreak{}PASSAGE\_\allowbreak{}2026\_\allowbreak{}09\_\allowbreak{}08} & {\scriptsize\ttfamily docs/\allowbreak{}research/\allowbreak{}spatial-\allowbreak{}continuum-\allowbreak{}passage-\allowbreak{}2026-\allowbreak{}09-\allowbreak{}08.md} \\
{\scriptsize\ttfamily SEPT\_\allowbreak{}DOC\_\allowbreak{}WILSON\_\allowbreak{}BACKGROUND\_\allowbreak{}2026\_\allowbreak{}09\_\allowbreak{}09} & {\scriptsize\ttfamily docs/\allowbreak{}validation/\allowbreak{}wilson-\allowbreak{}background-\allowbreak{}2026-\allowbreak{}09-\allowbreak{}09.json} \\
{\scriptsize\ttfamily SEPT\_\allowbreak{}DOC\_\allowbreak{}WILSON\_\allowbreak{}BACKGROUND\_\allowbreak{}INDEPENDENT\_\allowbreak{}REVIEW} & {\scriptsize\ttfamily docs/\allowbreak{}validation/\allowbreak{}wilson-\allowbreak{}background-\allowbreak{}independent-\allowbreak{}review.md} \\
{\scriptsize\ttfamily SEPT\_\allowbreak{}DOC\_\allowbreak{}WILSON\_\allowbreak{}G19\_\allowbreak{}CAUCHY\_\allowbreak{}REPAIR} & {\scriptsize\ttfamily docs/\allowbreak{}validation/\allowbreak{}wilson-\allowbreak{}g19-\allowbreak{}cauchy-\allowbreak{}repair.md} \\
{\scriptsize\ttfamily SEPT\_\allowbreak{}DOC\_\allowbreak{}WILSON\_\allowbreak{}G19\_\allowbreak{}CORPUS\_\allowbreak{}INVENTORY} & {\scriptsize\ttfamily docs/\allowbreak{}validation/\allowbreak{}wilson-\allowbreak{}g19-\allowbreak{}corpus-\allowbreak{}inventory.json} \\
{\scriptsize\ttfamily SEPT\_\allowbreak{}DOC\_\allowbreak{}WILSON\_\allowbreak{}G19\_\allowbreak{}CORPUS\_\allowbreak{}REVIEW\_\allowbreak{}MANIFEST} & {\scriptsize\ttfamily docs/\allowbreak{}validation/\allowbreak{}wilson-\allowbreak{}g19-\allowbreak{}corpus-\allowbreak{}review-\allowbreak{}manifest.json} \\
{\scriptsize\ttfamily SEPT\_\allowbreak{}DOC\_\allowbreak{}WILSON\_\allowbreak{}G19\_\allowbreak{}FIBERS\_\allowbreak{}REVIEW} & {\scriptsize\ttfamily docs/\allowbreak{}validation/\allowbreak{}wilson-\allowbreak{}g19-\allowbreak{}fibers-\allowbreak{}review.md} \\
{\scriptsize\ttfamily SEPT\_\allowbreak{}DOC\_\allowbreak{}WILSON\_\allowbreak{}G19\_\allowbreak{}GAUSSIAN\_\allowbreak{}REVIEW} & {\scriptsize\ttfamily docs/\allowbreak{}validation/\allowbreak{}wilson-\allowbreak{}g19-\allowbreak{}gaussian-\allowbreak{}review.md} \\
{\scriptsize\ttfamily SEPT\_\allowbreak{}DOC\_\allowbreak{}WILSON\_\allowbreak{}G19\_\allowbreak{}MARTINGALE\_\allowbreak{}REPAIR} & {\scriptsize\ttfamily docs/\allowbreak{}validation/\allowbreak{}wilson-\allowbreak{}g19-\allowbreak{}martingale-\allowbreak{}repair.md} \\
{\scriptsize\ttfamily SEPT\_\allowbreak{}DOC\_\allowbreak{}WILSON\_\allowbreak{}G19\_\allowbreak{}REVIEW\_\allowbreak{}2026\_\allowbreak{}09\_\allowbreak{}09} & {\scriptsize\ttfamily docs/\allowbreak{}validation/\allowbreak{}wilson-\allowbreak{}g19-\allowbreak{}review-\allowbreak{}2026-\allowbreak{}09-\allowbreak{}09.json} \\
{\scriptsize\ttfamily SEPT\_\allowbreak{}DOC\_\allowbreak{}WILSON\_\allowbreak{}G19\_\allowbreak{}SOURCES\_\allowbreak{}REVIEW} & {\scriptsize\ttfamily docs/\allowbreak{}validation/\allowbreak{}wilson-\allowbreak{}g19-\allowbreak{}sources-\allowbreak{}review.md} \\
{\scriptsize\ttfamily SEPT\_\allowbreak{}DOC\_\allowbreak{}WILSON\_\allowbreak{}MARKED\_\allowbreak{}2026\_\allowbreak{}09\_\allowbreak{}08} & {\scriptsize\ttfamily docs/\allowbreak{}validation/\allowbreak{}wilson-\allowbreak{}marked-\allowbreak{}2026-\allowbreak{}09-\allowbreak{}08.json} \\
{\scriptsize\ttfamily SEPT\_\allowbreak{}DOC\_\allowbreak{}WILSON\_\allowbreak{}MARKED\_\allowbreak{}AXIOMS} & {\scriptsize\ttfamily docs/\allowbreak{}validation/\allowbreak{}wilson-\allowbreak{}marked-\allowbreak{}axioms.lean} \\
{\scriptsize\ttfamily SEPT\_\allowbreak{}DOC\_\allowbreak{}WILSON\_\allowbreak{}SC17\_\allowbreak{}2026\_\allowbreak{}09\_\allowbreak{}09} & {\scriptsize\ttfamily docs/\allowbreak{}validation/\allowbreak{}wilson-\allowbreak{}sc17-\allowbreak{}2026-\allowbreak{}09-\allowbreak{}09.json} \\
{\scriptsize\ttfamily SEPT\_\allowbreak{}DOC\_\allowbreak{}WILSON\_\allowbreak{}SC17\_\allowbreak{}BC2\_\allowbreak{}UPDATE\_\allowbreak{}AUDIT} & {\scriptsize\ttfamily docs/\allowbreak{}validation/\allowbreak{}wilson-\allowbreak{}sc17-\allowbreak{}bc2-\allowbreak{}update-\allowbreak{}audit.md} \\
{\scriptsize\ttfamily SEPT\_\allowbreak{}DOC\_\allowbreak{}WILSON\_\allowbreak{}SC17\_\allowbreak{}INDEPENDENT\_\allowbreak{}REVIEW} & {\scriptsize\ttfamily docs/\allowbreak{}validation/\allowbreak{}wilson-\allowbreak{}sc17-\allowbreak{}independent-\allowbreak{}review.md} \\
{\scriptsize\ttfamily SEPT\_\allowbreak{}DOC\_\allowbreak{}WILSON\_\allowbreak{}SC17\_\allowbreak{}LATE\_\allowbreak{}LEDGER\_\allowbreak{}AUDIT} & {\scriptsize\ttfamily docs/\allowbreak{}validation/\allowbreak{}wilson-\allowbreak{}sc17-\allowbreak{}late-\allowbreak{}ledger-\allowbreak{}audit.md} \\
{\scriptsize\ttfamily SEPT\_\allowbreak{}DOC\_\allowbreak{}WILSON\_\allowbreak{}SELECTED\_\allowbreak{}2026\_\allowbreak{}09\_\allowbreak{}08} & {\scriptsize\ttfamily docs/\allowbreak{}validation/\allowbreak{}wilson-\allowbreak{}selected-\allowbreak{}2026-\allowbreak{}09-\allowbreak{}08.json} \\
{\scriptsize\ttfamily SEPT\_\allowbreak{}DOC\_\allowbreak{}WILSON\_\allowbreak{}SHELL\_\allowbreak{}2026\_\allowbreak{}09\_\allowbreak{}08} & {\scriptsize\ttfamily docs/\allowbreak{}validation/\allowbreak{}wilson-\allowbreak{}shell-\allowbreak{}2026-\allowbreak{}09-\allowbreak{}08.json} \\
{\scriptsize\ttfamily SEPT\_\allowbreak{}DOC\_\allowbreak{}WILSON\_\allowbreak{}SHELL\_\allowbreak{}AXIOMS} & {\scriptsize\ttfamily docs/\allowbreak{}validation/\allowbreak{}wilson-\allowbreak{}shell-\allowbreak{}axioms.lean} \\
{\scriptsize\ttfamily SEPT\_\allowbreak{}DOC\_\allowbreak{}WILSON\_\allowbreak{}SPATIAL\_\allowbreak{}2026\_\allowbreak{}09\_\allowbreak{}08} & {\scriptsize\ttfamily docs/\allowbreak{}validation/\allowbreak{}wilson-\allowbreak{}spatial-\allowbreak{}2026-\allowbreak{}09-\allowbreak{}08.json} \\
{\scriptsize\ttfamily SEPT\_\allowbreak{}DOC\_\allowbreak{}WILSON\_\allowbreak{}SPATIAL\_\allowbreak{}AXIOMS} & {\scriptsize\ttfamily docs/\allowbreak{}validation/\allowbreak{}wilson-\allowbreak{}spatial-\allowbreak{}axioms.lean} \\
{\scriptsize\ttfamily SEPT\_\allowbreak{}DOC\_\allowbreak{}WILSON\_\allowbreak{}TRUE\_\allowbreak{}BLOCKS\_\allowbreak{}2026\_\allowbreak{}09\_\allowbreak{}09} & {\scriptsize\ttfamily docs/\allowbreak{}validation/\allowbreak{}wilson-\allowbreak{}true-\allowbreak{}blocks-\allowbreak{}2026-\allowbreak{}09-\allowbreak{}09.json} \\
{\scriptsize\ttfamily SEPT\_\allowbreak{}DOC\_\allowbreak{}WILSON\_\allowbreak{}TRUE\_\allowbreak{}BLOCKS\_\allowbreak{}CHECK\_\allowbreak{}2026\_\allowbreak{}09\_\allowbreak{}09} & {\scriptsize\ttfamily docs/\allowbreak{}validation/\allowbreak{}wilson-\allowbreak{}true-\allowbreak{}blocks-\allowbreak{}check-\allowbreak{}2026-\allowbreak{}09-\allowbreak{}09.json} \\
{\scriptsize\ttfamily SEPT\_\allowbreak{}DOC\_\allowbreak{}WILSON\_\allowbreak{}VACUUM\_\allowbreak{}2026\_\allowbreak{}09\_\allowbreak{}09} & {\scriptsize\ttfamily docs/\allowbreak{}validation/\allowbreak{}wilson-\allowbreak{}vacuum-\allowbreak{}2026-\allowbreak{}09-\allowbreak{}09.json} \\
{\scriptsize\ttfamily SEPT\_\allowbreak{}DOC\_\allowbreak{}WILSON\_\allowbreak{}VACUUM\_\allowbreak{}LEDGER\_\allowbreak{}SNAPSHOT} & {\scriptsize\ttfamily docs/\allowbreak{}validation/\allowbreak{}wilson-\allowbreak{}vacuum-\allowbreak{}ledger-\allowbreak{}snapshot.yaml} \\
{\scriptsize\ttfamily SEPT\_\allowbreak{}DOC\_\allowbreak{}WILSON\_\allowbreak{}VACUUM\_\allowbreak{}THEOREMS\_\allowbreak{}SNAPSHOT} & {\scriptsize\ttfamily docs/\allowbreak{}validation/\allowbreak{}wilson-\allowbreak{}vacuum-\allowbreak{}theorems-\allowbreak{}snapshot.yaml} \\
{\scriptsize\ttfamily SEPT\_\allowbreak{}DOC\_\allowbreak{}WILSON\_\allowbreak{}WEIGHTED\_\allowbreak{}2026\_\allowbreak{}09\_\allowbreak{}08} & {\scriptsize\ttfamily docs/\allowbreak{}validation/\allowbreak{}wilson-\allowbreak{}weighted-\allowbreak{}2026-\allowbreak{}09-\allowbreak{}08.json} \\
{\scriptsize\ttfamily SEPT\_\allowbreak{}DOC\_\allowbreak{}YANGMILLS\_\allowbreak{}DERIVATIONS\_\allowbreak{}2026\_\allowbreak{}09\_\allowbreak{}08} & {\scriptsize\ttfamily docs/\allowbreak{}validation/\allowbreak{}yangmills-\allowbreak{}derivations-\allowbreak{}2026-\allowbreak{}09-\allowbreak{}08.json} \\
{\scriptsize\ttfamily SEPT\_\allowbreak{}DOC\_\allowbreak{}YANGMILLS\_\allowbreak{}FORMAL\_\allowbreak{}BRIDGES} & {\scriptsize\ttfamily docs/\allowbreak{}derivations/\allowbreak{}yangmills-\allowbreak{}formal-\allowbreak{}bridges.md} \\
{\scriptsize\ttfamily SEPT\_\allowbreak{}DOC\_\allowbreak{}YANGMILLS\_\allowbreak{}GPU\_\allowbreak{}AUDIT\_\allowbreak{}2026\_\allowbreak{}09\_\allowbreak{}09} & {\scriptsize\ttfamily docs/\allowbreak{}validation/\allowbreak{}yangmills-\allowbreak{}gpu-\allowbreak{}audit-\allowbreak{}2026-\allowbreak{}09-\allowbreak{}09.json} \\
{\scriptsize\ttfamily SEPT\_\allowbreak{}DOC\_\allowbreak{}YANGMILLS\_\allowbreak{}SOURCES\_\allowbreak{}2026\_\allowbreak{}09\_\allowbreak{}08} & {\scriptsize\ttfamily docs/\allowbreak{}validation/\allowbreak{}yangmills-\allowbreak{}sources-\allowbreak{}2026-\allowbreak{}09-\allowbreak{}08.json} \\
{\scriptsize\ttfamily TRACK\_\allowbreak{}A\_\allowbreak{}SUPPORTED} & {\scriptsize\ttfamily docs/\allowbreak{}derivations/\allowbreak{}track-\allowbreak{}a-\allowbreak{}supported-\allowbreak{}statements.md} \\
{\scriptsize\ttfamily UNIFIED} & {\scriptsize\ttfamily theory/\allowbreak{}MASTER\_\allowbreak{}THEORY\_\allowbreak{}UNIFIED\_\allowbreak{}2026-\allowbreak{}08-\allowbreak{}20\_\allowbreak{}v4\_\allowbreak{}3.md} \\
{\scriptsize\ttfamily WILSON\_\allowbreak{}ACTIVITY\_\allowbreak{}INPUT} & {\scriptsize\ttfamily docs/\allowbreak{}derivations/\allowbreak{}wilson-\allowbreak{}shell-\allowbreak{}inputs/\allowbreak{}G18\_\allowbreak{}WILSON\_\allowbreak{}WEIGHTED\_\allowbreak{}ACTIVITY\_\allowbreak{}BOUND\_\allowbreak{}20260905.md} \\
{\scriptsize\ttfamily WILSON\_\allowbreak{}BAND\_\allowbreak{}INPUT} & {\scriptsize\ttfamily docs/\allowbreak{}derivations/\allowbreak{}wilson-\allowbreak{}shell-\allowbreak{}inputs/\allowbreak{}G18\_\allowbreak{}WILSON\_\allowbreak{}INFINITE\_\allowbreak{}VOLUME\_\allowbreak{}PHYSICAL\_\allowbreak{}BAND\_\allowbreak{}20260905.md} \\
{\scriptsize\ttfamily WILSON\_\allowbreak{}CHART\_\allowbreak{}INPUT} & {\scriptsize\ttfamily docs/\allowbreak{}derivations/\allowbreak{}wilson-\allowbreak{}shell-\allowbreak{}inputs/\allowbreak{}G18\_\allowbreak{}WILSON\_\allowbreak{}CARDINALITY\_\allowbreak{}UNITARY\_\allowbreak{}CHART\_\allowbreak{}20260905.md} \\
{\scriptsize\ttfamily WILSON\_\allowbreak{}SELECTED} & {\scriptsize\ttfamily docs/\allowbreak{}derivations/\allowbreak{}wilson-\allowbreak{}selected-\allowbreak{}inverse-\allowbreak{}wall.md} \\
{\scriptsize\ttfamily WILSON\_\allowbreak{}SHELL} & {\scriptsize\ttfamily docs/\allowbreak{}derivations/\allowbreak{}wilson-\allowbreak{}marked-\allowbreak{}shell-\allowbreak{}transport.md} \\
{\scriptsize\ttfamily WILSON\_\allowbreak{}SPATIAL} & {\scriptsize\ttfamily docs/\allowbreak{}derivations/\allowbreak{}wilson-\allowbreak{}spatial-\allowbreak{}schur-\allowbreak{}excess.md} \\
{\scriptsize\ttfamily WILSON\_\allowbreak{}SPATIAL\_\allowbreak{}FLAT\_\allowbreak{}INPUT} & {\scriptsize\ttfamily docs/\allowbreak{}derivations/\allowbreak{}wilson-\allowbreak{}spatial-\allowbreak{}inputs/\allowbreak{}G19\_\allowbreak{}FLAT\_\allowbreak{}HOLONOMY\_\allowbreak{}SOURCES\_\allowbreak{}AND\_\allowbreak{}RANK\_\allowbreak{}REPAIR\_\allowbreak{}20260907.md} \\
{\scriptsize\ttfamily WILSON\_\allowbreak{}SPATIAL\_\allowbreak{}SCHUR\_\allowbreak{}INPUT} & {\scriptsize\ttfamily docs/\allowbreak{}derivations/\allowbreak{}wilson-\allowbreak{}spatial-\allowbreak{}inputs/\allowbreak{}G19\_\allowbreak{}FORM\_\allowbreak{}SCHUR\_\allowbreak{}SCALE\_\allowbreak{}COMPARISON\_\allowbreak{}20260905.md} \\
{\scriptsize\ttfamily WILSON\_\allowbreak{}TRUE\_\allowbreak{}BLOCK} & {\scriptsize\ttfamily docs/\allowbreak{}derivations/\allowbreak{}wilson-\allowbreak{}true-\allowbreak{}vacuum-\allowbreak{}block-\allowbreak{}estimates.md} \\
{\scriptsize\ttfamily WILSON\_\allowbreak{}VACUUM\_\allowbreak{}ASSEMBLY} & {\scriptsize\ttfamily docs/\allowbreak{}derivations/\allowbreak{}wilson-\allowbreak{}vacuum-\allowbreak{}aligned-\allowbreak{}assembly.md} \\
{\scriptsize\ttfamily WILSON\_\allowbreak{}WEIGHTED\_\allowbreak{}REPAIR} & {\scriptsize\ttfamily docs/\allowbreak{}derivations/\allowbreak{}wilson-\allowbreak{}weighted-\allowbreak{}repair-\allowbreak{}and-\allowbreak{}rotor-\allowbreak{}gap.md} \\
{\scriptsize\ttfamily YM\_\allowbreak{}GPU\_\allowbreak{}AUDIT} & {\scriptsize\ttfamily docs/\allowbreak{}derivations/\allowbreak{}yangmills-\allowbreak{}gpu-\allowbreak{}resolution-\allowbreak{}audit.md} \\
{\scriptsize\ttfamily v4.3} & {\scriptsize\ttfamily theory/\allowbreak{}MASTER\_\allowbreak{}THEORY\_\allowbreak{}UNIFIED\_\allowbreak{}2026-\allowbreak{}08-\allowbreak{}20\_\allowbreak{}v4\_\allowbreak{}3.md} \\
\end{longtable}

\subsection*{Standing: corpus (8)}
\addcontentsline{toc}{subsection}{Standing: corpus}
\begin{longtable}{@{}p{2.0in} p{4.2in}@{}}
\toprule Alias & Path \\ \midrule
\endfirsthead
\toprule Alias & Path \\ \midrule \endhead
\bottomrule\endfoot
{\scriptsize\ttfamily 818} & {\scriptsize\ttfamily corpus-\allowbreak{}import/\allowbreak{}records/\allowbreak{}transcripts/\allowbreak{}818.txt} \\
{\scriptsize\ttfamily GCSG} & {\scriptsize\ttfamily corpus-\allowbreak{}import/\allowbreak{}papers/\allowbreak{}gauge\_\allowbreak{}constrained\_\allowbreak{}spectral\_\allowbreak{}geometry/\allowbreak{}PAPER\_\allowbreak{}FLUX\_\allowbreak{}gauge\_\allowbreak{}constrained\_\allowbreak{}spectral\_\allowbreak{}geometry\_\allowbreak{}unified\_\allowbreak{}v1\_\allowbreak{}4\_\allowbreak{}2026-\allowbreak{}08-\allowbreak{}08.md} \\
{\scriptsize\ttfamily PAPER} & {\scriptsize\ttfamily corpus-\allowbreak{}import/\allowbreak{}papers/\allowbreak{}flat\_\allowbreak{}band/\allowbreak{}PAPER\_\allowbreak{}FLUX\_\allowbreak{}glueball\_\allowbreak{}flat\_\allowbreak{}band\_\allowbreak{}v1\_\allowbreak{}1.tex} \\
{\scriptsize\ttfamily THM\_\allowbreak{}FLUX} & {\scriptsize\ttfamily corpus-\allowbreak{}import/\allowbreak{}theory/\allowbreak{}notes/\allowbreak{}THM\_\allowbreak{}FLUX\_\allowbreak{}hodge\_\allowbreak{}cellular\_\allowbreak{}circuit\_\allowbreak{}mobility\_\allowbreak{}theorem.md} \\
{\scriptsize\ttfamily TRACK\_\allowbreak{}A\_\allowbreak{}RECEIVED\_\allowbreak{}G17} & {\scriptsize\ttfamily runs/\allowbreak{}track\_\allowbreak{}a\_\allowbreak{}formalization\_\allowbreak{}2026-\allowbreak{}09-\allowbreak{}09/\allowbreak{}sources/\allowbreak{}paper/\allowbreak{}research\_\allowbreak{}notes/\allowbreak{}G17\_\allowbreak{}HAMILTONIAN\_\allowbreak{}OSTERWALDER\_\allowbreak{}SEILER\_\allowbreak{}TRANSCRIPTION\_\allowbreak{}20260910.md} \\
{\scriptsize\ttfamily TRACK\_\allowbreak{}A\_\allowbreak{}RECEIVED\_\allowbreak{}SC17} & {\scriptsize\ttfamily runs/\allowbreak{}track\_\allowbreak{}a\_\allowbreak{}formalization\_\allowbreak{}2026-\allowbreak{}09-\allowbreak{}09/\allowbreak{}sources/\allowbreak{}docs/\allowbreak{}derivations/\allowbreak{}wilson-\allowbreak{}sc17-\allowbreak{}spatial-\allowbreak{}closure.md} \\
{\scriptsize\ttfamily TRACK\_\allowbreak{}A\_\allowbreak{}RECEIVED\_\allowbreak{}W6} & {\scriptsize\ttfamily runs/\allowbreak{}track\_\allowbreak{}a\_\allowbreak{}formalization\_\allowbreak{}2026-\allowbreak{}09-\allowbreak{}09/\allowbreak{}sources/\allowbreak{}paper/\allowbreak{}research\_\allowbreak{}notes/\allowbreak{}G19\_\allowbreak{}W6\_\allowbreak{}CUBIC\_\allowbreak{}RESIDUAL\_\allowbreak{}BOUNDS\_\allowbreak{}20260910.md} \\
{\scriptsize\ttfamily transcript} & {\scriptsize\ttfamily corpus-\allowbreak{}import/\allowbreak{}records/\allowbreak{}transcripts/\allowbreak{}\#-\allowbreak{}Final-\allowbreak{}unified-\allowbreak{}theory.txt} \\
\end{longtable}

\subsection*{Standing: unrecorded (3)}
\addcontentsline{toc}{subsection}{Standing: unrecorded}
\begin{longtable}{@{}p{2.0in} p{4.2in}@{}}
\toprule Alias & Path \\ \midrule
\endfirsthead
\toprule Alias & Path \\ \midrule \endhead
\bottomrule\endfoot
{\scriptsize\ttfamily CANON} & {\scriptsize\ttfamily } \\
{\scriptsize\ttfamily MOB} & {\scriptsize\ttfamily } \\
{\scriptsize\ttfamily PUB edition} & {\scriptsize\ttfamily } \\
\end{longtable}

\subsection*{Standing: quarantined (3)}
\addcontentsline{toc}{subsection}{Standing: quarantined}
\begin{longtable}{@{}p{2.0in} p{4.2in}@{}}
\toprule Alias & Path \\ \midrule
\endfirsthead
\toprule Alias & Path \\ \midrule \endhead
\bottomrule\endfoot
{\scriptsize\ttfamily SEPT\_\allowbreak{}DOC\_\allowbreak{}YANGMILLS\_\allowbreak{}RESOLVENT\_\allowbreak{}LOCALIZATION\_\allowbreak{}AND\_\allowbreak{}DIRICHLET\_\allowbreak{}GAP} & {\scriptsize\ttfamily docs/\allowbreak{}derivations/\allowbreak{}yangmills-\allowbreak{}resolvent-\allowbreak{}localization-\allowbreak{}and-\allowbreak{}dirichlet-\allowbreak{}gap.md} \\
{\scriptsize\ttfamily YM\_\allowbreak{}BALABAN\_\allowbreak{}MULTISCALE} & {\scriptsize\ttfamily docs/\allowbreak{}derivations/\allowbreak{}yangmills-\allowbreak{}continuum-\allowbreak{}balaban-\allowbreak{}multiscale-\allowbreak{}proof.md} \\
{\scriptsize\ttfamily quarantined master v3} & {\scriptsize\ttfamily corpus-\allowbreak{}import/\allowbreak{}corpus/\allowbreak{}MASTER\_\allowbreak{}THEORY\_\allowbreak{}UNIFIED\_\allowbreak{}2026-\allowbreak{}08-\allowbreak{}20\_\allowbreak{}v3.md} \\
\end{longtable}

\subsection*{Standing: superseded (1)}
\addcontentsline{toc}{subsection}{Standing: superseded}
\begin{longtable}{@{}p{2.0in} p{4.2in}@{}}
\toprule Alias & Path \\ \midrule
\endfirsthead
\toprule Alias & Path \\ \midrule \endhead
\bottomrule\endfoot
{\scriptsize\ttfamily MASTER\_\allowbreak{}THEORY} & {\scriptsize\ttfamily theory/\allowbreak{}superseded/\allowbreak{}MASTER\_\allowbreak{}THEORY.md} \\
\end{longtable}

% ---- end gen_documents.tex -----------------------------------------
% ---- end B5_documents.tex ------------------------------------------

% ---- begin inlined bibliography ----------------------------------
\begin{thebibliography}{100}

\bibitem{AIZENMAN_1982}
Michael Aizenman.
\newblock Geometric analysis of $\varphi$$^4$ fields and ising models. parts i
  and ii.
\newblock {\em Communications in Mathematical Physics 86, 1-48}, 1982.

\bibitem{AT_2020_SU3}
Andreas Athenodorou and Michael Teper.
\newblock The glueball spectrum of su(3) gauge theory in 3+1 dimensions.
\newblock {\em JHEP 11 (2020) 172}, 2020.

\bibitem{AT_2021_SUN}
Andreas Athenodorou and Michael Teper.
\newblock Su(n) gauge theories in 3+1 dimensions: glueball spectrum, string
  tensions and topology.
\newblock {\em JHEP 12 (2021) 082}, 2021.

\bibitem{BV_1981_GEOMETRY}
O.~Babelon and C.~M. Viallet.
\newblock The riemannian geometry of the configuration space of gauge theories.
\newblock {\em Communications in Mathematical Physics 81(4), 515--525}, 1981.

\bibitem{BALABAN_1985}
T.~Balaban.
\newblock Ultraviolet stability of three-dimensional lattice pure gauge field
  theories.
\newblock {\em Communications in Mathematical Physics 102, 255-275}, 1985.

\bibitem{BALABAN_1987}
Tadeusz Ba\l{}aban.
\newblock Renormalization group approach to lattice gauge field theories. i.
  generation of effective actions in a small field approximation and a coupling
  constant renormalization in four dimensions.
\newblock {\em Communications in Mathematical Physics 109, 249-301}, 1987.

\bibitem{BBIJ_1984}
Tadeusz Balaban, John Imbrie, Arthur Jaffe, and David Brydges.
\newblock The mass gap for higgs models on a unit lattice.
\newblock {\em Annals of Physics 158, 281-319}, 1984.

\bibitem{BALAJI_2026}
S.~Balaji and et~al.
\newblock Quantum simulation of su(3) lattice yang-mills theory: one-cube
  benchmarks and circuit resource estimates, 2026.
\newblock arXiv preprint.

\bibitem{BRS_1976}
C~Becchi, A~Rouet, and R~Stora.
\newblock Renormalization of gauge theories.
\newblock {\em Annals of Physics 98, 287-321}, 1976.

\bibitem{BB_1983}
B.~Berg and A.~Billoire.
\newblock Glueball spectroscopy in 4d su(3) lattice gauge theory (i).
\newblock {\em Nucl. Phys. B221, 109-140}, 1983.

\bibitem{BORGA_2024}
Jacopo Borga, Sky Cao, and Jasper Shogren-Knaak.
\newblock Surface sums for lattice yang-mills in the large-n limit, 2024.
\newblock arXiv:2411.11676 [math.PR].

\bibitem{BFS_1979_I}
David Brydges, J\"urg Fr\"ohlich, and Erhard Seiler.
\newblock On the construction of quantized gauge fields. i. general results.
\newblock {\em Annals of Physics 121, 227-284}, 1979.

\bibitem{BF_1980}
David~C. Brydges and Paul Federbush.
\newblock Debye screening.
\newblock {\em Communications in Mathematical Physics 73, 197-246}, 1980.

\bibitem{BFS_1980_II}
David~C. Brydges, J\"urg Fr\"ohlich, and Erhard Seiler.
\newblock Construction of quantised gauge fields. ii. convergence of the
  lattice approximation.
\newblock {\em Communications in Mathematical Physics 71, 159-205}, 1980.

\bibitem{BFS_1981_III}
David~C. Brydges, J\"urg Fr\"ohlich, and Erhard Seiler.
\newblock On the construction of quantized gauge fields. iii. the
  two-dimensional abelian higgs model without cutoffs.
\newblock {\em Communications in Mathematical Physics 79, 353-399}, 1981.

\bibitem{CAO_2023}
Sky Cao, Minjae Park, and Scott Sheffield.
\newblock Random surfaces and lattice yang-mills.
\newblock {\em Commun. Amer. Math. Soc. 5 (2025) 774-896}, 2023.

\bibitem{CM_2003}
Jesse Carlsson and Bruce H.~J. McKellar.
\newblock Su(n) lattice gauge theory on a single cube, 2003.
\newblock arXiv preprint (unpublished).

\bibitem{CHEN_2006}
Y.~Chen, A.~Alexandru, S.J. Dong, T.~Draper, I.~Horvath, F.X. Lee, K.F. Liu,
  N.~Mathur, C.~Morningstar, M.~Peardon, S.~Tamhankar, B.L. Young, and J.B.
  Zhang.
\newblock Glueball spectrum and matrix elements on anisotropic lattices.
\newblock {\em Phys. Rev. D 73, 014516}, 2006.

\bibitem{CKS_2021}
Anthony Ciavarella, Natalie Klco, and Martin~J. Savage.
\newblock Trailhead for quantum simulation of su(3) yang-mills lattice gauge
  theory in the local multiplet basis.
\newblock {\em Phys. Rev. D 103, 094501}, 2021.

\bibitem{CB_2024}
Anthony~N. Ciavarella and Christian~W. Bauer.
\newblock Quantum simulation of su(3) lattice yang-mills theory at leading
  order in large-n\_c expansion.
\newblock {\em Phys. Rev. Lett. 133, 111901}, 2024.

\bibitem{CBB_2026}
Anthony~N. Ciavarella, Ivan~M. Burbano, and Christian~W. Bauer.
\newblock Efficient truncations of su(n\_c) lattice gauge theory for quantum
  simulation, 2026.
\newblock arXiv preprint.

\bibitem{BESIII_2026}
BESIII Collaboration.
\newblock Lightest 0-+ glueball as dominant constituent of x(2370), 2026.
\newblock arXiv:2607.20366 [hep-ex].

\bibitem{CS_2006}
Benoit Collins and Piotr Sniady.
\newblock Integration with respect to the haar measure on unitary, orthogonal
  and symplectic group.
\newblock {\em Comm. Math. Phys. 264, 773-795}, 2006.

\bibitem{CREUTZ_1980}
Michael Creutz.
\newblock Monte carlo study of quantized su(2) gauge theory.
\newblock {\em Physical Review D 21, 2308-2315}, 1980.

\bibitem{DRS_2021}
Zohreh Davoudi, Indrakshi Raychowdhury, and Andrew Shaw.
\newblock Search for efficient formulations for hamiltonian simulation of
  non-abelian lattice gauge theories.
\newblock {\em Phys. Rev. D 104, 074505}, 2021.

\bibitem{FP_1967}
L.D. Faddeev and V.N. Popov.
\newblock Feynman diagrams for the yang-mills field.
\newblock {\em Physics Letters B 25, 29-30}, 1967.

\bibitem{FLPS_2026}
Andrea Falzetti, Matteo Lombardi, Mauro~Lucio Papinutto, and Francesco
  Scardino.
\newblock Update on the computation of the quenched su(6) yang-mills lattice
  spectrum, 2026.
\newblock PoS LATTICE2025 (to appear); arXiv:2603.13138.

\bibitem{FO_1976}
Joel~S Feldman and Konrad Osterwalder.
\newblock The wightman axioms and the mass gap for weakly coupled ($\phi$4)3
  quantum field theories.
\newblock {\em Annals of Physics 97, 80-135}, 1976.

\bibitem{FEYNMAN_1981}
Richard~P. Feynman.
\newblock The qualitative behavior of yang-mills theory in 2 + 1 dimensions.
\newblock {\em Nuclear Physics B 188, 479-512}, 1981.

\bibitem{FOS_1983}
J.~Frohlich, K.~Osterwalder, and E.~Seiler.
\newblock On virtual representations of symmetric spaces and their analytic
  continuation.
\newblock {\em Annals of Mathematics 118, 461--489}, 1983.

\bibitem{FROHLICH_1982}
J\"urg Fr\"ohlich.
\newblock On the triviality of $\lambda$$\varphi$$^4$\_d theories and the
  approach to the critical point in d $\geq$ 4 dimensions.
\newblock {\em Nuclear Physics B 200, 281-296}, 1982.

\bibitem{GAWEDZKI_1985}
K.~Gaw\k{e}dzki and A.~Kupiainen.
\newblock Renormalization of a non-renormalizable quantum field theory.
\newblock {\em Nuclear Physics B 262, 33-48}, 1985.

\bibitem{GINIBRE_1970}
J.~Ginibre.
\newblock General formulation of griffiths' inequalities.
\newblock {\em Commun. Math. Phys. 16, 310-328}, 1970.

\bibitem{GJ_1968_I}
James Glimm and Arthur Jaffe.
\newblock A $\lambda$$\varphi$$^4$ quantum field theory without cutoffs. i.
\newblock {\em Physical Review 176, 1945--1951}, 1968.

\bibitem{GJ_1970_II}
James Glimm and Arthur Jaffe.
\newblock The $\lambda$($\varphi$$^4$)$_2$ quantum field theory without
  cutoffs. ii. the field operators and the approximate vacuum.
\newblock {\em Annals of Mathematics 91, 362--401}, 1970.

\bibitem{GJ_1970_III}
James Glimm and Arthur Jaffe.
\newblock The $\lambda$($\varphi$$^4$)$_2$ quantum field theory without
  cutoffs. iii. the physical vacuum.
\newblock {\em Acta Mathematica 125, 203-267}, 1970.

\bibitem{GJ_1972_IV}
James Glimm and Arthur Jaffe.
\newblock The $\lambda$$\varphi$$^4$$_2$ quantum field theory without cutoffs.
  iv. perturbations of the hamiltonian.
\newblock {\em Journal of Mathematical Physics 13, 1568-1584}, 1972.

\bibitem{GJ_1973_POSITIVITY}
James Glimm and Arthur Jaffe.
\newblock Positivity of the $\varphi$$^4$$_3$ hamiltonian.
\newblock {\em Fortschritte der Physik 21, 327-376}, 1973.

\bibitem{GJ_1974_SINGLE_PHASE}
James Glimm and Arthur Jaffe.
\newblock $\varphi$$^4$$_2$ quantum field model in the single-phase region:
  Differentiability of the mass and bounds on critical exponents.
\newblock {\em Physical Review D 10, 536-539}, 1974.

\bibitem{GJ_1985_SELECTED_I}
James Glimm and Arthur Jaffe.
\newblock Collected papers vol. 1: Quantum field theory and statistical
  mechanics: Expositions.
\newblock {\em Birkh\"auser Boston, ISBN 978-0-8176-3271-7}, 1985.

\bibitem{GJ_1985_SELECTED_II}
James Glimm and Arthur Jaffe.
\newblock Collected papers vol. 2: Constructive quantum field theory: Selected
  papers.
\newblock {\em Birkh\"auser Boston, ISBN 978-0-8176-3272-4}, 1985.

\bibitem{GJ_1987_QP}
James Glimm and Arthur Jaffe.
\newblock Quantum physics: A functional integral point of view (second
  edition).
\newblock {\em Springer New York, second edition}, 1987.

\bibitem{GJS_1974}
James Glimm, Arthur Jaffe, and Thomas Spencer.
\newblock The wightman axioms and particle structure in the p($\varphi$)$_2$
  quantum field model.
\newblock {\em Annals of Mathematics 100, 585--632}, 1974.

\bibitem{GJS_1976_I}
James Glimm, Arthur Jaffe, and Thomas Spencer.
\newblock A convergent expansion about mean field theory. i. the expansion.
\newblock {\em Annals of Physics 101, 610-630}, 1976.

\bibitem{GJS_1976_II}
James Glimm, Arthur Jaffe, and Thomas Spencer.
\newblock A convergent expansion about mean field theory. ii. convergence of
  the expansion.
\newblock {\em Annals of Physics 101, 631--669}, 1976.

\bibitem{GW_1973}
David~J. Gross and Frank Wilczek.
\newblock Ultraviolet behavior of non-abelian gauge theories.
\newblock {\em Physical Review Letters 30, 1343-1346}, 1973.

\bibitem{GRS_1975_I}
F.~Guerra, L.~Rosen, and B.~Simon.
\newblock The p($\varphi$)$_2$ euclidean quantum field theory as classical
  statistical mechanics. i.
\newblock {\em Annals of Mathematics 101, 111--189}, 1975.

\bibitem{GRS_1975_II}
F.~Guerra, L.~Rosen, and B.~Simon.
\newblock The p($\varphi$)$_2$ euclidean quantum field theory as classical
  statistical mechanics. ii.
\newblock {\em Annals of Mathematics 101, 191--259}, 1975.

\bibitem{GRS_1975_GAP}
Francesco Guerra, Lon Rosen, and Barry Simon.
\newblock Correlation inequalities and the mass gap in p($\varphi$)$_2$. iii.
  mass gap for a class of strongly coupled theories with nonzero external
  field.
\newblock {\em Communications in Mathematical Physics 41, 19-32}, 1975.

\bibitem{HAAG_1992}
Rudolf Haag.
\newblock Local quantum physics: Fields, particles, algebras.
\newblock {\em Springer, first edition}, 1992.

\bibitem{HAMER_1989}
C.~J. Hamer.
\newblock Hamiltonian strong coupling expansions for glueball masses in su(3).
\newblock {\em Phys. Lett. B 224, 339-342}, 1989.

\bibitem{HIP_1986}
C.~J. Hamer, A.~C. Irving, and T.~E. Preece.
\newblock Cluster expansion approach to non-abelian lattice gauge theory in
  (3+1)d (ii). su(3).
\newblock {\em Nucl. Phys. B270, 553}, 1986.

\bibitem{HSB_2000}
C.~J. Hamer, M.~Samaras, and R.~J. Bursill.
\newblock Green's function monte carlo study of su(3) lattice gauge theory in
  (3+1)d.
\newblock {\em Phys. Rev. D 62, 074506}, 2000.

\bibitem{HAZRA_2026}
Tamaghna Hazra.
\newblock New class of exactly flat topological bands - compact localised
  states protected by local graph topology.
\newblock {\em SciPost Physics (submission)}, 2026.

\bibitem{THOOFT_1974}
G.'t Hooft.
\newblock A planar diagram theory for strong interactions.
\newblock {\em Nuclear Physics B 72, 461-473}, 1974.

\bibitem{IMBRIE_1981_I}
John~Z. Imbrie.
\newblock Phase diagrams and cluster expansions for low temperature p(phi)\_2
  models. i. the phase diagram.
\newblock {\em Communications in Mathematical Physics 82, 261-304}, 1981.

\bibitem{IMBRIE_1981_II}
John~Z. Imbrie.
\newblock Phase diagrams and cluster expansions for low temperature p(phi)\_2
  models. ii. the schwinger functions.
\newblock {\em Communications in Mathematical Physics 82, 305-343}, 1981.

\bibitem{IH_1984}
A.~C. Irving and C.~J. Hamer.
\newblock Methods in hamiltonian lattice field theory (ii). linked-cluster
  expansions.
\newblock {\em Nucl. Phys. B230, 361}, 1984.

\bibitem{IZ_1980}
Claude Itzykson and Jean-Bernard Zuber.
\newblock Quantum field theory.
\newblock {\em McGraw-Hill, New York}, 1980.

\bibitem{JAFFE_2000_CQFT}
Arthur Jaffe.
\newblock Constructive quantum field theory.
\newblock {\em Mathematical Physics 2000, pp.111-127, Imperial College Press /
  World Scientific}, 2000.

\bibitem{JW_2006}
Arthur Jaffe and Edward Witten.
\newblock Quantum yang-mills theory.
\newblock {\em Clay Mathematics Institute, official 14-page problem statement},
  2006.

\bibitem{KRS_2023}
Saurabh~V. Kadam, Indrakshi Raychowdhury, and Jesse~R. Stryker.
\newblock Loop-string-hadron formulation of an su(3) gauge theory with
  dynamical quarks.
\newblock {\em Phys. Rev. D 107, 094513}, 2023.

\bibitem{KLEIN_ROSENBERGER_2020_WKB}
Markus Klein and Elke Rosenberger.
\newblock The tunneling effect for schr\"odinger operators on a vector bundle,
  2020.
\newblock arXiv preprint arXiv:2005.13852v1.

\bibitem{KPS_1981}
J.~B. Kogut, R.~B. Pearson, and J.~Shigemitsu.
\newblock The string tension, confinement and roughening in su(3) hamiltonian
  lattice gauge theory.
\newblock {\em Phys. Lett. B 98, 63}, 1981.

\bibitem{KSS_1976}
J.~B. Kogut, D.~K. Sinclair, and L.~Susskind.
\newblock A quantitative approach to low-energy quantum chromodynamics.
\newblock {\em Nucl. Phys. B114, 199-236}, 1976.

\bibitem{KS_1975}
J.~B. Kogut and L.~Susskind.
\newblock Hamiltonian formulation of wilson's lattice gauge theories.
\newblock {\em Phys. Rev. D 11, 395-408}, 1975.

\bibitem{KP_1986}
R.~Kotecky and D.~Preiss.
\newblock Cluster expansion for abstract polymer models.
\newblock {\em Communications in Mathematical Physics 103, 491-498}, 1986.

\bibitem{LEMOINE_2026}
Thibaut Lemoine.
\newblock Universal dualities for wilson loops in lattice yang-mills, 2026.
\newblock arXiv:2604.16252 [math-ph].

\bibitem{LLL_2006}
Mushtaq Loan, Xiang-Qian Luo, and Zhi-Huan Luo.
\newblock Monte carlo study of glueball masses in the hamiltonian limit of
  su(3) lattice gauge theory.
\newblock {\em Int. J. Mod. Phys. A 21, 2905}, 2006.

\bibitem{MRS_1993}
Jacques Magnen, Vincent Rivasseau, and Roland S\'en\'eor.
\newblock Construction of ym4 with an infrared cutoff.
\newblock {\em Communications in Mathematical Physics 155, 325-383}, 1993.

\bibitem{MS_1980_Y3}
Jacques Magnen and Roland S\'en\'eor.
\newblock Yukawa quantum field theory in three dimensions (y3).
\newblock {\em Annals of the New York Academy of Sciences 337, 13-43; Third
  International Conference on Collective Phenomena}, 1980.

\bibitem{MALDACENA_1998}
Juan~M. Maldacena.
\newblock The large n limit of superconformal field theories and supergravity.
\newblock {\em Advances in Theoretical and Mathematical Physics 2, 231-252},
  1998.

\bibitem{MR_1976_CRITICAL}
Oliver~A. McBryan and Jay Rosen.
\newblock Existence of the critical point in phi4 field theory.
\newblock {\em Communications in Mathematical Physics 51, 97-105}, 1976.

\bibitem{MMM_2019_CURVATURE}
Vincent Moncrief, Antonella Marini, and Rachel Maitra.
\newblock Orbit space curvature as a source of mass in quantum gauge theory,
  2019.
\newblock arXiv:1809.06318v2; Annals of Mathematical Sciences and Applications
  4(2).

\bibitem{MONDAL_2023_GEOMETRIC}
Puskar Mondal.
\newblock A geometric approach to the yang-mills mass gap, 2023.
\newblock arXiv:2301.06996v7.

\bibitem{MORNINGSTAR_2025}
Colin Morningstar.
\newblock Update on glueballs.
\newblock {\em PoS LATTICE2024 (2025) 004}, 2025.

\bibitem{MP_1999}
Colin~J. Morningstar and Mike Peardon.
\newblock The glueball spectrum from an anisotropic lattice study.
\newblock {\em Phys. Rev. D 60, 034509}, 1999.

\bibitem{MUNSTER_1981}
G.~Munster.
\newblock Strong coupling expansions for the mass gap in lattice gauge
  theories.
\newblock {\em Nucl. Phys. B190, 439}, 1981.

\bibitem{MUNSTER_1985_TM}
G.~Munster.
\newblock Effective transfer matrix for low-lying glueball states in lattice
  gauge theory.
\newblock {\em Nucl. Phys. B256, 67-76}, 1985.

\bibitem{NSY_2019}
B.~Nachtergaele, R.~Sims, and A.~Young.
\newblock Quasi-locality bounds for quantum lattice systems. part i.
  lieb-robinson bounds, quasi-local maps, and spectral flow automorphisms.
\newblock {\em J. Math. Phys. 60, 061101}, 2019.

\bibitem{NELSON_1973_MARKOFF}
Edward Nelson.
\newblock Quantum fields and markoff fields.
\newblock {\em Proceedings of Symposia in Pure Mathematics 23, 413-420,
  American Mathematical Society}, 1973.

\bibitem{OBZ_1985}
K.~H. O'Brien and J.-B. Zuber.
\newblock Strong coupling expansion of large-n qcd and surfaces.
\newblock {\em Nucl. Phys. B253, 621-634}, 1985.

\bibitem{ORAIFEARTAIGH_1997}
Lochlainn O'Raifeartaigh.
\newblock The dawning of gauge theory.
\newblock {\em Princeton University Press}, 1997.

\bibitem{OS_1973}
Konrad Osterwalder and Robert Schrader.
\newblock Axioms for euclidean green's functions.
\newblock {\em Communications in Mathematical Physics 31, 83-112}, 1973.

\bibitem{OS_1975}
Konrad Osterwalder and Robert Schrader.
\newblock Axioms for euclidean green's functions ii.
\newblock {\em Communications in Mathematical Physics 42, 281-305; appendix by
  Stephen Summers}, 1975.

\bibitem{OS_SEILER_1978}
Konrad Osterwalder and Erhard Seiler.
\newblock Gauge field theories on a lattice.
\newblock {\em Annals of Physics 110, 440-471}, 1978.

\bibitem{OS_SENEOR_1976}
Konrad Osterwalder and Roland S\'en\'eor.
\newblock The scattering matrix is non-trivial for weakly coupled p(phi)\_2
  models.
\newblock {\em Helvetica Physica Acta 49, 525-535}, 1976.

\bibitem{POLITZER_1973}
H.~David Politzer.
\newblock Reliable perturbative results for strong interactions?
\newblock {\em Physical Review Letters 30, 1346-1349}, 1973.

\bibitem{RIVASSEAU_1991}
Vincent Rivasseau.
\newblock From perturbative to constructive renormalization.
\newblock {\em Princeton University Press}, 1991.

\bibitem{SALAM_1968}
Abdus Salam.
\newblock Weak and electromagnetic interactions.
\newblock {\em Elementary Particle Theory, Proceedings of the Eighth Nobel
  Symposium, Lerum, Sweden, pp.367-377}, 1968.

\bibitem{SCHIERHOLZ_1988}
G.~Schierholz.
\newblock Glueball masses in su(3) lattice gauge theory.
\newblock {\em Nucl. Phys. B (Proc. Suppl.) 4, 11-20}, 1988.

\bibitem{SZH_1997}
D.~Schutte, W.-H. Zheng, and C.~J. Hamer.
\newblock The coupled cluster method in hamiltonian lattice field theory.
\newblock {\em Phys. Rev. D 55, 2974}, 1997.

\bibitem{SW_1994}
Nathan Seiberg and Edward Witten.
\newblock Electric-magnetic duality, monopole condensation, and confinement in
  n=2 supersymmetric yang-mills theory.
\newblock {\em Nuclear Physics B426, 19-52}, 1994.

\bibitem{SEILER_1975_Y2}
Erhard Seiler.
\newblock Schwinger functions for the yukawa model in two dimensions with
  space-time cutoff.
\newblock {\em Communications in Mathematical Physics 42, 163-182}, 1975.

\bibitem{SEO_1982}
K.~Seo.
\newblock Glueball mass estimate by strong coupling expansion in lattice gauge
  theories.
\newblock {\em Nucl. Phys. B209, 200-216}, 1982.

\bibitem{SIMON_1983_DISCRETE}
Barry Simon.
\newblock Some quantum operators with discrete spectrum but classically
  continuous spectrum.
\newblock {\em Annals of Physics 146, 209-220}, 1983.

\bibitem{SINGER_1981_ORBIT}
Isadore~M. Singer.
\newblock The geometry of the orbit space for non-abelian gauge theories.
\newblock {\em Physica Scripta 24, 817-820}, 1981.

\bibitem{STREATER_WIGHTMAN_1964}
Raymond~F. Streater and Arthur~S. Wightman.
\newblock Pct, spin and statistics, and all that.
\newblock {\em W. A. Benjamin, New York}, 1964.

\bibitem{SYMANZIK_1969}
Kurt Symanzik.
\newblock Euclidean quantum field theory.
\newblock {\em Local Quantum Theory, R. Jost (ed.), Academic Press, New York,
  pp.152-226}, 1969.

\bibitem{T_HOOFT_1979}
G.~'t~Hooft.
\newblock A property of electric and magnetic flux in nonabelian gauge
  theories.
\newblock {\em Nucl. Phys. B153, 141-160}, 1979.

\bibitem{UELTSCHI_2004_CLUSTER}
Daniel Ueltschi.
\newblock Cluster expansions and correlation functions.
\newblock {\em Moscow Mathematical Journal 4, 511-522}, 2004.

\bibitem{URZUA_2019}
Alejandro~R. Urzua, Iran Ramos-Prieto, Manuel Fernandez-Guasti, and Hector~M.
  Moya-Cessa.
\newblock Solution to the time-dependent coupled harmonic oscillators
  hamiltonian with arbitrary interactions.
\newblock {\em Quantum Reports 1, 82-90}, 2019.

\bibitem{WEINBERG_1967}
Steven Weinberg.
\newblock A model of leptons.
\newblock {\em Physical Review Letters19, 1264-1266}, 1967.

\bibitem{WILSON_1977_LECTURES}
Kenneth~G. Wilson.
\newblock Quarks and strings on a lattice.
\newblock {\em New Phenomena in Subnuclear Physics, Part A, The Subnuclear
  Series13, pp.69-142, Plenum}, 1977.

\bibitem{WITTEN_1994_FOUR_MANIFOLD}
Edward Witten.
\newblock Supersymmetric yang-mills theory on a four-manifold.
\newblock {\em Journal of Mathematical Physics35, 5101-5135}, 1994.

\bibitem{YANG_MILLS_1954}
Chen~Ning Yang and Robert~L. Mills.
\newblock Conservation of isotopic spin and isotopic gauge invariance.
\newblock {\em Physical Review96, 191-195}, 1954.

\bibitem{YAROTSKY_2004_GS}
D.~A. Yarotsky.
\newblock Ground states in relatively bounded quantum perturbations of
  classical lattice systems, 2004.
\newblock arXiv preprint.

\bibitem{YAROTSKY_2004_QP}
D.~A. Yarotsky.
\newblock Quasi-particles in weak perturbations of non-interacting quantum
  lattice systems, 2004.
\newblock arXiv preprint.

\end{thebibliography}
% ---- end inlined bibliography ------------------------------------

\end{document}
